\documentclass[10pt,a4paper,twoside,bg=print,twocolumn,openany,nodeprecatedcode]{dndbook}
\usepackage[utf8]{inputenc}
\usepackage{tabulary}
\usepackage[normalem]{ulem}
\usepackage{color,soul}
\usepackage{float}
\usepackage{caption}
\usepackage{wrapfig}
\usepackage{tocloft}
\usepackage[hidelinks]{hyperref}
\raggedbottom
\extrafloats{2000}
\cftsetindents{section}{0em}{0em}
\cftsetindents{subsection}{0em}{0em}
\cftsetindents{subsubsection}{0.2em}{0em}
\setcounter{tocdepth}{1}
\definecolor {mytable} {HTML} {f4edd8}
\definecolor {mycomment} {HTML} {ffffff}
\definecolor {mysidebar} {HTML} {d0cfc2}
\definecolor {myreadaloud} {HTML} {d9e8f1}
\definecolor {myreadaloudborder} {HTML} {5a6173}
\definecolor {mypagenum} {HTML} {2a2d2e}
\DndSetComplexThemeColor[mytable][mycomment][mysidebar][myreadaloud][myreadaloudborder][titlered][titlegold][titlegold]\setlist[description]{nolistsep,listparindent=3pt,topsep=6pt}
\begin{document}
\tableofcontents
\newpage
\chapter{Table of Contents}
\textit{Tome of Beasts 3 Lairs} contains 23 short-form adventures, that feature new creatures from the Tome of Beasts 3. Each adventure is intended to take one or two sessions to complete and includes a map, adventure hooks, and treasures for your adventurers.

\DndDropCapLine{T}{}o run each of these adventures, you'll need the core rulebooks of the 5th edition of the world's first roleplaying game or the \textit{System Reference Document 5.1}. Creatures whose names appear in \textbf{bold} without a page or book reference can be found in the \textit{System Reference Document 5.1} or in those rulebooks. Similarly, spell and magic item names which appear in \textit{italics} without a page or book reference can be found in those rulebooks or in the \textit{System Reference Document 5.1}. In addition, you will need the Tome of Beasts 3 to run the encounters in these adventures.


Additional creature portraits and tokens can be found here.

\textit{Tome of Beasts 3 Lairs} includes the following adventures, listed by character level:

\subsection{Proudheart's Predator Problem}

\begin{DndTable}{XX}

\textbf{By Eytan Bernstein} & \textbf{(1st-level characters)}\\
\end{DndTable}

Farmer Grumbo Proudheart has a problem. A mysterious entity has been eating his prized sheep, winners of the Halfling Wool-Off the past ten years. At his wit's end, he searches for heroes to uncover the culprit and save his farm.

This adventure features the \textbf{herd skulker} from Tome of Beasts 3.

\subsection{Flight of the Dromedaries}

\begin{DndTable}{XX}

\textbf{By Benjamin Eastman} & \textbf{(1st-level characters)}\\
\end{DndTable}

A giant dragon soars overhead startling the camels of a caravan the heroes were hired to guard. The terrified beasts flee deep into a dangerous desert. Can the heroes survive the desert's dangers and round up the terrified mounts?

This adventure features the \textbf{desert slime} from Tome of Beasts 3.

\subsection{Furor in the Farmyard}

\begin{DndTable}{XX}

\textbf{By Kelly Pawlik} & \textbf{(2nd-level characters)}\\
\end{DndTable}

A notorious bandit and his crew have holed up at the Hundedel family farm following his daring escape from the local magistrate's cells. The criminals are holding the youngest Hundedel captive and have deceived the farm's dragonettes into helping them fortify and protect the farmhouse. Can the heroes free the young farm boy before he is harmed?

This adventure features the \textbf{barnyard dragonette}, \textbf{shovel dragonette}, and \textbf{razorback crab} from Tome of Beasts 3.

\subsection{Salvage at Lonely Cove}

\begin{DndTable}{XX}

\textbf{By Jeff Lee} & \textbf{(2nd-level characters)}\\
\end{DndTable}

A ship was sunk during a recent storm, and flotsam from the wreck has washed up in Lonely Cove, a hidden beach at the bottom of steep shoreline cliffs. Can the heroes collect the valuable salvage and survive the cove's inhabitants?

This adventure features the \textbf{beach weird}, \textbf{guardian archaeopteryx}, \textbf{moon weaver}, \textbf{pescavitus}, and \textbf{talus flow} from Tome of Beasts 3.

\subsection{Automated Banditry}

\begin{DndTable}{XX}

\textbf{By Basheer Ghouse} & \textbf{(3rd-level characters)}\\
\end{DndTable}

A group of bandits has stolen a small hoard of constructs, turning their new mechanical servants into a small army. The roads are no longer safe, and merchants are desperate for heroes who can end the bandits' reign.

This adventure features the \textbf{atavist}, \textbf{animated instrument}, \textbf{clockwork armadillo}, \textbf{clockwork pugilist}, and \textbf{harvest horse} from Tome of Beasts 3.

\subsection{Evil Rises}

\begin{DndTable}{XX}

\textbf{By Jeff Lee} & \textbf{(3rd-level characters)}\\
\end{DndTable}

The \textbf{bakery drake}, Amaranth, is the darling of the local nobility. Her baked goods are in high demand by the city's wealthy citizens, and it was a shock to all when she came under suspicion of poisoning several notable personages. Can the heroes brave the now-dangerous bakery, uncover the truth behind the poisoning, and bring justice to the drake?

This adventure features the \textbf{bakery drake} from Tome of Beasts 3.

\subsection{Four-Part Harmony}

\begin{DndTable}{XX}

\textbf{By Sebastian Rombach} & \textbf{(3rd-level characters)}\\
\end{DndTable}

Rumors of ghosts in the old haunted bardic college didn't stop the orphan Budanyek from breaking in. But now he's been gone for days, and his friends swear they hear terrifying sounds coming from the old college. Can a plucky group of heroes brave the tremulous terrors inside to find the missing Budanyek?

This adventure features the \textbf{atavist}, \textbf{animated instrument}, \textbf{animated quartet}, \textbf{clockwork conductor}, \textbf{clockwork pugilist}, and \textbf{origami golem} from Tome of Beasts 3.

\subsection{Harvest Fangs}

\begin{DndTable}{XX}

\textbf{By Benjamin Eastman} & \textbf{(3rd-level characters)}\\
\end{DndTable}

In the isolated farming village of Sycamore Crossing, this year's harvest is inexplicably failing. Faced with starvation, the villagers have fallen under the fell sway of a devious cult that promises a full harvest if the village regularly sacrifices people to the soil they till. The heroes must stop the cult and save the village before too many more travelers are abducted and sacrificed.

This adventure features the \textbf{tripwire patch} from Tome of Beasts 3.

\subsection{The Pavilion of Whispers and Wonderment}

\begin{DndTable}{XX}

\textbf{By Brian Suskind} & \textbf{(4th-level characters)}\\
\end{DndTable}

The Maestro, a fey with delusions of grandeur, operates out of a magnificent, extraplanar palace. An information hoarder, The Maestro collects uncomfortable secrets, scandals, and slander, using them to amass great power and wealth. The heroes must recover a particularly important secret kept by The Maestro, facing the odd inhabitants of the whimsical, dangerous pavilion along the way.

This adventure features the \textbf{doppelixir}, \textbf{lobe lemur}, \textbf{hippopotamus}, \textbf{moppet}, \textbf{swarm of penguins}, and \textbf{veritigibbet} from Tome of Beasts 3.

\subsection{Rainforest Reckoning}

\begin{DndTable}{XX}

\textbf{By Richard Green} & \textbf{(4th-level characters)}\\
\end{DndTable}

Several weeks ago, villagers from the jungle village of Bukofa fell victim to mild, mischievous pranks while clearing trees for farmland. Now, events have taken a sinister turn and something lurks in the jungle, killing Bukofa's inhabitants. The village elder needs heroes to travel to the fey's grove and end the threat to the village.

This adventure features the \textbf{aziza}, \textbf{giant walking stick}, \textbf{leavesrot ooze}, \textbf{monkey's bane vine}, \textbf{moon weaver}, \textbf{musk deer}, \textbf{rainforest ogre}, and \textbf{tripwire patch} from Tome of Beasts 3.

\subsection{Safe and Sound}

\begin{DndTable}{XX}

\textbf{By Sarah Madsen} & \textbf{(4th-level characters)}\\
\end{DndTable}

An unusually large and ferocious cat has recently claimed the area outside of the town of Bay's Burrow as its territory. The area is bisected by a major trade route, and unfortunately the territorial beast has been disrupting caravans, killing horses, and generally making a nuisance of itself. Heroes are needed to kill the beast or encourage it to find a new home away from Bay's Burrow.

This adventure features the \textbf{catamount} from Tome of Beasts 3.

\subsection{Lord Gorgo's Keep}

\begin{DndTable}{XX}

\textbf{By Jeff Lee} & \textbf{(5th-level characters)}\\
\end{DndTable}

Squatting atop a lonely peak surrounded by a dark wood, Spiteful Keep is the home of Lord Gorgo, a powerful ogre said to craft amazing magical items. Lord Gorgo's allies and servants fill the keep, discouraging trespassers and gathering materials for the artisan. Heroes willing to deal with the cunning ogre might be able to barter for magic items or take them by force, but the items might not be what they appear...

This adventure features the \textbf{cunning artisan ogre}, the \textbf{daeodon}, the \textbf{dokkaebi}, and the \textbf{kobold ettin} from Tome of Beasts 3.

\subsection{Redoubt of the Perfidious Guru}

\begin{DndTable}{XX}

\textbf{By Kelly Pawlik} & \textbf{(5th-level characters)}\\
\end{DndTable}

There is a sanctuary hidden in the mountains where it is said one can learn the secrets of inner peace and perfection. The master of this secret place is said to be able to see into a person's very soul and divine their future from it, but something isn't quite right. Can the heroes uncover the truth and put an end to the fiendish machinations of the sanctuary's elusive leader?

This adventure features the \textbf{alpine creeper}, \textbf{cloudhoof assassin}, and \textbf{lesser infernal tutor} from Tome of Beasts 3.

\subsection{Incident at Wrackwater}

\begin{DndTable}{XX}

\textbf{By Robert Fairbanks} & \textbf{(6th-level says 5th-level.} characters)}\\
\end{DndTable}

The fog-bound village of Wrackwater Cove thrives on the trade brought into its small harbor by the light of a lighthouse on a nearby spit of land. Recently, the beacon went out and the village's two search parties haven't returned. A ship with much-needed supplies is due any day, and the villagers seek heroes to find the missing patrols and light the beacon.

This adventure features the \textbf{bilge gremlin}, \textbf{bilge gremlin bosun}, \textbf{giant flea}, \textbf{gullkin}, \textbf{gullkin hunter}, and \textbf{muraenid} from Tome of Beasts 3.

\subsection{The Lost Forge}

\begin{DndTable}{XX}

\textbf{By Jonathan Miley} & \textbf{(6th-level says 5th-level.} characters)}\\
\end{DndTable}

A dwarven fort high in the mountains was abandoned decades ago when a dragon's attack left much of it in ruins. The once-mighty forge stayed dormant for years, but now the forge blazes to life under the care of twin giants and their followers, who raid dwarven caravans carrying ore and ale. Locals turn to traveling heroes for help in dethroning the giants and bringing peace back to their area.

This adventure features the dwarf \textbf{pike guard}, dwarf \textbf{pike guard captain}, \textbf{thursir armorer} giant, \textbf{thursir hearth priestess} giant, and \textbf{alleybasher ogre} from Tome of Beasts 3.

\subsection{The Monstrous Machinations of Margrave Millard}

\begin{DndTable}{XX}

\textbf{By Basher Ghouse} & \textbf{(7th-level characters)}\\
\end{DndTable}

The eccentric Margrave Millard has long considered himself a pioneer in magically created beasts, often parading his creations as short-lived show pets. However, he has recently begun creating more monstrous creatures and testing their murderous prowess on kidnapped citizens. Can the heroes stop the Margrave before he releases his creations to wreak havoc on the surrounding lands?

This adventure features the \textbf{capybear}, \textbf{dire owlbear}, \textbf{grolar bear}, \textbf{grolar bear alpha}, \textbf{razorfeather raptor}, and \textbf{slithy tove} from Tome of Beasts 3.

\subsection{The Green Sanctum}

\begin{DndTable}{XX}

\textbf{By Mike Welham} & \textbf{(7th-level characters)}\\
\end{DndTable}

An eccentric scholar spent her life studying exotic plants in an arboretum on a remote island. The scholar has since died and other scholars now want a particularly important plant she studied. Can the heroes survive the abandoned arboretum's deadly plants and convince its guardian to part with the exotic plant?

This adventure features the \textbf{alpine creeper}, \textbf{ice willow}, and \textbf{ice golem} from Tome of Beasts 3.

\subsection{The Salons of Mother Celeste}

\begin{DndTable}{XX}

\textbf{By Robert Fairbanks} & \textbf{(8th-level characters)}\\
\end{DndTable}

Days ago, the crew of a local fishing boat pulled up the remains of a sailor thought to belong to an overdue, missing merchant vessel. The nearby town's fishermen have since found the location of the wreck, and the local authorities put out a call for heroes willing to brave the water's depths and rescue what remains of the vessel's cargo.

This adventure features the \textbf{bilge gremlin}, \textbf{bilge gremlin bosun}, \textbf{breakwater troll}, \textbf{brine hag}, \textbf{dire lionfish}, and \textbf{wrackwraith} from Tome of Beasts 3.

\subsection{Chaos at the Caldera}

\begin{DndTable}{XX}

\textbf{By Eytan Bernstein} & \textbf{(8th-level characters)}\\
\end{DndTable}

The town of Dorma sits on the fertile soil next to Mount Magra, a dormant volcano. Recently, portals have opened in the heart of the mountain, expelling violent creatures from other planes. Can the heroes vanquish these threats before the the volcano erupts and destroys the town?

This adventure features the \textbf{ibexian} and \textbf{obsidian ophidian} from Tome of Beasts 3.

\subsection{The Eclipsed Chapel}

\begin{DndTable}{XX}

\textbf{By Sebastian Rombach} & \textbf{(9th-level characters)}\\
\end{DndTable}

A painter named Pol seeks escort to an abandoned chapel. Convinced he's met an angel who wants him to paint a magic mural there in just one day, Pol's request may require a leap of faith. Is there a group of heroes who will stand vigil while Pol works his miracle?

This adventure features the \textbf{apostle}, \textbf{copperkill slime}, \textbf{gearmass}, \textbf{leavesrot ooze}, \textbf{sinoper ooze}, \textbf{stained-glass moth}, and \textbf{gryllus swarm} from Tome of Beasts 3.

\subsection{The Twisted Sanctuary}

\begin{DndTable}{XX}

\textbf{By Brian Suskind} & \textbf{(10th-level characters)}\\
\end{DndTable}

Down the street, among the shops and inns, sits the Veiled Mannequin, a tailor's shop for discerning customers. Unbeknownst to most, the tailor and her assistant guard the secret entrance to the Sanctuary of the Twisted Claw, a rakshasa safehouse hidden in the heart of the city. Can the heroes uncover the secret and stop the rakshasa before they cause irreparable damage to the city's denizens?

This adventure features the \textbf{apostle}, \textbf{atavist}, \textbf{myrmidon rakshasa}, \textbf{pustakam rakshasa}, \textbf{servitor rakshasa}, and \textbf{slayer rakshasa} from Tome of Beasts 3.

\subsection{A Midnight Ride}

\begin{DndTable}{XX}

\textbf{By Sarah Madsen} & \textbf{(11th-level characters)}\\
\end{DndTable}

A strange, black horse draped in silver chains has been seen near a small farming village. The locals claim it is a fey creature that will run off with their sons and daughters. More concerning, however, is the group of cultists who recently arrived to capture the creature and use it to gain access to other planes. The farmers need heroes to chase off the cultists, and they promise a bonus if the heroes can do something about the fey as well.

This adventure features the \textbf{leashed lesion}, \textbf{psychophant} cultist, and \textbf{púca} from Tome of Beasts 3.

\subsection{Down in the Dark of The Deep King's Domain}

\begin{DndTable}{XX}

\textbf{By Brian Suskind} & \textbf{(12th-level characters)}\\
\end{DndTable}

Down in the depths of the world, among twisting natural caverns where bioluminescent fungi shine as motes of light in the darkness, lies the lair of the Deep King. It hunts the world below, feeding on the spiritual essence of its victims and dragging their bodies back to its lair. Many have fallen in an attempt to put an end to this fell alien intelligence; can the heroes succeed where so many others have failed?

This adventure features the \textbf{cave sovereign} and \textbf{duskwilt} from Tome of Beasts 3.

\chapter{Proudheart's Predator Problem}

\begin{DndTable}{c}

\textbf{An adventure for four to five characters of 1st level}\\
\end{DndTable}

\section{Adventure Background}
\DndDropCapLine{E}{}very year, halfling farmer Grumbo Proudheart enters his animals into the Halfling Wool-Off, a competition that attracts fellow farmers from throughout the region. The competition includes five categories: long wool, medium wool, fine wool, goat's wool, and lamb's wool. Grumbo has won the lamb's wool competition for the last three years. Recently, he's developed his goat breeding program and could now also win in that category.


The Wool-Off is in four days, and Grumbo is frantic, grooming his animals, sharpening his sheers, and loading his wagons. All was running smoothly until three nights ago when an unseen predator began menacing his flock.

Grumbo has had the bad luck of attracting the attention of a \textbf{herd skulker}, a dangerous and crafty wolflike predator that stalks herd animals, shapeshifting into one of their kind and hiding amongst them until it can get close enough to feast. It's been stalking the sheep, goats, chickens, and horses of Grumbo's farm, radiating an aura of calm that lulls the animals into a false sense of security.

At his wit's end, Grumbo seeks the help of the PCs in rooting out the menace and protecting his animals until the end of the week when he leaves for the Wool-Off. It would pain him greatly if harm came to any of his animals.

\section{Adventure Hook}
The following can entice the PCs into helping Grumbo:


\begin{itemize}
\item \textbf{\textit{Market Meeting.}} The PCs run into Grumbo as he's purchasing supplies at the local market. Seeing their weapons and armor, Grumbo begs them to come help him. He didn't want to leave his animals alone for the day, as only his dogs remain to protect them, but he had little choice.
\item \textbf{\textit{Fair Competition.}} Lucy Porridge, a rival competitor in the Wool-Off (and Grumbo's cousin), heard about Grumbo's troubles. She wants to win the Wool-Off fairly and that's only possible if Grumbo brings his best. Lucy runs into the PCs on the road and explains the situation, offering the same reward as explained below.
\end{itemize}

Whichever hook you choose, if the PCs ask for reward, they are offered 50 gp if they can rid the farm of whatever is hunting Grumbo's animals, plus another 50 if they keep all of his animals alive and uninjured in the process.

\section{Grumbo and His Farm}
\begin{DndReadAloud}{}
Grumbo's farm occupies an acre of land, situated amid rolling pastures, fields, and trees. Every building, plot, and pen is pristine, recently painted and meticulously maintained. Carefully tended flower boxes adorn the farmstead with sunflowers, tulips, and daisies. Grumbo beams as he admires his farm, but then his eyes fill with tears as he looks at the baby sheep and goats playing in their pens. He pulls out a handkerchief, wiping his eyes. "Please, you're my only hope. It would kill me if that monster ate my animals. I couldn't bear it."

\end{DndReadAloud}

Grumbo Proudheart (LG halfling \textbf{commoner}) is a perpetually nervous halfling in his early 40s with thinning gray hair and sideburns, a substantial belly, a strong musculature from decades of farm labor, and a quick laugh. His pride and joy are his sheep and goats, and he frets constantly about their safety. He is a lifelong bachelor, and he sent all of his farmhands away shortly before the PCs arrived for fear of the creature harming them. He is alone now, except for his dogs and livestock.

\subsubsection{The Predator}
The creature plaguing Grumbo's farm is a \textbf{herd skulker}.

\subsubsection{The Farmstead}
Grumbo's vast farmstead is comprised of animal pens, a barn, a chicken coop, a vegetable garden, his home, and his pasture. Grumbo gives the PCs a tour and grants them complete access to all the buildings and lands, including his cottage. (See Areas 1–6 for more information.)

\subsubsection{Grumbo's Prize Animals}
Grumbo cares deeply for all of his flock, but he has four prized animals he loves more than life itself. He has three baby sheep (lambs): Lady Lanolyn of Lambshire, Duke Fuzzington von Felt, and Prince Lambert, plus a baby goat (kid) named Sir Cashimir of the Order of the Fleece.

\subsubsection{Grumbo's Dogs}
Grumbo has an eager, vigilant wolfhound named Fancy. She's massive, easily triple Grumbo's own weight. She could scare off any normal predator, but the herd skulker is far from normal. If the PCs get into in trouble while fighting the creature, Fancy comes to their aid.

\textbf{Fancy} uses the statistics of a \textbf{wolf}, but replace Pack Tactics with the following feature:

\subparagraph{Nimble Protector}
Fancy can take the Dash or Disengage action as a bonus action on each of her turns.

Grumbo also has a trio of sheepdogs—Lucy, Dart, and Priscilla—that herd his sheep and goats. Grumbo has been keeping them in the house or by his side for their safety.

\subsubsection{Overnight Accommodations}
Grumbo doesn't have any extra beds available, but he welcomes the PCs to stay with him, allowing them to sleep on the floor, using the many pillows and rugs scattered about his living space.

If they prefer, they are welcome to stay in Grumbo's barn instead; if so, he puts down fresh hay in the many empty stalls, though the barn was essentially immaculate already. The PCs will have to share the barn with a cow, two draft horses, and a pony. (See 4.)

Wherever the PCs decide to spend the night, whenever Grumbo cooks a meal for himself (which is early and often), he brings some for each of the PCs as they stand guard.

\subsection{What Grumbo Knows}
If the PCs question Grumbo about the creature and/or its predations, he'll gladly share anything he knows:


\begin{itemize}
\item All of Grumbo's preparations for the Wool-Off were running smoothly until three nights ago when something began stalking one of his prized lambs, Duke Fuzzington.
\item Fancy scared it off before he could get a look at it.
\item Grumbo searched the farm afterward and even hired some locals to investigate, but they found nothing.
\item Last night, the creature made a play for another lamb, Prince Lambert. Grumbo heard a strange sound outside his window and rushed to investigate. The creature had cornered Lambert and was about to pounce when Grumbo blew a bullhorn, startling it, and the creature fled.
\item This time, Grumbo caught sight of it, but he couldn't believe his eyes; it looked just like his other lambs, but its movements were aggressive and not-at-all lamblike.
\item Grumbo doesn't know if the creature is a threat during the day or not, but both incidents happened at night.
\item Despite the threat of the creature, Grumbo has no choice but to continue opening the gate each morning because the animals need to eat. He refuses any suggestion that he do otherwise.
\end{itemize}

\begin{DndSidebar}[float=!b]{Grumbo's Menagerie}
Since this adventure has the PCs protecting a flock from a predator, they might wish to know how many animals Grumbo has in total. For ease of reference, his total number of animals is included here, with any names of particularly prized animals in parentheses:




\begin{itemize}
\item Chickens: 9 (Penelope)
\item Roosters: 1 (Cody)
\item Goats: 30 (Sir Cashimir of the Order of the Fleece [a prized kid])
\item Sheep: 45 (Lady Lanolyn of Lambshire, Duke Fuzzington von Felt, and Prince Lambert [all prized lambs])
\item Cows: 1 (Pinkle)
\item Horses: 2 (Binny and Winny)
\item Ponies: 1 (Theodora)
\item Sheepdogs: 3 (Lucy, Dart, and Priscilla)
\item Wolfhounds: 1 (\textbf{Fancy})
\end{itemize}



The totals above include all of the prized animals and any babies (i.e., lambs or kids).



\end{DndSidebar}

\begin{DndSidebar}[float=!b]{Avoiding Cruelty to Animals}
This adventure depicts a very real struggle facing farmers—predators hunting their livestock. While it's important to get the PCs to feel for Grumbo, the goal is never to graphically describe the death of animals. Give the PCs hints and chances to catch the herd skulker in the act before it harms the animals. If the PCs make a mistake or fail crucial checks, give them the opportunity to confront the herd skulker and scare it off rather than having it simply murder an animal. If that fails and the herd skulker captures or kills an animal, leave details to a minimum, so as to avoid distressing players. Perhaps the PCs find a bit of wool, a feather, or a single drop of blood. Check with your players about their comfort level with this and adjust accordingly.



This advice also applies to \textbf{Fancy} and the three sheepdogs, Lucy, Dart, and Priscilla. They can help the PCs if they are struggling, but it is strongly advised that you avoid allowing the \textbf{herd skulker} to harm them. For his part, Grumbo refuses to send the sheepdogs into danger; if there's a threat, he sends \textbf{Fancy}, but her Nimble Protector feature should allow her to help but avoid harm. In any case, the \textbf{herd skulker} does not go after her or the sheepdogs, focusing instead on the PCs or escape.



The \textbf{herd skulker} natural form bears similarities to that of wolves, but it should not be depicted as anything but the horrible monster that it is. If harming a creature that resembles a real animal would distress your players, modify the \textbf{herd skulker} true form to something altogether alien.



\end{DndSidebar}

An unlabelled version of the map can be found here.

\subsection{1. Sheep Pen}
\begin{DndReadAloud}{}
A short, wooden fence surrounds a large sheep pen that occupies the northwest quarter of Grumbo's farmstead. Lazing within are thirty sheep, though others wander in and out of the open twenty-foot-wide gate on the western side. More sheep graze in the pastures surrounding the farmstead. Several fuzzy lambs chase each other enthusiastically, the older sheep observing disinterestedly. One of the lambs approaches the fence and baas at you in greeting.

\end{DndReadAloud}

Grumbo keeps the gate open during the day and closes it at night, though a determined predator can easily leap the low fence. The sheep, too, will leap the fence at the first sign of a predator, which may necessitate the PCs locating them. The lamb that greets the PCs is Lady Lanolyn of Lambshire.

\subsection{2. Goat Pen}
\begin{DndReadAloud}{}
A large goat pen surrounded by a low, wooden fence occupies the northeast quarter of the farmstead. Milling about inside are twenty goats, and you can see more goats grazing in the pasture beyond the pen as well. The wide gate on the eastern side of the fence is open, allowing the goats the freedom to go in or out. Several downy kids gambol about within and are occasionally joined by the older goats. One baby goat approaches the fence when you near, sticking his nose out for a scritch or a treat.

\end{DndReadAloud}

Like the sheep pen, Grumbo keeps the gate open during the day and closes it at night. A determined predator, such as the \textbf{herd skulker}, can leap the easily low fence, as can all but the smallest of the goats. The kid seeking scritches is Sir Cashimir of the Order of the Fleece.

\subsection{3. Grumbo's House}
\begin{DndReadAloud}{}
Grumbo's small house is located in the southwestern portion of the farmstead. He welcomes you inside, insisting on providing tea and cookies for all. The interior of the cottage is immaculate and homey. Overstuffed pillows and plush rugs cover the floor, a comfy bed occupies the northwest corner, a table for six sits in the center, and a wood-burning stove graces the southeast corner. Custom cabinets and closets are built into the walls. Three merry sheepdogs lie in a tangle on a fluffy dog bed near the stove.

\end{DndReadAloud}

Grumbo painstakingly built this home. The door doesn't have a lock, as Grumbo has little fear of intruders. Even if the door is left wide open, the \textbf{herd skulker} does not enter.

\subsection{4. Barn}
\begin{DndReadAloud}{}
Grumbo leads you to a neatly painted red barn and opens the doors. Most of the stalls are empty, save for the four nearest the door, which contain a cow, two draft horses, and a pony. Grumbo points to the cow. "Pinkle here is a good cow. She's very patient." He waves his hand in the direction of the horses. "Binny and Winny are a tad lazy, but they do their jobs." Finally, he stops at the stall of the pony, who noses his hand, and hands her a carrot. "Theodora is the best pony that ever lived. I'm riding her to the fair."

\end{DndReadAloud}

Like everything else on Grumbo's property, the barn's been kept immaculate.

\subsection{5. Chicken Coop}
\begin{DndReadAloud}{}
Grumbo points out a substantial chicken coop east of the barn. "Happy chickens lay tasty eggs. That what I always say." He waves at a hen squawking at several others. "Penelope keeps the other hens in line." He nods at the lone rooster. "Cody mostly keeps to himself, though a couple of the hens like him quite a bit."

\end{DndReadAloud}

\subsection{6. Vegetable Garden}
\begin{DndReadAloud}{}
A large vegetable garden lines the southern edge of the farmstead, featuring neat rows of cucumbers, radishes, tomatoes, onions, zucchini, and turnips.

\end{DndReadAloud}

Grumbo loves gardening almost as much as raising animals. He doesn't need to say so; this bounteous garden and the numerous flower boxes scattered around the property make that clear.

\section{Catching the Herd Skulker}
Each night the PCs stay at Grumbo's farm, they have a chance to catch the \textbf{herd skulker} in the act, but it's a clever and cautious foe with an uncanny ability to disappear within herds, so it will be difficult. Though the \textbf{herd skulker} possesses a dangerous bite, it's a coward and runs if it feels threatened.

\subsection{Guard Duty}
Any PCs on guard duty must succeed on a DC 16 check using one of the following skills: Wisdom (Animal Handling, Perception, or Survival) or Intelligence (Nature).

\subsubsection{Success}
On a success, the PCs notice the creature entering the chicken coop and arrive in time to stop it before it attacks. The \textbf{herd skulker} flees (using its Spring Attack ability to avoid opportunity attack) when the PCs arrive and attempts to hide amongst one of the other groups of animals. If a PC blocks the entrance to the chicken coop, the \textbf{herd skulker} leaps at them, attempting to shove them out of the way so it can escape.

\subsubsection{Failure}
On a failure, the PCs hear a commotion of squawking chickens. When they arrive at the coop, the chickens are running around wildly, screeching, the \textbf{herd skulker} hidden among them. If the PCs hesitated for more than a round before rushing to investigate, the \textbf{herd skulker} assumes its true form, eats one of the chickens, then shifts back into a chicken before they arrive.

\subsection{Creature Considerations}
The animals are all on edge, as if they can sense the danger in the air.

\subsubsection{Animal Handlers}
PCs proficient in Animal Handling have a well-honed sense of animal behavior. They can observe the animals and succeed on a Wisdom (Animal Handling) check contested by the \textbf{herd skulker} Charisma (Deception) to notice that one of them is behaving suspiciously.

\subsubsection{Answers from Animals}
If the animals are questioned using magic, they communicate that there's a new animal in the pen that doesn't smell right. If the \textbf{herd skulker} is still present, the animals can direct the PCs to it. Being not terribly intelligent, they have little useful information, but if the PCs spend too long trying to talk with them they might allow the \textbf{herd skulker} to slip away.

\subsection{Running the Herd Skulker}
The \textbf{herd skulker} is a cunning, vicious predator that behaves cautiously and attacks only when advantageous.

\subsubsection{Monstrosity Not Beast}
It is important to remember that the herd skulker is a monstrosity, not a beast—even though it much of the time looks and acts like a beast. Accordingly, spells that only affect beasts will not work on the \textbf{herd skulker}. For instance, if the PCs use \textit{speak with animals} to attempt to talk to the \textbf{herd skulker} (regardless of what form it is in), it doesn't respond; a successful DC 12 Intelligence (Nature) or Wisdom (Animal Handling or Survival) check determines that the \textbf{herd skulker} doesn't understand the words being spoken, despite the spell.

\subsubsection{Patient Pounce}
The \textbf{herd skulker} waits for the PCs to relax before it risks attacking one of the goats or sheep. If the PCs remain vigilant, it spends the rest of the evening among the animals, then attempts to slip away when Grumbo opens the gate for the goats to graze at dawn, wandering off to sleep until nightfall. Vigilant PCs can attempt a Wisdom (Perception) check opposed by the \textbf{herd skulker} Stealth; on a success, the PCs can spot the \textbf{herd skulker} attempting to slink away.

\subsubsection{True Form Revealed}
If the \textbf{herd skulker} reveals itself to any of Grumbo's animals, each PC may attempt a DC 13 check to hear them crying or otherwise sense the disturbance, using one of the following skills: Wisdom (Animal Handling, Perception, or Survival) or Intelligence (Nature).

On a success, if the PCs immediately rush to investigate, they catch the \textbf{herd skulker} before it attacks. It fights if pressed, but, if it's reduced below half its hit points, it flees.

On a failure, the \textbf{herd skulker} heads back to the farm to hide, either in the stables or in the sheep or goat pen, where Grumbo's animals still remain. If any of Grumbo's animals are presently immune to the \textbf{herd skulker} Herd-Hidden ability, it avoids those specific enclosures until said immunity has expired.

\subsubsection{The Cost of Delay}
If the PCs detect the animals' distress but delay rushing to their aid for a round, they arrive at the scene to discover frantic sheep and goats running about the pasture. At the edge of the goat pen, they find one dazed-looking goat with wet, mussed fur on the back of its neck, but otherwise unharmed; the \textbf{herd skulker} abandoned this goat when the PCs approached and hid somewhere nearby.

If the PCs delay longer than one round, the herd skulker kills one of the goats and carries it off. They find a few bits of goat fuzz and signs of a struggle in the dirt.

\subsubsection{Escape}
If the \textbf{herd skulker} escapes and the PCs track it, they can follow the trail of the creature into the pasture and confront it there.

\subsubsection{Prey Play}
The \textbf{herd skulker} continues to play these games with the animals and PCs for the up to three nights; after the third night, if the PCs do not drive it off or kill it, the \textbf{herd skulker} will wander off in search of easier prey.

\section{Concluding the Adventure}
The adventure concludes when the PCs either kill the \textbf{herd skulker} or defend (or attempt to defend) the animals for three nights.

If the PCs protect all the animals for three nights and kill (or drive off) the \textbf{herd skulker}, Grumbo is eternally grateful and happily pays the PCs the full promised reward. The following day, he leaves for the Wool-Off both relieved and exhausted.

If the \textbf{herd skulker} kills any of Grumbo's animals—but none of his four prized babies (Fuzzington, Lanolyn, Lambert, and Cashimir)—Grumbo will give the PCs the initial 50 gp reward and will still pay the 50 gp bonus if a PC succeeds on a DC 15 Charisma (Persuasion) check or if there is a fellow halfling in the party who is kind (or at least not unkind) to Grumbo. Grumbo sadly continues his plans to attend the Wool-Off, but it's clear that a substantial amount of wind has been taken out of his sails.

If the PCs dispatch the \textbf{herd skulker} after it has killed any of Grumbo's prize animals, he grudgingly gives them the 50 gp reward but refuses to pay any more than that. He is absolutely despondent, as if he had lost a member of his own family and is not shy about expressing his extreme disappointment in the PCs' failure to follow through on their promise to protect his flock. Grumbo cancels his plans to attend the Wool-Off this year, too overcome with grief to make the journey.

\chapter{Flight of the Dromedaries}

\begin{DndTable}{c}

\textbf{An adventure for four to five characters of 2nd level}\\
\end{DndTable}

\section{Adventure Background}
\DndDropCapLine{T}{}asnim Azul (LN dwarf \textbf{veteran}), a tall, eagle-eyed woman with braided black hair, is looking to hire capable guards to defend the caravan she is leading across the desert to the city of Onist.


\section{Adventure Hook}
The PCs encounter the caravan on the road while traveling to a new locale, or perhaps they run into Tasnim in the market square of a city, where she recognizes their kind (i.e., people who look like they can handle themselves in a fight) and inquires about engaging their services.

\subsubsection{Getting Hired}
Tasnim pays well for caravan guards skilled in steel and spell—but she demands the best. She offers the PCs 300 gp each in exchange for their services protecting the caravan along its journey to the city of Onist. The PCs must defend the caravan from bandits, protect the cargo, and keep the ten camels from escaping. Any failure to deliver all of the promised cargo to the city of Onist results in significant financial and reputational loss for her. When offered the job, it is made abundantly clear to the PCs that her losses come out of their own pay if she finds them to be at fault. Tasnim expects the journey will take approximately two months.

The caravan is transporting one or two items of valuable cargo. Choose from the list below, or choose your own:


\begin{itemize}
\item Fresh limes bound for a seawater port;
\item Fragrant citrons for manufacturing fine perfumes;
\item Whole coffee beans for sale in faraway souks;
\item Ornate jewelry inlayed with precious gems, coveted by concubine and courtesan alike;
\item Finely crafted steel helmets promised to the commander of heavy infantry.
\end{itemize}

\section{Dragon Sighting}
Shortly after the caravan begins its early morning trek, an \textbf{adult blue dragon} soars low overhead. While the dragon is oblivious to the humanoids below, it's flying low enough that everyone in the caravan must make a saving throw against its Frightful Presence.

While the PCs deal with their own fears, a number of the caravan's camels equal to the number of PCs break their bindings and bound down the desert trail. Due to their movement speed (50 feet), the camels likely leave any pursuing PCs in the literal dust.

Tasnim is quick to remind the PCs that they were in charge of securing the camels. If the PCs don't recover the camels (50 gp each) and the cargo they carried (worth a total of 50 gp × the number of PCs), those losses come out of their pay.

Luckily, the camels leave a trail in the sands—the expensive cargo that falls from their saddlebags. A successful DC 15 Wisdom (Survival) check allows a PC to track the camels to a sandy valley while it's still midmorning. On a failure, the PCs arrive at the sandy valley after the sun reaches high in the midday sky.

\section{The Sandy Valley}
Before their fear waned, the escaped camels fled down the desert trail and into a small valley surrounded by steep escarpments. This out-of-the-way valley houses the ruins of an ancient shrine to a forgotten god, as well as a primordial pool of rock and minerals that guards the valley.

An unlabelled version of the map can be found here.

\subsection{1. Valley Entrance}
\begin{DndReadAloud}{}
The camel tracks lead off the beaten path and into a small valley. Steep escarpments shoot up at sharp angles, rising to the flat butte that meets the skyline.

\end{DndReadAloud}

The eastern butte blocks the rising sun from beating down on the sands—even if the PCs arrived midday.

If the PCs brought the caravan in their pursuit, there is adequate space for the remaining camels to graze while the PCs locate the lost dromedaries. Tasnim and her other workers focus on caring for—and keeping a close eye on—the camels that didn't escape, and so they are not in a position to help the PCs recover those that fled.

\subsection{2. The Oasis}
\begin{DndReadAloud}{}
Toward the northern end of the valley entrance, another small caravan is resting by a small oasis. Their camels block access to the fresh water.

\end{DndReadAloud}

The PCs can approach members of the other caravan and inquire if anyone saw the fleeing camels.

\subsubsection{Friendly Approach}
If the PCs don't appear threatening, the caravan workers explain they saw the fleeing camels gallop into the sandy valley and up the narrow incline no less than 15 minutes ago. If asked, they have no knowledge of a blue dragon flying through this area recently.

\subsubsection{Unfriendly Approach}
If the PCs appear threatening, the caravan workers try to avoid the dangerous adventurers. The PCs can put the caravan workers at ease with a successful DC 13 Charisma (Deception or Persuasion) check, thereby convincing them to reveal the information about the fleeing camels.

\subsubsection{Notable Travelers}
In addition to the caravan workers, two notables lounge at the oasis. The PCs may speak with them before the other caravan resumes its journey.

Tadalesh Tesfay (NG human \textbf{noble}) is a wealthy gambler from a distant nation. He is traveling to Onist and needs bodyguards for once he arrives in the city. After noticing the PCs' weapons and capable nature, he asks them to meet him at the Azure Flame Gambling House. Once they arrive, Tadalesh asks them to help him pull one last heist—recovering treasured artifacts first stolen from his people.

Zehava Rasason (N high elf \textbf{mage}) is a transmuter traveling to Onist, where she'll teach at the wizarding college. Before she can begin classes, she must clear thugs from the transmutation tower. She is considering hiring adventurers to deal with this nuisance below her station.

\subsubsection{The Pool}
While there's no space for Tasnim's camels to drink from the pool, the PCs can squeeze between the animals to take sips from the pool or even fill canteens with fresh water.

Any creature that drinks from the pool has advantage on Constitution saving throws against the heat for 1 hour; any water taken from the oasis loses this beneficial property 1 hour after its removal from the pool.

\subsection{3. The Narrow Incline}
\begin{DndReadAloud}{}
Just past the oasis, the path narrows, as if some primordial giant had tried to pinch shut the valley eons ago. The slope rises at a steady incline, leading to the northern half of the sandy valley.

\end{DndReadAloud}

\subsubsection{Pit Trap}
As the camels rushed through this area, they greatly disturbed the sand, which had long rested precariously above an underground pit. PCs who are actively searching for traps or hazards identify the danger with a successful DC 13 Wisdom (Perception or Survival) check. Any PC who succeeds on a DC 10 Dexterity (Acrobatics) check can carefully skirt around the edge of the narrow incline, thereby avoiding triggering the hazard.

If the PCs don't notice the hazard (or fail to avoid it), the sand beneath their feet gives way. As the sand collapses, the PCs fall into an underground pit that is 15 feet in diameter and 15 feet deep. Fortunately, the sand below cushions the PCs' fall, and they suffer no falling damage.

\subsubsection{The Western Wind}
Though the PCs survive the fall without difficulty, the cursed western wind suddenly whips up, and the pit begins filling with sand. The pit fills in 1 minute, burying any PCs still inside after 10 rounds.

\subsubsection{Escaping the Pit}
To climb out of the pit, each PC must succeed on three DC 13 Strength (Athletics) checks; using ropes or other climbing gear grants advantage on the rolls. If the d20 roll is a natural 20, it counts as two successes; if the d20 roll is a natural 1, the PC falls to the bottom of the pit, losing any accrued successes.

The PCs shouldn't tarry, as the pit quickly fills with sand. After five rounds, the PCs have disadvantage on all ability checks they make while in the pit.

Everyone in Tasnim's caravan is too far away to reach the PCs before the pit is filled. If the PCs find themselves in desperate straits, one of the notable travelers from 2 is close enough to help, arriving with rope at the end of the third turn. (They refuse to proceed any further into the valley—the pit is already risky enough.)

\subsubsection{Buried PCs}
PCs buried in the pit are restrained and unable to speak. The number of successes determines the level at which the PCs are buried:


\begin{itemize}
\item 2 successes: 5 feet
\item 1 success: 10 feet
\item No successes: 15 feet
\end{itemize}

Buried PCs are likewise suffocating. It may be helpful for you to reference those rules and calculate how long each PC can hold their breath before this event begins.

\subsubsection{Rescuing Trapped PCs}
If PCs not trapped in the pit have appropriate equipment, such as a shovel, they can dig out their buried companions. It takes 1 minute to dig down 5 feet in the loose sand. Once a PC reaches the depth of a buried companion, that PC can breathe and stops suffocating.

\subsubsection{Development}
After 1 minute, the sand completely fills the pit, and the PCs can walk over it without issue. If the PCs avoided triggering the trap while traveling up the narrow incline, they cannot avoid the trap when returning with the camels. Camels falling into the pit are buried in the sand after 1 minute, perishing soon thereafter.

\subsection{4. The Eastern Dunes}
\begin{DndReadAloud}{}
To the east of the incline, flat sands stretch to the edge of the valley. Drought-resistant trees line the edge of these dunes, while the late morning sun glitters from the arid ground.

\end{DndReadAloud}

If the PCs arrive in the morning, the eastern escarpments block the sun from shining down on the dunes.

\subsubsection{Hot Sands}
At midday, the dunes are fiery hot, forcing the PCs each to make a DC 10 Dexterity saving throw when they begin traveling through this area; on a failure, they take 1d4 fire damage as the hot sands burn their feet. (Barefoot creatures have disadvantage on this saving throw.) This is intended to be a one-time effect, but if the PCs linger in the area you might require them to succeed on the saving throw again or suffer additional damage.

\subsubsection{Creatures}
A PC who succeeds on a DC 12 Wisdom (Perception) check notices five \textbf{vulture} fly into the sky at the PCs' approach. If the PCs follow the birds' flight path, they witness the vultures meet up with a \textbf{giant vulture}. These scavengers view the PCs as potential meals and swoop down to attack the PCs if they proceed further into this area.

\subsubsection{Tactics}
The vultures exploit their Pack Tactics feature, splitting into two groups of three. They target anyone who looks weaker or who has lagged behind the group.

\subsubsection{Development}
The \textbf{giant vulture} disengages from combat and flees to its roost if it starts its turn with less than 8 hit points. The other \textbf{vulture} attempt to fly away if the \textbf{giant vulture} dies or has disengaged, or if at the start of their turn there are two or fewer other \textbf{vulture} remaining.

\subsection{5. The Bushes}
\begin{DndReadAloud}{}
Thick bushes line a tapered sliver of the sandy valley. A number of broken branches are scattered about the sand near the bushes, and some half-broken boughs still cling to their branches.

\end{DndReadAloud}

If the PCs arrive in the morning, there is a faint odor in the air that they can't quite place. If the PCs arrive midday, the unmistakable smell of a rotting corpse fill their nostrils.

\subsubsection{In the Bushes}
If the PCs inspect the bushes, they locate the bloated corpse of a slender, redheaded halfling. A successful DC 13 Wisdom (Insight) check deduces that the halfling was trying to climb the steep escarpment when she fell to her death.

\subsubsection{Gear}
A PC who spends 10 minutes searching the area and succeeds on a DC 10 Intelligence (Investigation) check finds the components of a \textit{climber's kit} near the steep escarpment.

\subsubsection{Treasure}
On the body is an ornate, ruby-crusted \textit{dagger} (worth 10 gp); the shards from two shattered potions of healing; and a fine wooden case containing 5 engraved rings (worth 20 gp). A PC who succeeds on a DC 17 Intelligence (History) check recognizes the engravings of the ruling family of Onist—and that the ruling family are not halfling.

\subsection{6. The Gnoll God's Temple}
\begin{DndReadAloud}{}
Aging, cream-colored pillars rise from a marble foundation, covered with ancient crimson stains. Underneath the crumbling canopy sits the gigantic, crumbling effigy of a gnoll god, clad in armor and wielding a vicious glaive.

\end{DndReadAloud}

Long ago, this temple served as the center of worship for the gnolls in this region. Great packs would gather before this giant idol, offering praise and sacrifices to their progenitor god before embarking on great marauding campaigns across the realm.

A PC who succeeds on a DC 17 Intelligence (History) check recalls that, ages ago, desert druids drove the gnolls out of the sandy valley. These druids left some kind of hidden guardians to protect the valley from gnolls and any other interlopers (see 8).

Time has not been kind to the temple. Over time, the sand has chipped away at the pillars and robbers have stripped the valuables from the once ornate statue—all save for the gnoll god's right eye.

\subsubsection{The Gnoll God's Eye}
The gnoll god statue stands 20 feet tall. A PC who wishes to remove the eye gem can do so with 5 minutes of careful work with a \textit{dagger} or similar implement; time has worn away most of what held the stone in place so it doesn't require much skill to remove it. Climbing the statue requires two successful DC 13 Strength (Athletics) checks; any creature wearing medium or heavy armor makes this check with disadvantage.

\subsubsection{Treasure}
The gnoll god's eye is a crimson almandine garnet the size of a large grapefruit; most buyers would gladly pay 100 gp for it, but, if the PCs find a buyer interested in gnoll culture or religion, they could sell the gem for as much as 500 gp.

\subsubsection{Creatures}
If the PCs remove the eye, they disturb a \textbf{swarm of wasps}, which attacks anyone holding the eye. Development. While the temple is potentially a good place to rest, resting there allows the sun to rise higher in the midday sky. Additionally, if the PCs linger, their presence draws the \textbf{desert slime} to attack at the temple (see 8).

\subsection{7. The Western Dunes}
\begin{DndReadAloud}{}
Camel prints head off along the western dunes, heat gleaming off the sand.

\end{DndReadAloud}

The harsh sun already sits high enough in the sky that it has already been beating down on the western dunes for several hours. From this vantage point, the PCs can see the camels resting in 8.

\subsubsection{Hot Sands}
PCs who arrived in the late morning must make DC 10 Constitution saving throws the first time they cross this area. On a failure, the PCs' movement speed is reduced by 10 feet. PCs who fail by 5 or more also suffer one level of exhaustion.

The desert sun is even hotter for PCs who arrive midday; this increases the Constitution saving throw DC to 15.

PCs who drank water from the oasis have advantage on this saving throw. If they brought water from the oasis, you may suggest they feel the urge to take a sip at this time.

\subsection{8. The Primordial Pool}
\begin{DndReadAloud}{}
At the northern edge of the western dunes, the camels rest, tired from their escape. Beyond the dromedaries, a bubble of gas rises up from a pit of slightly undulating sand.

\end{DndReadAloud}

The steep escarpments cast a shadow over this area. Desert druids created this primordial pool of silt, fine rock, and minerals to birth their desert slime guardians to keep the gnolls from returning to the sandy valley.

\subsubsection{Creatures}
Two \textbf{desert slime} emerge from the primordial pool, startling the camels and sending the beasts running.

\subsubsection{Tactics}
The first \textbf{desert slime} pursues the PCs. The other slime chases after the camels, but struggles to catch the beasts even at a slow trot. Once the PCs defeat the first slime, the second slime stops chasing the camels and targets the PCs instead.

\subsubsection{Dangerous Foe}
Having the PCs fight both slimes at the same would be extremely challenging for PCs of this level and could easily end with the death of multiple PCs.

\subsubsection{Treasure}
PCs who spend 5 minutes watching the pit eventually notice a gnomish hand poke above the surface. A rope can be used to fish the corpse out of the roiling sand. Strapped to the gnome's waist is a \textit{+1 shortsword} with the name "Howler" engraved into its pommel. Tucked into her shirt pocket is a flask filled with a fig liqueur, and she wears a necklace with a large platinum pendant (worth 50 gp).

A PC who succeeds on a DC 15 Intelligence (History) check recognizes the pendant as the symbol of the Onyx Sabers, an organized crime group in Onist.

\section{Concluding the Adventure}
Once the PCs lead the camels from the pool and back to the caravan, a grateful (and perhaps a little surprised) Tasnim shoots them an impressed and appreciative look. Along the way to Onist, the PC with the highest passive Perception spies some bandits lurking along the caravan's path, but the presence of armed guards is enough to scare them off without a fight. The PCs and the caravan travel the rest of the way to the city without incident, and when they arrive, Tasnim gladly pays them the agreed-upon amount.

\chapter{Furor in the Farmyard}

\begin{DndTable}{c}

\textbf{An adventure for four to five characters of 2nd level}\\
\end{DndTable}

\section{Adventure Background}
\DndDropCapLine{F}{}ollowing his escape from the local magistrate's cells, the bandit Sben Albacht and his gang holed up in the Hundedel family's farmhouse while they wait to escape the area under cover of darkness. The family fled to the nearby village of Vandlewin before the bandits moved in, but in the chaos of their rush to escape, they realized too late that their oldest child, Kristoph, had somehow gotten off their cart and had been left behind. Kristoph is unharmed and is being kept prisoner, under guard, in his bedroom. Sben hopes he can trade the boy's life for his freedom should all other avenues of escape be lost.


\section{Adventure Hook}
The following hooks start the PCs on this adventure:


\begin{itemize}
\item \textbf{\textit{Harried Rider.}} A traveler on horseback who is pushing the animal to its limits waves down the PCs, then stops, panting. He breathlessly explains that he's from the village of Vandlewin, just a half-day's ride up the road. A group of blackguards broke their leader out of the magistrate's cells and then took some poor farmer's boy hostage and are holed up at the family's farmstead.
\item \textbf{\textit{Tavern Talk.}} At a tavern in the city the PCs are in, they overhear a merchant talking about the events in Vandlewin. The people he's speaking to aren't capable fighters and are afraid to help. "Shame about that boy, but I'm not messing about with the likes of the Albacht gang," one of them says. The merchant explains that Vandlewin is just about a half-day's ride from here.
\end{itemize}

\section{Hundedel Hostage}
Once the PCs arrive in Vandlewin, they are brought to the elderly local magistrate, Milliam Volstacker. When the PCs are brought before him, he says, too loudly, "Adventurers! Thank the gods!" then conveys the following details:


\begin{itemize}
\item After the Albacht Gang freed their leader, Sben Albacht, from the magistrate's cell, the bandits fled on stolen horses toward the Hundedel family's farmhouse located a half-day's ride from town.
\item Mr. Hundedel was out in the fields when he saw the bandits riding hard in their direction. He recognized Albacht and knew he needed to get his family to safety. He hastily gathered his family into their carriage and rushed to town.
\item They'd put their oldest child, Kristoph (age 5) inside the carriage and closed him in; Mrs. Hundedel insists that he didn't know how to get the carriage door open. They'd loaded several of their small animals into the cart with him, but he'd been crying because they couldn't find his favorite, Ogyss.
\item Unbeknownst to the family, the boy ran off when Mr. Hundedel was tethering the horses, and Mrs. Hundedel was getting their baby, Raimeson, ready for the trip. Both mother and father rode outside on the driver's seat, and so in the chaos of their rush to escape the bandits, they didn't notice Kristoph had run off. Mrs. Hundedel thinks the boy must have seen Ogyss and ran after him.
\end{itemize}

If asked what kind of animals the Hundedels raise, Milliam says, "Oh, I don't know. All kinds, really. They've got some hogs, some chickens and goats. Some kind of lizard, I think." He doesn't know what kind of animal Ogyss is.

If the PCs ask to speak to the Hundedels directly, Milliam says that the parents are extremely distraught, and he would prefer not to question them right now—and time is of the essence. If the PCs insist, a DC 15 Charisma (Persuasion) check convinces him to allow it. The parents don't have any additional information relevant to the matter, but they can tell the PCs that Ogyss is a \textbf{barnyard dragonette}, a kind of friendly draconic creature the size of a house cat.

\subsubsection{Complication}
If the PCs need any convincing, Milliam notes that four members of the town watch were killed during Sben's escape, so he doesn't trust the remaining members of the watch to get the job done without getting the boy killed.

\subsubsection{Reward}
Milliam offers to pay the PCs: (1) 100 gp for the rescue of Kristoph Hundedel; (2) 50 gp for each bandit captured or killed; and (3) for Albacht himself, 100 gp if killed or 500 gp if the notorious bandit is recaptured so that he can be properly tried.

\section{Razorback Crab Burrow}
If the PCs spend an hour exploring the area around the Hundedel family farm they can find an access to the burrow of the \textbf{razorback crab} by making a DC 16 Intelligence (Investigation) check. Access points can be found up to a half-mile from the farm. If the PCs attempt to travel through the burrow, they can reach the southernmost portion of the \textbf{razorback crab} burrow within 10 minutes by making a successful DC 13 Intelligence (Investigation) or Wisdom (Survival) check.

The pungent scent of ammonia, fainter in the tunnels and sharper in the chambers, fills the burrow. Scuttling can be heard to the northeast.

\subsubsection{Ceilings}
The burrow's tunnels are 5 feet high, forcing Medium creatures to crouch to move through them.

\subsubsection{Lighting}
The burrow is dark and has no natural light sources.

\subsubsection{Smell}
The burrow smells of ammonia from the crabs' excretions. Upon entry, a creature must succeed on a DC 10 Constitution saving throw or suffer from the poisoned condition while within the burrow. A poisoned creature can attempt a new saving throw every minute; a successful save ends the effect. A creature that succeeds on its saving throw cannot be affected by the burrow again until 24 hours have passed.

\subsubsection{Sound}
The crabs can be heard burrowing to the north. Their hard carapaces make a dull scratching sound as they scrape against the walls.

\subsubsection{Creatures}
For every 10 minutes the PCs spend in the burrow, roll 1d6 and consult the Razorback Crab Burrow table. Adult \textbf{razorback crab} are aggressive and defend the burrow to the death.

\begin{DndTable}[header=Razorback Crab Burrow]{XX}

d6 & Encounter\\
1 & 2 adult \textbf{razorback crab}\\
2 & 1 adult \textbf{razorback crab}\\
3–6 & 2d10 harmless Tiny \textbf{juvenile razorback crab}. (Juvenile crabs have AC 14, 1 hit point, and can take no actions.)\\
\end{DndTable}

If the PCs reach 11 before the \textbf{razorback crab} finish burrowing into the farmhouse, they can complete the tunnel—and thus burrow into the farmhouse—in 1d6 minutes. If they do not have proper tools, they can use the severed scoop-claw limbs of the crabs to do so.

The PC doing the digging must make a Strength (Athletics) check (with disadvantage if using the razorback claw as their tool); if they surpass a DC equal to 10 + the number of minutes rolled on the d6 above, the bandits assume that the scratching is coming from the farm's \textbf{shovel dragonette} and thus are surprised when the PCs burst into 1. If the d20 roll is a 1 on the Strength (Athletics) check, the digging PC suffers one level of exhaustion.

\section{Hundedel Farmyard}
An unlabelled version of the map can be found here.

\begin{DndReadAloud}{}
The low grass of the farmyard is worn to bare dirt paths leading from the farmhouse to the pastures and outbuildings. Figures can occasionally be seen moving furtively near the fenced perimeter and inside the farmhouse itself.

\end{DndReadAloud}

The farmyard is neatly bordered by a wide-slat wooden fence which has an ungated 30-foot-wide entryway to the south. All of the fences are 5 feet tall and can be vaulted during a creature's movement with a successful DC 12 Dexterity (Acrobatics) check.

\subsubsection{Creatures}
Twelve \textbf{bandit} have been posted in the yard to keep watch. They have advantage on Wisdom (Perception) checks against attempts to approach the farm from the south, east, or west. The tall grass to the north of the farm makes the area lightly obscured. The bandits use the fence and buildings for cover as best as possible and try to stay out of sight of observers.

\subsection{1. Hundedel Great Room}
\begin{DndReadAloud}{}
This large room serves as kitchen, dining area, and recreation room. A large bearskin sits next to the wood-burning stove on the northern wall, and a prep table can be found to the east. A worn but solid wooden table with sturdy chairs occupies much of the room's center, and dirty dishes litter its surface. The floor is covered with muddy footprints, both from boots and what looks like some sort of reptile.

\end{DndReadAloud}

Sben's gang have made themselves at home here. The sound of digging is audible in this room; the gang ignores it, thinking it's the farm's dragonettes. The digging is actually being done by a colony of \textbf{razorback crab} that has been burrowing beneath the farmhouse for weeks. (See the "Razorback Crab Burrow" section.)

\subsubsection{Creatures}
Sben (NE \textbf{bandit captain}) and five \textbf{bandit} occupy this room. The \textbf{bandit} are not actively watching for intruders and will parley if they are not immediately attacked.

\subsubsection{Developments}
If combat ensues, once two or more of the \textbf{bandit} are defeated, Sben flees to 2; once inside, he knocks over the wardrobe inside to block the door (requiring a DC 15 Strength Athletics] check to break through). After he flees, any remaining \textbf{bandit} surrender if injured; if not under guarded or tied up, any \textbf{bandit} that surrender will attempt to run at the most opportune moment.

At the beginning of the second or third round of combat—or after the PCs have been parleying with the \textbf{bandit} for a minute or more—four \textbf{razorback crab} burst into the room from below. The crabs immediately attack the closest humanoid to them.

\subsubsection{Treasure}
Mrs. Hundedel's mandolin was gifted to her by a master luthier and is worth 200 gp. The muddy grizzly bearskin on the floor near the fireplace, if cleaned up, is worth 75 gp.

\subsection{2. Mr. and Mrs. Hundedel's Room}
\begin{DndReadAloud}{}
This once tidy room is now tracked with muddy boot prints. The room has obviously been ransacked for valuables. The wardrobe has been opened and rifled through. All of the common woolen clothes that once hung within it now in a heap on the floor. On the bed is a bloody, brutally injured woman who writhes in pain.

\end{DndReadAloud}

On the bed is Sben's sister, Marcella Albacht (N \textbf{bandit captain}, with only 1 hit point remaining), who was injured by the \textbf{boar} in 5. She only moves from the bed if forced. Under her body, the bedding is a muddy, bloody mess.

\subsubsection{Development}
If Sben flees into this room, after using the wardrobe to block the door he then exits the farmhouse through the window and heads north into the tall grass where he attempts to hide until nightfall.

Sben trusts the good nature of the local villagers and would sooner leave his sister so she can recover rather than risk her life—and his capture—by taking her with him.

\subsubsection{Treasure}
An ornately carved wooden case with silver hinges sits on a small table by the bed (worth 15 gp). Inside is a \textit{silvered dagger} with a mother of pearl hilt (worth 75 gp).

\subsection{3. Kristoph Hundedel's Room}
\begin{DndReadAloud}{}
Finely crafted wooden toys are scattered in a chaotic jumble around the floor of this room. On the floor, beside the unmade bed, is a plate of half-eaten food. Sitting on top of the bed is a young boy whose eyes light up with hope upon seeing you. As you enter, leaping to his feet from a chair next to the bed is a black-haired ruffian.

\end{DndReadAloud}

Kristoph is being held hostage in his bedroom by the ruffian, Yosuf. The toys littering the floor, make the room difficult terrain.

\subsubsection{Creatures}
In this room are Yosuf (NE \textbf{bandit}), plus one \textbf{barnyard dragonette} that is hiding under the bed; this dragonette is Ogyss, which Kristoph left the carriage to chase after.

\subsubsection{Development}
Yosuf has been feeding the \textbf{barnyard dragonette} from the plate of food on the floor, so if the PCs attack him, the dragonette emerges from beneath the bed and strikes at them; if that happens, Kristoph shouts \textit{"Oggie, no!"} at the creature.

If he's not surprised, Yosuf attempts to grapple Kristoph on his first turn, then pulls a \textit{dagger} and threatens to kill the boy. If attacked, he shoves Kristoph at his attacker and attempts to escape. If two or more PCs enter the room, the \textbf{bandit} surrenders with a successful DC 13 Charisma (Intimidation) check; if only one PC enters, he fights.

\subsection{4. Pantry}
\begin{DndReadAloud}{}
Jars of pickles and preserves have been knocked off the shelves and onto the floor of this pantry, soaking the bags of barley, potatoes, and carrots below in fruit and pickling brine. A few hams, links of smoked sausages, and sacks of onions hang from hooks on the ceiling. It looks like there was enough food in the pantry to carry the Hundedel family through the coming winter, but the bandits' depredations have reduced their stores considerably.

\end{DndReadAloud}

\subsubsection{Creatures}
A \textbf{barnyard dragonette} is greedily eating pickled asparagus and cherries off the floor. If the door is opened, it uses its breath weapon to drive the interlopers away and then returns to its meal. If further disturbed, it becomes angry and attacks. If reduced to 10 hit points, the dragonette disengages and attempts to escape the farmhouse to lick its wounds behind the barn at 6.

\subsection{5. Pigsty}
\begin{DndReadAloud}{}
Wide, horizontal planks of wood fence in a pig sty. Filthy straw is scattered across the muddy ground and an empty trough runs the length of the sty's northern border. A door in the eastern barrier leads into the pigs' stall in the barn.

\end{DndReadAloud}

A sow and a \textbf{boar} pace the sty in agitation. Sben's gang gives the boar a wide berth since it gored Marcella and their leader ordered them to leave the beast alone.

\subsubsection{Development}
The \textbf{boar} can be freed with a successful DC 13 Wisdom (Animal Handling) check, provided no creature enters the sty. The \textbf{boar} can't be commanded, but the successful check prevents it from attacking the PCs or their companions as it rages through the farmyard. The freed \textbf{boar} roams the farm looking for food and luxuriating in its freedom; it attacks any \textbf{bandit} it encounters on its way toward the farm gate in the south fence.

\subsubsection{Treasure}
Marcella's gold and ivory cameo locket (worth 45 gp) was torn from its chain and trampled into the mud when she was attacked by the \textbf{boar}. A creature with a passive Wisdom (Perception) score of 14 or higher notices the sun reflecting off an exposed piece of the jewelry.

\subsection{6. Barn}
\begin{DndReadAloud}{}
This small barn smells of soil and hay. Clean straw is heaped in the two horse stables and the interior pigsty along the northern portion of the structure, and stacked bales are piled neatly along the west wall. A variety of farm tools hang tidily from hooks and brackets on the south wall, and an empty barrow rests in the northwest corner. The sound of snoring can be heard from a stall to the northeast.

\end{DndReadAloud}

The exterior of the barn has recently been painted a bright red. The bandits have mostly avoided the dirt-floored building, due to the excitable \textbf{shovel dragonette} and \textbf{barnyard dragonette} lairing in it.

\subsubsection{Development}
If the PCs make it to the barn without drawing attention to themselves, roll a d20. On a 1-4, two of the \textbf{bandit}, Kel and Sora, are quietly mid-tryst in the center stable.

\subsubsection{Creatures}
Two \textbf{shovel dragonette} and one \textbf{barnyard dragonette} are napping in the stall to the northeast. If they are awoken by the PCs, they excitedly yip and jump at their ankles, trying to initiate play. If this occurs, any \textbf{bandit} on watch near the barn investigate within 3 rounds. If the PCs misconstrue the dragonettes' playfulness as aggression, they will defend themselves; if any of them are killed or knocked unconscious, the remaining dragonettes will attempt to flee, making a lot of noise in the process.

\subsection{7. Well}
\begin{DndReadAloud}{}
A large wooden pail hangs from the well's free-spinning crossbeam. It can be lowered and raised using the attached hand crank.

\end{DndReadAloud}

\subsubsection{Development}
Father Hundedel has long been terrified that Kristoph might tumble down the well. Thus, he has secured rungs on the southern shaft of the well that allow one to easily climb up or down the 80-foot shaft. From above, a successful DC 12 Wisdom (Perception) check is required to notice them. Creatures at the bottom of the well have advantage on Dexterity (Stealth) checks when attempting to hide from being seen from ground level.

\subsection{8. Chicken Coop}
\begin{DndReadAloud}{}
A stout post at each corner suspends the floor of the chicken coop. An occasional cluck can be heard from one of the hens roosting inside.

\end{DndReadAloud}

Inside, dirty straw litters the floor, while the straw filling the nesting beds is clean; twenty-four chickens amble about, some of them taking umbrage at the intrusion. The \textbf{bandit} ignore noise from the coop, as any loud sound sends the hens clucking loudly, but they investigate if any of them escape. Like the barn, the coop has recently been painted bright red.

\subsection{9. Horse Paddock}
\begin{DndReadAloud}{}
The gate to this paddock is open. Mounds of soil mar the pasture, indicating the presence of some small burrowing creature.

\end{DndReadAloud}

The paddock is empty as the Hundedels took their horses when they escaped the farm. There's loose soil in one corner that indicates a very small creature had been digging here; a shovel dragonette burrowed down here and emerged into 13.

\subsubsection{Development}
If the PCs have the equipment, they can expand the dragonette's tunnel wide enough to squeeze through to 13 with 30 minutes of work.

\subsection{10. Sheep and Goat Pen}
\begin{DndReadAloud}{}
Bleating ewes and lambs graze alongside stammering nanny goats and their rambunctious kids.

\end{DndReadAloud}

The grass here is kept short by the twenty sheep and sixteen \textbf{goat} grazing in it. While an observer can see clearly through the pen fence from one side of it to the other, the beasts within provide cover from view.

The corpse of Nitti, the Hundedels' collie, has been left unceremoniously by the bandits just outside of the fence to the northeast.

\subsubsection{Development}
The PCs can move undetected through the animal enclosure with a DC 10 Dexterity (Stealth) or Wisdom (Animal Handling) check.

The \textbf{barnyard dragonette} from 6 venture out to this pen at sunup and sundown every day to check on the livestock. If they notice unfamiliar humanoids in the pen, they excitedly speak to the sheep or goats (via \textit{speak with animals}) to find out what is going on. If this happens, any checks made to move through the pasture without being noticed are made with disadvantage.

1d4 \textbf{bandit} located in the farmyard investigate within 6 rounds if any of the sheep or goats escape the pen.

\subsection{11. Farmhouse Tunnel}
\begin{DndReadAloud}{}
A mild but unpleasant odor of ammonia lingers in this low, recently dug tunnel.

\end{DndReadAloud}

As noted in 1, \textbf{razorback crab} are tunneling into the Hundedel farmhouse from here.

\subsection{12. Crab Nursery}
\begin{DndReadAloud}{}
The walls of this rough-hewn cave are stacked with roe-like crab eggs. The odor of ammonia is almost overpowering.

\end{DndReadAloud}

\subsubsection{Creatures}
Three \textbf{razorback crab} defend the nursery, moving to intercept and dispatch any threats to the eggs. They fight to the death.

\subsection{13. Food Storage}
\begin{DndReadAloud}{}
The smell of ammonia is strong in here. Piles of unidentifiable organic matter are moldering next to stalks of pale gray fungus. A dead shovel dragonette lays atop one of the piles.

\end{DndReadAloud}

A PC that makes a DC 13 Intelligence (Investigation) check can determine from the soil scattered among the piles and the state of the chamber's ceiling that the dragonette must have burrowed in from above.

\subsubsection{Development}
If the PCs have the equipment, they can expand the dragonette's tunnel wide enough to squeeze through to 9 with 30 minutes of work.

\subsection{14. Sleep Niche}
\begin{DndReadAloud}{}
Hundreds of dormant young razorback crabs are resting on the many earthen shelves of this nook. The smell of ammonia is stronger here than anywhere else in the burrow.

\end{DndReadAloud}

If the PCs enter this area with a light source, the \textbf{juvenile razorback crab} wake and tumble off the shelves toward 12 like a wave. Any PC in the niche or the tunnel leading to 12 must succeed on a DC 12 Constitution saving throw or be poisoned by the odor (and revolted by the sensation of scuttling legs) for 1 minute.

\subsubsection{Development}
When the PCs arrive in this area, roll 1d6. On a result of 1 or 2, there is one \textbf{razorback crab} watching the niche. The crab fights to the death to allow the \textbf{juvenile razorback crab} to escape to 12.

\section{Concluding the Adventure}
If the PCs rescue Kristoph Hundedel, his family gladly gifts them with the mandolin in 1 and the dagger in 2. Likewise, the magistrate gladly pays the agreed upon reward, and afterward the townsfolk insist on throwing the PCs a big celebratory feast. Any bandits the PCs didn't kill are successfully captured and arrested.

\chapter{Salvage at Lonely Cove}

\begin{DndTable}{c}

\textbf{An adventure for four to five characters of 2nd level}\\
\end{DndTable}

\section{Adventure Background}
\DndDropCapLine{F}{}inlay Barrett was a notorious pirate and smuggler who sails the seas in his ship, the \textit{Dead Reckoning}. Barrett was an old salt, but even the most seasoned sailor runs afoul of nature. The \textit{Dead Reckoning} was wrecked during a terrible storm, and part of the ship and its cargo are now washed ashore in Lonely Cove, an isolated strip of beach at the base of towering sandstone cliffs. Word spreads quickly about the wreck. Ambitious souls move quickly to gather what they can of whatever valuables Barrett was carrying. Time is of the essence, because the next high tide or storm could carry away any treasures currently in Lonely Cove.


\section{Adventure Hook}
The following can be used to start the PCs on this adventure:


\begin{itemize}
\item \textbf{\textit{Chance Encounter.}} The PCs are aboard a passing ship and spy the wreckage. They convince the captain to drop anchor so they can row ashore and investigate, or perhaps the captain commissions them to do so, in order for the ship to claim salvage rights.
\item \textbf{\textit{Rumor Has It.}} A survivor of the \textit{Dead Reckoning} was brought into port aboard a naval cutter. Rumors are flying around town about the wreck and the possibility that Finlay Barrett's treasure is out there for the taking.
\item \textbf{\textit{Hired Hands.}} A merchant named Deverick Stanmore who was recently preyed on by Barrett wishes to recover a family heirloom the pirate plundered from his ship; it's a jewel-encrusted dagger with the words "We Were, We Are, We Will Be" engraved on the blade. He's willing to pay 200 gp for its safe return.
\end{itemize}

\section{Arrival at Lonely Cove}
\begin{DndReadAloud}{}
You arrive at Lonely Cove at low tide. This section of coast is comprised of sheer, towering sandstone cliffs. The interior of the cove is a crescent of rocky beach surrounded by tidal pools. You can see the front half of the Dead Reckoning run aground on the rocks just offshore of the northern end of the beach. To the south, a dark opening of a cave can be clearly seen in the cliff wall. Debris from the ship, barrels, crates, and other containers, float in the pools or lie washed up on the beach.

\end{DndReadAloud}

Within the confines of the tidal pools and beach, the PCs only have to contend with the inhabitants of Lonely Cove as indicated in the numbered encounter areas. However, should they spend time in the deeper waters outside the cove, they find that one of the reasons the cove bears its name is that this area is known for its predators, which keep local fishermen and others away from the locale. For every hour the PCs spend afloat, there is a 25\% chance of an encounter with 1d3 \textbf{reef shark} or one \textbf{swarm of quippers}. If anyone is actually in the water, the chance increases to 40\%; if anyone is wounded and in the water, the chance increases to 75\%.

An unlabelled version of the map can be found here.

\subsection{1. Rum!}
\begin{DndReadAloud}{}
Several barrels float in a tide pool at this end of the beach.

\end{DndReadAloud}

\subsubsection{Creature}
The tidal pools of the cove are the domain of a \textbf{beach weird}. The creature moves freely along the length of the cove and can be encountered anywhere in the tidal pools on the map, arriving in \textbf{1d4} rounds at any of the numbered areas if the PCs linger at a location for longer than 1 minute. It is upset with the intrusion of the wreckage and other creatures that the recent storm brought into the area, attacking any creatures it notices that do not normally reside in its domain. It pursues fleeing foes, but not out into the deeper waters or up onto the beach itself.

\subsubsection{Treasure}
Among the other cargo items that were washed into Lonely Cove are the remains of ship's rum supply. Each of these barrels can be sold for 20 gp.

\subsection{2. Shipwreck}
\begin{DndReadAloud}{}
The front end of a sloop is run aground on the rocks here, the wreck leaning hard to starboard. The mast has snapped and fallen, its crossbeam and furled sail soaking in the water below. What sounds like mournful birdsong can be heard coming from the inside of the shattered hull, as well as the sounds of scratching and skittering, like claws on wood.

\end{DndReadAloud}

This is most of what remains of the \textit{Dead Reckoning}. The back half of the ship is sunk beneath the waters off the cove, and the vast majority of its cargo sunk with it or is scattered here in the cove. However, a section of the cargo remains lashed down and in place within the bow of the cargo deck. Barrett often made money capturing rare and exotic creatures from far-off lands and selling them as pets or private zoo exhibits to unscrupulous nobles and other wellto-do patrons. Two of his recent finds are here in the wreck.

\subsubsection{Creatures}
Inside the hull, a stack of crates, lashed and secured to the deck, stands against the port side. At the top of the stack sits a cage containing the source of the sad singing, a bird that looks like an eagle-sized finch—a \textbf{moon weaver}. Also present—but uncaged—is a \textbf{guardian archaeopteryx}; it's the size of a small dog, and this particular archaeopteryx has had its wings clipped, thus it cannot sustain flight for more than 25 feet and must land at the end of its movement.

The archaeopteryx is hungry and would rather feast on the \textbf{moon weaver} than fend for itself against the \textbf{giant crab} that roam the beach. It has been attempting, unsuccessfully thus far, to reach the \textbf{moon weaver}, which is still locked in its cage, wedged atop the stack of crates. The dinosaur defends itself against any intruders, but attempts to flee elsewhere if brought below half its hit points.

The \textbf{moon weaver} cage is locked, but it can be picked with a successful DC 12 Dexterity check using \textit{thieves' tools}. The creature has also had its wings clipped and suffers the same flight penalties as the dinosaur. It takes at least six months for either creature to grow enough new plumage to regain its full flight speed.

\subsubsection{Treasure}
The crates still in the hull contain mundane goods, most of which have been ruined by exposure to seawater. However, with a successful DC 12 Intelligence (Investigation) check, searching the deck of the wreck uncovers a polished wooden case that has been wedged under the fallen mast. A successful DC 12 Strength (Athletics) check is needed to pull the case loose.

The lock on the case was damaged during the wreck, but inside is a red velvet-lined interior, inside of which is a jewel-encrusted \textit{dagger} with the words "We Were, We Are, We Will Be" engraved on the blade.

\subsubsection{Developments}
If the archaeopteryx is killed or chased off, the \textbf{moon weaver} attempts to communicate its wishes to be freed from its cage. While intelligent, it does not speak Common, and it is unlikely the PCs know its language. It has a passing familiarity with Common if spoken to, though it only recognizes common words and phrases. The \textbf{moon weaver} also readily responds to simple codes, such as "Whistle once for yes, two for no." If treated well, the \textbf{moon weaver} might become a companion to one or more of the PCs, until it has an opportunity to return home.

\subsection{3. Flotsam on the Beach}
\begin{DndReadAloud}{}
Several crates have washed up on the rocky beach here among a scattering of planks and other debris.

\end{DndReadAloud}

\subsubsection{Creature}
The beach is the haunt of a \textbf{talus flow}. The elemental subsists mainly on whatever \textbf{giant crab} it can catch. The \textbf{talus flow} can be found anywhere on the beach, but its constant hunger for the bodily fluids of living creatures means an attack on anything that lingers for too long on any part of the beach.

\subsubsection{Treasure}
The crates here are still intact, but can be opened with a successful DC 10 Strength (Athletics) check (with disadvantage if no \textit{crowbar} or similar tool at hand). Most of the crates have been inundated with seawater, ruining their contents, but four still contain items of value. The first crate holds various tools and sealed containers of herbs, enough to comprise a total of 10 herbalism kits (worth 5 gp each). The second contains sealed jugs of rare dyes (worth 50 gp total). A third crate holds waterlogged, ruined books, but also a small box containing 4 potions of healing. The fourth crate contains 10 lbs. of incense (worth 25 gp).

\subsection{4. A Big Fish in a Small Pond}
\begin{DndReadAloud}{}
The bottom of this deep tidal pool is thick with growing seaweed that waves in the current. A shiny form can be seen moving about amidst the foliage.

\end{DndReadAloud}

\subsubsection{Creature}
A \textbf{pescavitus} was another of the rare creatures Finlay Barrett captured and had planned on selling. The box containing the creature spilled into this tide pool as the remains of the ship beached. The creature was able to free itself, but by then the tide had gone out and it found itself trapped here. Should it notice the presence of the PCs, it comes cautiously to the surface to converse. It would like to be freed from the pool and released into the ocean so it can escape. It offers to use any of its magical abilities on the party's behalf in exchange for its release. If attacked, the \textbf{pescavitus} flees to the bottom of the tidal pool, which is about 12 feet deep, defending itself if pursued.

\subsection{5. Sodden Textiles}
\begin{DndReadAloud}{}
Beyond the various detritus floating on the surface here, a number of colorful objects can be seen on the sand at the bottom of this tidal pool.

\end{DndReadAloud}

\subsubsection{Treasure}
The \textit{Dead Reckoning} was carrying a haul of expensive cloth and furs taken from a merchant ship. Most were lost in the deep water, but some of the smaller bolts made it to this tidal pool. The PCs can find the following: a 5-yard bolt of gold and silver brocade (worth 40 gp); an 8-yard bolt of white muslin (worth 32 gp); a 10-yard bolt of wool dyed midnight blue (worth 80 gp); and a small bundle of assorted furs (worth 34 gp).

\subsection{6. Sunken Treasure}
\begin{DndReadAloud}{}
A large collection of flotsam—sundered planks, rope, sailcloth, and other debris—cover the top of this tidal pool.

\end{DndReadAloud}

\subsubsection{Treasure}
A chest sits at the bottom of this 4-foot-deep pool. A successful DC 14 Wisdom (Perception) check notices it without moving the debris or entering the water. The chest is locked (but can be picked with a successful DC 15 Dexterity check using \textit{thieves' tools}); inside are 50 1-lb. ingots of copper (worth 25 gp total).

\subsection{7. Quarreling Crabs}
\begin{DndReadAloud}{}
A corpse lies on the beach here, not far from the waterline. A pair of crabs, each easily as wide as a human is tall, scuttle around the corpse, snapping at one another with their claws.

\end{DndReadAloud}

\subsubsection{Creatures}
The two \textbf{giant crab} here discovered the body of one of the pirates washed ashore and are currently fighting over the food. They don't attempt to injure one another; they are simply posturing in an attempt to drive the other off and claim the prize. If the PCs approach, the crabs abandon their argument in order to go after the new, fresh sources of food.

\subsection{8. Crab Graveyard}
\begin{DndReadAloud}{}
This section of beach is littered with what at first looks like shards of wood, bleached white from exposure to sun and salt. Upon closer inspection, however, they are revealed to be jagged, broken pieces of the shells of giant crabs. Among the smaller pieces, you can see sections of legs and even a couple of claws half-buried in the rocks. Nearby, to the east, a cave mouth looms in the side of the cliff.

\end{DndReadAloud}

The \textbf{giant crab} that reside and breed in the sea cave here are the main prey of the \textbf{talus flow} that dominates the beach (see 3). It catches them here as they make their way from the cave out to the water, hence the signs of carnage. This is also why this location is the most likely spot for PCs to encounter the \textbf{talus flow} if they have not already done so elsewhere.

\subsubsection{Creature}
A \textbf{guardian archaeopteryx} wanders this section of beach, cautiously scavenging among the crab shells for any remaining scraps. This creature is the mate of the other archaeopteryx aboard the \textit{Dead Reckoning} (see 2) that that Barrett captured in the hopes of forming a breeding pair in order to sell the offspring to collectors of exotic creatures. It attacks despite being outnumbered by larger opponents, as it is desperately hungry.

\subsection{9. Unlucky Treasure Hunter}
\begin{DndReadAloud}{}
A path here has been dragged through the sand and seaweed on the cave floor. At its end sits a chest, with a corpse draped over its lid.

\end{DndReadAloud}

A local fisherman unafraid to ply these waters spied the wreck of the ship and came to see what he could salvage. Unfortunately for him, the \textbf{beach weird} (see 1) dislodged his boat from the rocks and pushed it out into the deeper water as he made his way ashore. Then he ran afoul of the \textbf{talus flow} (see 3) while on the beach. Wounded, he made for the sea cave with a chest he'd found, only to run afoul of one of its inhabitants and met his end there. The fisherman's corpse has been decapitated, and the man's head is currently being snacked on by the cave's occupant.

\subsubsection{Creature}
A \textbf{giant crab} lurks in the back of the cave, eating the head of the corpse. It considers any creatures advancing on it or the chest (where the body is located) a threat, causing the crab to attack.

\subsubsection{Treasure}
The chest is locked (but can be picked with a successful DC 15 Dexterity check using \textit{thieves' tools}); it contains 40 1-lb. silver ingots (worth 200 gp total).

\subsection{10. Crab Den}
\begin{DndReadAloud}{}
Judging from the seaweed and sand, this cave is at least partially flooded at high tide. The smell of brine is strong, and the walls echo with the sound of chitin scraping on stone.

\end{DndReadAloud}

\subsubsection{Creatures}
There are currently four \textbf{giant crab} lounging in the cave. They aggressively defend their home against intruders.

\subsubsection{Treasure}
Among the seaweed and other debris in the cave is an ebony walking stick, capped and inlaid with silver, with a large, faceted quartz crystal set on top (worth 40 gp). It was brought in by one of the crabs that was attracted by its shine.

\section{Concluding the Adventure}
Once the PCs have overcome the various creatures in the cove and collected what treasures they can find, all that remains is for them to return to civilization and receive whatever compensation they're due (or to find buyers for the merchandise if working independently).

\chapter{Automated Banditry}

\begin{DndTable}{c}

\textbf{An adventure for four to five characters of 3rd level}\\
\end{DndTable}

\section{Adventure Background}
\DndDropCapLine{T}{}he Makville Artisan Academy is a respected institution, but a harsh one to attend. Its tuition inevitably beggars students, while recent graduates find that all of their creations are confiscated by the university upon graduation for reasons unknown. Three recent graduates, without job prospects, crippling debts, and none of the creations they'd toiled on, opted to turn to banditry. They stole their creations back from the university and fled into the countryside, using their constructs to rob merchants in order to pay off their debts. Though they haven't harmed their victims so far, their actions have drawn the attention of the town of River's End, whose economy relies on such merchants passing through.


\section{Adventure Hook}
The following hooks can start the PCs on this adventure:


\begin{itemize}
\item \textbf{\textit{Frightened Merchant.}} On their way to another location, the PCs encounter a passing merchant who is at first wary of them as if expecting to be robbed. When he realizes they're no threat, he informs them about a plague of banditry in the area specifically targeting merchants. He suggests they go speak with the mayor of River's End.
\item \textbf{\textit{Makville's Reputation.}} A representative of the Makville Artisan Academy seeks out the PCs due to their reputation as problem solvers and engages them to track down the students-turned-bandits. The Academy is willing to pay 150 gp to each PC in order to keep their name from being further sullied by the miscreants. She suggests the PCs speak with the authority in the town of River's End, as it seems to be a locus for the thuggish activity.
\end{itemize}

\section{River's End}
The adventure starts in the small, wooded town of River's End, noted for its apples and as a welcoming stop along long trade routes. The town militia comprises a trio of well-mannered but ineffectual young human men, whose primary job is breaking up bar fights and scaring foxes. The PCs are approached by Mayor Agar, who wishes to hire them to deal with his bandit problem.

\section{Mayor Agar}
Agar is an elderly dwarf (LN dwarf \textbf{commoner}), tall for his species, with tanned skin and a long, curly beard that comes down to his waist. He's understanding, empathetic, and very talkative in his old age. If the conversation drifts, he is likely to ask characters about themselves and offer well-meaning advice.

\subsubsection{The Job}
Mysterious bandits have been setting constructs on local traders. Mayor Agar wants the PCs to track down the bandits and stop the attacks.

\subsubsection{Reward}
The mayor is offering 100 gp to each PC if the attacks stop, plus an additional 50 gold if they bring him evidence of the destruction of the constructs used in the attacks. He doesn't care too much about the bandits, though as they've been nonviolent so far he'd prefer to avoid their deaths.

\subsubsection{Additional Information}
Agar shares the following information with the PCs:


\begin{itemize}
\item There are many caves in the woods east of town, Agar suspects the bandits are hiding in one of them.
\item The bandits haven't actually hurt anyone so far. The constructs jump out of the woods, scare folks into running and/or disarm their guards, and then run off with the most valuable cargo and some food, leaving the rest for whenever the merchants return.
\item Agar doesn't think the bandits are local and doesn't know who they could be. He suspects they have at least one spellcaster.
\end{itemize}

\section{The Hideout}
The bandit's hideout is a natural cave a half-day's ride east of town. The entrance is obscured by a sheet of hanging vines but isn't hard to find. The PCs easily find it after spending an hour looking for caves.

\subsubsection{Illumination}
While the entrance is left unlit to ward off suspicion, the interior is illuminated throughout with magical lights, providing bright light in all indoor areas unless otherwise noted.

\subsubsection{Interior}
The hideout is entirely unworked stone and contains furniture made from repurposed barrels or dismantled carts. There are no doors. The ceiling is 10 feet high unless otherwise noted.

An unlabelled version of the map can be found here.

\subsection{1. Entryway}
\begin{DndReadAloud}{}
Beyond the sheet of vines, a wide cavern opens before you. Within, natural rock formations rise up from the cave floor, blocking light and obscuring its true size, making the area feel claustrophobic and small despite its vastness. The entrance continues for 50 feet, leading into a larger chamber from which you can see light filtering out.

\end{DndReadAloud}

\subsubsection{Lighting}
This area has no magical lights to maintain the façade that it's an uninhabited cave. During the day, dim light filters in from outside. During the night, it is entirely dark.

\subsubsection{Creatures}
One \textbf{animated armor} and three \textbf{clockwork armadillo} wait in a hidden alcove that isn't visible from the entrance. Once the PCs move more than 10 feet into the cave, the constructs emerge; PCs with a passive Perception of 14 or lower are surprised. The constructs initially stalk and lunge toward the PCs, posturing, trying to scare them off. If the PC don't leave, the constructs attack.

\subsubsection{Tactics}
The constructs each engage the nearest PC. They have been ordered to engage nonlethally and to not engage retreating targets, so any PC reduced to 0 hit points is stable. If the PCs retreat or are defeated, the constructs head to Area 2 with any prisoners to show that the hideout has been discovered.

\subsubsection{Treasure}
When defeated, the \textbf{animated armor} becomes a fully functional suit of \textit{plate armor}. It has a maker's mark, "M I," emblazoned in the right shoulder; it was crafted from scratch by one of the bandits, Mahfouz, as his senior project.

\subsection{2. Living Area}
\begin{DndReadAloud}{}
A low table and bench large enough for three people sits in the center of this room; both the table and bench look as if they were built from a deconstructed cart. There is a small stove in the north of the chamber, and modest decorations and furnishings fill the rest of the space.

\end{DndReadAloud}

\subsubsection{Creatures}
Mira (N \textbf{bugbear}) and Muhfaz (N human \textbf{priest}) are sitting at the table and talking when the PCs enter. One \textbf{harvest horse} and one \textbf{clockwork pugilist} are awaiting orders in the center of the room. When Mira and Muhfaz notice the PCs, the duo do their best to convince them to leave (see the "Playing the Bandits" sidebar), attacking only if the PCs attack first or negotiations break down completely.

\subsubsection{Tactics}
The \textbf{harvest horse} and \textbf{clockwork pugilist} engage casters and archers first. Mira protects Muhfaz, who heals and supports her with spells. The constructs fight to the death, but Mira and Muhfaz surrender once both are reduced below half of their hit points or if either is reduced below 5 hit points. If a PC or Mira is reduced to 0 hit points, Muhfaz casts \textit{spare the dying} on his next turn. If Mira or Muhfaz is killed, the other fights to the death and stops trying to spare the PCs.

\subsubsection{Diplomacy}
The PCs can resolve this encounter without a fight by talking the pair down. With a successful DC 16 Charisma (Persuasion) check, Mira and Muhfaz are convinced to surrender peacefully and agree to pursue a different path to getting out from under their debts; appealing to their safety, the terror and losses suffered by their victims, or suggesting more legitimate ways to pay off their debts provides advantage on this check. On a success, the pair also offer to help talk down the third member of their cohort, Hadad. On a failure, they insist that the PCs vacate the premises, but they offer to leave the area and continue their banditry somewhere else instead.

If the PCs offer to personally help pay off their debt and make a convincing case that they might be able to do it, no check is necessary.

\begin{DndSidebar}[float=!b]{Playing the Bandits}
Mira and Muhfaz are desperate college kids trying to handle a bad situation with some semblance of their morals intact. They've made peace with fighting and banditry, but want to avoid actually hurting or killing anyone. They're inclined to play to the sympathy of the PCs and try to convince them that what they're doing isn't a big deal. They'll freely share the following information:




\begin{itemize}
\item There are three bandits: Mira, Muhfaz, and Hadad.
\item They're all graduates from the Makville Artisan Academy. All are capable of magic, though only Muhfaz can actually cast spells in combat.
\item They turned to banditry due to their inability to pay off their student loans with the jobs available to them. They haven't seriously hurt anyone and claim they leave most of the cargo behind so that their victims can still make a living.
\item Each of the bandits owes at least 2,500 gp to various debt collectors. See 5 for additional details about their debts.
\end{itemize}



\end{DndSidebar}

\subsection{3. Cart Storage}
\begin{DndReadAloud}{}
A group of small carts are arrayed along the south wall here, their harnesses modified so that they can be tied to various constructs. The floor is scoured with deep drag marks, but is otherwise unremarkable. A number of musical instruments have been placed in a pile on the ground in front of the carts; a tall clockwork guardian stands beside them.

\end{DndReadAloud}

\subsubsection{Creatures}
One \textbf{clockwork pugilist} and four \textbf{animated instrument} (a trombone, a dulcimer, a flute, and a shawm) animate and attack immediately if the PCs aren't accompanied by one of the bandits.

\subsubsection{Tactics}
On the first round of combat, the \textbf{animated instrument} bleats a long, low note and a \textbf{brown bear} emerges from it into a space within 20 feet of the trombone. The bear immediately takes its turn, moving to attack the closest PC; thereafter, it takes its turn immediately after the trombone's.

\subsubsection{Treasure}
The trombone is a magical instrument, created by Hadad for her sophomore year's band class. It functions as a rust bag of tricks, but is activated by playing notes rather than pulling items out of it.

\subsection{4. Supply Room}
Barrels and crates are stacked haphazardly around this room; some are empty, but many contain provisions, such as food, paper, charcoal, spare clothes, or water.

On the southern wall, stacked barrels hide the corridor leading to 5, but can be located with a successful DC 14 Intelligence (Investigation) check.

\subsubsection{Treasure}
45 days of rations, 112 gallons of clean water, 100 sheets of paper (worth 20 gp total), 12 sets of \textit{common clothes} (worth 6 gp total), 10 lbs. of steel ingots (worth 10 gp total), and 15 charcoal styluses (worth 3 sp total).

\subsection{5. Treasury}
\begin{DndReadAloud}{}
Assorted loot is stored in chests and secured sacks in this room. Pinned to the southern wall are three sheets of paper with words and figures written on them; upon closer inspection, they're revealed to be lists detailing each bandit's debt.

\end{DndReadAloud}

Each bandit owes 2,500 gp in tuition fees, plus an additional amount as follows:


\begin{itemize}
\item Muhfaz: \textit{500 gp for the components for a \textit{raise dead} spell to use on his brother;}
\item Mira: \textit{120 gp for reagents she had to replace after a magical accident at the Academy;}
\item Hadad: \textit{1,700 gp for an unspecified "magical treatment."}
\end{itemize}

They've made very little progress paying down their debt thus far, but they've paid off 78 cp each via legitimate work, plus an additional 120 gp each with their banditry.

\subsubsection{Treasure}
In the various chests and sacks, there is a collection of pottery (worth 50 gp total), a barrel containing 10 lbs. of saffron (worth 150 gp total), and crates containing 10 sq. yd. of linen (worth 50 gp total).

\begin{DndSidebar}[float=!b]{Playing Hadad}
Hadad has similar morals to that of her companions but is more committed to banditry as the solution. She's significantly more bitter than Mira or Muhfaz, critical not merely of the Academy's exorbitant tuition, but the entire system that allowed such fees to be levied and for the unscrupulous loan sharks to enforce it. If the PCs talk to her, she freely shares the following information:




\begin{itemize}
\item Her debt is larger than that of the others, as she was exposed to blood magic due to a professor's negligence, and the school charged her for the treatments she needed to survive.
\item They can't sell the constructs because they technically stole them from the Academy.
\item Banditry was Hadad's idea. She thought adventuring was too likely to get one of them killed, and the skills and abilities they learned at the Academy didn't offer them enough consistent work opportunities to outpace the debt collectors.
\end{itemize}



\end{DndSidebar}

\subsection{6. Workshop}
\begin{DndReadAloud}{}
A massive furnace has been carved into one wall of the cave here, while rock formations have been flattened and beaten into anvils and working surfaces. Hammers, tongs, barrels of oil, and spare parts are scattered throughout the room.

\end{DndReadAloud}

\subsubsection{Creatures}
Three \textbf{clockwork pugilist} are awaiting maintenance in this room. They animate and attack immediately unless the PCs are accompanied by one of the bandits.

\subsubsection{Treasure}
Mira's \textit{smith's tools} rest atop an anvil; a creature proficient in smith's tools doubles their proficiency bonus when using them to repair items (worth 60 gp). Two barrels of oil, used for quenching hot metal, are also present (worth 4 sp). Scrap metal (worth 4 gp) is scattered throughout the room.

\subsection{7. Study Room}
\begin{DndReadAloud}{}
Handwritten notes and stashed tomes on a variety of subjects, both mundane and magical, are stored in stone-carved bookshelves set into the walls here. The center of the room has been cleared out and fashioned into an impromptu sparring ring lined with stones.

\end{DndReadAloud}

\subsubsection{Creatures}
Hadad (CN human \textbf{atavist}) and one \textbf{flying sword} are in this room. If the PCs are not accompanied by Mahfouz and/or Mira, Hadad demands to know why they're there (see "Playing Hadad" sidebar); she attacks if discussions break down.

\subsubsection{Tactics}
The \textbf{flying sword} focus their attacks on any spellcasters in the group (preferably unarmored ones), while Hadad goes for the most dangerous melee threat. Hadad stabilizes unconscious PCs if the fight ends badly for them; she surrenders if reduced to 10 or fewer hit points.

\subsubsection{Diplomacy}
Hadad only surrenders without a fight if Mira and Muhfaz have already yielded to the PCs. If they have, a successful DC 14 Charisma (Persuasion) check convinces Hadad to take whatever deal they offered Mira and Muhfaz.

\subsubsection{Treasure}
The shelves contain books on construct repair and enchanting lore (worth 85 gp total). Any more valuable tomes the trio had have already been sold to pay off debts.

\subsection{8. Sleeping Area}
\begin{DndReadAloud}{}
Three crude sleeping mats, a layer of straw bedding, and several threadbare blankets are strewn about this room in a haphazard communal sleeping arrangement. A single chest sits, open, pushed up against the western wall.

\end{DndReadAloud}

\subsubsection{Treasure}
Three diplomas from the Makville Artisan Academy; a steel locket containing a picture of a large bugbear family (worth 5 cp); Hadad's diary, which expresses the hope that she and her friends will be debt-free in a year or two; 1 lb. of soap; and two potions of healing.

\section{Concluding the Adventure}
If the PCs let the students go or recruit them as aides, they are determined to repay their debts legally and quickly. The trio can repair equipment, analyze magic items and constructs, or perform research for the PCs, but, because of their financial situation, they always charge for such services.

If the PCs bring the students to justice, Agar listens to their story and has them repay their debt via community service while protecting them from debt collectors. If the PCs ever return to the town in the future, they find the students have settled into life as smiths and enchanters for the local community.

\chapter{Evil Rises}

\begin{DndTable}{c}

\textbf{An adventure for four to five characters of 3rd level}\\
\end{DndTable}

\section{Adventure Background}
\DndDropCapLine{T}{}he \textbf{bakery drake}, Amaranth, is the darling of the local upper crust. Since her arrival in the city, the baked goods she crafts have sold well among the nobility and other wealthy citizens. Her culinary craft has been in high demand, and the drake herself is known to be an excellent conversationalist and an appealing figure to invite to high-profile events.


Recently, several notable figures among the city's elite have died under mysterious circumstances. Initial investigations discovered they had been poisoned, but the method of delivery remained a mystery until a bright, young alchemist made the connection that all the victims had eaten foods baked at Amaranth's establishment. The poison had been delivered as two or three separate components which, when combined, created the toxin. Shocked officials now have a suspect, though her motive remains unclear.

What they do not realize is that the \textbf{bakery drake} is under duress. During renovations to her bakery, a hidden room was discovered, previously walled off from the rest of the bakery. When investigating this find, Amaranth was attacked by the ghost of \textbf{Panifex}, the former owner and another bakery drake. A decade ago, \textbf{Panifex} baked goods were the talk of the town. However, unlike Amaranth, he was not simply a baker, but also a spy, using his place in high society to glean sensitive information to pass on to his masters. He was eventually discovered, and one night agents of the nobility sneaked into the bakery, subdued \textbf{Panifex}, and walled him up in his own basement as punishment for his treason. \textbf{Panifex} died a slow death, cursing his captors with his final breath. But the drake's spirit lingered, bent on revenge . . . yet he could not leave the rooms in which he died. Years later, when Amaranth took over the bakery, she discovered the hidden rooms and accidentally unleashed \textbf{Panifex} ghost. The spirit possessed Amaranth and began plotting his revenge, using his baking skills—and Amaranth's fame—to learn which of the city's elite had caused his demise, then bake his way to sweet, sweet revenge.

\section{Adventure Hook}
The following can start the PCs on this adventure:


\begin{itemize}
\item \textbf{\textit{Deputized.}} The PCs are hired to arrest Amaranth, as capturing a \textbf{bakery drake} is not something typically in the purview of the average city guard.
\item \textbf{\textit{Mercenaries.}} The next of kin of one of the murdered elites hires the PCs to bring Amaranth to them before the city officials can get ahold of her, believing the drake to be too popular to face true justice without their intervention.
\item \textbf{\textit{To the Rescue.}} The PCs are friends and confidants of Amaranth. Having heard the news that she is the prime suspect in the recent murders, they head to the bakery to warn her and/or get more information about the situation.
\end{itemize}

In the "Deputized" and "Mercenaries" hooks, the hiring party emphasizes that Amaranth should be brought in alive. Either of these two hooks offer a reward of 150 gp per PC; with the "To the Rescue" hook, Amaranth will offer the same amount as a reward if she is rescued from \textbf{Panifex}.

\section{Draconic Confections}
Amaranth's bakery sits just off a main thoroughfare in the mercantile quarter of the city, near its border to the noble quarter. Of stone construction, the building has a partial second story, and the roof of its first story sports several chimneys, which are almost always sending streams of woodsmoke into the sky. The front of the shop has several large windows to display goods, and above the door hangs a wooden sign with the shop's name, as well as an image of a lithe, draconic form coiled around a hoard of baked goods.

Much gossip is discussed while patrons wait in line at the bakery, so \textbf{Panifex} is well aware of the general situation and expects the law to descend upon the bakery soon. The staff has been sent home, and the place is locked up tight.

An unlabelled version of the map can be found here.

\subsection{1. Shop Front}
\begin{DndReadAloud}{}
Tables of delicious-looking confections line the front of the shop under the windows. Flowering shrubs planted in large, glazed pots decorate the waiting area. A polished, wooden countertop divides the room, behind which are shelves holding various baked goods for sale. A bright red curtain covers the doorway to the back areas of the bakery.

\end{DndReadAloud}

Though the place looks ready for business, there is no one present, and a hand-painted sign reading "closed until further notice" has been placed in the window just to the right of the doors. The doors themselves are locked, but can be opened with the key (which Amaranth holds) or can be picked with a successful DC 15 Dexterity check using \textit{thieves' tools}.

\subsubsection{Treasure}
Behind the counter, on the floor, is an iron lockbox (which is locked, but can be picked with a successful DC 15 Dexterity check using \textit{thieves' tools}) containing 100 cp, 40 sp, and 40 gp.

\subsection{2. Bakery}
\begin{DndReadAloud}{}
This warm room smells of woodsmoke and baked bread. The center of the room is dominated by ten brick ovens that still crackle with flame. Around the room are several prep tables, their surfaces covered with flour and the remnants of dough and icing. Casks, crates, and sacks of ingredients line the back wall between two doors, alongside stacks of firewood. A large trapdoor in the floor is visible in the southeast corner.

\end{DndReadAloud}

The ovens in the main work area of the bakery remain hot, the coals having been banked but kept lit when the workers were released by Amaranth/\textbf{Panifex}. The sacks and crates contain dried fruits, nuts, and other ingredients; the casks contain various flours, water, honey, molasses, and salt.

\subsubsection{Discovery}
Searching the area reveals the kitchen manager's ledger on the table on the west side of the room, in which he keeps track of things such as the bakers' hours and pay, ingredient stock on hand, etc. A DC 12 Intelligence (Investigation) check finds a slip of paper in between some pages near the back of the book with the words "three dozen cabbage buns, day old" written with care on it. This phrase can be used to bypass the magical trap on the door leading to 5.

While looking around the room, the PCs notice a sheet of paper has been pinned to the trapdoor. In large letters the words "do not enter" have been written on it, with each word underlined three times; below that, in smaller letters, it reads: "by order of the boss!" If the ledger has been located, a successful DC 10 Intelligence check determines that the handwriting on the sign is the same as in the ledger; a creature proficient with \textit{calligrapher's supplies} makes the check with advantage.

\subsubsection{Trapped Trapdoor}
If a creature other than Amaranth attempts to open the trapdoor to 7, the entire room shakes, knocking things off shelves, toppling glass bottles, and most importantly, causing six bags of flour that were perched on the rafters above to fall to the floor and burst open, spreading flour into the air. This makes the air visibly dusty, rendering anything more than 20 feet away lightly obscured to the observer. It takes 5 minutes for the flour to settle.

If the trapdoor is opened, a spurt of flame bursts out of one of the ovens, igniting the flour in the air in a sudden flash of roaring flame, inflicting 11 (2d10) fire damage to all creatures in the room and igniting any unattended flammable materials, such as paper or cloth. Anyone standing in the space directly in front of the middle oven on the south side of the building also takes an additional 5 (1d10) points of fire damage from the spurt of flame it releases.

\subsubsection{Detection}
With a successful DC 15 Wisdom (Perception) check, a creature examining the area around the trapdoor (or the trapdoor itself after the first part of the trap is triggered) detects a sense of menace and foreboding hinting at a magical trap or that there is a hostile supernatural effect at play.

\subsubsection{Prevention}
Casting \textit{gust of wind} or using similar magic to remove the flour from the air prevents the trap's flame from igniting it, as does waiting 5 minutes or more to open the trapdoor. Otherwise, a creature who succeeds on a DC 15 Intelligence (Arcana) or DC 17 Dexterity check using \textit{thieves' tools} can disable the trapdoor for 1 minute, preventing the second part of the trap from triggering.

\subsubsection{Creatures}
If the flame spurt is triggered, four \textbf{flour mephit} (use statistics for \textbf{dust mephit}) emerge from flour casks after the fire subsides and attack any creatures they see that are not wearing Draconic Confections uniforms. For their innate spell, the \textbf{flour mephit} can cast \textit{grease} instead of \textit{sleep}; the effect created by the spell is clearly butter.

After 1 round of combat, the three \textbf{wolf} from 3 come investigate the noise and join the fight.

\subsubsection{Magical Wards}
Everything in this room has been magically warded to be immune to fire.

\subsubsection{Treasure}
Among the other ingredients in the room are containers of gold foil and dust, used to decorate the bakery's fancier confections (worth 50 gp total).

\subsection{3. Delivery Room and Storage}
\begin{DndReadAloud}{}
A pair of carts, their beds covered by canvas, occupy most of this wide room. A staircase leads to the second floor.

\end{DndReadAloud}

If approached from outside, there is a third set of double doors that open to a ramp down to the cellar, 7. All the doors to and from this room are locked (but can be picked with a successful DC 15 Dexterity check using \textit{thieves' tools}). The carts are used by bakery employees to make deliveries and pick up supplies for the bakery. The room also serves as a temporary storage area until supplies can be moved down to the cellar. The stairs lead up to 4.

\subsubsection{Creatures}
Three guard dogs (use statistics for \textbf{wolf}) patrol this area, here to guard the bakery against thieves during the off-hours. Two are currently resting on one of the carts beneath the canvas. If they detect any intruders moving through this area, they attack. If combat ensues, after the first round, the third guard dog comes down the stairs and joins the fray. A successful DC 12 Wisdom (Perception) check spots a bowl full of water under the wagon, possibly providing a clue to the dogs' presence prior to an attack.

After 1 round, the \textbf{flour mephit} from 2 join the fight.

\subsection{4. Entry Hall}
\begin{DndReadAloud}{}
The room at the top of the stairs is spacious and airy. To the north, several windows overlook the street that runs outside past the bakery, while a single window to the south offers a view of the smokestacks on the building's rooftop. A comfortable divan and a pair of plush chairs occupy the southeast corner, situated around a low table of polished wood. A door to the east and a door leading out to the rooftops are the only apparent exits other than the stairs.

\end{DndReadAloud}

Amaranth uses this room to entertain guests and meets with potential clients to discuss orders. The door to the roof is unlocked. The door to 5 is locked (but can be picked with a successful DC 15 Dexterity check using \textit{thieves' tools}); it is also trapped.

\subsubsection{Magical Trap}
If an attempt is made to open this door without Amaranth's key, a blast of golden magical energy lashes out at the creature, forcing them to make a DC 15 Dexterity saving throw. On a failure, the creature is paralyzed for 1 minute; on a success, the creature is not paralyzed, but is instead under the effect of the \textit{slow} spell for 1 minute.

Touching the door also triggers the \textit{alarm} spell that was placed on it, causing it to start ringing loudly. After the alarm rings for 1 round, the gargoyle in 6 comes to investigate and attacks any creatures not wearing Draconic Confections uniforms it sees in the Entry Hall.

\subsubsection{Detection}
With a successful DC 15 Wisdom (Perception) check, a creature examining the area around the door detects a sense of menace and foreboding hinting at a magical trap or that there is a hostile supernatural effect at play.

\subsubsection{Prevention}
If Amaranth's key is used to unlock the door, the trap is not triggered. Alternately, the passphrase found in the ledger in Area 2 deactivates the trap for 1 minute. This is to allow Amaranth's foreman—who has a copy of the key not attuned to the trap—access to the room for emergencies or unusual circumstances.

\subsection{5. Amaranth's Room}
\begin{DndReadAloud}{}
This is a large, tidy living space. To the south is a fireplace and a large, four-poster bed that's piled with cushions and thick blankets. A comfortable-looking stuffed sofa faces the window to the north, and a low table occupies the center of the room. A bookshelf stands adjacent to the fireplace, holding a number of books.

\end{DndReadAloud}

The fireplace is connected to the chimney of one of the ovens in the kitchen below, allowing the smell of baking to rise up into the room when the oven is used. The bookshelf is filled with cookbooks, containing recipes for a variety of baked goods from a wide range of lands and cultures. The table is very low to the ground, a little awkward for a Medium bipedal creature, but comfortable for a Small-sized one. The bed is comfortable, but full of crumbs.

\subsubsection{Treasure}
Under the bed is an iron lockbox (which is locked, but can be picked with a successful DC 15 Dexterity check using \textit{thieves' tools}); inside is 100 cp, 400 sp, and 100 gp, as well as a set of silver measuring spoons (worth 25 gp).

\subsection{6. Rooftop}
\begin{DndReadAloud}{}
This flat roof is featureless save for a number of smokestacks and a stone rainwater cistern with pipes leading from it down into the kitchen.

\end{DndReadAloud}

\subsubsection{Creature}
Perched atop the cistern is a \textbf{gargoyle}, an old ally of \textbf{Panifex} that has been recruited to aid him. It can be easily seen by any creature that ventures to the southern edge of the rooftop. It remains motionless unless attacked, or it realizes that there are unauthorized creatures inside the building—either because they open the door to the roof or because they set off the alarm on the door. The \textbf{gargoyle} is brash but not completely foolhardy—nor overly loyal to \textbf{Panifex}. It attempts to flee if reduced to a quarter of its hit points.

\subsection{7. Cellar}
\begin{DndReadAloud}{}
This cellar is stocked with crates, casks, and sacks, as well as stacks of firewood. A wheelbarrow sits in the southwest corner. A ramp leads up to a pair of double doors to the northwest, and another to a trapdoor in the ceiling to the southeast. This room is stocked with various goods used in the bakery—flour, sugar, molasses, nuts, dried fruits, and other foodstuffs, as well as cords of wood of the sort that would be used to feed ovens.

\end{DndReadAloud}

Specifically looking under the ramp leading up to the trapdoor reveals a hole broken into the wall that leads to Area 8; otherwise, a successful DC 12 Wisdom (Perception) check locates it.

\subsection{8. Hidden Room}
\begin{DndReadAloud}{}
Bricks and mortar from the hole in the wall are scattered across the floor on the northern side of this room. Bookshelves line the walls here, sagging with age and the weight of their contents. One has fallen to the floor, its books scattered around it. There's an archway in the middle of the western wall, leading to another room. The smell of old books is in the air, as is a slight odor of dampness or mildew—but, strangely, the predominant smell that fills this area is that of fresh-baked pastries.

\end{DndReadAloud}

This was once \textbf{Panifex} library. All of the books here have succumbed to rot and an infestation of silverfish. There are hundreds of rotted tomes.

If the PCs linger in this area, Amaranth/\textbf{Panifex} calls out from 9: \textit{"No confections today, I'm afraid. We're closed."}

\subsection{9. Death Chamber of Panifex}
\begin{DndReadAloud}{}
The floor of this room is covered with shredded paper and leather, as well as several books that have been torn apart by claws and teeth. In the southwest corner are the skeletal remains of a reptilian creature about the size of a jackal. The scent of pastry is much stronger here, and it is easy to determine the source: The bakery drake in the white vest with "Draconic Confections" stitched on one of the breast pockets; the creature smiles at you as you enter, and somehow the sweet aroma in the air becomes even stronger.

\end{DndReadAloud}

\subsubsection{Creatures}
The \textbf{bakery drake}, Amaranth, is lurking here, curled up in the northeast corner of the room (and thus is not visible from 8). She was driven to linger in this place by the ghost of \textbf{Panifex}. In her possessed state, Amaranth fights until she dies or is knocked unconscious; when either eventuality occurs, \textbf{Panifex} emerges from her fallen form as a \textbf{ghost} to continue the combat, fighting until destroyed.

\subsubsection{Panifex}
\textbf{Panifex} can use an action to summon a \textbf{swarm of insects} (silverfish) to attack his foes; once he has used this ability, he cannot do so again until the next dawn. \textbf{Panifex} Possession ability only works on creatures of the dragon type, and he cannot move beyond Areas 8 and 9 unless he has possessed a corporeal body.

\subsubsection{Tactics}
On his first turn after he emerges as a ghost, \textbf{Panifex} uses his action to summon the \textbf{swarm of insects}. The silverfish swarm out of the old books in 8 and move to attack the nearest creature hostile to \textbf{Panifex}.

\subsubsection{Treasure}
While searching the debris in the room, a successful DC 13 Intelligence (Investigation) check discovers a scroll tube of carved ebony, inlaid with and capped in silver (worth 75 gp). Inside the scroll tube are the collected recipes created by \textbf{Panifex}. Bakers far and wide would pay dearly for the secret to such confections, and the recipes can be sold for 200 gp to the right buyer. In one of the pockets of Amaranth's vest is a key that unlocks any of (and bypasses any traps on) the locked doors in Draconic Confections.

\subsubsection{Drake Bones}
The skeletal remains in the southwest corner belong to \textbf{Panifex}. A successful DC 10 Wisdom (Medicine) or Intelligence (Investigation) check determines that \textbf{Panifex} was likely also a \textbf{bakery drake}.

\section{Concluding the Adventure}
If Amaranth still lives, once \textbf{Panifex} has been defeated, she is safe from his influence and can explain what happened to her, as noted in the Adventure Background.

If the PCs were deputized or hired to deal with the drake, they can take her in to plead her case. Sympathetic PCs may testify on her behalf about the ghost and her possession.

If the PCs came to rescue her, they can still provide testimony, but may decide to not turn her over to authorities—though if given a choice, Amaranth would prefer to turn herself in and be given a chance to plead her case. Many of the city authorities and nobility remember \textbf{Panifex}, and those directly involved in his demise are grateful to be spared any more of his vengeance—so they are likely to believe any truthful testimony from the PCs and be lenient toward Amaranth as a result.

The PCs are duly rewarded for their efforts; depending on which adventure hook you chose, their reward comes from either the city's ruler, a noble/merchant, or Amaranth herself. Provide them with what seems a fitting reward, or whatever was promised, if an amount was already set.

Amaranth herself is insistent on providing the PCs with regular gifts of goods from her bakery in addition to any monetary sum agreed upon, and she would certainly also be willing to purchase \textbf{Panifex} recipes if the PCs found them.

\chapter{Four-Part Harmony}

\begin{DndTable}{c}

\textbf{An adventure for four to five characters of 3rd level}\\
\end{DndTable}

\section{Adventure Background}
\DndDropCapLine{T}{}hree days ago, an orphan named Budanyek entered the haunted Heinrosch Bardic College, goaded by the gang he was with to sneak inside. When Budanyek didn't come back, most of the other youths blew off his absence, save for a few friends who grew worried for his safety. Held back by superstition, they have not risked going in after him, believing Budanyek was taken by a spirit of the college. And they're right—the ghost of the college's founder, Mikolaj Heinrosch has captured and possessed Budanyek. Now, Heinrosch's ghost works Budanyek to the bone to finish his life's work. If the ghost succeeds, he believes he shall have eternal rest, despite the cost—Budanyek's life.


\section{Adventure Hook}
The following can used to start the PCs on this adventure:


\begin{itemize}
\item \textbf{\textit{Gang Prank.}} A group of teens, part of an unnamed, mostly harmless street gang, approaches the PCs, asking for their help in rescuing their friend, Budanyek. Led by his sibling, Rabbit, they explain how they lost their friend; loitering around the old college, daring each other and teasing him with his childhood nickname, Buttanyek. Penniless, Rabbit and the orphans try to appeal to the PCs' good nature as a reward for Budanyek's safe return. Failing that, they share that Budanyek often talked about someday cracking a safe he'd glimpsed inside the college.
\item \textbf{\textit{Collector.}} A collector hires the PCs to recover any Heinrosch memorabilia from the college grounds. He knows that much of it has likely been pilfered by bandits, but he hopes that something of interest remains. He offers 50 gp for each piece of Heinrosch memorabilia delivered. He doesn't tell the PCs what his total budget is—because he wants to be able to pick and choose from whatever they find—but he has a soft cap of 300 gp.
\end{itemize}

\section{Heinrosch Bardic College}
The college, built more than sixty years ago, is as much a monument to the ambitions and ego of Mikolaj Heinrosch as it is an academy. With wealth earned through years of adventuring, he retired and turned his hand to composing music—and indeed achieved some minor repute in that realm, with a handful of his compositions occasionally (okay, once) mentioned among those of the great composers. Not content with the "minor" part of "minor repute," or the fact that he was only once mentioned among the greats, Heinrosch hired gnomish architects and inventors to design an institute that would elevate his bardic legacy. Alas, Heinrosch died before the college ever became successful, and the institution folded abruptly after his demise.

Now, within the college's walls, musical constructs have awakened by Heinrosch's ghost, assisting him while he works towards his goal of completing his unfinished magnum opus.

The building is clearly abandoned, with all windows boarded up except where otherwise noted. All of the doors have likewise been boarded, but the eastern entrance door is ajar, left so when Budanyek used a crowbar to break inside.

An unlabelled version of the map can be found here.

\subsection{1. Entrance}
\begin{DndReadAloud}{}
The door opens to reveal a vast space with vaulted ceilings that reach 60 feet high. A tattered banner of purple cloth hangs from the ceiling, with yellow stitching spelling out the words "Welcome Students." To the north, a dusty, glass-faced cabinet stands in the gloom. Beyond the cabinet, and around the corner to the north, a long corridor stretches into the dark as the rest of the room to the west continues forward into another large, open area.

\end{DndReadAloud}

\subsubsection{Display Cabinet}
This cabinet once held trophies for the school's triumphs and Heinrosch's own accomplishments. Now, it's almost completely empty; the only remaining sign of a trophy is on the top shelf, a brass plaque with the heading "The Gilded Quill" in large letters, with the following written below in smaller text: "The Pen With Which Heinrosch's Symphony No. 3, Oboe Concerto No. 2, and Violin Sonata No. 7 Were Written." There is a mounting bracket where an item was clearly once displayed, but there is presently no sign of the quill in question.

\subsection{2. Hall of Revelry}
\begin{DndReadAloud}{}
Four boarded-up, curtained windows on the south wall allow in slivers of dim light that lance across a disarray of cocktail tables. Opposite the windows, posters depicting bygone performances plaster a wall between three closed doors.

\end{DndReadAloud}

Music can be heard coming from inside 10. A successful DC 11 Wisdom (Insight) or Charisma (Performance) check makes out five distinct sounds coming from the room: four woodwind instruments playing in practiced harmony while a young man's voice fluctuates between cries of pain and frustrated vocalizations. A success that surpasses the DC by 5 or more determines that as the music plays on and off, it subtly changes each time.

Regardless of the skill check result, any creature in this area hears the voice's cries.

\subsubsection{Performance Hall Doors}
Heinrosch's influence magically holds shut all six doors leading to 10. The doors can be opened with a successful DC 23 Strength (Athletics) or Intelligence (Arcana) check.

For each encounter resolved in Areas 5–8, reduce the DC by 1. By defeating the \textbf{origami golem} in 9, reduce the DC by 3. If all encounters in Areas Areas 5–9 are completed, Heinrosch's influence over the Performance Hall is broken, and the doors open normally thereafter.

\subsection{3. Hall of Muses}
\begin{DndReadAloud}{}
Three heavy wooden tables and benches are arrayed here, and light floods the area, shining through the pair of partially boarded windows. A trio of statues are lined up along a wall to the north.

\end{DndReadAloud}

\subsubsection{Statues}
The three stone statues—of The Harpist, The Painter, and The Warrior-Poet—depict the three bardic muses; students would lay offerings at the foot of the statues for luck and creative inspiration. Examining the statues reveals that they are in various states of disrepair.

A successful DC 13 Intelligence (Investigation) check notes the following: All of The Harpist's strings have snapped; The Painter's easel holds no canvas; and The Warrior-Poet holds a quill made of stone in one hand, poised as if to write something in a book held in his other hand, but his other hand is empty. There are bits of broken string still attached to the harp, indicating that some elements of the statues were originally not stone.

The intent of the statues being only partially stone was so that the students could interact with them—play a song on The Harpist's instrument, place a painting on The Painter's easel, or add a poem to The Warrior-Poet's book.

\subsubsection{Development}
If all three statues are repaired, the muses offer a blessing. Any creature that assisted at all in the repair has advantage on Intelligence, Wisdom, and Charisma ability checks for the next eight hours.

Suitable materials to repair the statues can be found within the college or can possibly be carried or crafted by well-prepared PCs.

\subsection{4. Dean's Office}
\begin{DndReadAloud}{}
Thick dust coats every surface in this richly appointed office. The window on the western wall lets in soft light, and a crack in one of its panes allows in a breeze that flutters yellowed stationary atop a handsome desk. On the walls are paintings depicting heroic deeds, and light glints off a basket-hilted saber mounted on the wall behind the desk. A black safe looms in the corner.

\end{DndReadAloud}

\subsubsection{Locked Door}
The door to this room is locked, but can be picked with a successful DC 13 Dexterity check using \textit{thieves' tools}. A spare key is hidden above the door, which can be spotted with a successful DC 16 Wisdom (Perception) or Intelligence (Investigation) check. Budanyek (who is in the Performance Hall) also has a key.

\subsubsection{Paintings}
A successful DC 13 Intelligence (Investigation) check notices that the same bardic figure in a jaunty hat is present in all of the paintings (though he is one of many figures depicted); proficiency with \textit{painter's supplies} provides advantage on this check.

\subsubsection{Treasure}
What's left of the late founder's wealth is contained within this room. Spending an hour to gather the room's contents accrues the following: six paintings (worth 300 gp total), rare books (worth 100 gp total), and 3d4 small mundane objects of dubious value or utility. The mounted saber is a \textit{silvered longsword} (worth 150 gp).

\subsubsection{The Desk}
A nameplate on the desk reads "Dean Heinrosch." While searching the desk, a successful DC 10 Intelligence (Investigation) check discovers a drawer with a false bottom. If the check surpasses the DC by 5 or more, it also reveals that the drawer is trapped; if it doesn't, and the PC specifically looks for traps, a successful DC 15 Wisdom (Perception) detects it. A successful DC 13 Dexterity check using \textit{thieves' tools} disarms the trap. Failure to disarm the trap triggers a small poisoned needle when the false bottom is accessed, forcing one creature interacting with the drawer to make a DC 13 Constitution save. The creature takes 2 piercing damage and 16 (3d10) poison damage on a failed save or half as much damage on a successful one.

Inside the false bottom are a collection of blueprints for a massive auditorium, arcane schematics for the creation of animated quartets and symphonies, a scroll of \textit{detect thoughts}, and a scroll of \textit{locate object}.

\subsubsection{The Safe}
A combination lock secures the safe; it has three rotating dials, each numbered 0-9. The code is 3, 2, 7—the numbers of Heinrosch's three best-regarded compositions; a clue to both the numbers and the order in which they appear in the combination can be found in the display cabinet in 1. Without the code, the safe code can be cracked with a successful DC 19 Intelligence check; a creature with proficiency in \textit{thieves' tools} can add their proficiency bonus to the check.

If the PCs can't unlock the safe and are stumped about the code, assuming they examined the display cabinet in 1, a successful DC 13 Intelligence (Investigation) or Wisdom (Perception) check recalls that they previously encountered three numbers extremely significant to the (egotistical) owner of this safe.

\subsubsection{Treasure}
The safe contains a \textit{+1 hand crossbow}, a \textit{whip}, Heinrosch's iconic hat, 20 pp, a \textit{potion of climbing}, and a \textit{potion of greater healing}.

\subsection{5–8. Practice Rooms}
Each one of these rooms are functionally the same, with 20-foot ceilings and thick, alchemically concocted insulation that isolates the occupants from hearing activities in the other rooms. The rooms are setup with the standard equipment one would expect to find in each type of space.

While combat might ensue in these rooms, each presents opportunities for quick-thinking PCs to find non-combat solutions. If PCs choose to resolve an encounter through a skill challenge, calculate the DC as 10 plus the number of creatures in the room. If the PCs succeed three times before failing three times they successfully resolve the encounter peacefully. If the PCs fail three times as a group, the creatures within become hostile.

\subsubsection{Area 5}
A music rehearsal space. 1 \textbf{clockwork conductor} leads 1 \textbf{animated piano} (use statistics for \textbf{mimic}, with the changes listed below) in practicing a piece, the piano playing its own keys.

\subsubsection{Creature Change}
The \textbf{animated piano} is a construct and doesn't have the Shapechanger or Adhesive abilities. Its Pseudopod (Piano Wire) attack is the same, but change the last line to read: If the target is a creature, it is grappled (escape DC 13). Ability checks made to escape this grapple have disadvantage.

\subsubsection{Area 6}
A music rehearsal space. 2 \textbf{animated instrument} (cellos, Medium instrument) engage in a musical duel while a panel of 3 \textbf{clockwork conductor} judge them.

\subsubsection{Area 7}
A dance rehearsal space. Ballet barres are mounted on each wall; above and below the barre on the east wall are mirrors that go floor to barre and barre to ceiling. 2 \textbf{clockwork pugilist} rehearse ballet movements, with 1 \textbf{animated instrument} (violin, Small instrument) providing the music.

\subsubsection{Area 8}
A music rehearsal space. 1d4+2 \textbf{animated instrument} (a harp, Medium instrument, with various woodwinds, Small instruments) practice syncopated melodies, led by 1 \textbf{clockwork conductor}. A small supply cabinet in the southeast corner contains a variety of instrument maintenance gear, such as replacement strings, reeds, rosin, spare string instrument bows, etc.

\subsection{9. Professor's Study}
\begin{DndReadAloud}{}
Light filters through frosted windows, casting a dim haze over this large room. Loose, half-graded pages litter the surface of the office furniture while tools and art supplies lie scattered across a wall-mounted workbench. An enormous model swan the size of a bear, made from folded parchment, sits on a platform amid sheets of colorful paper and a pair of silver shears that have fallen apart. Several stepladders stand open behind the swan.

\end{DndReadAloud}

\subsubsection{Treasure}
Books and art supplies are scattered throughout this room, all associated with the teachings of over a dozen different disciplines of art and performance.

In a journal open on the desk is written a dissertation's first draft on the ethics of giving inanimate objects life and sentience. A nameplate on the desk reads Professor Ravenshollow." Inside a desk drawer are 5 blank journals, 20 quill ink pens, a tarnished silver necklace (worth 15 gp), and 1d4+2 small mundane objects of dubious value or utility.

PCs who spend at least 10 minutes at the workbench can assemble a set of \textit{painter's supplies} or \textit{calligrapher's supplies}. Under the workbench are 10 blank canvases and 5 folded easels.

\subsubsection{Creatures}
The paper swan is not just a piece of art, but an \textbf{origami golem}. The \textbf{weakened origami golem} activates when anything in the room is picked up. It follows its last directive, "Stop any student from interfering with my work." The \textbf{weakened origami golem} takes this directive very literally and pursues any creature that activates it until it, or the creature, is destroyed, though it will not leave the college or enter room 10. The \textbf{weakened origami golem} was left unfinished when its creator vanished or perished; in its weakened state, it is CR 3 (700 exp) with the following changes:


\begin{itemize}
\item It has 32 hit points.
\item Its Trumpeting Blast recharges on a 6.
\item Any damage it deals to the PCs is halved.
\end{itemize}

\subsubsection{Silvered Weapons}
Each arm of the broken silver shears near the \textbf{weakened origami golem} platform can be used as a silver dagger.

\subsection{10. Performance Hall}
\begin{DndReadAloud}{}
This enormous auditorium tapers in toward an impressive stage. Rows of seating littered with old playbills sit underneath a dark, sixty-foot-high vaulted ceiling while torches hover overhead and dimly glow with ghostly fire. You see a young man hunching over a sheet of paper pressed flat on the stage as he furiously writes something on it with a quill pen with golden barbs. Suddenly, he shouts "No, no, no!" and sweeps the paper into the air. An \textbf{animated quartet} plays a sorrowful melody while the man—no, just a teen boy—breaks down into mutters and sobs.

\end{DndReadAloud}

\subsubsection{Lights}
Sixteen torches float throughout the room and glow with ghostly green fire, casting dim light in a 20-foot radius around each one.

\subsubsection{Doors}
The six doors to the area aren't traditionally locked, but are held fast by the will of Heinrosch's ghost, which is currently possessing the young man by the stage, Budanyek. Once the PCs open at least one of the doors, Heinrosch loses his hold on them, and all doors work normally thereafter.

\subsubsection{Creatures}
Opening the doors attracts the attention of both Heinrosch's \textbf{ghost} (currently possessing Budanyek) and the \textbf{animated quartet}.

\subsubsection{Tactics}
Heinrosch does not engage, and instead tries to get back to work. The \textbf{animated quartet}, however, attacks immediately after a door opens.

If the \textbf{animated quartet} is defeated, or at the end of the third round of combat, with Budanyek's voice, Heinrosch shouts:

\begin{DndReadAloud}{}
\DndQuote%
{}
{''Stop, stop, stop! It is impossible to concentrate with all this racket! How am I supposed to craft my magnum opus amidst this . . . this cacophony?''}
{}
\end{DndReadAloud}

At Heinrosch's words, the \textbf{animated quartet} (if still engaged in combat) stops attacking and moves to hover just above the young man's shoulder. Heinrosch continues:

\begin{DndReadAloud}{}
\DndQuote%
{}
{''Listen to me. Listen! I need just a little more time. Just a little more time. Please! I'm so close. I've spent so many years searching for the right melody. And . . . I've almost got it. I know it! I can't let you ruin this for me. My entire legacy hangs upon this achievement.''}
{}
\end{DndReadAloud}

\section{Ghostwriting Heinrosch's Final Sonata}
Heinrosch is obsessed with completing his magnum opus, a final sonata his death left unfinished. As the PCs speak to him, the \textbf{ghost} maintains perfect control over Budanyek's movement and actions, though the boy still intermittently wrests back control of his own body to cry out in pain before Heinrosch regains control.

If the PCs question Heinrosch, other than his work, he doesn't remember anything they might ask about, such as how, when, or how long ago he died. If they ask about anything else, he angrily shouts that he doesn't have time for idle chatter, but he's willing to listen to anything the PCs offer that might help him finish his opus.

Here are three likely scenarios via which the PCs can resolve the situation:


\begin{itemize}
\item \textbf{\textit{Have a PC Finish It.}} Heinrosch can be convinced to relinquish control of Budanyek if a PC offers to be possessed instead; Heinrosch only agrees if the PC has a Performance modifier of 3 or higher, or with a successful DC 17 Charisma (Persuasion) check (made with advantage if the PC is wearing Heinrosch's iconic hat found in 4). If the PCs choose this option, Heinrosch's \textbf{ghost} leaves Budanyek, possesses the PC, and immediately returns to his work. The possessed PC requires 2d8+2 hours to finish Heinrosch's work. The PC must make a DC 10 Constitution saving throw each hour; if they fail three times, they suffer a level of exhaustion and then suffer an additional level each additional failure.
\item \textbf{\textit{Let Budanyek Finish It.}} Heinrosch wishes to keep possessing Budanyek until his work is finished. If the PCs choose this option, it takes Heinrosch 1d8+1 hours to finish his work; this process requires Budanyek to succeed on a DC 13 Constitution saving throw every hour or else suffer a level of exhaustion. This could very easily end with Budanyek's death without intervention, as the teen is a \textbf{commoner} who is already suffering from three levels of exhaustion due to the rigors he has already endured. The PCs are free to assist Budanyek however they wish, whether it's via the Help action and/or spells or other means; allow a player's creative thinking to be beneficial to this process regardless of normal game mechanics.
\item \textbf{\textit{Stop the Ghost.}} The PCs can choose to exorcise the \textbf{ghost} from Budanyek's body, via the means detailed in the \textbf{ghost} Possession ability. Choosing this action always results in combat as the \textbf{animated quartet} quickly moves to defend the \textbf{ghost}, fighting until it is destroyed. Because Budanyek is in no shape to fight and Heinrosch is obsessed with his legacy, Heinrosch only takes the Dodge action in the hopes of finishing the sonata amid the fracas.
\end{itemize}

\subsubsection{Treasure}
Heinrosch/Budanyek holds "The Gilded Quill" pen (see 1). It's worth 10 gp purely for the novelty of it being a gold quill, but a savvy negotiator might be able to get someone to pay up to 25 gp for it as a curiosity if they tell the tale of Heinrosch or invent a compelling fiction.

If you used the Collector adventure hook, however, the collector will excitedly pay triple the fee he promised the PCs for this one item—150 gp—and will go above and beyond any budget limits in order to acquire this piece.

\subsection{Evaluating the Composition}
If Budanyek completes the composition, he must make a Charisma (Performance) check (with a +2 modifier); when you make this roll, do it publicly so the players can see the die roll result. However, unless restorative magic has been used on Budanyek, do note that he has disadvantage on ability checks due to exhaustion.

If a PC completes the composition, they must make their own Charisma (Performance) check.

Once you have the result, consult the "Magnum Opus?" table to determine how well the composition turned out.

\begin{DndTable}[header=Magnum Opus?]{XX}

DC & Result\\
10 & Middling. It's . . . fine? Children who are just learning to play their instruments seem to enjoy performing it.\\
15 & Good. A work that will likely find great appreciation among the masses, but often does not excite an expert's ear. Heinrosch's previous "masterworks" were of this level.\\
20 & Incredible. A masterpiece that will surely be remembered forever.\\
25 & Astounding. A work of unsurpassed genius that only one-in-a-million composers could even attempt.\\
35 & Legendary. Undoubtedly among the finest works ever produced. Truly a magnum opus of magnum opuses.\\
\end{DndTable}

\section{Concluding the Adventure}
If Budanyek is not saved, the PCs find his gang waiting outside when they leave the college. They're heartbroken at the news (explaining who they are to the PCs if they haven't met before). If the PCs brokered a deal with the gang, the teens now refuse to fulfill their end of any bargain, citing Budanyek's death as a deal-breaker. Resentful of the PCs' failure, Rabbit and the gang disparage the PCs in whatever locale they're going to next; while in that city/town, the PCs have disadvantage on any Charisma ability checks until they can repair their reputation.

If Budanyek is saved, he gets halfway through offering his thanks before fainting from exhaustion. Budanyek's gang waits just outside the college and are beside themselves with joy when they're reunited with their friend (perhaps confusing the PCs if they haven't met before); they gladly pay any reward that was agreed to. If the PCs refuse the reward (or never asked for one in the first place), Rabbit and the gang talk up the PCs in whatever locale they're going to next; while in that city/town, the PCs have advantage on any Charisma ability checks for one week.

If Heinrosch completes his final composition and the result on the "Magnum Opus?" table is 15–24, Heinrosch is thrilled and asks the PCs to give it an impressive title and share it with the world. If the result is 25 or higher, Heinrosch reacts the same as above, but afterward breaks down, weeping with joy and whispering, \textit{"I've done it. I've truly done it."} The \textbf{ghost} then fades into the afterlife and disappears for good. If the PCs do share the composition with the world, and the result is 25 or higher, they hear it being played any time such music might be played, such as at royal's or noble's party.

If Heinrosch completes his final composition and the result on the "Magnum Opus?" table is 14 or lower, Heinrosch is at first furious, naming the PCs incompetent hacks, but quickly turns morose, lamenting his legacy. As he continues muttering, before long, he fades away into the afterlife and disappears for good—\textit{or so the PCs think}. Heinrosch is enraged at the failure to cement his legacy and, in the future, seeks revenge on the PCs, possessing the leader of a group of bandits that comes to the college to loot whatever isn't nailed down; in this new warrior body, Heinrosch bides his time, striking when it's most opportune, with the goal of possessing either the PC that he previously possessed or the PC with the highest Performance modifier.

\chapter{Harvest Fangs}

\begin{DndTable}{c}

\textbf{An adventure for four to five characters of 3rd level}\\
\end{DndTable}

\section{Adventure Background}
\DndDropCapLine{I}{}n the small farming village of Sycamore Crossing, this year's harvest is inexplicably failing. Faced with starvation (or worse), the villagers have fallen under the fell sway of Sheriff Ialos, who promises a full harvest if the village sacrifices living subjects to the earth they till.


The farmers have already kidnapped two travelers and fed them to a strange, living plant creature in the woods. The positive results on the crops were noticeable immediately, and so now even the local druid and the village priest have agreed to go along with Ialos' dark solution. Last night, a traveler named Rinus Aldebrand came to the village and was taken prisoner by the growing cult.

\section{Adventure Hook}
Tosef Smallfoot, the sheriff of Yrunby, hires the PCs to travel to Sycamore Crossing to look for a man named Rinus Aldebrand. The sheriff offers 150 gp per PC if they can return Rinus to Yrunby.

Using the above hook, you can choose one of the following identities for Rinus Aldebrand, which best suits your PCs' motivations:


\begin{itemize}
\item \textbf{\textit{Rust Runner.}} Rinus Aldebrand, a thug who recently ran with a gang called the Rust Wolves, fled north when their operation was destroyed. The sheriff hires the PCs to capture Rinus so he can stand trial. Two days ago, he was seen heading in the direction of the village of Sycamore Crossing.
\item \textbf{\textit{Tome Hunting.}} Rinus Aldebrand is a young scholar (and aspiring wizard) who has gleaned the location of a lost tome relevant to his arcane research. Yendra Aldebrand, his mother (and a minor noble), comes home from a trip to find a note from her son saying that he traveled to Sycamore Crossing to confirm the location of the tome. Yendra goes to the sheriff of Yrunby in a panic, afraid that her son would go after "that damned book of his" on his own and get himself killed. The sheriff hires the PCs (on Yendra's behalf) to go to Sycamore Crossing to make sure Rinus stays safe. One of the Aldebrand servants told Yendra that Rinus left two days ago.
\end{itemize}

The sheriff describes Rinus' horse as a dappled chestnut mare named Vossi, and he notes that the horse knows her name. Sycamore Crossing is a day's travel from Yrunby (or whatever starting location you choose).

\section{Sycamore Crossing}
The PCs arrive in Sycamore Crossing searching for Rinus. Unbeknownst to them, the village's cultists have kidnapped the man and are currently taking him to the woods north of town.

\subsubsection{Talking With Villagers}
All villagers are reluctant to speak with the PCs—and refuse if they ask about recent visitors, indicating that they don't appreciate strangers poking their noses into their business. While only a few villagers participated in the kidnapping, the rest don't want to risk the cultists' ire. If threatened, the villagers tell the PCs that if they want to know more they'll need to talk to the sheriff.

\subsubsection{Searching the Village}
With a successful DC 11 Intelligence (Investigation) check, the PCs find Rinus' dappled chestnut steed comfortably stabled by the village's tavern. If the horse's name (Vossi) is spoken aloud, the mare turns to look at the speaker and flicks her head upward as if saying hello.

If the PCs fail that check and don't otherwise make their way to the tavern, a young woman holding a handful of flowers offers one to a PC and discretely informs them that a visitor came to Sycamore Crossing last night and booked lodging at the tavern; she moves away quickly before she can answer any questions (whispering "Please, I can't!") and refuses to say more if followed.

\subsubsection{The Tavern}
Sycamore Crossing's sole tavern doesn't have a name. It serves green anise liquor and red wine sweetened with beet sugar. Patrons can order a thick garlicky parsnip stew with occasional chunks of mutton that become stuck in their teeth.

The bartender, Lenka (CG human \textbf{commoner}), is a large brunette woman who looks like she once knew her way around a donnybrook. Though she's tough, she's in her middle years now and no hero; she's frightened of the sheriff and hasn't been brave enough to cross him, even though she finds what he's doing repugnant.

When the PCs come in and start asking questions, she initially thinks that they're secretly working for Ialos, to see if any in the town are disloyal to the sheriff. Whenever anyone voiced any objection to the sheriff 's plan, he warned that everyone had better stay in line if they know what's good for them, and that if anyone dared betray him, he'd know who was responsible.

As a result, when questioned, Lenka attempts to deflect, claiming to know nothing, having never heard of any Rinus. She's not sure who that mare belongs to, telling them they would have to ask the stableboy, but he's out trying to fetch a horse that ran off.

A successful DC 10 Wisdom (Insight) check sees through Lenka's deceptions. If the PCs get angry with her or attempt any kind of intimidation, she panics and leans forward, saying quietly, "Please, please. I'm loyal, I swear. I just . . . I'm no good at . . . this kind of thing. You can tell the sheriff I'm steadfast. I promise."

At this point, if the PCs are honest with her about why they're there, her eyes get big, and she looks like her salvation is at hand. She asks them to come into the backroom with her so they can talk freely. She calls over to her barmaid and asks her to watch the bar. "Hurry," she whispers, "\textit{we might still be able to save him}."

\subsection{What Lenna Knows}
Lenka can provide the following information:


\begin{itemize}
\item The local farmers all fear that this year's crop will fail, which would spell disaster for Sycamore Crossing.
\item Sheriff Ialos said they could improve the crop by "offering sacrifices to the land." She's not sure what he meant by "sacrifices to the land" exactly. But there were a couple of folks who came to town recently who disappeared and never came back.
\item A month ago, a fast-talking gnome came to town. He talked a big game all night in the bar, then never came down from his room. She went up to find it empty, but with his belongings still there. No one ever saw him again. Shortly after that, the crop began to improve.
\item Then, about two weeks ago, there was a dragonborn with cold eyes who looked like he slit your throat sooner than he'd give a smile. After that one disappeared, the crop rebounded even more strongly.
\item All along, the main resistance against the sheriff came from Lenka; the local druid, Broshod; and the priest, Andel. After the crops started doing well, Broshod convinced herself that this was all part of the natural order, and Andel said that no matter how vile, ultimately it would save more lives than it cost. Leaving Lenka seemingly as the lone voice of dissent.
\item In the time that's passed since the last sacrifice, the crop has begun to falter again.
\item Last night, Ialos concocted some charge as a pretext to arrest Rinus shortly after his arrival in town. She saw the sheriff dragging the young man away with his hands bound.
\item Lenka heard Broshod talking about the old standing stones in the woods north of town and thinks that must be where the sheriff is doing his foul deeds.
\item Only one path leads to the woods in the north; Lenka can direct the PCs to it.
\item If asked about the sheriff, Lenka says that he's been sheriff for twenty years or more. Always been a good sort until recently, when all this unpleasantness started. She doesn't have any idea what might have happened to change him.
\end{itemize}

\section{Sheriff Ialos's Awakening}
Ialos takes his prisoners to a sycamore grove in the forest north of town. While chasing a bandit out of town, he discovered a \textbf{tripwire patch} in the center of the standing stones that have stood there in that grove for as long as anyone could remember. Everyone always avoided the stones out of superstition, but when he was chasing the bandit the tripwire patch swallowed them both. For some reason, however, the plant later regurgitated Ialos alive.

While in the patch's gullet, Ialos had a religious awakening and came to believe that the plant had chosen him as its servant. His fervent belief soon manifested itself in the ability to cast divine spells.

Over the last several years, Ialos has returned regularly to the grove to commune with his master, bringing with him another sacrifice, such as a common criminal or other ne'er-do-well. He noticed long ago that everything in the area—from birches to beets to barley—grows better after a sacrifice.

Ialos viewed the village's recent agricultural issues as the opportunity to go public about his beliefs. The sheriff has subsequently recruited several villagers into his now-burgeoning cult.

If the PCs are to save Rinus, time is of the essence; so tarrying on the trail—or taking any kind of rest—almost certainly results in his demise.

An unlabelled version of the map can be found here.

\subsection{1. Entrance}
\begin{DndReadAloud}{}
Well-worn ruts in the road cut a clear path north from Sycamore Crossing, weaving between the gentle sloping hills outside the village. The path ends near a narrow, treelined trail. By the entrance, a small hickory wagon rests, seemingly abandoned.

\end{DndReadAloud}

Ialos brought Rinus north in a hickory wagon that is drawn by a small \textbf{draft horse}. Ialos parked the wagon here, inviting Andel and six devoted cultists to witness the ceremony. The \textbf{draft horse} is now devouring a bag of oats hung around its neck.

PCs who succeed on a DC 10 Wisdom (Animal Handling) check can approach the \textbf{draft horse} without spooking it. If the check succeeds by 2 or more, upon examining the animal, the PC also determines that the \textbf{draft horse} stopped pulling the wagon within the last fifteen minutes.

\subsection{2. Open Trail}
\begin{DndReadAloud}{}
On the trail, the foliage begins to grow thicker, casting ominous shadows. Birds sing a throaty, legato melody, breaking the near-silence of any footfalls. Oak and sycamore trees line the entrance to the lair. The deep wagon wheel ruts in the road come to an end here, though a single narrow wheel rut continues ahead along a new path.

\end{DndReadAloud}

\subsubsection{Creatures}
PCs who succeed on DC 13 Wisdom (Perception) checks notice three ravens sitting high in an oak tree. The ravens are friends of Broshod and watch for trespassers for the druid. If the birds spot the PCs, the ravens alert Broshod about the incoming intruders by emitting a series of loud caws.

A successful group DC 14 Dexterity (Stealth) check allows the PCs to move through the area without detection.

\subsection{3. The Druid's Grotto}
\begin{DndReadAloud}{}
The path ahead turns briefly east before quickly curving north once more, bending around a low, rocky hill; as it does so, the mouth of a small, earthen grotto dug into a hillside comes into view, opening along the western edge of the path.

\end{DndReadAloud}

After learning of Ialos' scheme, Broshod took up residence in this grotto so that she could monitor his activities; while the druid has rationalized the sheriff's actions as aligning with the natural order, Broshod fears that such fervor will lead to unnecessary sacrifices.

A successful group DC 15 Dexterity (Stealth) check allows the PCs to move without being detected through this area; this check is made with disadvantage if the ravens in 2 noticed the PCs and called out their alarm.

\subsubsection{Creatures}
Broshod (N human \textbf{druid}) and two \textbf{giant badger} are in the grotto.

\subsubsection{Tactics}
The badgers move to confront the PCs outside the grotto, while Broshod stays back. She tries to remain at range, but if the PCs engage her in melee combat, she casts \textit{thunderwave} to try to push them away. Because the grotto doesn't have any other exits, the druid and her badgers fight to the death.

\subsubsection{Treasure}
Searching the grotto uncovers an \textit{herbalism kit}, two potions of healing, and a \textit{potion of animal friendship}.

\subsubsection{Development}
If Broshod casts \textit{thunderwave}, it causes the sheriff and his cultists in 5 to go on high alert; if they see the PCs, they attack them on sight.

\subsection{4. The Divide}
\begin{DndReadAloud}{}
Beyond the grotto, the path continues north; 70 feet ahead, it continues further north, but also branches off to the east.

\end{DndReadAloud}

Upon reaching the divide, the PCs must decide whether to continue north or veer to the east. Any PC may make a Wisdom (Survival) or Intelligence (Investigation) check; consult the "Inspecting the Tracks" table to determine what they learn:

\begin{DndTable}[header=Inspecting the Tracks]{XX}

DC & Information Learned\\
8 & A number of boot prints have gone through this area recently, to the north, the east, and back south.\\
13 & Nine or ten sets of boot prints came through this intersection recently. More than half of them headed east.\\
18 & The boot prints heading north look like they're sunk deeper into the earth, suggesting the people might have been carrying something heavy.\\
\end{DndTable}

\subsection{5. The Shrine}
\begin{DndReadAloud}{}
The path here opens into a wide, oblong clearing approximately 50 feet in diameter. To the southwest, near the edge of the woods, is a small shrine. As you look around, five figures rush into the glade from the southeast. Four are dressed as farmers and carry pitchforks, while the fifth walks with the air of legal authority.

\end{DndReadAloud}

Perhaps as a way to assuage his guilt, Andel built a shrine to the god of life in honor of those sacrificed for the greater good.

\subsubsection{Creatures}
Ialos (NE human \textbf{cult fanatic}) is the sheriff of Sycamore Crossing. He is accompanied by his deputy, Edmon (NE human \textbf{bandit captain}), and four human \textbf{cultist}. The cultists wield pitchforks instead of scimitars; use the same statblock for \textit{scimitar}, except change the damage type to piercing.

Ialos knows the PCs are outsiders who must not interrupt the sacrifice. If he was alerted by Broshod's \textit{thunderwave} from 3, he attacks on sight. If he wasn't alerted by the \textit{thunderwave}, he doesn't attack immediately; instead, he tries to get the PCs to leave, saying "This isn't your concern, strangers. If you leave now, no harm will come to you." A successful DC 14 Wisdom (Insight) check reveals that this is a lie; in truth, if the PCs leave, Ialos sends Edmon and four \textbf{cultist} to follow after them, as they've already seen too much.

\subsubsection{Tactics}
Ialos orders Edmon and the \textbf{cultist} into melee and uses \textit{hold person} on his first turn to incapacitate the most dangerous-looking PC. Here in the clearing, Ialos, Edmon, and the \textbf{cultist} fight to the death. A single \textbf{cultist} may surrender—but only after Ialos and their other companions are dead.

\subsubsection{Treasure}
Inspecting Ialos' body discovers a \textit{periapt of wound closure}. The sheriff took this from Rinus and hasn't yet attuned to the item. If he survives, Rinus gifts this item to the PCs.

\subsubsection{Development}
If the PCs take a \textbf{cultist} prisoner, they warn that their fellows are about to sacrifice Rinus to the land. If the \textbf{cultist} has 4 or fewer hit points remaining, they also warn them of the random nature of the rune in 6.

\subsection{6. The Rune}
\begin{DndReadAloud}{}
A soft, glowing blue-colored light looms ahead on the path. As you approach, you see, drawn into the ground, a large, strange rune. The lines of the rune look to be filled with some kind of powder, from which the glimmer emanates. It shifts and changes shape as you look at it, transforming from one character to another.

\end{DndReadAloud}

Years ago, Ialos had a wizard etch this rune into the ground outside the grove. The sheriff hoped the rune would discourage any trespassers from proceeding further.

\subsubsection{Compulsion}
A creature that moves within 15 feet of the rune must succeed on a DC 12 Wisdom saving throw or be forced to examine the rune more closely. On a success, the creature is immune to this effect for 10 minutes; the creature may still choose to examine the rune.

\subsubsection{The Rune's Magic}
If a creature examines the rune closely, it becomes subject to a random effect that is determined by rolling a d6 on the "Rune Effects" table.

\begin{DndTable}[header=Rune Effects]{XX}

d6 & Effect\\
1 & \subsection{Faerie Mark}A creature that reads the rune must succeed on a DC 12 Dexterity saving throw or become subject to the effects of the \textit{faerie fire} spell, but only to creatures with the Plant type; a creature of any other type confers no benefit from the \textit{faerie fire}. This effect lasts for 1 hour.\\
2–5 & \subsection{Unsettling Mark}A creature that reads the rune must succeed on a DC 12 Wisdom saving throw. On a failure, the creature is turned for 1 minute or until it takes any damage. This effect works like the cleric's Channel Divinity: Turn Undead class feature, except it can affect any creature type, and it treats the \textbf{tripwire patch} in 10 as the origin of the effect.\\
6 & \subsection{Reviving Mark}A creature that reads the rune gains 5 (1d10) temporary hit points; at the same moment, a nearby plant withers and dies. A creature can only be benefit from this mark once per long rest.\\
\end{DndTable}

\subsection{7. The Boulder}
\begin{DndReadAloud}{}
The trail north curves east after 80 feet. A large gray boulder sits in the center of the path.

\end{DndReadAloud}

There are no signs of recent disturbance of the trees nearby, so the PCs can easily determine that the boulder has been where it is for a very long time; it perhaps once rolled down from some not-too-distant mountain and has remained there ever since. It's completely innocuous—it's just a big rock—but PCs may be suspicious of it, and this distraction may cause them to be less cautious when approaching 8.

\subsection{8. The Trap}
To thwart or slow any do-gooders like the PCs from interfering with their sacrifices, the cultists placed a pit here where the path narrows just past 7.

\subsubsection{Pit Trap}
This simple pit trap is a 10-foot-square, 10-foot-deep hole dug in the ground. The hole is covered by canvas anchored on the pit's edge and camouflaged with dirt, leaves, and other debris. Spotting the trap requires a successful DC 12 Wisdom (Perception) check. When a creature enters the trap area, it must succeed on a DC 12 Dexterity saving throw to avoid falling into the pit.

A creature falling into the pit takes 3 (1d6) piercing damage from the few hastily carved wooden spikes at the bottom, plus 3 (1d6) bludgeoning damage from the 10-foot fall. To climb out of the pit, a PC must succeed on a DC 12 Strength (Athletics) check; using ropes or other climbing gear grants advantage on the roll.

If the PCs detect the trap, they can easily move around it, either by skirting the edge of it on the path or moving through the trees.

\subsection{9. Trail's End}
\begin{DndReadAloud}{}
As you continue south, the foliage grows thicker and thicker the farther you progress. About halfway down this spur, the ground becomes covered with thorns, burs, and pricker bushes. From the south, you hear the faint sounds of chanting, which is then followed by what can only be the muffled cry of a person in distress.

\end{DndReadAloud}

\subsubsection{Bramble Hazard}
The last 70 feet of this southern part of the trail is covered in thorny plants and brambles. If a creature moves at half speed, it can traverse the path normally.

To move at full speed, a creature must succeed on a DC 12 Dexterity saving throw every 10 feet or take 1 piercing damage from the thorns and brambles. A creature wearing no armor and/or boots has disadvantage on the saving throw; a creature wearing boots and heavy armor automatically succeeds.

\subsection{10. The Standing Stones}
\begin{DndReadAloud}{}
Six sycamore trees are interspersed throughout this broad clearing, and the trees lining the outer edge slant away from the center, as if they are giving something within a wide berth. Seven standing stones are arranged in the center of the grove. There, a priest stands, in the midst of offering a muttered prayer as he lowers an unconscious body to the ground. Two cultists stand near the body, one of them cutting the palm of their right hand with a knife. As the blood drips onto the ground at their feet, thick vines on the ground start to rustle and writhe.

\end{DndReadAloud}

By the time the PCs arrive at 10, the \textbf{cultist} have carried Rinus' unconscious form to the center of the standing stones.

\subsubsection{Creatures}
Upon seeing the PCs enter the area, \textbf{Andel} (NG human \textbf{acolyte}, with \textit{command} prepared instead of \textit{sanctuary}) raises his hands in a "hold" gesture—in the hopes of finding a peaceful solution to the current situation.

The two \textbf{cultist}, meanwhile, raise their pitchfork to fight. The \textbf{cultist} wield pitchforks instead of scimitars; use the same statblock for \textit{scimitar}, except change the damage type to piercing.

The \textbf{tripwire patch}—which is located at the center of the standing stones—is alert and ready to eat.

\subsubsection{The Priest}
\textbf{Andel} attempts to parley with the PCs, trying to convince them that the people of Sycamore Crossing had no choice and that what they've done is the lesser of two evils; he truly believes this, and that his god—the god of life—has endorsed this action (in this, however, he is misguided). If the PCs are not persuaded by his words, he offers to give them the treasure noted below.

\subsubsection{Tactics}
Once combat ensues, the \textbf{tripwire patch} uses one Tripwire Vine attack to grapple Rinus and flails the other toward a creature at random (choosing from any creature within range, not just the PCs). It uses Drag to pull any creature it grapples toward its maw so that it can use its Bite to devour its prey.

If the patch bites an NPC, they are killed instantly. If that happens, unless the PCs continue to attack, the patch returns to its slumber and digests its food. \textbf{Andel} won't attack with violence, but if a PC draws a weapon and/or makes a hostile move toward him or one of the cultists, he casts \textit{command}; the command he gives is "flee."

If \textbf{Andel} or either of the \textbf{cultist} take damage, the wounded creature attempts to flee.

\subsubsection{Treasure}
\textbf{Andel} carries two potions of greater healing and a \textit{robe of useful items} tucked into his satchel.

\subsubsection{Lore}
If a PC spends 10 minutes examining the area, a successful DC 14 Intelligence (History or Religion) check determines that the standing stones have been in this location for hundreds—perhaps thousands—of years; there are also a number of markings (unreadable now due to the ravages of time) that suggest an older culture may have worshipped this creature if anything the villagers believed about it is true.

\section{Concluding the Adventure}
If the PCs save Rinus, he thanks them profusely. From there, the PCs can return him to the sheriff of Yrunby, and then otherwise proceed as the adventure hook you used suggests.

If the PCs didn't encounter the sheriff in 5, the PCs find Ialos and the four cultists in 1, waiting for \textbf{Andel} signal that the sacrifice has begun. Ialos and the four \textbf{cultist} attack Rinus and the PCs on sight. The sheriff orders his minions to make nonlethal attacks to ensure replacement sacrifices; if the PCs point out that the \textbf{tripwire patch} is dead, he says "Then you've doomed us all" and attempts to run. If he's captured, the PCs can reasonably conclude that (also) taking him to the sheriff of Yrunby is the proper action. Ialos also flees the battle if he starts his turn at 11 hit points or fewer, commandeering the wagon if possible. The \textbf{cultist} fight to the death unless Ialos flees; if he does, they too attempt to run.

\chapter{The Pavilion of Whispers and Wonderment}

\begin{DndTable}{c}

\textbf{An adventure for four to five characters of 4th level}\\
\end{DndTable}

\section{Adventure Background}
\DndDropCapLine{T}{}ucked into the feylands is The Pavilion of Whispers and Wonderment, a magical building formed of softly glowing amber bubbles resting upon titanic lily pads, where it floats on the surface of a wide lake of unmatched beauty. An arching bridge connects The Pavilion to the mistshrouded shoreline.


The Pavilion is the home of the Maestro, a connoisseur of secrets, scandals, and slander. The Maestro has numerous contacts among planar sages and other lore hungry creatures. He archives his accumulated information in an extraplanar library and guards his home with a variety of strange creatures and enchantments.

Over the years, the fame of the Maestro has grown far and wide, as have the stories of his great power and wealth. But the Maestro seldom lowers himself to trade information, unless tidbits of lore of greater value are offered. He prefers, instead, to hoard his information like a dragon sitting upon an enormous pile of gold coins.

Recently, the Maestro acquired a tome of forbidden knowledge called \textit{The Mardom Codex}. While investigating its pages, he fell victim to the book's magical curse and now exists in a paranoid nightmare of terrifying, chaotic delusions. This magical malady has also spread to infect the entire Pavilion. An unknowing visitor might expect to meet a savvy, reasonable information broker, but what they encounter may instead make them the next victim of the Pavilion of Whispers and Wonderment.

\section{Adventure Hook}
The following can used to start the PCs on this adventure:


\begin{itemize}
\item \textbf{\textit{Information Seekers.}} The PCs need a particular bit of lore to defeat a powerful foe, and the only entity possessing this information is the information broker known as the Maestro. Arriving at The Pavilion, the PCs can try to negotiate with the Maestro's minions or attempt to gain the information they need by means of a heist or raid.
\item \textbf{\textit{Noble Quest.}} A noble or patron of the PCs engages their services to collect a scroll containing information embarrassing to them (or implicating them—perhaps falsely—in some wrongdoing), which has fallen into the hands of the information broker known as the Maestro. The PCs are offered 400 gp each to acquire the scroll from the Pavilion of Whispers and Wonderment.
\item \textbf{\textit{Mercy Mission.}} A noble or patron of the PCs has fallen ill, struck down by a mysterious malady that has vexed the finest healers. The noble/patron's majordomo offers the PCs 400 gp each to travel to the Pavilion of Whispers and Wonderment to see if the information broker known as the Maestro might know of a cure.
\end{itemize}

\section{The Pavilion}
The dome-like buildings that comprise the Pavilion look like a collection of huge orange and yellow opaque bubbles resting on gargantuan lily pads. The surface of the bubbles is harder than steel, perfectly smooth, and seems to contain no windows or doors when viewed from the outside.

The only visible entrance is the one connected to the bridge in 1. The other doorways are hidden by magic; a successful DC 15 Intelligence (Investigation) check sees through the illusion. Alabaster bridges, lavishly adorned with gold and ivory, connect the bubbles.

The interior of the Pavilion is brightly lit by magical lights and richly appointed, though the style of the decorations seems odd, even horrific, to the common eye. The walls, floors, furniture, and other accoutrements are formed of circles and gracefully flowing curves. The circular doorways do not have traditional doors, but some have magical barriers in place to keep creatures out.

\begin{DndSidebar}[float=!b]{A peaceful home or dangerous lair?}
If you wish to use the Pavilion as a recurring location to get lore or specialized information to the PCs, you can present this material as if the Maestro never acquired \textit{The Mardom Codex}. In that circumstance, the creatures in this adventure do not attack first and treat with the PCs politely.



\end{DndSidebar}

An unlabelled version of the map can be found here.

\subsection{1. The Welcoming Bridge of Welcome}
\begin{DndReadAloud}{}
A ten-foot-wide bridge of ivory and gold leads up from the shore to an arched opening set into the side of the towering amber bubble. A field of energy fills the archway like a shimmering door. A small green shrub, with an orange bow tied on its topmost branch, sits in a decorative ceramic pot to one side of the archway. An ornate sign on the bridge reads "The Welcoming Bridge of Welcome."

\end{DndReadAloud}

When the PCs approach, the shrub greets them and introduces itself as Quince. It quickly adds that the Maestro isn't receiving visitors at this time.

Convincing Quince to allow the PCs entry requires a successful DC 12 Charisma (Deception, Intimidation, or Persuasion) check—with advantage if the PCs legitimately have rare, secret, or important information that they can trade (Persuasion) or say that they have such information (Deception).

\subsubsection{Energy Field}
The magical barrier in the archway was created by the \textit{wall of force} spell; it can be deactivated by the wand hidden in Quince's pot.

\subsubsection{Creatures}
Quince (N \textbf{awakened shrub}) defends itself if attacked but poses no real threat to the PCs. However, if it feels threatened, it alerts the household via an \textit{alarm} spell set into its decorative pot.

\subsubsection{Treasure}
A small wand hidden at the bottom of Quince's pot can be used by any creature to deactivate the energy field in the archway to 2.

\subsection{2. The Restful Hall of Lounging}
\begin{DndReadAloud}{}
The walls of this large, multileveled chamber arch upward to a domed apex 50-feet-high. The center of the room is an ornamental and elegant foyer with stone benches lining the walls. Throughout the entire chamber, 10 feet overhead small globes of magical light hang from ropes of braided silk like an ornamental lattice. A metal plaque on the wall reads "The Restful Hall of Lounging." To the east, a curving platform stretches from a doorway, and silver baskets of fruit hang on the wall. To the north, a short flight of stairs descends 5 feet to a recessed area that branches toward other chambers.

\end{DndReadAloud}

When the PCs enter the chamber, the Butler, the Maestro's majordomo, immediately notices their arrival.

\subsubsection{Eastern Kitchen}
The baskets on the walls magically produce a wide variety of fresh fruit, but cease functioning if removed from the wall.

\subsubsection{Creatures}
\textbf{The Butler} is a \textbf{lobe lemur} dressed in a yellow vest with brass buttons. While under the effects of \textit{The Mardom Codex}, he lurks among the ropes above the PCs heads and attempts to ambush them with hit-and-run attacks.

Add the following action to \textbf{the butler} statblock:

\subparagraph{Lightfall (Recharge 5\textendash 6)}
The Butler shakes the ropes hanging above, dislodging a rain of the glowing lights. All creatures in the central area of the room must make a DC 12 Dexterity saving throw; a creature takes 7 (2d6) acid damage on a failed save, or half as much damage on a successful one.

\subsubsection{Treasure}
In a pocket of \textbf{the butler} vest is a small wand that deactivates the barrier door to 9.

\subsection{3. The Exuberant Treasury of Wonderful Things}
\begin{DndReadAloud}{}
A short flight of stairs leads into a twenty-five-foot-wide circular dome. A vast array of items hangs on the walls, displayed and organized with care. A glass cloche rests atop a stone plinth that stands on the end of a raised dais on the western side of the room. A metal plaque near the entrance reads, "The Exuberant Treasury of Wonderful Things."

\end{DndReadAloud}

\subsubsection{Displayed Items}
The objects on the walls are a mixture of oddities and trinkets of no immediately apparent value, such as a ribbon from a child's bow, half a broken teacup, or a drawing of a somewhat ferocious dragon, etc. Each is accompanied by a small, handwritten card with names like "Last Chess Piece of the Forgotten King" or "Broken Sword of the Nearly Famous."

\subsubsection{The Glass Cloche}
The glass cloche atop the plinth initially seems to be empty, but when the PCs approach it they can see that inside is what, from afar, looks to be a small (about 3" x 5") sheet of paper.

\subsubsection{Animate Objects Trap}
A trap guards the glass dome. One or more creatures entering the 5-foot area around the plinth triggers the casting of the \textit{animate objects} spell. 10 \textbf{animated object (small)} around the room come to life. Detecting the trap requires a successful DC 15 Wisdom (Perception) check before any creature enters the 5-foot area around the plinth. The \textbf{animated object (small)} have 15 hit points, but otherwise have the same statistics as described by the spell.

\subsubsection{Disarming the Trap}
Hidden in one of the steps of the dais is a hidden panel; a successful DC 15 Wisdom (Perception) check notices it. Inside is a small lever that deactivates the Animate Objects trap.

\subsubsection{Treasure}
The objects on the walls are a mixture of nonmagical (and worthless) mementos and minor magical trinkets. But among the other items are a blue sapphire elemental gem marked as "The Eye of Hava" and a \textit{bag of holding} mislabeled as "Gamda's \textit{Bag of Devouring}."

Beneath the glass cloche is a single card, facedown. A small metal plaque reads, "Card from the \textit{Deck of Many Things}."

This is indeed a card from a \textit{deck of many things}. When a creature turns the card over—or in any way attempts to determine which card it is—the magic of a randomly selected card from the deck of many things takes effect. If you wish to control the chaos presented by the \textit{deck of many things}, you could choose the card instead of it being randomly selected; Comet, throne, or vizier cards are all beneficial and relatively innocuous (at least in comparison to the other cards in the deck).

Or, if you wish to avoid the deck entirely, replace the card with a magic item associated with knowledge or secrets, such as: a \textit{helm of comprehending languages}, a \textit{ring of mind shielding}, or a \textit{medallion of thoughts}.

\subsection{4. The Poppet Playhouse}
\begin{DndReadAloud}{}
Toys are strewn across a decorative rug in this chamber; that, coupled with its brightly-colored walls indicate its function is as a child's playroom. Near the entrance, a metal plaque on the wall near reads, "The Poppet's Playhouse." Seemingly out-of-place, a glowing, rune-carved \textit{greatsword} hangs on the western wall. In the center of the room, three intricately carved, articulated wooden dolls play, serving imaginary tea to each other around a small table with a child-sized ceramic tea set.

\end{DndReadAloud}

The three living constructs in this chamber were the childhood playthings of the Maestro's sister, and he keeps them for sentimental purposes having long ago become immune to their nightmarish powers. The sword on the wall was once wielded by the man (now deceased) who murdered the Maestro's sister.

\subsubsection{Creatures}
Three \textbf{moppet} make this room their home. When they become aware of the PCs, they invite one or more of them to have pretend tea with them. Creatures who play along with the moppets become their friends. If befriended, the moppets gladly allow the PCs to take the glowing sword.

They are very forgetful, but reveal the following if asked:


\begin{itemize}
\item The toys in the Treasury (3) are "mean."
\item \textbf{The Butler} has the key to the Maestro's bedroom (7).
\item "Fluffy" (in 6) just wants to eat snacks.
\item The Maestro doesn't let them drink anything in the Decantory (5).
\item They think the Maestro is very sad and lonely.
\end{itemize}

Attacking the moppets, attempting to take the sword without permission, or refusing to have tea causes all three \textbf{moppet} to attack.

\subsubsection{Treasure}
The glowing greatsword is a luck blade with no wishes remaining.

\subsection{5. The Decantory}
\begin{DndReadAloud}{}
Three brass, wheeled carts stand against both the east and west walls of this comfortable and well-appointed chamber. Each cart holds an array of strange bottles, glasses, and other items suited to mixing and tasting wines, liqueurs, and spirits. The floor is covered with a thick, ornate rug. A metal plaque near the entrance reads "The Decantory."

\end{DndReadAloud}

The Maestro stores his collection of fine spirits, wines, and liqueurs in this chamber.

\subsubsection{Creature}
Two \textbf{doppelixir} hide in bottles on opposite sides of the room. They are bitter foes and use the appearance of the PCs to compete to be the first to get one of them to succumb to their Telepathic Urge ability. If attacked, the \textbf{doppelixir} eagerly team up to combat intruders.

\subsubsection{Treasure}
The carts hold 12 bottles of exquisite wine (worth 300 gp total), 6 bottles of strong spirits (worth 90 gp total), and assorted bottles of fruit juice and sparkling water. A successful DC 10 Wisdom (Perception) check also reveals what appears to be 3 potions of greater healing (two of them are, but one is actually a \textit{potion of poison}), a \textit{potion of water breathing}, and a \textit{philter of love}.

\subsection{6. Fluffy's Room}
\begin{DndReadAloud}{}
The center of this room is open to the water below, with a five-foot-wide walkway that encircles it just above the waterline. A plaque near the entrance reads "Fluffy's Room."

\end{DndReadAloud}

Years ago, the Maestro adopted a strange creature who had wandered into his lake. This is her home.

\subsubsection{Creatures}
"Fluffy" is a \textbf{hippopotamus} with red bows tied to her ears and a bejeweled leather collar around her neck.

When a PC enters the room, Fluffy surfaces, expecting snacks. She is very aggressive about being fed and refuses to let anyone pass without giving her some food.

\subsubsection{Treasure}
Fluffy's collar—which has her name set into it in gemstones—is worth 1200 gp.

\subsection{7. The Garden of Resting Tranquility}
\begin{DndReadAloud}{}
Flowering vines climb the walls of this chamber, and lush, green grass carpets the floor. A bed to the east seems to be naturally formed from a tree stump extruding from the wall. In the air, fireflies flicker and butterflies flitter. A wooden plaque near the entrance reads "The Garden of Resting Tranquility."

\end{DndReadAloud}

This is the Maestro's bedroom, but he doesn't spend much time here. It contains nothing of value.

\subsection{8. The Periculosum Viaduct}
\begin{DndReadAloud}{}
A forked bridge of alabaster, gold, and ivory connects the main section of The Pavilion to a pair of outbuildings floating on their own lily pads. A small metal plaque near the foot of the bridge reads "The Periculosum Viaduct." At the far end of the bridge to the northwest, a magical barrier fills the entrance to the building beyond, blocking any means of ingress.

\end{DndReadAloud}

The magical barrier was created by the \textit{wall of force} spell; it can be deactivated by the wand carried by \textbf{the Butler} (2).

\subsection{9. The Dizzying Library of Rumors and Whispers}
\begin{DndReadAloud}{}
This chamber seems to extend upward infinitely, with no ceiling in sight, and endlessly tall bookcases rise up toward the heavens. Other than their incredible height, the other very strange thing about these bookcases is that they move on their own, rotating around the periphery of the room, constantly repositioning themselves. Dozens of rolling ladders at different heights allow access to the higher levels, though only one ten-foot-high length actually reaches to the floor. Another shimmering barrier of energy fills the doorway to the northeast. A bronze plaque by the doorway reads "The Dizzying Library of Rumors and Whispers."

\end{DndReadAloud}

When the PCs enter the chamber, one of the bookcases starts asking them general questions such as their names, affiliations, race, occupation, and so forth. This is the Maestro using his Throw Voice feature. The Maestro can hear anything uttered in the library.

\subsubsection{Bookcases}
It is nearly impossible to find anything specific in the sprawling hoard of information until the bookshelves stop rotating. If \textit{The Mardom Codex} is destroyed or The Maestro is killed, the shelves stop moving, and specific information can then be recovered with a few hours of searching.

What other information the PCs discover in the library (and the repercussions of finding it) is beyond the scope of this adventure and is left for you to decide.

\subsubsection{Creatures}
The Maestro—a \textbf{veritigibbet}—is in 10, using his Throw Voice feature to pretend that the bookcase is talking to the PCs. He can hear anything uttered in the library.

\subsubsection{Development}
Once the Maestro receives three true answers to his questions, he says, "Well, you seem interesting at least. You may enter." He then opens the \textit{wall of force} that otherwise blocks the way to 10.

If the PCs lie to him (which he can unerringly detect), he calls them on it and courteously requests they don't lie to him again. If they lie a second time, he admonishes them for doing so and again politely asks for the truth. If they lie a third time, he says "Oh, don't be so boorish." and summons a \textbf{magma mephit} into the chamber to attack them. Once it's dispatched, The Maestro says, "Well, that was interesting at least. I guess you might as well come in," then opens the barrier to 10.

\subsection{10. The Icecarved Babbletorium}
\begin{DndReadAloud}{}
The floors, the curving walls, and the thirty-foot-high ceiling of this chamber seem to be fashioned from carved ice. Despite that, though, the temperature in the room is comfortable. Faintly glowing runes shine through the layer of ice on the floor. To the northeast, a pair of curved staircases lead up a twenty-foot-high, circular stone dais, atop which is a large desk covered with scrolls, quills and inkwells, and other scribal equipment. A small, narrow-framed creature with a disproportionally large head dressed in what looks like well-tailored finery sits in a massive chair behind the desk, facing the entrance. He stares unblinking at a huge book resting open in front of him. Arrayed in front of the desk on the dais is a flock of black and white birds standing unnervingly motionless.

\end{DndReadAloud}

\subsubsection{Eternal Ice}
The ice in this room is "eternal ice," a magical substance that never melts.

\subsubsection{The Codex}
The book on the desk is \textit{The Mardom Codex}. It is magical, cursed, and currently attuned to the Maestro; the only way to free the fey is to destroy the book. While the Maestro is within 30 feet of the book, it casts \textit{counterspell} as a reaction to any spell directly targeting the Maestro.

The book can be destroyed with any slashing or bludgeoning melee weapon. It has AC 10 and 15 hit points.

\subsubsection{Creatures}
On the dais are the Maestro—a \textbf{veritigibbet}—and a \textbf{swarm of penguins}; the swarm is weakened and only has 90 hit points.

The Maestro attempts to parley with the PCs, trying to extract any information or knowledge from them that he doesn't already know. Whether the PCs provide such information or not, the paranoia instilled in the Maestro by the \textit{Codex} eventually convinces him that the PCs are there to destroy him and/or steal his knowledge; at that point, he initiates combat if the PCs haven't already.

\subsubsection{Tactics}
The Maestro attacks from the dais, using his desk and chair for cover whenever possible; if any PCs near the stairs, he takes to the air, flying up to the ceiling and to the other end of the chamber.

The \textbf{swarm of penguins} toboggan down the stairs and toward the PCs to engage them in melee. The \textit{Codex} forces the Maestro to fight to the death, though the PCs can choose to simply knock him unconscious or otherwise incapacitate him, then deal with the cursed book.

\subsubsection{Treasure}
In one of the drawers of the desk is a \textit{pearl of power} and a \textit{tome of clear thought}, whose magic has been expended.

To save the tracking of the passage of years in order to determine when the tome's magic returns, you can instead have it restored at a certain PC level. Once the tome is identified, have a player roll a d20 and consult the "Tome's Magic Returns" table; the result determines what level the PCs need to achieve in order for the magic to return:

\begin{DndTable}[header=The Tome's Magic Returns]{XX}

d20 & Information Learned\\
1–12 & Level 11\\
13–17 & Level 10\\
18–19 & Level 9\\
20 & Level 8\\
\end{DndTable}

\subsection{Concluding the Adventure}
If the Maestro is freed from the clutches of \textit{The Mardom Codex}, the fey is grateful and rewards the PCs with the information they came to the Pavilion seeking, as well as 5 spell scrolls of their choice (up to 4th level). The Maestro can also become a powerful ally and valuable source of information going forward, provided the PCs can meet his price.

If the Maestro is killed, the PCs are free to search through the library and can discover the information they seek with a few hours of searching. If you don't wish the PCs to have access to this vast trove of information in perpetuity, the Maestro's library could have been warded against unauthorized access, making all of the books magically unintelligible to the reader, or the Pavilion could sink shortly after the Maestro's demise (but after the PCs get the information they seek).

\chapter{Rainforest Reckoning}

\begin{DndTable}{c}

\textbf{An adventure for four to five characters of 4th level}\\
\end{DndTable}

\section{Adventure Background}
\DndDropCapLine{T}{}he small, riverside jungle village of Bukofa is being menaced by hostile fey. At first, it was just a herd of tiny rainforest fey playing mischievous pranks on the villagers who were felling trees in the jungle. Now, something much worse is killing—and eating—Bukofa's inhabitants.


The people of Bukofa have long relied on the jungle for their livelihoods. They hunt deer and other forest animals for food and clothing, harvest fruits, nuts, and plants to eat and for medicinal purposes, and gather wood to build their homes and to make crafts.

In the past few months, as Bukofa's population swelled, the inhabitants began to chop down trees and burn back the jungle to grow crops and graze cattle. One of these trees belonged to a dryad named Delande. She became so distraught when her mahogany tree was cut down by the villagers that she threw herself into the river and drowned.

Delande's sister, Aminata, was enraged by her tragic death and sent a message to King Kashama, telling him that the jungle was being threatened and beseeching him to intervene. In response, the fey lord sent a herd of \textbf{aziza} to make life difficult for the loggers by befuddling them with magic and stealing their tools. This worked for a while, but when the logging resumed it became clear to Aminata that more drastic measures were needed, so she appealed to King Kashama again. This time, he commanded a rainforest ogre named Korug to travel there to frighten the villagers.

Since then, the fearsome ogre has been regularly sneaking into Bukofa to murder and consume its people. Unsurprisingly, the villagers are now too scared to venture into the jungle and are desperate to find a solution to the situation they find themselves in.

\section{Adventure Hook}
The following can used to start the PCs on this adventure:


\begin{itemize}
\item \textbf{\textit{Vexed Villagers.}} The PCs hear of the plight of Bukofa when they're traveling in the area or when word of the ogre attacks reaches whatever city they're in.
\item \textbf{\textit{Jungle Search.}} The PCs need—or are hired to retrieve—a rare plant needed for a spell component or alchemical project. The plant only grows in the jungle near the village of Bukofa. On their way there, they either pass through Bukofa, or someone from the village sees them and flags them down. (This hook could also pair with the "Mercy Mission" hook from "The Pavilion of Whispers and Wonderment," if you ran that adventure with your PCs.)
\item \textbf{\textit{Explorer Escort.}} The PCs are hired to search for a group of missing explorers, who went searching for a rumored lost dungeon in the jungle near Bukofa.
\end{itemize}

Once in Bukofa, Mayifa, the leader of the village elders, explains what has been happening, from the fey pranks to the recent slaughter of her people. The villagers are understandably now too frightened to hunt or gather food in the rainforest, and, as a result, their supplies are running dangerously low. She begs the adventurers to put a stop to the depredations of the fey before Bukofa's inhabitants starve. Mayifa believes the hostile fey are living at Emerald Cascade, a waterfall a few miles from Bukofa.

\section{Into the Jungle}
To reach Emerald Cascade, the PCs must head north into the hot and humid tropical rainforest, following a small stream. As they trek through the steamy jungle, mosquitos and other biting insects buzz around their heads, colorful parrots squawk at them as they fly past, and monkeys chatter to each other in the canopy overhead.

After an hour's walk, the PCs spot a tiny \textbf{musk deer} on the opposite side of the stream. The musk deer has been sent as an \textit{animal messenger} (as the spell) by the \textbf{aziza} living at Emerald Cascade with instructions to warn two-legged intruders away from the grove.

The animal has reddish-brown fur marked with white spots and lacks antlers; it stands barely one-and-a-half feet tall at the shoulder. The deer opens its mouth to speak to the adventurers, revealing a pair of four-inch-long canine teeth. In a high-pitched voice it says, "I'd turn back now if I were you. Only pain and suffering lie ahead."

Once it has delivered its message, it speeds off into the undergrowth. The PCs can travel the last two miles to Emerald Cascade without further incident.

\section{Emerald Cascade}
Located five miles upstream from Bukofa, Emerald Cascade is a grove of tall mahogany and kapok (silk cotton) trees grouped around the beautiful waterfall which gives the glade its name. The borders between the material plane and the fey realms are thin at Emerald Cascade, and visitors that cross over are common.

The fey grove has long been home to the \textbf{dryad} Aminata who lives in a mahogany tree growing above the waterfall. Following her appeal to King Kashama—a \textbf{rainforest king} fey lord—the \textbf{dryad} now shares Emerald Cascade with a herd of \textbf{aziza} and a \textbf{rainforest ogre}, as well as several other dangerous creatures.

\subsubsection{Trees}
The towering mahogany and kapok trees are between 100 and 150 feet tall and have 10-foot-diameter trunks. Smaller palm trees are 30 feet tall with 3-foot-diameter trunks. Trees can be climbed with a successful DC 12 Strength (Athletics) check.

\subsubsection{Foliage}
A huge variety of tropical bushes and other plants grow in the grove; many have colorful flowers or fruits and provide concealment. Roll 1d6 any time a creature touches one of the bushes; on a 1, the plant has poisonous leaves, and the creature must succeed on a DC 11 Constitution saving throw or take 3 (1d6) poison damage.

\subsubsection{Jungle Floor}
Vines and roots trail along the jungle floor, offering an increased risk of tripping. Creatures that move at their normal speed must succeed on a DC 8 Dexterity saving throw or fall prone; creatures moving at half speed automatically succeed on the saving throw.

\subsubsection{Stream}
Fed by the waterfall, the slow-moving stream that flows through the grove from north to south is 10 feet wide and 10 feet deep. Creatures who swim or fall into the stream attract the attention of a \textbf{swarm of quippers}, arriving in 1d3 rounds.

An unlabelled version of the map can be found here.

\subsection{1. Cascade Approach}
\begin{DndReadAloud}{}
Following the stream has brought you to an enchanting forest glade of tall mahogany and kapok trees. Colorful blossoms and ripe fruits grow on the bushes, the sweet smell of flowers permeates the air, and the sound of a waterfall can be heard through the trees.

\end{DndReadAloud}

This is where the PCs arrive in the grove if they've been walking upstream from Bukofa.

\subsubsection{Creatures}
An \textbf{aziza} with wings is hiding in the branches of a mahogany tree on the eastern side of the stream. She watches the PCs' approach, then gives the call of a tree frog to alert her companions in the anthills and treehouses (Areas 2–4). If she is spotted—which can be done with a successful DC 16 Wisdom (Perception) check—she retreats to 3 to warn \textbf{Kabou}, the \textbf{aziza} leader.

\subsection{2. Giant Anthills}
\begin{DndReadAloud}{}
Three roughly shaped pyramids of earth stand in a row between two tall silk-cotton trees. Each mound is around ten feet in diameter at the base and nearly fifteen feet tall. Several openings—each a foot and a half high—lead inside the mounds at different heights.

\end{DndReadAloud}

These large anthills have been abandoned by their original occupants and are now inhabited by \textbf{aziza}. Inside each mound, a network of 1-foot-wide tunnels connects several tiny earthen chambers, each 3 to 5 feet in diameter. The three anthills are connected underground by more tunnels, allowing the \textbf{aziza} to travel between them unseen. The anthills have AC 17 and 40 hit points.

\subsubsection{Creatures}
There are five \textbf{aziza}—two with wings (females) and three without (males)—in the anthills. With a successful DC 12 Wisdom (Perception) check, as the PCs approach, they catch sight of a tiny fey peering out at them from one of the tunnels.

The shy fey are curious (rather than hostile) about the PCs' presence in the grove, but if a \textit{find familiar}, a Wild Shape druid, or a reduced-size PC enters one of the anthills—or a PC pokes a sword or the like inside or otherwise acts aggressively—the \textbf{aziza} react angrily and attack.

PCs who can speak Sylvan can encourage the \textbf{aziza} to emerge from their hiding places by succeeding on a DC 15 Charisma (Persuasion) check. The \textbf{aziza} insist that the adventurers speak to their leader, \textbf{Kabou}, who lives in one of the treehouses in 3.

\subsection{3. Silkhome}
\begin{DndReadAloud}{}
This huge kapok tree towers over a hundred feet above the jungle floor and has large buttress-like roots that extend thirty feet up the trunk and radiate out in all directions. The trunk and its major branches are covered in large conical thorns. Miniature treehouses are nestled in the branches sixty feet from the ground.

\end{DndReadAloud}

The \textbf{aziza} have built three tiny treehouses in the branches of this tree, which they call Silkhome. Each treehouse is 4 feet square and 2 feet tall and surrounded by a 2-foot-wide wooden platform. Aziza-sized rope bridges connect the three dwellings.

Climbing up to the treehouses requires a successful DC 12 Strength (Athletics) check. For every 10 feet a creature climbs along the trees, it must succeed on a DC 10 Dexterity saving throw or take 1 (1d3) piercing damage from the sharp thorns on the trunk. \textbf{Aziza} and their \textbf{gliding frog} are used to climbing kapok trees and can avoid the thorns automatically.

\subsubsection{Creatures}
Six \textbf{aziza}—three with wings (females) and three without (males)—live in the treehouses; each male aziza rides a \textbf{gliding frog}.

Their leader is \textbf{Kabou}; she uses the same statistics as the other \textbf{aziza}, but she has 21 hit points.

\subsubsection{Tactics}
The \textbf{aziza} seek to confound the PCs and drive them away from the grove, but do not shy away from using deadly force if the PCs kill one or more of their number. If combat ensues, the three females take to the air, preferring to stay out of the range of the PCs and use their shortbows to rain their poisoned arrows down on them. The males each drop one \textit{terror melon} (see below) on their first turn of combat, then mount their \textbf{gliding frog} and use their glide ability to remain aloft (also using their shortbows from above) for as long as they can.

If the battle goes badly for the \textbf{aziza}, their comrades in the anthills (2) and the other treehouses (4) can come to their aid.

\subsubsection{Terror Melons}
In addition to their standard weapons, these \textbf{aziza} have four terror melons, 5-inch-diameter psychotropic fruits from the fey realms. They drop them onto intruders beneath the tree by making a ranged weapon attack with a +4 bonus.

On a hit, the melon explodes, spraying mind-warping juice in a 5-foot radius. Each creature in the area must succeed on a DC 11 Constitution saving throw or drop whatever it is holding and become frightened for 1 minute. While frightened, a creature must take the Dash action and move away from the point of impact on each of its turns. A frightened creature can repeat the save at the end of each of its turns, ending the effect on a success.

\subsubsection{Treasure}
Inside the treehouse are any remaining terror melons and two shiny red berries; one acts as \textit{potion of diminution}, the other acts as a \textit{potion of growth}.

\subsubsection{Development}
If the PCs attempt to parley with the \textbf{aziza}, \textbf{Kabou} is willing to talk but insists on doing this from the wooden platform in front of her treehouse, which requires the PCs to climb, levitate, or fly up the tree to have the conversation. \textbf{Kabou} (who can speak Common) outlines the events described in the Adventure Background, admitting that things may have gone too far. The \textbf{aziza} hate the murderous Korug—the \textbf{rainforest ogre}—but are not prepared to go against the will of \textbf{rainforest king}. She requests that the PCs persuade the \textbf{dryad} to call off the ogre.

\subsection{4. Breezebough}
This 100-foot-tall kapok tree (called Breezebough by the \textbf{aziza}) has two treehouses built in its branches, 50 feet above the ground.

\subsubsection{Creatures}
Four \textbf{aziza}—two with wings (females) and two without (males)—make their home in this tree. Two \textbf{gliding frog} cling to the trunk nearby. The \textbf{aziza} do not attack unless provoked but reinforce the rest of the herd in 3 if they are having trouble driving the PCs away.

\subsection{5. Fallen Trees}
\begin{DndReadAloud}{}
This open area is littered with fallen trees; their rotting trunks are covered with lichen and vines. Here and there, colorful orchids have sprouted from the remains of these fallen giants.

\end{DndReadAloud}

\subsubsection{Creatures}
A \textbf{giant walking stick} hides among the fallen trees. It leaps up onto its six long, spindly legs to attack when the PCs approach.

\subsection{6. Overgrown Ruins}
\begin{DndReadAloud}{}
Ancient stone ruins covered in vines, moss, and other foliage lie here in a roughly semicircular pattern, half-buried in the undergrowth.

\end{DndReadAloud}

The ruins stand from 3–4 feet tall. Cutting back the foliage enables the PCs to uncover a series of worn pictographs depicting tiny fey creatures (\textbf{aziza}) teaching human jungle dwellers how to make fire and which plants to harvest for food and medicine.

\subsubsection{Creatures}
A \textbf{tripwire patch} lurks among the tangle of vines and other plants inside the semicircle of stone blocks, attacking any PC that comes within range of its Tripwire Vine (50 feet).

\subsubsection{Treasure}
A successful DC 13 Wisdom (Perception) check spots a pointed, oval \textit{+1 shield} covered in zebra hide which lies among the foliage around the southernmost stone block.

\subsection{7. Log Bridge}
\begin{DndReadAloud}{}
A once-mighty fallen tree covered in moss and lichen forms a makeshift bridge across the stream.

\end{DndReadAloud}

The log bridge is 50 feet long and 10 feet wide and is treacherous to cross. A creature attempting to use the slippery bridge must succeed on a DC 15 Dexterity saving throw or fall into the water below.

\subsubsection{Creatures}
One \textbf{swarm of quippers} always stays close to the bridge; they have learned this is a good place to find food. If a creature falls in, a second swarm arrives in 1d3 rounds.

\subsection{8. Deadly Clearing}
\begin{DndReadAloud}{}
The forest floor here is strewn with dead leaves in a riot of autumnal colors: reds, oranges, browns, and yellows.

\end{DndReadAloud}

\subsubsection{Creatures}
What looks like a pile of dead leaves here is actually a \textbf{leavesrot ooze}, which lurks in this area, ready to ambush any creature who ventures too close.

\subsection{9. Waterfall and Pool}
\begin{DndReadAloud}{}
A waterfall tumbles down forty feet from a rocky outcropping covered in lush vegetation into a deep pool below. A sparkling rainbow forms in the spray where it catches the sunlight overhead.

\end{DndReadAloud}

A successful DC 12 Wisdom (Perception) check spots the opening to a cave (12) behind the rushing water.

\subsubsection{Pool}
The pool is 15 feet deep and is home to two \textbf{swarm of quippers}. The hungry fish attack any creature that enters the water.

\subsubsection{Treasure}
A woodcutter's axe and a set of \textit{carpenter's tools} lie at the bottom of the pool, stolen from the loggers and dropped there by the \textbf{aziza}.

\subsection{10. Rocky Outcropping}
\begin{DndReadAloud}{}
Thick, creeping vines decorated with dozens of brightly-colored flowers trail down the sides of this large rock formation. Ripe orange fruit grows on some of the vines hanging down the western side of the outcropping.

\end{DndReadAloud}

The vines are strong enough to support the weight of PCs climbing to the top of the rock formation. The climb is 40 feet and requires a successful DC 10 Strength (Athletics) check, made with advantage.

The top of the outcropping is covered in foliage, which is considered difficult terrain. A 5-foot-diameter hole in the ground on the western side of the stream looks down into the Hidden Cave (12).

\subsubsection{Creatures}
The orange fruit serves to lure prey toward the \textbf{monkey's bane vine} that lies in wait amid the foliage. The vine lashes out with its tendrils at any creature attempting to climb within 20 feet of it. Several monkey skeletons tangled in the foliage—previous victims of the plant monster—offer a warning to climbers.

\subsection{11. Aminata's Tree}
\begin{DndReadAloud}{}
A huge mahogany tree, over a hundred feet tall, with great buttress-like roots, towers over the outcropping. An eaglesized finch with mottled black, gray, and white plumage perches in its branches.

\end{DndReadAloud}

\subsubsection{Creatures}
This mighty tree is home to the \textbf{dryad} Aminata who watches over Emerald Cascade (see Adventure Background). The bird is a \textbf{moon weaver} named Tayo and an ally of Aminata's. Tayo calls out to the \textbf{dryad} as soon as he spots the PCs.

\subsubsection{Development}
When the PCs approach, Aminata appears from her tree and ask them why they have come to Emerald Cascade. She is prepared to listen to what they have to say, as long as they are sympathetic to the destruction of the forest by the villagers and Delande's tragic death.

The PCs can persuade Aminata to meet with Mayifa, the leader of Bukofa's elders, with a successful DC 15 Charisma (Deception or Persuasion) check; consider granting advantage on the check (or waiving the need for a roll) if the players roleplay the conversation well.

If the PCs win Aminata over, she tells them she doesn't think Korug, the \textbf{rainforest ogre}, will listen to reason. The adventurers must slay the monstrous creature if they wish to put a stop to its murderous ways.

If the PCs fail to convince Aminata to negotiate, or upset her, the \textbf{dryad} orders them to leave her grove. If they don't acquiesce, she attacks, calling upon Tayo and two \textbf{awakened tree} to defend her.

\subsection{12. Hidden Cave}
\begin{DndReadAloud}{}
Dank and foul-smelling, this cave is lit only by a beam of sunlight coming through the five-foot-diameter hole in the ceiling forty feet above. A fire pit smolders in the center of the chamber and a heap of thick animal skins is piled up near the wall opposite the entrance. Numerous human bones lie strewn across the floor.

\end{DndReadAloud}

\subsubsection{Creatures}
This cave is home to Korug, the fearsome \textbf{rainforest ogre}, and his \textbf{giant boar} companion. The \textbf{rainforest ogre} wears gruesome jewelry made from the bones of his victims. Both creatures attack when they catch sight of the PCs.

\subsubsection{Treasure}
Inside an unlocked wooden chest in the eastern part of the cavern is: 227 sp, 90 gp, an ebony headdress (worth 25 gp), a leopard skin cloak (worth 50 gp), an ivory bangle carved with a dozen elephants (worth 100 gp), a gold toe ring (worth 75 gp), and a silver amulet studded with blue topazes (worth 150 gp).

These items all once belonged to the slain villagers, but, if the PCs kill the \textbf{rainforest ogre}, the mourning families all insist that the PCs keep it all in return for slaying the monster.

\section{Concluding the Adventure}
If the PCs convince Aminata to meet with Mayifa, they return to Bukofa to fetch the elder and escort her to Emerald Cascade. Provided the \textbf{rainforest ogre} has been dealt with, Aminata and Mayifa discuss matters long into the night before eventually coming to an agreement on how the fey and the villagers can live alongside each other in peace. \textbf{Kabou} and the other \textbf{aziza} elect to stay at Emerald Cascade, where they live happily, continuing to play tricks on the villagers and occasionally helping them out with their magic.

If the PCs kill the \textbf{rainforest ogre}, the villagers are greatly appreciative, but if the PCs are also unable to convince Aminata to meet with Mayifa (or they kill the \textbf{dryad}), the threat to Bukofa won't end—because these actions infuriate \textbf{rainforest king}, who then continues to send more and more fey until the last few surviving villagers are driven away.

\chapter{Safe and Sound}

\begin{DndTable}{c}

\textbf{An adventure for four to five characters of 4th level}\\
\end{DndTable}

\section{Adventure Background}
\DndDropCapLine{B}{}ay's Burrow is a stopping-point for merchants and others traveling along the major trade route that runs through the city. Its economy is reliant on the constant flow of people coming and going, but that is currently threatened by a creature that has begun terrorizing travelers in the area.


A \textbf{catamount} has claimed the area north and east of Bay's Burrow as its territory—including a portion of the trade route—and has carved a lair into the rock and earth of the foothills. While a \textbf{catamount} would normally ignore the townsfolk and travelers in favor of game, it recently birthed a pair of kits, causing it to become far more aggressive and territorial.

This new danger has disrupted the travel and trade through Bay's Burrow, as the \textbf{catamount} has taken to attacking caravans, killing horses, and generally making it incredibly hazardous to move about the region.

\section{Adventure Hook}
Sarra Voz hires the PCs to travel to Bay's Burrow, the last known location of a shipment she has been eagerly awaiting. The shipment never arrived, and her caravanner, Perinigan, has not been heard from since. She offers the PCs 200 gp each if they return with news of her shipment's whereabouts—or 400 gp each if they return with the shipment and/or Perinigan. She describes Perinigan as a half-elf with long red hair and a long scar along the left side of his face.

Using the above hook, you can choose one of the following identities for Sarra Voz, which best suits your PCs' motivations:


\begin{itemize}
\item \textbf{\textit{Exotic Clothier.}} Sarra Voz is a clothier of fine and exotic fabrics. The shipment she's expecting is \textbf{Ettercap} silks, which is highly fashionable this year.
\item \textbf{\textit{An Herbalist in Need.}} Sarra Voz is an herbalist awaiting a shipment of rare plants and herbs that she needs to treat an illness in her town. The disease has resisted magical curing, but she's figured out a medicinal solution and desperately needs that shipment. The town council provides the funds to hire the PCs.
\item \textbf{\textit{Noble Love.}} Sarra Voz is a minor noble who has fallen in love with a caravanner she met while traveling. In this case, the shipment is all of Perinigan's worldly goods, which he's transporting from his original home to Sarra's city so that they can be together. For this scenario, Sarra only offers the 400 gp option and Perinigan's return is not optional.
\end{itemize}

\section{Bay's Burrow}
The walled city of Bay's Burrow is a bustling and busy place nestled in the foothills of an arid mountain range, full of shops, restaurants, and inns.

\subsection{On the Outskirts}
While heading south toward Bay's Burrow, a mile north of the city the PCs come upon an overturned cart. Two horses, still attached to the cart, are dead, with long, bloody gashes in their sides. But most striking is an enormous rift in hard-packed dirt and gravel of the road.

The yawning fissure stretches 30 feet long by 15 feet wide by 10 feet deep, starting somewhere off-road and coming to an end in the middle of the thoroughfare.

Any PC who wishes to investigate the area may make a Wisdom (Survival) check; consult the "The Cart and the Fissure" table to determine what they learn:

\begin{DndTable}[header=The Cart and the Fissure]{XX}

DC & Information Learned\\
8 & The cart's tracks indicate that the cart was heading north, out of town. Bootprints lead away from the cart, south toward town. The prints are small, as if made by a child or perhaps one of the smaller races.\\
10 & The horses' wounds were made by a creature with vicious claws and teeth, like a mountain lion—but larger.\\
12 & The wounds were inflicted within the last day or so.\\
14 & The horses were running in the direction of the fissure, but turned sharply to avoid falling into it and tipped the cart (and themselves) over in their panic and haste. The fissure is clearly impossible to miss, so, strange as it may be, it seems as if it must have appeared suddenly.\\
20 & In addition to the bootprints and other tracks, there are also very faint signs of paw prints that seem to belong to some sort of large cat creature, and judging from the evidence on display it seems likely that this large cat was the culprit. (Having found the \textbf{catamount} tracks already, the PCs can skip ahead directly to the "Tracking the Creature" section.)\\
\end{DndTable}

\subsubsection{The Cart's Contents}
The cart contains five sacks of flour, three broken casks of wine, and three crates filled with common clothes and linen fabric.

\subsection{Fearne's Table}
Most of the inns in the city are currently full due to the backlog of travelers, but Fearne's Table still has some rooms available. If the PCs begin asking around about Perinigan, a young woman says she noticed a man fitting his description drowning his sorrows there.

A tavern and inn of decent size and repute, Fearne's Table is filled nearly to the brim with patrons. When the PCs enter, the bartender, Gideon, apologizes for the wait and explains that, with the roads being unsafe, quite a lot of visitors are stranded in Bay's Burrow while they adjust their travel plans.

\subsubsection{Perinigan}
Perinigan, the caravanner, is one of the many guests at Fearne's Table. He is sitting at a table alone, nursing an ale and feeling sorry for himself. He's easily recognizable by his red hair and large scar that pulls the corner of his mouth down in a perpetual half-scowl.

If the PCs tell Perinigan they've been sent by Sarra, he's happy to have them join him; otherwise, offering to buy him a drink or a successful DC 13 Charisma (Persuasion) check makes him willing to talk. He has been in Bay's Burrow for several days now, waiting for the roads to become safe to travel, and he's losing his metaphorical shirt hiring round-the-clock security to guard the caravan while he waits. He's heard that a giant cat has been attacking caravans and travelers on the road, but doesn't know much more than that. He suggests, however, that the PCs should talk to Terran Teakettle—and gestures to a halfling sitting at the bar; his cart was attacked yesterday.

If the PCs propose that they escort him along his route, Perinigan refuses; he says he's not going anywhere until that cat has been dealt with, and he also refuses to leave his goods behind. Likewise, he won't allow the PCs to take the cart to its destination without him.

\subsubsection{Terran Teakettle}
Terran Teakettle is a halfling merchant with a mop of dark hair and piercing green eyes. His cart was attacked and his horses were killed yesterday when he attempted to leave Bay's Burrow despite the risk (see the "On the Outskirts" section). He got a good look at the creature, seeing as it attacked his horses right in front of him. It was an enormous cat, like a cougar, but far bigger than any cougar he's ever seen, though it had similar coloring and appearance.

Even more strangely, the earth around the creature seemed to react to its presence, shaking as if from an earthquake and splitting open in a huge fissure right in front of his panicked horses. He thanks every god great and small that he was able to walk away with but a few scrapes and bruises—his life isn't worth the goods in the cart, though he mourns the loss of his horses, Cinder and Blaze.

\section{Tracking the Creature}
The PCs can return to the site of Terran's overturned cart and attempt to find any tracks the creature might have left. Any PC who wishes to look for tracks may make a Wisdom (Survival) check; consult the "Cat Tracks" table to determine how successful they are at following the trail.

The DCs required to follow the tracks has been reduced to reflect the fact that the PCs now know what they're looking for.

\begin{DndTable}[header=Cat Tracks]{XX}

DC & Tracking Success\\
12 & The tracker finds, then loses, the trail. It takes 1d6 hours to search for the trail again, during which all of the PCs suffer one level of exhaustion. After the 1d6 hours pass, the tracker must make another Wisdom (Survival) check, this time with disadvantage. If the result is 16 or higher, they find the trail again; if the result is lower, the trail is cold. Proceed to the "Failure to Track" section.\\
14 & The tracker finds, then loses, the trail. The PC must make another Wisdom (Survival) check, this time with disadvantage. If the result is 16 or higher, they find the trail again and can follow it all the way to the lair; if the result is lower, the trail is cold. Proceed to the "Failure to Track" section.\\
20+ & The tracker expertly guides the PCs to the creature's lair. Along the way, they also find the bootprints of several small humanoid creatures. (The goblins in 6.)\\
\end{DndTable}

\subsection{Failure to Track}
If the PCs lose the trail and cannot get back on track (or find it in the first place), bring the trail to them instead.

\subsubsection{Creatures}
Once the PCs' efforts to follow the trail stall out, a pack of five \textbf{wolf} attacks. After three rounds—or after one wolf has been killed (whichever comes first)—the \textbf{catamount} joins the fray, attracted by the sounds of predators in its territory.

When the \textbf{catamount} approaches, any PCs with a passive Perception of 14 or higher notice the broad-shouldered beast before it uses its Fissure ability to open a rift in the ground beneath their feet; PCs who don't notice the \textbf{catamount} must make the Dexterity saving throw against falling into the fissure with disadvantage.

\subsubsection{Tactics}
The \textbf{catamount} initially focuses on attacking the \textbf{wolf}, thinking them the bigger threat. Once the \textbf{wolf} are killed, the \textbf{catamount} turns on the PCs. If only one PC remains above the fissure, the \textbf{catamount} attacks them; if more than one PC remains aboveground, or if the PCs manage to escape the fissure before the \textbf{catamount} has downed its current target, the \textbf{catamount} flees back to its den, preferring to fight the PCs on its turf; it uses its Fissure and Stone Wall abilities to allow itself to escape.

\subsubsection{Fresh Tracks}
After the encounter, the \textbf{catamount} tracks are easily followed back to its lair, requiring no ability check.

\subsection{Adventurer in Caravanner's Clothing}
The PCs may think to set a trap for the \textbf{catamount}, using a fake caravan or the like to lure the creature into coming out of hiding to attack. If they try this gambit, once they get a half-mile or so past the location of the attack on Terran's cart, run the same encounter as described above in "Failure to Track," but without the \textbf{wolf}. Or, if the PCs go to extensive lengths to ensure they attract the \textbf{catamount} attention, such as feigning injury, producing the smells of a fresh kill, or the sounds of an injured prey animal, keep the \textbf{wolf} in the encounter.

\begin{DndSidebar}[float=!b]{Save the Boss Fight}
If the PCs manage to reduce the \textbf{catamount} to 0 hit points when they encounter it outside its lair, have the \textbf{catamount} stabilize at 1 hit point instead of dying and allow it to flee back to its den to lick its wounds; if pursued, it aggressively uses its Fissure and Stone Wall abilities to allow itself to escape. If the PCs pursue it immediately, they find it in 10, at half-health. If the PCs took a short or long rest during their pursuit, however, the \textbf{catamount} does as well—and thus when the PCs find it in its den, it is at full health.



\end{DndSidebar}

\section{The Catamount's Den}
The den is a series of tunnels carved into the southern face of a hillside an hour's walk from the main road. It consists of a series of passages and cavern-like chambers, crafted by the \textbf{catamount} itself. Most of the areas were only used by the \textbf{catamount} temporarily as it expanded its lair and now lie empty and abandoned.

The den has several common features:


\begin{itemize}
\item \textbf{\textit{Pitch Black.}} The interior has no lighting of any kind; as such, it is considered complete darkness.
\item \textbf{\textit{Strangely Smooth.}} The lair was created by the \textbf{catamount} innate ability to carve rock and earth. The floor, ceiling, and walls are all smooth and nearly tubular and show no sign of being hewn by hand.
\item \textbf{\textit{Tall Ceilings.}} The ceilings are all 15 feet tall.
\end{itemize}

An unlabelled version of the map can be found here.

\subsection{1. Main Tunnel}
\begin{DndReadAloud}{}
The scent of musk assaults you as you enter the mouth of the cave. Ahead, the tunnel continues on into darkness. It is completely, eerily silent—and utterly dark and black.

\end{DndReadAloud}

The main tunnel is approximately 15 feet wide at its narrowest. It winds into the darkness of the mountainside.

\subsection{2. Original Chamber}
\begin{DndReadAloud}{}
This large chamber smells of stale musk and a slight hint of rot. Some old bones are scattered about, long-ago stripped of their flesh, and there are tufts of gold and brown fur that have collected at the edges of the room.

\end{DndReadAloud}

This oblong chamber, which was the original cave carved by the \textbf{catamount}, measures roughly 30 feet by 55 feet.

\subsection{3. Pit Trap}
\begin{DndReadAloud}{}
This portion of the hallway widens slightly, like a bulge in a weakened vein. You see more detritus of bones and fur, and, ahead, you can see that the cavern continues north then bends westward.

\end{DndReadAloud}

The floor in the center of the room is thin and easily broken. Creatures weighing more than 100 lbs. (including their gear) must make a DC 14 Dexterity saving throw when they enter the space; on a failure, the false floor breaks beneath them in a 5-foot square, and they fall 30 feet into the pit below, taking 10 (3d6) bludgeoning damage from the fall and landing \textbf{prone}.

To climb out of the pit, a creature must succeed on a DC 18 Strength (Athletics) check; using a rope and \textit{grappling hook} or other climbing gear grants advantage on the roll. If more than one PC avoids falling into the pit and can lower a rope, together, with 10 minutes of work, they can help any fallen compatriots escape the trap.

The pit spans a 20-foot square in the center of the room. The \textbf{catamount} can pass through this area without triggering the trap.

\subsection{4. Main Tunnel Continuation}
When the PCs pass the pit trap and enter this portion of the tunnel, they can make a DC 15 Wisdom (Perception) check. On a success, they hear the sound of hushed bickering from up ahead, and any PC that speaks goblin recognizes it as that language, though they can't quite make out what's being said beyond the words "gonna wait"; on a result of 18 or higher, they are able to determine that the bickering is coming from the goblin war band (see 6) arguing about how or when to proceed against the \textbf{catamount}, since their scout Brux has not returned.

\subsection{5. Another Old Den}
The room holds nothing of interest and is almost identical to 2.

\subsection{6. False Den}
This room is also very similar to 2, except for its \textbf{goblin} occupants.

\subsubsection{Creatures}
A war band of eight \textbf{goblin} huddle together in this room, arguing about whether or not they should attack the \textbf{catamount} now or wait. They are not immediately hostile, but fight if attacked.

When the PCs enter this room, the \textbf{goblin} draw their weapons and back against the southern wall as far from the PCs as they can get. A successful DC 14 Charisma (Deception, Intimidation, or Persuasion) check convinces the \textbf{goblin} to put away their weapons and talk. The leader of the band identifies herself as Herx.

\subsubsection{Development}
The \textbf{goblin} are part of a hunting party there to kill the \textbf{catamount} and claim the lair for themselves. One of their member, Brux, went ahead to scout the den (see 10), and they are waiting for him to return before making their final push. A successful DC 18 Charisma (Deception, Intimidation, or Persuasion) check is necessary to convince the \textbf{goblin} to fight alongside the PCs; otherwise they are content—now that the brave adventurers are here—to wait for them to kill the \textbf{catamount} and then claim the lair for themselves.

\subsubsection{Treasure}
Along with the gear listed in their stat block, the \textbf{goblin} each carry 15 cp.

\subsection{7. Waste Room}
\begin{DndReadAloud}{}
While the ground has been mostly firm, hard-packed earth elsewhere in the den, in this chamber the dirt is loose and clearly has been disturbed with some frequency. A slight acrid scent fills the air, sharper than the scent of musk you've detected thus far.

\end{DndReadAloud}

This chamber is simply the \textbf{catamount} "litter box," though since it can literally move earth at will, it buries its urine and feces thoroughly.

\subsection{8. Yet Another Old Den}
Another old den, mostly identical to 2, but the southwestern corner of the chamber shows sign of collapse, as if the \textbf{catamount} had been trying to tunnel further westward, but the earth was too unstable.

This empty cavern once served as the \textbf{catamount} primary nesting area, but has since been abandoned.

\subsection{9. Funnel}
The tunnel widens to a bulb, with the only exit leading to the west to 10.

\subsection{10. Nest}
\begin{DndReadAloud}{}
The cavern widens ahead, opening into an enormous chamber. At the far western end sits a large stone, bowlshaped nest. The bowl is formed of the same stone as the floor, in one continuous piece. The bowl looks to be lined with fur, and two newly-born kits rest within, so young that their eyes are still shut. Lying in front of the stone nest and feasting on the carcass of a goblin is an enormous catlike creature with gold and brown fur.

\end{DndReadAloud}

The lip of the stone nest rises 5 feet off the ground, and the interior depression is 2 feet deep. The kits within are newborns and defenseless; at this young age, they're Small creatures, just over 2 feet in length.

The goblin carcass is that of Brux, the scout mentioned in 6.

\subsubsection{Creatures}
As soon as the PCs enter this room, the \textbf{catamount} poises itself to attack.

\subsubsection{Tactics}
The \textbf{catamount} primary concern is the safety of its kits. It attacks whatever creature moves closest to the nest and uses its Control Earth ability to inhibit any intruders from reaching the kits.

If the \textbf{catamount} health is reduced by half (55 hit points or lower if it started the combat with full health, 25 hit points or lower if it started the combat with half health), it uses its Fissure ability to create a rift between the intruders and the kits, then leaps across it and attempts to lure the PCs out of the den and back out into the wilderness.

\subsubsection{Diplomacy}
The \textbf{catamount} is a beast, so PCs may attempt to cast spells such as \textit{speak with animals} to attempt to resolve the situation peacefully. It is intelligent enough to be able to reason, but its instincts are currently driving it to protect its young at all costs. Because of this, any attempt to reason with the \textbf{catamount} requires a successful DC 18 Charisma (Deception or Persuasion) or Wisdom (Animal Handling) check, made with disadvantage—unless the PCs demonstrate any magic or abilities that show that they too are capable of manipulating earth. If there is any attempt to use Charisma (Intimidation), the \textbf{catamount} attacks. If the PCs don't offer a reasonable solution as part of their entreaty, the attempt at diplomacy automatically fails, and the \textbf{catamount} attacks.

\section{Concluding the Adventure}
The adventure concludes when the \textbf{catamount} is killed or is otherwise no longer a threat to the area.

If the situation is resolved peacefully and the PCs convince the \textbf{catamount} to relocate and/or provide it with ways to protect its kits and survive without harming townsfolk, the \textbf{catamount} considers the PCs friends for life. If the PCs ever encounter the creature again, it remembers them and greets them (over)enthusiastically like a lost pet being reunited with its person.

If any of the \textbf{goblin} remain alive, they move into the den if the \textbf{catamount} is either killed or removed, claiming the tunnels as their new home. After a week, a clan of 30 \textbf{goblin} takes up residence there. The \textbf{goblin} are wary of being too much of a threat to Bay's Burrow, as they don't want the wrath of the city—or adventurers—brought down upon them, but they attack lone travelers when the opportunity arises, and occasionally raid caravans on the trade route.

The \textbf{catamount} kits are too young to survive on their own; if left behind in the den, they starve to death in a manner of days—if the \textbf{goblin} don't kill them first. If taken from the den, the kits can be hand-reared and relocated elsewhere once they are old enough to hunt on their own. Particularly bold adventurers may attempt to keep the \textbf{catamount} as pets or animal companions. If kept as animal companions, they are completely helpless for a week, after which they use the statistics for a \textbf{badger}; at 3 months, they use the statistics for a \textbf{giant badger}; at 6 months, they use the statistics for a \textbf{brown bear}. Once the kits are roughly 1 year of age, they mature into adult \textbf{catamount} (and then use the full \textbf{catamount} stat block) and become too large to keep—and their instincts drive them to hunt, establish territory, and build a den. Alternatively, there is a market for two such kits, if the PCs know where to look.

When the \textbf{catamount} has been dealt with, Perinigan is thankful to the PCs for making the roads safe for travel again. If they offer to accompany him back to Sarra, he happily accepts the company and starts off the next morning. Terran also accompanies Perinigan—the two have hit it off over the past few days of sharing drinks and commiserating and have been discussing the possibility of merging their two businesses.

\chapter{Lord Gorgo's Keep}

\begin{DndTable}{c}

\textbf{An adventure for four to five characters of 5th level}\\
\end{DndTable}

\section{Adventure Background}
\DndDropCapLine{S}{}piteful Keep sits atop a small peak near a trade rode, amid a small forest. It is the home of an ogre artisan of the arcane who self-styles himself as "Lord" Gorgo. Only the bravest or most desperate dare to seek audience with the ogre


\section{Adventure Hook}
The following can used to start the PCs on this adventure:


\begin{itemize}
\item \textbf{\textit{To Lift a Curse.}} A local knight named Veladra is cursed. She purchased a sword from Gorgo, years past, when she was an adventurer. Gorgo recently sought her aid and was rebuffed. The ogre activated the curse, to make the knight give in to his demands. Instead, she hires the PCs to take care of the problem, by killing Gorgo or bringing him the sword and getting him to remove the curse.
\item \textbf{\textit{Vanishing Goods.}} Yintalyn, the head of a town council at a nearby settlement, is concerned about a series of mysterious recent thefts. Rare and precious materials simply vanish, disappearing from locked strongboxes in merchant wagons. There are no witnesses or clues; things simply vanish. The authorities eventually realize that all of the goods stolen are useful in magical crafting. Gorgo immediately comes under suspicion. The PCs are hired to investigate Gorgo's keep and put a stop to these thefts.
\item \textbf{\textit{In the Market for Magic.}} A wizard by the name of Gelarien is in need of a certain magic item and knows that Gorgo is the only artisan skillful enough to craft such a thing. She hires the PCs to make the journey to Gorgo's keep and convince the ogre to craft the item. In truth, Gelarien is a fake name, and the wizard was once wronged by Gorgo, and her true wish is that the PCs go there and will kill the ogre, who will not view an uninvited visit kindly.
\end{itemize}

\section{Spiteful Keep}
This keep is the home of Lord Gorgo, a \textbf{cunning artisan ogre}. The fortress is quite old, but well-maintained and well-fortified, with a curtain wall bolstered by five round towers and a gatehouse.

\subsubsection{Approaching the Keep}
As the PCs approach the keep within sight of any of its towers, they hear the sound of a distinctive whistle from above, clearly a signal of some kind alerting the keep that visitors are approaching.

\subsubsection{Keep Denizens}
The keep's denizens are utterly under the thrall of Lord Gorgo, and, indeed, they truly believe he's a great man and a master artisan worthy of following. The keep denizens include: 16 \textbf{hobgoblin}, a \textbf{bugbear}, a \textbf{daeodon}, 4 \textbf{goblin}, \textbf{kobold ettin}, a \textbf{dokkaebi} , and Lord Gorgo himself. All of the denizens (except the \textbf{goblin}) are loyal to Gorgo and won't act contrary to his orders, even to save the life of one of their fellows.

\subsubsection{Gorgo's Jape}
One day, years ago, Gorgo demanded that his followers always refer to him as "Lord Gorgo, my mighty master"; he'd intended it as a jape that wouldn't last more than a day or two, but it pleased him so much he never rescinded the order. As a result, any time any of the keep denizens refer to Gorgo, they specifically use the words "\textit{Lord Gorgo, my mighty master}."

An unlabelled version of the map can be found here.

\subsection{1. Gatehouse}
\begin{DndReadAloud}{}
As you approach the keep's iron-banded, oaken gates of the gatehouse, a slit slides open behind a small metal grate in the door and a pair of yellow eyes peer out from behind it. "What's your business here?" a voice demands.

\end{DndReadAloud}

\subsubsection{Polite Approach}
If the PCs are polite and present a reasonable story for why they're there, the guard tells them to wait there, and one of them will go see if Lord Gorgo is willing to meet with them. After a long time—long enough for the PCs to start to get impatient—the guard comes back and says, "Lord Gorgo, my mighty master, has agreed to an audience. Follow me." The guard leads the PCs to 9.

\subsubsection{Aggressive Approach}
If the PCs are overtly aggressive or insulting, the guards refuse them entry; once insulted, the guards can be convinced to hear the PCs out with a successful DC 17 Charisma (Deception, Intimidation, or Persuasion) check; if Charisma (Intimidation) is attempted, the check is made with disadvantage. On a failure, the guards demand they leave.

If the PCs refuse to leave or try to attack the guards, they sound the alarm, which results in a volley of arrows from the guards in 15 and in the two southernmost towers (3).

\subsubsection{Inside the Gatehouse}
The keep's outer gate opens into a small chamber with arrowslits on the eastern and western walls and murder holes in the ceiling above.

\subsubsection{Creatures}
Four \textbf{hobgoblin} are stationed here at any given time—two in the main gatehouse entryway and one each in the chambers to the east and west, manning the arrowslits; the latter are armed with pikes and heavy crossbows, which they use to attack hostile creatures through the arrowslits.

If combat ensues here, the two hobgoblins in 15 provide archery support from above, raining arrows down through the murder holes.

\subsubsection{Arrowslits}
The arrowslits provide the guards within with three-quarters cover. The chamber behind each arrowslit contains a pair of chairs and an iron ladder leading to 15.

\subsubsection{Inner Gate}
On the northern wall of the entryway are the inner gates that lead to the keep proper. Like the outer gates, these are iron-banded and made of solid oak. They have AC 19, 27 hit points, and a damage threshold of 10. The inner gates are barred on the courtyard side, requiring someone inside the keep—normally one of the arrowslit guards—to open them.

\begin{DndSidebar}[float=!b]{Sound the Alarm!}
If the keep's denizens are alerted to their compatriots being attacked—either by the alarm being sounded or by the clamor of battle—they converge on the threat and all fight to the death (except the \textbf{goblin}, who surrender or flee at the first opportunity).



If this occurs, it is likely to go very poorly for the PCs, as they're considerably outnumbered and would be fighting foes with superior fortified positioning.



\end{DndSidebar}

\subsection{2. Courtyard}
\begin{DndReadAloud}{}
This large, open area is mostly packed earth with a few patches of scrubby grass. Outbuildings sit in the shelter of the keep's walls, while a covered well stands in the center of the courtyard. A pile of refuse has been dumped beside the north wall between the main building and the tower.

\end{DndReadAloud}

The courtyard is often unoccupied, except for Gorgo's pet \textbf{daeodon}, Gouger. The beast can usually be found in a wallow it has dug between the smithy (4) and the westernmost tower, or it is sometimes found feeding in the refuse pile on the north side of the courtyard.

Occasionally two or more \textbf{hobgoblin} from the barracks (4) can be found here training.

\subsubsection{Creatures}
Lord Gorgo's pet \textbf{daeodon} Gouger has the run of the courtyard most of the time. The \textbf{daeodon} is still half wild, and only Gorgo can command the beast. He readily attacks any creature he does not recognize, unless commanded otherwise by Gorgo—or unless that creature is in the company of one or more \textbf{hobgoblin}.

\subsection{3. Tower, Ground Floor}
\begin{DndReadAloud}{}
This small, round room is empty except for a narrow staircase that winds its way up the side of the tower to the floor above.

\end{DndReadAloud}

All five of the keep's towers are identical. The stairs wind up to the tower's upper level (16).

\subsection{4. Smithy}
\begin{DndReadAloud}{}
A roofed, outdoor area on the north side of this stone structure contains a forge and anvil.

\end{DndReadAloud}

This building houses the keep's smithy, consisting of the forge, a small storage room, and a room containing a bed and small table, where the smith sleeps and takes his meals.

\subsubsection{Creatures}
Khorr, the \textbf{bugbear} blacksmith, is typically found here at work in front of the forge.

\subsubsection{Treasure}
A set of smith's tools, 50 5-pound ingots of iron (worth 25 gp total), a sack containing 3 1-pound silver ingots (worth 15 gp total), and a pouch containing 20 sp and 13 gp.

\subsection{5. Stable}
\begin{DndReadAloud}{}
This wooden building has a large set of double doors facing the center of the courtyard. Above the doors, you can see a block and tackle in front of an opening into a loft above.

\end{DndReadAloud}

\subsubsection{Creatures}
This is the keep's stables. Inside are 8 stalls, half of them occupied by \textbf{riding horse}. The other stalls are unoccupied.

\subsubsection{Treasure}
The stable holds saddles, saddlebags, and bits and bridles for four horses.

\subsection{6. Barracks}
\begin{DndReadAloud}{}
This one-story stone building holds four bunk beds, each with footlockers beneath them. A long table and four chairs take up one side of the room. Two cast iron wood stoves sit in the northeast and southeast corners.

\end{DndReadAloud}

\subsubsection{Creatures}
There are up to 8 \textbf{hobgoblin} here. If alerted by an alarm or the sound of fighting, they rush to reinforce their allies. There is a 50\% chance that any \textbf{hobgoblin} located here is asleep. In that case, their AC is only 13.

\subsubsection{Treasure}
In the footlockers is a total of 700 sp and 400 gp.

\subsection{7. Kitchen}
\begin{DndReadAloud}{}
Iron pots bubble away in the large fireplace in the southwest corner of this room. A large, metal sink, work table and tall stools, and a set of shelves take up most of the space. An iron rack hung with various cookware, utensils, and bundles of drying herbs is suspended over the table.

\end{DndReadAloud}

\subsubsection{Creatures}
Four \textbf{goblin} work here, preparing food for Gorgo and his followers. They won't attack the PCs, though will defend themselves; at the first opportunity, they attempt to flee through 8.

\subsection{8. Pantry/Storage}
\begin{DndReadAloud}{}
One side of this room is stacked with crates, barrels, and sacks, while the other is occupied by coils of rope, stacks of lumber, and more assorted crates and barrels.

\end{DndReadAloud}

This room serves as a pantry for the kitchen, containing general staples and various beers and wines.

\subsection{9. Main Hall}
\begin{DndReadAloud}{}
A large fireplace dominates the wall opposite the door. A long table stretches nearly the length of the room, surrounded by tall-backed chairs. At one end of the table sits a large, throne-like chair of wrought iron, set with leather cushions. A massive chandelier festooned with candles hangs above the table.

\end{DndReadAloud}

This room serves as a dining hall and Gorgo's audience chamber. When Gorgo is receiving visitors, \textbf{goblin} minions bring in drinks and appetizers, then return to the kitchen (7).

\subsubsection{Secret Door}
A successful DC 17 Wisdom (Perception) check notices a secret door at the southern end of the eastern wall that opens into the stairwell.

\subsubsection{Audience with Gorgo}
If the PCs attempt to reason with Gorgo, he quietly listens, politely nodding. The conversation plays out differently depending on which adventure hook you used, but in all cases almost everything he says is a lie; the PCs can discern this with a successful DC 17 Wisdom (Insight) check.

\subsubsection{To Lift a Curse}
Gorgo explains that he made some youthful mistakes, this unpleasant situation with Veladra being one of them. He says he will gladly accede to this most reasonable request. He calls out to a guard in the hall, "Kludek, have Shovar bring me my toolbox."

\subsubsection{Vanishing Goods}
Gorgo seems surprised and distressed at this news and says that he'll question his majordomo about this immediately, as it is clearly unacceptable. He calls out to a guard in the hall, "Kludek, have Shovar come down here."

\subsubsection{A Wanting Wand}
Gorgo is dismayed at this news. He pulls on a glove and asks to see the wand. After examining it a moment, he says "Ah yes, this is my apprentice's work. I do apologize. He still has a lot to learn, it seems." He calls out to a guard in the hall, "Kludek, have Shovar come down here."

\subsubsection{Ambush Code}
"Kludek, have Shovar . . ." is code to prepare an attack. Gorgo makes pleasant small talk while waiting. After 5 minutes pass, the \textbf{kobold ettin} Squib-Jib (see 10) comes barreling through the secret door at the same time as two \textbf{hobgoblin} burst in through the regular door—and combat ensues; if it does, see the "Encountering Gorgo" sidebar.

\begin{DndSidebar}[float=!b]{Encountering Gorgo}
The upper level of the keep is the private domain of Lord Gorgo. Only his servants and select visitors ever see this level of the keep. Gorgo is equally likely to be found in Areas 11, 12, or 13. Less frequently, he is in 14, consulting with his prisoner, or escorting Kwon to his workshop for assistance with something he's working on.



\end{DndSidebar}

\subsection{10. Upper Hallway}
A set of stone steps leads to the upper level of the keep. The upper hallway stretches the length of the building. Heavy wooden shutters cover the four windows on the northern wall.

\subsubsection{Magical Trap}
If a creature steps into the hallway from the stairs without first saying the words "Lord Gorgo, my mighty master," a bolt of lightning streaks down the hall. Each creature in the path of the bolt must make a DC 14 Dexterity saving throw. A creature takes 22 (4d10) lightning damage on a failed save or half as much damage on a successful one.

Once triggered, the trap is expended and takes 10 minutes to recharge.

\subsubsection{Trap Prevention}
On the face of the top step, a successful DC 12 Wisdom (Perception) check notices that the words "Say Hail" have been scratched into the stone. Other than uttering the passphrase, the only way to bypass the trap is to sense the magical trigger using \textit{detect magic}; if discovered in this way, a successful DC 15 Intelligence (Arcana) check determines how to dissipate the trap's magical charge harmlessly, nullifying it for 10 minutes while it recharges.

\subsubsection{Creatures}
Squib-Jib, a \textbf{kobold ettin}, is around the corner, in the hall leading to 14. If the trap is triggered, they are alerted to intruders and come to investigate.

\subsubsection{Gorgo's Demeanor}
 If confronted with intruders, Gorgo's response depends on their numbers, posture, and attitude. He attempts to parley with intruders in the hopes of lulling them into a false sense of safety. If they are openly hostile, Gorgo attacks and/or defends himself.

\subsubsection{Alertness}
If the trap in 10 is triggered, Gorgo is alerted to the PCs' presence. If combat with Squib-Jib then ensues, Gorgo joins the fray after 1 round, throwing open the door of 11 and activating his \textit{gem of brightness}, even if Squib-Jib is in the line of fire.

\subsubsection{Combat Tactics}
Gorgo fights to the death, unless he sees an opportunity to escape to another location where his minions can join the battle. He uses his \textit{gem of brightness} every turn (since he can activate it with a bonus action), with no regard for saving any of its charges.

\subsubsection{Treasure}
Gorgo wears a fine \textit{breastplate}, though it only fits a Large creature. He also carries a \textit{battleaxe}, a \textit{gem of brightness}, the keys to the prison cell in 14 and the chest located in 12.

\subsection{11. Parlor}
\begin{DndReadAloud}{}
The walls of this room are hung with thick tapestries, and the floor covered in fur rugs. A large, overstuffed chair with a matching footstool sits in one corner next to a carved wooden cabinet.

\end{DndReadAloud}

\subsubsection{Treasure}
Inside the cabinet are a silver tray and goblet set (worth 50 gp), 2 crystal decanters (worth 40 gp total), a carved teakwood pipe (worth 15 gp), and a pouch of fine, apple-flavored tobacco (worth 35 gp). On the chair is a large tome about curses (worth 10 gp).

\subsection{12. Bedroom}
\begin{DndReadAloud}{}
A massive bed, piled with furs, takes up most of the room. A round table stands on one side of it, a chest on the other. A tall wardrobe stands in one corner, a freestanding, ten-foot-tall ovoid mirror in the other. There is a single, shuttered window in the center of the southern wall and a closed door on the eastern wall.

\end{DndReadAloud}

The wardrobe contains fine clothing (but very large). The table holds a ceramic washbasin and an ewer full of water.

\subsubsection{Treasure}
The chest is locked, but can be picked with a successful DC 15 Dexterity check using thieves' tools. Gorgo holds the key. It contains 10,000 sp, 1,000 gp, and a large gold bracelet worth 250 gp.

\subsection{13. Magical Workshop}
\begin{DndReadAloud}{}
Tables cluttered with tools, alchemical ingredients, and piles of notes line the walls of this room. A shelf full of old tomes stands against one wall. A round table is located in the southwest corner, a large chair next to it. On the table are a number of engraving tools, as well as an adjustable stand holding a magnifying lens.

\end{DndReadAloud}

If Gorgo is encountered in this room, there is a 50\% chance he has Kwon (14) with him.

\subsubsection{Treasure}
On the tables are sets of \textit{alchemist's supplies}, \textit{calligrapher's supplies}, \textit{jeweler's tools}, and \textit{woodcarver's tools}. There's also an assortment of gemstones, as well as gold and silver dust and wire (worth 500 gp total), and four potions of healing.

There are also several finished magic items here that function but also bear Gorgo's signature curses. Consult the "Gorgo's Cursed Items" table to see how each item is cursed; aside from the stated effects of the curses, the items work normally.

\begin{DndTable}[header=Gorgo's Cursed Items]{XX}

Magic Item & Curse\\
\textit{ring of jumping} & Curse of Cowardice: On any round in which you would otherwise be surprised, you can still act, but you must use your action to activate the ring and jump away from the nearest enemy.\\
\textit{wand of magic detection} & Curse of Falsehood: There is a 50\\% chance that one non-magical item in range is identified as magical.\\
\textit{+1 battleaxe} & Curse of Reciprocation: When you hit a creature with an attack using his magic weapon, the next attack the target makes against you has advantage.\\
\end{DndTable}

\subsection{14. Prison Cell}
\begin{DndReadAloud}{}
This spartan chamber contains only a small cot and a chamber pot.

\end{DndReadAloud}

The door to this room is locked; Gorgo holds the key.

\subsubsection{Creatures}
Inside the room is Gorgo's prisoner, a \textbf{dokkaebi} named Kwon. The ogre forces the fey to summon items used in the creation of his magic items. Kwon is desperate to be free and is grateful to anyone that rescues him. He would rather flee immediately, but stays and helps his rescuers if he must. If attacked, he turns \textbf{invisible} and attempts to flee.

\subsection{15. Upper Gatehouse}
\begin{DndReadAloud}{}
The walls of this room feature several arrowslits. The room is empty other than a pair of small casks filled with arrows and the two hobgoblins standing watch.

\end{DndReadAloud}

\subsubsection{Creatures}
Two \textbf{hobgoblin} are on watch in this room at all times. Each \textbf{hobgoblin} has a cask of 60 arrows to use in defense of the keep. If combat ensues in the courtyard, these guards use the inner arrowslits to rain arrows down on the intruders below.

\subsection{16. Upper Tower}
\begin{DndReadAloud}{}
The staircase emerges into a round room with arrowslits in the walls. A small cask of arrows sits in the center of the room.

\end{DndReadAloud}

All five of the keep's upper towers are identical. The stairs wind down to the tower's lower level (3).

\subsubsection{Creatures}
Two \textbf{hobgoblin} stand watch at the top of each tower. Each \textbf{hobgoblin} has a cask of 60 arrows to use in defense of the keep. If combat ensues in the courtyard, these guards use the inner arrowslits to rain arrows down on the intruders below.

\begin{DndSidebar}[float=!b]{Escaping the Keep}
If the PCs wish to escape the keep without going through the gatehouse, they can do so—and get beyond the keep's curtain wall—by leaping or rappelling down via: (1) the windows on the northern wall in 10 (30 feet to the ground); (2) the windows on the eastern wall in 13 (20 feet to the ground); and (3) 16, the Upper Towers (70 feet to the ground).



\end{DndSidebar}

\section{Concluding the Adventure}
The adventure comes to an end after the PCs conclude their visit to Spiteful Keep.

If Gorgo is killed, any other remaining denizens of the keep abandon the place for greener pastures.

If the PCs assault and/or raid the keep, but leave Gorgo alive, the ogre becomes an enemy who avenges himself at the most inopportune moment for the PCs. Any magic items taken from the keep may have their curses activated by him at a later date, possibly forcing another confrontation with the ogre to remove them.

\chapter{Redoubt of the Perfidious Guru}

\begin{DndTable}{c}

\textbf{An adventure for four to five characters of 5th level}\\
\end{DndTable}

\section{Adventure Background}
\DndDropCapLine{H}{}igh in the mountains, tunneled directly into one of the peaks, is the Abbey of the Serene Lily. The temple was founded by ascetic scholars from disparate nations who sought to overcome the differences between their respective peoples and live in peaceful contemplation. When acolytes leave the abbey, they often find positions as advisors and councilors to political and religious leaders in their home nations.


For the last several years, visitors to the abbey have returned to their homelands with tales of a mentor with powerful precognitive abilities. This teacher, the temple's leader, Abbot Qudos, is frequently sought after for advice on matters ranging from which suitor a young noblewoman should marry to how a trade mogul can maximize his profits on a given venture.

In truth, "Qudos" is a fiend who slayed the true Abbot Qudos and assumed his identity. He seeks to corrupt everyone he comes into contact with, and, over time, he has slain, supplanted, or sent away all of the noble monks and adepts of the abbey and replaced them with his own sycophantic followers.

\section{Adventure Hook}
The PCs have been asked to gain Abbot Qudos' counsel by a ruler whose nation suffered a cowardly attack by another nation. He doesn't wish to reignite the hostilities between the nations that has led to the loss of so many lives, but he also cannot afford to do nothing, lest it give his enemies free rein to do as they please against his people. It seems as if there is no good answer, and so Qudos' wisdom is needed.

\section{Mountain Hazards}
The abbey has been carved directly into the summit of a mountain (at 10,000 feet) located at the northwestern portion of the range. A narrow, steep—and sometimes precarious—trail winds up the mountain. A direct path, with no compilations, takes 12 hours, though attempting the trek all in one day incurs one level of exhaustion; if the trip is split into two days of travel, the PCs can reach the summit without incurring levels of exhaustion.

\subsubsection{Elevation Sickness}
Once the PCs reach the summit, whenever they complete a long rest, they must make a DC 13 Constitution saving throw; on a failed save, they gain a level of exhaustion from the effects of high altitude and fail to recover any levels of exhaustion they'd previously accrued. A PC accustomed to high altitudes automatically succeeds on this saving throw. After seven days at the summit, the PCs' bodies become inured to the rigors presented by the thinner air at the abbey and no longer have to make saving throws to avoid exhaustion.

\subsubsection{Mountain Encounters}
Twice while the PCs are scaling the mountain, roll a d6. On a result of 1 or 2, run one of the following encounters:

\subsection{Encounter 1: Alpine Creeper Ambush}
\begin{DndReadAloud}{}
Up the mountainside, you see the rocky escarpment is dotted with patches of blue and green. As you get closer, you see it is a colorful moss or lichen that has adhered itself to the landscape.

\end{DndReadAloud}

As the PCs approach this hazard, allow them to spot the stray, partly submerged bones littering the ground with a successful DC 14 Wisdom (Perception) check, or if they have passive Perception scores of 14 or higher. The bones are suspiciously humanoid and might alert the PCs to the impending threat.

\subsubsection{Creatures}
The lichen is comprised of two \textbf{alpine creeper}. They can be evaded entirely by moving out of range. If the creepers knock all of the PCs unconscious, some of the cultists from the abbey rescue them, and they later wake in 5.

\subsection{Encounter 2: Cloudhoof Assassination}
\subsubsection{Creatures}
As the PCs move across a narrow shelf, they are charged by six \textbf{cloudhoof assassin}. The goat-like creatures seek to knock the PCs from the shelf so they can rest there themselves. Creatures knocked from the ledge fall 40 feet to another ledge below, taking 4d6 bludgeoning damage and falling prone. The \textbf{cloudhoof assassin} cannot be reasoned with as long as the PCs are on the plateau, but each one retreats from the fight once it is reduced to 10 hit points or less.

\section{Abbey of the Serene Lily}
The abbey is built into the mountaintop and has no external profile other than the doors leading into it.

\subsubsection{Structure}
The floor of the abbey is comprised of 1-foot-square sandstone tiles. The walls are smooth stone and the ceiling is intermittently engraved with bas relief hummingbirds and lilies.

\subsubsection{Ceilings}
The corridors are 5 feet wide and 8 feet tall. The chambers have 15-foot ceilings.

\subsubsection{Lighting}
The abbey is dimly lit at all times through some form of magical ambient illumination.

\subsubsection{Incense}
The abbey smells of frankincense, which can be found burning in thuribles in each chamber. The pervasiveness of the odor is such that creatures making Wisdom (Perception) checks based on smell have disadvantage.

The incense smoke masks the odor of the narcotic burning with it. When a creature in the abbey makes a saving throw against being charmed or frightened—or makes a Wisdom (Insight) check to determine if a creature is telling the truth—it must subtract 1d4 from the number rolled. Creatures that are wearing some type of face covering—or who don't need to breathe—are not affected by the narcotic.

\subsubsection{Sound}
During daylight hours, the sound of someone playing the sitar can be heard. It is not uncommon to hear rhythmic chanting as well.

\subsubsection{Aspirants and Adepts}
There are twenty-seven \textbf{serene lily aspirant} (use statistics for \textbf{cultist}) and seven \textbf{serene lily adept} (use statistics for \textbf{cult fanatic}) residing at the abbey. All of the masters in the abbey are \textbf{barbed devil} or \textbf{bearded devil} in disguise.

\subsubsection{Creature Change}
Replace the \textbf{serene lily aspirant} Scimitar action and the \textbf{serene lily adept} Dagger action, with the following:

\subparagraph{Unarmed Strike}
\textsl{Melee Weapon Attack:} +4 to hit, reach 5 ft., one creature. Hit: 5 (1d6 + 2) bludgeoning damage.

\subsubsection{Color Coding}
All of the residents of the abbey wear pale blue robes. Aspirants wear red sashes, adepts wear dark blue sashes, and masters wear black sashes.

\subsubsection{Monk Distribution}
The \textbf{serene lily adept} and \textbf{serene lily aspirant} should be distributed throughout the abbey, with the following parameters depending on time of day:


\begin{itemize}
\item In the morning, the majority are outside sparring in 1.
\item In the afternoon, they are praying in 6.
\item At night, they rest in 5.
\end{itemize}

\subsubsection{Lies and Lies}
The cultists frequently lie about the purpose of the abbey and the behavior of its inhabitants. Detecting this deception requires a successful DC 14 Wisdom (Insight) check.

\subsubsection{Devils}
Each \textbf{barbed devil} or \textbf{bearded devil} encountered in the abbey is wearing a \textit{hat of disguise}, using it to appear to be a master in pale blue robes with a black sash. A creature must make a successful DC 13 Intelligence (Investigation) check to detect the disguise.

An unlabelled version of the map can be found here.

\subsection{1. Sparring Ground}
If the PCs approach in the morning, read or paraphrase the following:

\begin{DndReadAloud}{}
Several pairs of robed people are sparring against each other in the large flat area before the abbey's entrance. Four calm-faced monks walk among them, correcting their stances and offering advice. A short walkway descends slightly to the carved double doors leading into the abbey itself.

\end{DndReadAloud}

If the PCs approach at any other time of day, read or paraphrase the following:

\begin{DndReadAloud}{}
A large area in front of the entrance to the abbey has been worn flat by the repeated passage of many feet. It is curiously quiet here, with the ice-cold wind creating the only sound. A short walkway descends slightly to the carved double doors leading into the abbey itself.

\end{DndReadAloud}

This area is an expanse of dirt packed hard by years of sparring and treading upon it. A chill breeze blows dust and grit in eddies across the ground.

\subsubsection{Creatures}
In the morning, twenty \textbf{serene lily aspirant} are sparring here. Four \textbf{serene lily adept} (the monks referenced in the boxed text above) walk among them. They are all barefoot, though some have wrapped their feet in white cloth strips.

The people in the area are polite, but don't have much to say and won't answer questions in any helpful way. They're likely to utter the phrase, "Abbot Qudos knows all."

The double doors that lead into the abbey have been skillfully carved in bas-relief with a scene of lilies blowing in the wind while hummingbirds sip nectar from the flowers.

\subsection{2. Entrance}
\begin{DndReadAloud}{}
Upon entering the abbey, your eyes are immediately drawn to a fountain featuring a jade statue of a robed female monk wrestling a giant carp. The smell of frankincense is strong, and the sound of a distant sitar being played can be heard. Two white-goateed monks approach to greet you.

\end{DndReadAloud}

Frankincense is burning in the censers that flank each of the passage entrances leading further into the abbey.

\subsubsection{Creatures}
The aged monks are Adept Toub and Adept Yestel (both LE \textbf{serene lily adept}). They can answer the PCs questions and are willing to arrange for them to meet with Qudos the morning after they arrive.

Answers to some questions the PCs might ask over the course of their stay at the abbey follow. For questions that aren't presented here, the residents of the abbey answer deceptively about the purpose of the abbey, who Qudos is, and any activities that don't seem in-character for monks who seek serenity. If they feel that answering a question won't compromise them or the abbey, they answer truthfully.


\begin{itemize}
\item \textbf{Q: Why can't we see Qudos sooner?}
\item A: "The abbot is on an Astral retreat searching for the true source of serenity." (This is a lie. Qudos and the cultists want the narcotic in the air to make the PCs as pliant as possible).
\item \textbf{Q: What is going on with the torture chamber?}
\item A: "Some of our adherents feel soothed following a brief but intense encounter with pain. The excruciation chamber allows them to reach bliss in a safe manner." (Technically true, though the abbey's current residents use the chamber to torture reluctant \textbf{serene lily aspirant} for information).
\end{itemize}

\subsubsection{Audience with Qudos}
If the PCs do have an audience with Abbot Qudos, it occurs in the morning in 6. During the brief meeting, the elderly, white-haired monk is polite and listens attentively before claiming he needs time to look into their concerns. He directs other inquiries to Master Ghi (LE \textbf{barbed devil} in disguise).

Qudos is dressed as the other residents, but wears a white sash. His true nature is that he is a LE \textbf{lesser infernal tutor}; when he appears as a human, he does so using his Change Shape ability, and so is not using illusion magic to disguise himself as the other devils are.

\subsection{3. Relaxation Room}
\begin{DndReadAloud}{}
The scent of frankincense is overpowering in this chamber, and the air swirls with a pall of pale smoke. Aspirants sitting or lying on floor mats come here to take a moment to relax. A monk sits against the southwest wall idly playing a sitar.

\end{DndReadAloud}

Rectangular 3-foot by 6-foot jute mats are evenly spaced throughout this chamber, and it is rare for more than half of them to be occupied at any given time. Two large thuribles are always kept burning.

The PCs are welcome to rest in this chamber, and if they spend the night in the abbey, they are given quarters here rather than in the dormitory at 5.

\subsubsection{Creatures}
The sitar player is Master Kethi (LE \textbf{barbed devil}). A successful DC 19 Intelligence (Religion) check recognizes the music being played as being from \textit{Hymns for the Acquisition of Glorious Wealth}, a prominent hymnal of a god of greed and wealth.

\subsubsection{Development}
If the PCs make trouble in Areas 2–7, the \textbf{barbed devil} investigates and assists the monks.

\subsection{4. Kitchen and Dining Area}
\begin{DndReadAloud}{}
The odor of frankincense, so prominent in the other areas of the abbey, is diminished beneath the smells of stir-frying vegetables and roasting barley flour. The grim-faced cook bustles silently around the area while preparing meals. The chamber has a slightly steamy haze from the lack of ventilation, and beads of condensation stand on the walls and ceiling.

\end{DndReadAloud}

During the day, a cauldron of water is kept at a low boil over an open flame, and a stone oven is kept heated. The cook, Jaitna, prepares all the food for the abbey. Most meals consist of a porridge of roasted barley, stir-fried roots and legumes, and butter tea.

\subsubsection{Creatures}
Aspirant Jaitna (N human \textbf{acolyte}) is the only remaining holdout of the abbey's original residents. He refuses to break the vow of silence he took as a young monk, but a PC can intuit the intent and meaning of his expressions and gestures with a successful DC 15 Wisdom (Insight) check. The result of one check is sufficient for each conversation. Jaitna is uncertain of the true nature of the changes at the abbey, but he knows Abbot Qudos and the black-sashed monks never eat, and he is certain they don't sleep either. He assists the PCs as he can, but he doesn't take violent action against any creature.

\subsection{5. Dormitory}
\begin{DndReadAloud}{}
Orderly rows of bunks fill this chamber. Each monk's belongings are rolled into a small bundle that sits at the foot of their cot. Censers set on low round tables burn the abbey's incense through the day and night.

\end{DndReadAloud}

The abbey's aspirants and adepts rest here. Occasionally, this chamber can be found empty during daylight hours.

\subsubsection{Creatures}
At night there are 15 + 1d6 \textbf{serene lily aspirant} and 1d6 \textbf{serene lily adept} present here.

\subsubsection{Development}
A PC that succeeds on a DC 17 Wisdom (Perception) check while looking for secret doors finds the secret door leading to Qudos' chamber (8). A successful DC 17 Intelligence (Investigation) check is then needed to determine how to open the secret door.

If the PCs are detected searching or stealing in this room, the aspirants and adepts attempt to detain and secure them in 7.

\subsection{6. Central Chamber}
\begin{DndReadAloud}{}
Murmured chanting and singing can be heard as you approach the central chamber. The vast room is some 50 feet in diameter. At the northern side is a short staircase leading up to a wide dais with four burning braziers upon it.

\end{DndReadAloud}

The monks of the abbey spend their afternoons in peaceful meditation in the large central chamber. After lunch, most of the aspirants and adepts kneel or sit crosslegged on jute mats and chant or sing their prayers. PCs that understand Infernal recognize that the prayers are a supplication for Mammon's beneficence.

\subsubsection{Creatures}
Prayers are led by Master Miloch (LE \textbf{bearded devil} in disguise) and Master Ghi. During prayer time in the afternoon, there are 15 + 1d6 \textbf{serene lily aspirant} and 1d6 \textbf{serene lily adept} present in addition to the prayer leaders.

If the PCs disrupt a prayer session, they are asked to leave. If they persist, Master Miloch along with Adept Toub and Yestel from 1 attempt to coerce them away.

If the PCs are detected searching or stealing in this room, the \textbf{serene lily aspirant} and \textbf{serene lily adept} attempt to detain and secure them in 7.

\subsubsection{Development}
If the PCs start a fight in this chamber, all abbey residents (except Jaitna) turn against them and fight to the death to protect their compatriots.

\subsection{7. Excruciation Chamber}
\begin{DndReadAloud}{}
A rack and a spiked table occupy much of the space in this chamber, though an iron chair with built-in manacles to secure a humanoid's wrists and ankles is also secured to the floor. Six humanoids can be secured to the wall by manacles attached to iron rings.

\end{DndReadAloud}

A creature who is strapped to the rack or the spiked table can wriggle free by succeeding on a DC 17 Dexterity (Acrobatics) check or can break free with a successful DC 20 Strength (Athletics) check. Creatures secured with manacles can be freed by picking the locks with a successful DC 16 Dexterity check using \textit{thieves' tools}, or they can break free by making a DC 20 Strength (Athletics) check.

A creature strapped to the spiked table is restrained and takes 1d4 piercing damage each minute it remains restrained; it also takes the same amount of damage any time it attempts to move or take any actions while restrained.

\subsubsection{Creatures}
The excruciation chamber is overseen by the ever-present Master Atha (LE \textbf{barbed devil} in disguise), who is almost frighteningly cheerful. She answers most questions asked with a lie, and if her deception is discovered, she feigns ignorance.

\subsubsection{Development}
A successful DC 25 Wisdom (Perception) check while looking for secret doors discovers the secret door leading to Qudos' redoubt; a successful DC 20 Intelligence (Investigation) check is then needed to determine how to open it.

If Master Atha is free, she won't allow the PCs to search the chamber. If she is secured before being questioned, she knows where the secret door is and how to open it, and knows the passphrase to enter 9 from 8 without injury.

\subsection{8. False Lair}
\begin{DndReadAloud}{}
This spartan chamber holds a cot, two tables, some floor pillows, and a wardrobe. It's surprisingly dusty for a living area that is in use. A pair of circles has been etched in the eastern portion of the wall.

\end{DndReadAloud}

Most of the abbey residents believe this dusty chamber is where Abbot Qudos rests, but it's clear no one has stayed here in weeks.

\subsubsection{Tracks}
A track in the dust on the floor runs from 5 to the secret door here in 8, through the hidden passage to 9.

\subsubsection{Development}
The location of the secret door in this area is easy to find. A successful DC 15 Intelligence (Religion) check determines that the circles etched into the wall form two-thirds of Mammon's holy symbol. If a third circle is traced above the two existing circles, the door opens. If the phrase, "\textit{This is the path to profit}," is spoken in Infernal as the circle is traced, the \textit{glyph of warding} inscribed on the opposite side of the wall is not set off.

If, after the glyph is set off, the PCs attempt a short or long rest in 8 or the passageway leading to 9, Qudos opens the secret door from 9 to interrupt them, insisting they leave or attacking if they refuse.

\subsection{9. True Redoubt}
\begin{DndReadAloud}{}
A shimmering portal rimmed in crimson and black flames dominates the center of this chamber. On the floor in the center of the northern wall is a small altar surrounded with a hoard of coins and gems.

\end{DndReadAloud}

\subsubsection{Creatures}
Qudos spends most of his time in this chamber, where he confers with Master Ghi and Master Atha or with other devils via the portal. If he is encountered here, he urges the PCs to leave but fights if they refuse to do so. If reduced to 15 hit points or less, he uses his action to Disengage and exits through the portal; if he's able to flee through the portal, it closes behind him instantly.

\subsubsection{Development}
The portal can be closed by destroying the altar, which has AC 14, 27 hit points, and a damage threshold of 10. If the altar is not destroyed, the portal opens again if items and coin worth more than 1,000 gp are piled on or near it.

If a creature without the fiend type attempts to enter the portal, it bumps into it as if it was a solid object; inanimate objects also cannot pass through it if they are not carried by a fiend.

\subsubsection{Treasure}
The following can be found in this chamber: 3,345 sp, 1,096 gp, a jade and ivory statuette of an elephant (worth 65 gp), a scroll of \textit{silence}, and a pair of \textit{goggles of night}.

\section{Concluding the Adventure}
If Qudos escapes the PCs, he may return to have his vengeance at a future date, bringing Masters Ghi and Atha if he can. If he is killed, but Ghi and/or Atha escapes, they similarly seek vengeance. Once the abbey is cleared of its infernal residents, it inevitably becomes the lair of some dangerous creature or other group if the PCs don't make an effort to find new residents or render the lair unlivable.

\chapter{Incident at Wrackwater}

\begin{DndTable}{c}

\textbf{An adventure for four to five characters of 5th level says 6th level.} characters}\\
\end{DndTable}

\section{Adventure Background}
\DndDropCapLine{F}{}ar along the cold, rocky coastline lay the oft-fogbound port of Wrackwater Cove. Once a tiny frontier furrier outpost, Wrackwater has grown to become a burgeoning trade port. So when the town's lighthouse beacon inexplicably goes out during the month in which the season's first merchant ships are due to arrive, the townsfolk deployed maintenance teams immediately . . . who never returned.


Now, a fortnight later, three search parties and multiple individuals remain missing or confirmed dead. Moreover, due to the unlit beacon, a merchant ship is believed to have already been lost to the infamous fog and rocks outside Wrackwater Cove.

The sounds of shrieks and crashes, aggressive flocks of unfamiliar seabirds, and strange lights and inhuman shapes cavorting atop the darkened lighthouse are being reported regularly. This, combined with the townsfolk already dead or missing, has made the locals understandably fearful.

With local authorities and a citizenry demoralized, ill-equipped, and too thinly stretched to do anything more, the hue and cry for help went out. Across the region, messengers, signal fires, and courier birds carry news of a desperate populace . . . and of the urgent danger and exotic rewards that await at Wrackwater Cove.

\section{Adventure Hook}
Wrackwater can be located in any rugged, isolated coastal region prone to significant trade and travel hazards like treacherous currents reefs and/or fog.

The following hooks can used to start the PCs on this adventure. The PCs:


\begin{itemize}
\item Are wearily traveling back to civilization from exploits even further afield than Wrackwater Cove and stop there for supplies.
\item Have associates, relatives, or business interests aboard the missing ship or elsewhere in Wrackwater's vicinity.
\item Heard word of the incident via couriers sent seeking aid, local rumors, or other means.
\item Hear of a "haunted" lighthouse in a desolate area and of the town now threatened by it.
\end{itemize}

\section{Wrackwater Cove}
Following news of the PCs arrival, town council members approach them regarding Wrackwater's plight. The councilmembers wish to employ the PCs to search the lighthouse for evidence of the missing townsfolk and to ensure that the lighthouse is cleared of whatever "spooks and specters" are now occupying it. The beacon must also be relit immediately; the town relies on trade to survive, and with the lighthouse not functioning, the very continued existence of the town is threatened.

If the PCs are willing to take the job, the council offers them 1,200 gp in ermine fox and beaver pelts, plus another 900 gp in raw copper abalone shells and coldwater pearls.

Additionally, the council offers each PC: 1 \textit{sealskin coldsuit} (which grants the wearer a swimming speed of 20 feet) and 2 glowmold globes.

\section{The Lighthouse}
The mile-long approach to the lighthouse winds along a rocky promontory. As the PCs draw near, they discern sundry piles of garbage, broken furniture and tableware, wine casks, and much more of the same. It all appears to have been scattered about the structure's base, seemingly thrown willy-nilly from broken-out windows above. The lighthouse exterior also appears somewhat worse for wear, its outer walls streaked with bright stains of unknown origin.

Once the PCs approach within 300 feet of the lighthouse, they easily spot several malformed, winged shapes taking flight from one of the lighthouse windows (4).

\subsubsection{Creatures}
Flying toward the PCs are six \textbf{gullkin}. The \textbf{gullkin} don't reply if any attempt is made to converse. Once the PCs get to about 150 feet of the island, the \textbf{gullkin} attack. Their shrieking calls summon two swarms of maddened sea birds (use statistics for \textbf{swarm of ravens}) en route. The swarms fight to the death. Once there are only two \textbf{gullkin} remaining, they retreat to 4.

\subsubsection{Making an Entrance}
Upon arriving at the lighthouse, the PCs find it has been barricaded, battered down, and then re-barricaded. Breaking through the boarded-up doorway (or the western window) requires a successful DC 12 Strength (Athletics) check.

Carefully removing the boards and entering quietly requires the PCs to work at it for 5 minutes and to make a successful DC 14 Dexterity (Stealth) check to do so without making noise; a PC proficient with \textit{carpenter's tools} makes the check with advantage.

\subsubsection{Dozing Dangers}
On the building's western face, a single salt-crusted window that has been nailed shut looks into a darkened, recently ransacked bunkroom. Forced entry via this window or the main entrance awakens the creatures in 3, alerting them to intruders.

\subsubsection{Down to the Depths}
Diving off the western side of the island and swimming 30 feet down leads to Sea Caves Area 9, though the cave opening is in no way visible from land.

An unlabelled version of the map can be found here.

\subsection{1. Main Storeroom}
\begin{DndReadAloud}{}
This supply room has been thoroughly ransacked and is in a complete shambles. Heaps of shattered crates and barrels, rotting foodstuffs, broken lenses, tools, and other refuse litter the floor. The northern wall is dominated by a narrow staircase leading upward. The door to the west stands ajar and has clearly been forced open. The closed door to the south appears to be the only thing in the area that has escaped vandalism.

\end{DndReadAloud}

\subsubsection{Difficult Terrain}
The heaps of refuse in this room make it difficult terrain.

\subsubsection{Out of Place}
A successful DC 14 Intelligence (Investigation) check reveals several items—such as a damaged block and tackle, stacks of ruined sail cloths, tar barrels, etc.—that are far more nautical in nature than would be required in any lighthouse. All of these items show signs of lengthy (and recent) exposure to sea water.

\subsection{2. Lighthouse Keeper's Quarters}
\begin{DndReadAloud}{}
This apparently vacant living area shares the storeroom's pillaged and vandalized condition. Overturned beds, broken dishes, and other day-to-day items lay strewn about. A peat stove squats near a west-facing window.

\end{DndReadAloud}

\subsubsection{Secret Cache}
A successful DC 16 Wisdom (Perception) check notices a bloody hand smear on the wall behind the stove that reveals the location of a hidden compartment.

\subsubsection{Treasure}
Inside the compartment is a small lockbox, which can be picked open with a successful DC 17 Dexterity check using \textit{thieves' tools}; it contains 13 cp, 8 sp, 22 gp. In a leather case on top of the lockbox is a small, waterproof magnetic navigators' compass (worth 175 gp). A chest on the floor next to the stove has had its contents strewn about the room and is empty.

\subsection{3. Workroom}
\begin{DndReadAloud}{}
A workbench stands against the southern wall of what is obviously a workroom. Visible here are tool chests, spare window panes, beacon mirrors, painting supplies, and similar maintenance and repair items as one might expect to find in a lighthouse. While it has been well-plundered, the room has suffered comparatively little of the wanton defacement seen elsewhere.

\end{DndReadAloud}

\subsubsection{Creatures}
If the PCs have not already encountered them, two \textbf{bilge gremlin} and two \textbf{bilge gremlin bosun} are napping here, hidden inside a crude hunter's blind formed of stacked crates, rope coils, and barrels. While in the blind, the gremlins have full cover, but the soft sounds of them snoring can be detected with a successful DC 14 Wisdom (Perception) check.

\subsubsection{Sea Cave Access}
The blind also conceals access to the sea caves below, from which the gremlins tunneled up into the lighthouse; it can easily be located with even a cursory investigation of the room. Descending down into the sea caves takes the PCs to Sea Caves Area 1.

\begin{DndReadAloud}{}
Behind the blind, you see a rough-hewn floor breach that peers down into a dark, dripping tidal tunnel running eastwest beneath the lighthouse. Five small grappling hooks ring the aperture, their knotted hemp cords dangling into the damp darkness below.

\end{DndReadAloud}

\begin{DndSidebar}[float=!b]{Bilge Gremlin Tactics}
\textbf{bilge gremlin} (and \textbf{bilge gremlin bosun}) are wily and slippery foes. They rely on their evasiveness to be effective, and so they take full advantage of their ability to cast \textit{misty step at will} nor \textbf{bilge gremlin bosuns} can do this.}, rarely ending a turn adjacent to an enemy if at all possible. Likewise, whenever they detect danger, they make frequent use of their at-will \textit{invisibility} nor \textbf{bilge gremlin bosuns} can do this.} to spring ambushes.



\end{DndSidebar}

\subsection{4. Lower Landing}
\begin{DndReadAloud}{}
As you ascend the stairs and cross onto the next level of the lighthouse, you see that the floor here is similarly strewn with detritus—a squalid, one-foot-deep mixture of kelp, trash, tree branches, torn clothing, and soiled bedding. As you take it all in, the dull white of bone catches your eye, and a quick assessment flags it as a human femur.

\end{DndReadAloud}

\subsubsection{Creatures}
Three \textbf{giant flea} have infested this batch of detritus. They immediately attack any living, warm-blooded creatures that move more than 10 feet into the chamber.

\subsubsection{Treasure}
If the PCs spend 5 minutes or more searching the offal and debris in this area, determine what they find by having one PC make an Intelligence (Investigation) check. Once you have the result, consult the "Trash or Treasure?" table to determine what they find; a higher result finds the items from lower results as well.

\begin{DndTable}[header=Trash of Treasure?]{XX}

DC & Treasure Found\\
12 & A miniature, antique spyglass (worth 400 gp).\\
14 & A half-eaten human arm, its bloody sleeve still bearing the insignia of a Wrackwater sergeant-at-arms.\\
16 & A collection of international coinage and broken jewelry (worth 80 gp total).\\
18 & A bolt case with 3 +2 crossbow bolts in it.\\
\end{DndTable}

\subsubsection{Stairs}
The stairs continue upward to a second dimlylit landing. The landing and nesting area above Area 4 is currently unoccupied; however, its normal attendants (1d4 + 1 \textbf{gullkin}) won't be gone long. Roll 1d6 and consult the "Gullkin Return" table to determine when the creatures return.

\begin{DndTable}[header=Gullkin Return]{XX}

d6 & Result\\
1–2 & Gullkin arrive at 5 while the PCs are fighting the \textbf{giant flea} in Area 4; the \textbf{gullkin} screech and join the combat.\\
3–4 & \textbf{Gullkin} arrive at Area 5's window while the PCs are ascending/descending from Area 4.\\
5–6 & \textbf{Gullkin} arrive while the PCs are searching Area 4.\\
\end{DndTable}

\subsection{5. Upper Landing}
\begin{DndReadAloud}{}
This area resembles the lower landing, but contains a crowded nest for what looks like a half-dozen creatures. A glint of light catches your eye on the seaward windowsill, and you see two handheld sailor's signal mirrors resting there.

\end{DndReadAloud}

The next flight of stairs winds another 40 feet up, terminating at an unlocked, wooden ceiling hatch.

\subsubsection{Treasure}
If the PCs search the nest area, a successful DC 14 Intelligence (Investigation) check discovers a bag containing 4 rings of gold and silver and a broken necklace studded with fine gems (worth 140 gp total). There is also a matching pair of 2-inch-wide bracelets that are magical and function as \textit{bracers of defense}.

\subsection{6. Widow's Walk}
To enter this area, the PCs must open the wooden ceiling hatch where the stairs end. It is likely that lifting the hatch immediately initiates combat with the creatures in this area, but if the PCs have been particularly stealthy, it's possible that the creatures are unaware of their presence and could be surprised.

\begin{DndReadAloud}{}
As you open the hatch, you can see a strange gathering. A gremlin stands beside a sahuagin, both in conversation with a fierce-looking gullkin.

\end{DndReadAloud}

\subsubsection{Creatures}
Stubbe (a \textbf{bilge gremlin bosun}) stands beside a \textbf{sahuagin}, speaking with Laridaena (a \textbf{gullkin hunter}); the \textbf{gullkin hunter} clutches a nautical signal mirror and a logbook in her claws. The \textbf{bilge gremlin bosun} and \textbf{sahuagin} both carry crude spyglasses. Behind Stubbe stand two bilge gremlins, and beside Laridaena stands one \textbf{gullkin}.

\subsubsection{Argument}
A successful DC 13 Wisdom (Insight) check determines an argument is underway in the primordial language. If the PCs understand primordial, they deduce that the \textbf{sahuagin} is an ambassador, up from the sea caves, and seems to be arguing about some sort of loot-sharing agreement with Laridaena, who is queen mother of the gullkin tribe.

\subsubsection{The Hatch}
Climbing up through the hatch and onto the widow's walk for Medium or larger creatures is considered difficult terrain due to the awkwardness of the narrow staircase and tight quarters.

\subsubsection{Tactics}
The \textbf{gullkin} take flight and attack from the air. The \textbf{bilge gremlin bosun} and \textbf{sahuagin} charge the hatch. If Laridaena is reduced to below half her hit point maximum, she and her \textbf{gullkin} compatriot flee the area and do not return.

\subsubsection{Logbook}
The logbook is a standardized instruction manual on sending and translating maritime code. Scrawled on the inside cover is a crude diagram detailing routes to here (Lighthouse Area 6) from Sea Caves Area 8.

\subsection{7. Beacon House}
\begin{DndReadAloud}{}
Through the windows of the beacon house, you can see Wrackwater's extinguished beacon—and all around its base, an enormous nesting pile.

\end{DndReadAloud}

\subsubsection{Restoring the Beacon}
If there no enemies harrying the PCs, it is a simple matter to relight the beacon brazier, but one of the mirrors of the mirror array is cracked. It can be repaired with a spell such as \textit{mending} or using one of the beacon mirrors in 3. Replacing the mirror manually requires a successful DC 14 Dexterity check using \textit{thieves' tools}, \textit{tinker's tools}, \textit{jeweler's tools}, or \textit{glassblower's tools}; a creature using \textit{glassblower's tools} makes the roll with advantage.

\subsubsection{Treasure}
With a successful DC 12 Intelligence (Investigation) check, searching the refuse discovers a coral ocarina adorned in pearls and agate (worth 375 gp) amid the bones and offal of what must be some of the missing Wrackwater townsfolk.

\section{The Sea Caves}
The submerged grottoes beneath the lighthouse are tidally cleansed (except 5), maintaining water clarity and visibility if sufficiently illuminated. Natural surfaces throughout the sea caves are encrusted with countless forms of sea life (nudibranchs, crustaceans, seaweeds, etc.). It is chilly and dark, and swimming is occasionally difficult (due to tidal flow).

\subsection{1. Access Tunnels}
\begin{DndReadAloud}{}
You drop down through the hole in the lighthouse floor and find yourselves inside a naturally formed cavern—and in almost complete darkness save for the faintly glimmering light provided by a distant bioluminescent glow coming from the east. At the western terminus is a sinkhole filled with seawater that heaves and ebbs in time with the distantly crashing waves.

\end{DndReadAloud}

This tunnel runs east, widening after 40 feet into 2. Diving down into the sinkhole to the west (3) leads to 7.

\subsection{2. Cargo Grotto}
\begin{DndReadAloud}{}
The tunnel widens into a large water-filled grotto. Bioluminescent plants and invertebrates thrive along the walls above and below the waterline. You stand on a gravelly beach littered with mostly empty or ruined barrels and crates—and a crude, but sturdily built raft equipped with homemade oars and loaded with water-logged cargo that's been pulled up onto the beach of this enormous subterranean tide pool. A small rock islet, also piled with rubbish and damaged cargo, lies fifteen feet ahead.

\end{DndReadAloud}

The raft measures 10 feet by 10 feet and is seaworthy; it safely accommodates 4 Medium creatures with gear.

\subsection{3. Sinkhole Tube}
This twisting sinkhole tube descends 15 feet before widening into an underwater cavern.

\subsection{4. Rock Islet}
Piled with more cast-off cargo (making it difficult terrain), this stone islet affords PCs some level solid footing, partial cover, and a dry spot to rest.

\subsubsection{Treasure}
Spending 10 minutes or more searching through the refuse uncovers 3 cp, 4 sp, and a small leather bag containing 13 fangs and 11 claws from assorted beasts.

\subsection{5. Bilge Gremlin Berthing}
\subsubsection{Ambush}
If the PCs attempt to pass the island (either by raft or swimming), moving within 5 feet of its eastern edge triggers an ambush from below.

\subsubsection{Creatures}
Six \textbf{bilge gremlin} surge up from the shallows and attack. If the PCs slay three or more of the gremlins, they call out for help. The following round, one \textbf{bilge gremlin} arrives (on the other \textbf{bilge gremlin} initiative count), riding an \textbf{aquatic basilisk} (use statistics for \textbf{basilisk}, adding the Amphibious trait and a 20-foot swimming speed).

\subsubsection{Sea Bottom}
The bottom of this gigantic tide pool is akin to an underwater junkyard, with piles of cargo junk and debris everywhere.

\subsection{6. Cargo Hold}
\begin{DndReadAloud}{}
Clumps of garbage float nearby the corpse of a skinned and decapitated harbor seal that's half-beached on the ground here. Numerous crates and barrels line the walls.

\end{DndReadAloud}

\subsubsection{Treasure}
Various crates contain 4 sets of \textit{painter's supplies}, 10 yards of silk (worth 100 gp), and 50 lbs. of salt (worth 3 gp).

\subsection{7. Lower Depths}
The sinkhole in 3 descends 15 feet through a narrow channel before opening up into a fully submerged cavern that branches off to the east.

\subsection{8. Brig}
\begin{DndReadAloud}{}
Ahead, a large cavern opens before you. A gremlin swims beside two merrow and an enormous, serpentine yellow eel covered in dark spots.

\end{DndReadAloud}

\subsubsection{Creatures}
A \textbf{bilge gremlin bosun} leads two \textbf{merrow} and a \textbf{muraenid} around this area.

\subsubsection{Prisoner}
With a successful DC 14 Wisdom (Perception) check a PCs detects a captured \textbf{merfolk} scout named Minn'Vyr who is shackled and concealed inside an alcove to the southeast. She can be freed with a successful DC 16 Dexterity check using \textit{thieves' tools}. If freed, Minn'Vyr tells the PCs that she was captured while witnessing a ship go down weeks ago. She has limited information but says that a sudden radical uptick in \textbf{bilge gremlin} salvage operations—and there has been an increasingly strange stream of assorted creatures passing through the region that normally don't consort with each other.

\subsection{9. Tidal Cave}
\begin{DndReadAloud}{}
This chamber contains more pilfered, piled-up cargo. To the west, the cave mouth opens up to the sea.

\end{DndReadAloud}

\subsubsection{Creatures}
Chief Bosun Chark (a \textbf{bilge gremlin bosun}) is here separating out metal and valuables from common cargo and wreckage. He gives direction to a \textbf{monstrous crab} with tentacles coming out of its maw (use statistics for \textbf{chuul} with the changes noted below); it displays keen intelligence for a crab, and its pincers root through and tear into piles of scavenged wreckage with alarming efficiency. Two \textbf{sahuagin} swim alongside the \textbf{bilge gremlin bosun}, observing operations.

If combat ensues in 8, all but the \textbf{monstrous crab} move into the northern alcove behind piled cargo to ambush the PCs. The \textbf{monstrous crab} moves to the southeastern part of the chamber and burrows partially down into the sand, camouflaging itself with the copious kelp in the area; it can be spotted with a successful DC 16 Wisdom (Perception) check.

\subsubsection{Creature Change}
The \textbf{monstrous crab} creature type is monstrosity; it only understands primordial; its Intelligence is 8; and it does not have the Sense Magic trait or the Tentacles action. Any creature grappled by its Pincer attack is also restrained.

\subsubsection{Treasure}
Three casks of rum (worth 100 gp total); various gems and jewelry, holy symbols (worth 295 gp total); and a chest filled with bars of silver (worth 300 gp total).

\subsection{10. Aquatic Stable}
\begin{DndReadAloud}{}
As you swim down this winding cavern heading southeast, you notice the abundance of seaweed increasing as you descend the twisting tunnel. Flourishing kelp soon stretches from floor to ceiling, thickening to a jungle-like density, reducing visibility to five feet. You can't see it due to the low visibility, but, after traveling twenty to twenty-five feet, you sense the tunnel widening into a larger cavern.

\end{DndReadAloud}

\subsubsection{Creatures}
Two \textbf{aquatic basilisk} (use statistics for \textbf{basilisk}, adding the Amphibious trait and a 20-foot swim speed) and three \textbf{hunter shark} used for mounts are stabled here. There are also 2 \textbf{giant sea horse} here, but they are non-hostile and tame.

\section{Concluding the Adventure}
Upon the PCs' triumphant return to Wrackwater, a reclamation/repair team is immediately assembled and sent to the lighthouse. The PCs can confirm that no missing Wrackwater citizens remain alive (near the lighthouse) and that a ship indeed must have gone down somewhere outside the cove that was being scavenged from by gremlins below the lighthouse.

If you wish to further explore any of the plot threads of this adventure, "Incident at Wrackwater" can be played as a direct lead-in to "The Salons of Mother Celeste".

\chapter{The Lost Forge}

\begin{DndTable}{c}

\textbf{An adventure for four to five characters of 5th level says 6th level.} characters}\\
\end{DndTable}

\section{Adventure Background}
\DndDropCapLine{H}{}igh in an alpine forest lie the remains of a temple that was dedicated to dwarven gods of fire and forge. The site was abandoned by its creators after a horrific dragon attack which left most of the stone walls devastated by dragon fire. The dwarves felt the place had been cursed, and so this site became known as "The Lost Forge" by the local villagers who have long spun stories of the dwarven ghosts that haunt it.


More recently the twin warlords Kadal and Valgred have taken over the ruins with their band of orcs and ogres and made it the base of operations for their conquest of the region. Kadal quickly awoke the Lost Forge and began using it to make armor and weapons for their followers, while Valgred communed with the spirits of the ruins to gather what secrets they held. The twins then began raiding small villages in the mountains and attacking trade caravans. The local villagers pray that their homes won't be next.

\section{Adventure Hook}
The following can used to start the PCs on this adventure:


\begin{itemize}
\item \textbf{\textit{Ale Aid.}} While traveling through a mountain forest, the PCs come across a frantic and desperate dwarf guard named Jaster (N dwarf guard), who is a survivor of a dwarven ale caravan that was recently attacked. He's tracked the culprits and found where the survivors of the attack who were captured have been taken. He pleads for the PCs to aid them.
\item \textbf{\textit{Village Protectors.}} While stopping to rest at a mountain village, the PCs are approached by the folk who live there and offered 100 gp each to put a stop to the twin warlords before their village is pillaged.
\end{itemize}

\section{Sibling Disagreements}
Twins Kadal (LE \textbf{thursir armorer}) and Valgred (LE \textbf{thursir hearth priestess}) have been frequently arguing as of late—over how to proceed with their continuing conquest of the mountains. Kadal believes they should push forward with more raids, while Valgred believes they need to be more reserved and strategic with their attacks. The continual arguments have split their followers into two factions. Those who follow Valgred wear blue headbands and those that follow Kadal decorate their armor with animal skulls.

PCs who spend some time observing the occupants of the encampment quickly pick up on the tension between the two groups, though if none of them speak orc or Giant they might not learn the specifics.

\section{Ruins of the Lost Forge}
The encampment is surrounded by a fifteen-foot-tall palisade that collapsed in a few places in the past and has been hastily patched more recently. Dark scorch marks and slashes from powerful claws are obvious on most defensive surfaces that haven't collapsed or been covered over by vines. The sounds of a loud argument in 6—which might eventually erupt into a fight—cover much of the noise the PCs might make while infiltrating the camp.

Aside from 6 (which is lit by a campfire) and 15 (which is lit by the forge), all other areas of the camp are quite dark.

For creatures who are noted as being asleep, have the PCs roll a group Dexterity (Stealth) check against the creature's passive Perception.

\begin{DndSidebar}[float=!b]{Kadal's Craftsmanship}
Any of the weapons and armor crafted by Kadal are of extremely high quality and can be sold for three times their normal listed price. The \textbf{orc} all carry greataxes, and the \textbf{veteran} carry a \textit{battleaxe} and a \textit{handaxe}, all crafted by Kadal.



\end{DndSidebar}

An unlabelled version of the map can be found here.

\subsection{1. South Wall}
\begin{DndReadAloud}{}
Vines rise from the forest floor to cover much of a fifteen-foot-high wooden palisade. A stone building marred by large swaths of scorch marks stands at one corner of the wall. The center of the palisade in this area collapsed sometime in the distant past and has been more recently filled with stone boulders that are stacked ten feet high. Two orcs wearing blue headbands crouch near the pile of boulders playing a game of dice.

\end{DndReadAloud}

\subsubsection{Getting Through}
The PCs can attempt to get through this section of the structure by scaling the palisade with a successful DC 14 Strength (Athletics) check or squeezing past the boulder barrier with a successful DC 14 Dexterity (Acrobatics) check; a failure by 5 or more alerts the creatures in 6 who come to investigate the source of the noise.

\subsubsection{Creatures}
Two \textbf{orc} have been assigned as lookouts in this area outside the palisade, but both are currently playing a game of dice by the boulder barrier to pass the time. They wear blue headbands. If they see any trouble, they attempt to run to the main gate in 3 to warn the camp.

\subsubsection{Hostilities}
An orc with skulls adorning the pauldrons of his armor approaches after the most recent round of dice has completed to ask (in orc) for a report of any sightings. During the interaction, PCs can make a successful DC 12 Wisdom (Insight) check to discern that the new orc is not welcomed by the others and there is a sense of hostility between them.

\subsection{2. Ogre Sleeping Area}
\begin{DndReadAloud}{}
A 25-foot-tall stone building forms one side of the gate complex of this encampment. The floor is covered in with hay and detritus, and the air is stagnant with the strong odors of rotting food. The tattered remains of dwarven banners still adorn the walls, now too damaged and worn with age to read.

\end{DndReadAloud}

The building can only be entered from within the encampment. A successful DC 12 Intelligence check using \textit{mason's tools} or dwarf Stonecunning determines that though the structure presently shows significant signs of wear and neglect, it was once very well-crafted and suffered an attack sometime in the past.

\subsubsection{Creatures}
Two \textbf{ogre} are currently napping in this chamber, one wears a blue headband and the other has a single pauldron decorated with the skull of a large bird. They have a passive Perception of 8.

\subsection{3. Front Gate and Ruined Building}
\begin{DndReadAloud}{}
The front wooden gate is 15 feet tall and shows signs of recent repair work done by unskilled hands. The connected stone building north of the gate is missing a corner and shows signs of being burned. The inside is a mess of blackened, damaged rubble. A set of double wooden doors is the only exit from this ruined building, which leads into the encampment.

\end{DndReadAloud}

\subsubsection{Getting Through}
The front gate and the ruined building's double doors are both barricaded from the inside, requiring a successful DC 20 Strength (Athletics) check to break through.

\subsubsection{Hunting Traps}
A poison-laced hunting trap has been hidden in the rubble of the ruined tower. A creature that steps onto the trap must succeed on a DC 13 Dexterity saving throw or take 1d4 piercing damage and 1d6 poison damage. A successful DC 15 Wisdom (Perception) check reveals the hidden trap.

\subsubsection{The Ancient Attacker}
A successful DC 15 Intelligence (Investigation or Nature) check reveals the damage marks on the inside are claw marks. A success by 5 or more also reveals the claw marks and burnt stones were caused by a dragon (probably red, given the burn marks).

\subsection{4. Fallen Watch Tower}
\begin{DndReadAloud}{}
The remains of a stone tower are strewn across this area. The wooden palisade to the east of the once-imposing tower has a large gap that has been plugged by a wagon flipped onto its side.

\end{DndReadAloud}

\subsubsection{Difficult Sneaky Terrain}
Walking through the fallen tower remains is considered difficult terrain and due to the stone debris, all Dexterity (Stealth) checks are made with disadvantage.

\subsubsection{Creatures}
Two \textbf{orc} wearing armor decorated with animal skulls stand guard between the wagon and fallen tower on the outside of the palisade. They talk (in orc) about how Valgred needs to stop getting in the way of Kadal's plans.

\subsection{5. Ruined Building}
\begin{DndReadAloud}{}
This large stone building is mostly destroyed, with one side completely collapsed. Dark scorch marks mar what is left of the structure. To the west of the remains there is a gap in the wooden palisade that has been hastily repaired with an overturned wagon.

\end{DndReadAloud}

\subsubsection{Hunting Traps}
There is a poisoned hunting trap (see 3) hidden in the debris of the ruined building.

\subsection{6. Firepit Lounge}
\begin{DndReadAloud}{}
Three large logs surround a firepit which is the only light source in the encampment. Nearby rests a haphazard stack of wooden barrels, some open, some intact, and others smashed.

\end{DndReadAloud}

If the creatures in this area have not been already alerted to investigate 1, also read the following:

\begin{DndReadAloud}{}
The occupants of the firepit are embroiled in a loud argument. Four orcs, each wearing armor decorated with animal skulls, square off against one towering ogre in a blue headband who looks to be on the edge of violence.

\end{DndReadAloud}

\subsubsection{Creatures}
Three \textbf{orc} and one \textbf{orc veteran}—all followers of Kadal—are arguing loudly with an \textbf{ogre} (obviously drunk) who is aligned with Valgred. The \textbf{ogre} is suffering from the poisoned condition due to his overindulgence of dwarven ale.

\subsubsection{Creature Change}
For the \textbf{veteran}, replace all mentions of "\textit{longsword}" with "\textit{battleaxe}" and "\textit{shortsword}" with "\textit{handaxe}."

\subsubsection{They Just Keep Coming}
If battle begins in this area and lasts longer than 2 rounds, Nogg from 7 is alerted and joins the combat.

After 4 rounds of combat, Valgred hears the commotion and comes outside of her chamber in 8 to join the battle. Any creatures with blue headbands (potentially two \textbf{ogre} and one \textbf{orc}) within range act as her allies for the purposes of her special ability Call for Aid.

After 1 round of combat, the orc in 9 (loyal to Valgred) awakens and joins the fray.

\subsubsection{Hostilities}
Roll a d20 when the PCs move within sight of this area. On a 1–10 the argument continues as it is. On an 11–20 the \textbf{ogre} bellows with rage and attacks the \textbf{orc}. If the \textbf{ogre} attacks and the PCs do not intervene, Nogg emerges from 7 to break up the fight after 2 rounds.

\subsection{7. Prisoner Holding Area}
\begin{DndReadAloud}{}
This stone building is one of the least damaged in the encampment. The structure has been converted to a holding area for prisoners, with chains anchored into the western wall and the eastern portion being used as storage for recently acquired supplies.

\end{DndReadAloud}

\subsubsection{Creatures}
Nogg (CE \textbf{alleybasher ogre}), is guarding the dwarven prisoners and boxing with some empty wooden crates in the middle of the building. Nogg wears a blue headband and is itching for a fight. If she notices the PCs approach, she consumes a \textit{potion of speed} that Valgred gave her.

\subsubsection{Rallying the Troops}
Ruby Runemaker (LN \textbf{pike guard captain}) and six \textbf{pike guard} are chained together on the western side of the building. If the \textbf{pike guard} are freed, Ruby and her people are more than happy to bring the fight to the forces of the twin warlords. If the GM and/or the PCs prefer, the dwarves can head to other areas of the camp to clean up groups of enemies the PCs have not engaged yet rather than having them travel with the PCs.

\subsubsection{Stolen Supplies}
The crates in this area contain 3 potions of healing, a crate of salted jerky that counts as thirty days of rations, and various spices and cooking herbs (worth 100 gp total).

\subsection{8. Valgred's Chamber}
A squat stone building stands relatively intact at the center of the encampment, though it too bears weathered scorch marks.

\begin{DndReadAloud}{}
The walls of this room are decorated with tapestries, and a bed piled high with animal skins of various types takes up the southeast corner of the room. On the east side of the room, there are stacks of tomes and maps of the region spread over a table.

\end{DndReadAloud}

\subsubsection{Creatures}
Valgred, a \textbf{thursir hearth priestess}, is here studying maps of the region and plotting out the locations of nearby villages and towns they might raid next.

If a fight ensues in Area 6 and lasts more than 4 rounds, Valgred realizes it might be more than just the usual tussling amongst the troops and joins the battle when she sees the PCs.

\subsubsection{Bonded in Loyalty}
Any of Valgred's allies that wear blue headbands count as thursir giants for her Call for Aid ability.

\subsubsection{Treasure}
Valgred wears \textit{hide armor} (sized for a Large creature) and carries a \textit{wand of binding}. Her \textit{runic staff} is a magic quarterstaff that deals an additional 1d8 lightning damage on a hit once per short rest.

\subsection{9. Outhouses}
\begin{DndReadAloud}{}
Two stone outhouses stand in a state of disrepair along the southern palisade with the wooden doors of both displaying indications of advanced rot.

\end{DndReadAloud}

\subsubsection{Creatures}
One \textbf{orc} wearing a blue headband is currently napping on the western side of the outhouses. The \textbf{orc} has a passive Perception of 10. If a battle takes place in 6, he wakes and joins after 1 round.

\subsection{10. Converted Stables}
\begin{DndReadAloud}{}
An expansive stone building with two solid wooden doors still stands strong despite the scars from the destructive attack in the past. The floor is covered in hay, and it's clear it was once a stable, but now it's been converted into a makeshift living space. A desk stands against the southern wall, and a worn wooden sits on the floor between a pair of bunk beds against the eastern wall.

\end{DndReadAloud}

This large stone building was used as a stable by the dwarves, but it's now being used as sleeping quarters by followers of Kadal. Once the occupants are dealt with, this room can be used by the PCs to safely have a short rest if the doors are barricaded.

\subsubsection{Creatures}
Four \textbf{orc} and one \textbf{orc veteran} are currently slumbering.

\subsubsection{Creature Change}
For the \textbf{veteran}, replace all mentions of "\textit{longsword}" with "\textit{battleaxe}" and "\textit{shortsword}" with "\textit{handaxe}."

\subsubsection{Entering the Room}
The \textbf{orc} have a passive Perception of 10; the \textbf{orc veteran} is 12. If a PC attempts to open the door quietly, their Dexterity (Stealth) check is made with disadvantage due to the doors creaking loudly when they open. If the PCs just open the door without any concern for quietness, the guards awaken, but are surprised for the first round.

\subsubsection{Treasure}
The chest contains a pouch of eight rubies (worth 200 gp total) and a silver chalice (worth 45 gp).

\subsection{11. Lost Forge Entry}
\begin{DndReadAloud}{}
The entrance of this towering stone building has two heavy stone doors etched with depictions of dwarven blacksmiths kneeling before the sky. Two stone statues flank the doorway. The statue on the left is an armored dwarf with a maul raised to the sky; the head of this statue has been crudely removed and lies on the ground nearby. The statue on the right has been partially melted by an intense source of heat.

\end{DndReadAloud}

The two statues once bore the likenesses of the dwarven deities of the forge; the one to the north was destroyed during the dragon attack long ago, and the one to the south has been defaced by the followers of the twin warlords. The double doors are currently shut, but not locked.

\subsection{12. Chamber of Armor}
\begin{DndReadAloud}{}
Unlike what you've seen elsewhere in the encampment, this chamber is spotless and kept in a state of impeccable repair. Eight stands holding armor forged of dark metal and adorned with skulls line the walls in rows as straight as the finest honed edge. There are no light sources in this darkened chamber, except for a glowing light—and oppressive heat—straight ahead coming from the forge. You also hear the ringing sounds of someone hammering on metal.

\end{DndReadAloud}

Kadal prefers a clean workplace and has his followers routinely clean the building to his standards.

\subsubsection{Forging in Progress}
When the PCs enter this area, Kadal is standing north of the forge, by the anvil (see 15), hammering out a \textit{greatsword} blade. He is initially just out of view, though the PCs can see the end of the anvil and the sparks flying off of it as the hammer comes down. Due to the hammering, it's possible Kadal won't notice the PCs until they get quite close.

\subsubsection{Treasure}
The suits of armor lining the hall are 8 suits of splint mail adorned with skulls, forged by Kadal and sized for 9-foot-tall (Large) creatures; it weighs twice as much as regular splint mail.

\subsection{13. Kadal's Chamber}
\begin{DndReadAloud}{}
This room contains a large bed piled with bear furs in the western corner. Two heavy iron chests stand against the northern wall.

\end{DndReadAloud}

\subsubsection{Treasure}
The iron chests are filled with plunder from recent raids and require either a successful DC 18 Dexterity check with \textit{thieves' tools} or a successful DC 20 Strength (Athletics) check to break them open. One contains 1,100 gp and 4 bloodstone gems (worth 200 gp total). The other contains 175 cp, 1,700 sp, and 2 moonstones (worth 100 gp total).

\subsection{14. Blacksmith Storage}
\begin{DndReadAloud}{}
This is a simple room occupied mostly by crates of copper ore. A stack of dusty blacksmithing molds and an impressive collection of blacksmithing tools clutter a sturdy workbench along the eastern wall.

\end{DndReadAloud}

\subsubsection{Treasure}
There are enough tools in this room to make three complete sets of \textit{smith's tools}. There are 7 crates of copper ore (worth 35 gp total), but each one is extremely heavy, weighing several hundred pounds each.

\subsection{15. The Lost Forge}
\begin{DndReadAloud}{}
This large chamber with a vaulted ceiling is lit by the Lost Forge. A smelting area, a massive dark-iron anvil, and a huge worktable dominate much of the room. Atop the tables are immense blacksmithing tools that look as if made for the hands of giants.

\end{DndReadAloud}

\subsubsection{Creatures}
Kadal, a \textbf{thursir armorer}, is currently working angrily at the forge after a recent argument with his sister. As the PCs approach, four small creatures emerge from the forge—these four \textbf{magma mephit} follow Kadal and aid him with his work. Kadal stops working when he discovers intruders in the forge and immediately grabs his \textit{battleaxe} and \textit{shield} to deal with them.

\subsubsection{Power of the Lost Forge}
The magical forge emits an aura in the chamber that powers the heating of metal, which effects the \textbf{magma mephit} by replacing their Innate Spellcasting ability from (1/Day) to (Recharge 5\textendash 6).

\subsubsection{Treasure}
Kadal's \textit{battleaxe} and \textit{shield}, which he crafted himself, are a vicious axe and a \textit{+1 shield}. His armor is non-magical, but crafted with his expert smithing hand.

\section{Concluding the Adventure}
With the twin warlords defeated, any remaining followers flee into the mountains to find other employment. If Ruby Runemaker survives being prisoner, she offers a favor to the PCs if they should ever need it, along with a barrel of their finest dwarven ale, as thanks. The local villages see the PCs as their saviors, and they are offered free lodging and meals whenever they travel through.

After seeing that the magic of The Lost Forge has been rekindled, Ruby Runemaker brings word to her people, and the dwarves begin the process of rebuilding the walls to once again use the forge for making armor and weapons for the defense of the surrounding lands.

The thursir giant conclave where the twins came from hear of their defeat will someday come seeking vengeance for their fallen kin.

\chapter{The Monstrous Machinations of Margrave Millard}

\begin{DndTable}{c}

\textbf{An adventure for four to five characters of 7th level}\\
\end{DndTable}

\section{Adventure Background}
\DndDropCapLine{M}{}argrave Millard always had a fascination with monsters, even as a child. This was a harmless fascination until he inherited his estate and met Researcher Voss. The pair shared an interest, and while Millard loved the power and primal strength of monstrosities, Voss loved the act of creation, idolizing the wizards who created creatures like the \textbf{owlbear}. The two believed each was capable of achieving their goals—that Voss would create great and terrible creatures to unleash upon the world and that Millard would \textit{become} one of the monsters he so idolized.


Eventually, the pair escalated from plotting and research to experimentation. Disfavored servants and travelers began to disappear, rare and exotic monsters were imported, and Millard's mansion was converted into a laboratory so that the pair might create their monsters and unleash them upon unwitting victims to see how effectively they fought. As they ran out of convenient victims, they began to branch out, kidnapping merchant caravans, passing mercenaries, and unruly peasants—flinging them into an arena with their newest creations.

This was bad enough, but in the last month things have escalated. Margrave Millard is close to a breakthrough. He just needs a worthy foe to test himself against. He has begun to act openly, simply storming villages and towns he is lord of and kidnapping dozens to test against himself and his beast. The people don't know what's happening, why their lord has begun to terrorize them, what he plans in his isolated manse, or what his ultimate goal is.

\section{Adventure Hook}
This adventure can take place in any rural area, where news travels slowly and a powerful noble can get away with abusing his own populace. The PCs are hired in the small mining town of Iron Hill, where the local lord has been terrorizing his populace for months. Their contact is Jasmine, a local miner who has begged their urgent assistance saving the townsfolk from the local lord.

Jasmine is a respected mining forewoman (N human \textbf{scout}) in Iron Hill. She has sun-damaged skin, the strong build of a miner, and has lived in Iron Hill her whole life. She's gruff, serious, and self-assured but knows that she's out of her league with regard to the margrave's recent rampage.

\subsubsection{The Job}
Margrave Millard has been kidnapping citizens for unknown experiments for weeks, and no one he's taken has emerged from his secluded manor. Jasmine wants the PCs to stop the experiments and kill Millard for his crimes.

\subsubsection{Reward}
400 gp for each PC, as well as everything they can carry from the margrave's home.

\subsubsection{Additional Information}
Jasmine also shares the following:


\begin{itemize}
\item Margrave Millard lives in an isolated manor in the center of the woods, only an hour's ride from Iron Hill.
\item Millard has always been an odd man with an obsession with magical creatures, but was largely considered an acceptably neglectful lord until the experiments started, as he'd ignore his subjects and occasionally forget to collect taxes.
\item Millard has at least a dozen servants—thugs of various sorts—and a trained owlbear.
\item Before the kidnappings started, a shipment of a dozen or more caged animals was delivered to his manor.
\end{itemize}

\section{Millard's Manse}
Millard's Manse is a squat, overgrown gothic manor surrounded by looming willows and covered in creeping vines. Jasmine's directions bring the PCs straight to it without incident.

\subsubsection{Doors}
In addition to a large front door, there are 15-foot-wide barndoors on the eastern and western sides of the house, toward the back of the building. Unless otherwise noted, all doors—including the front door—are wooden and unlocked. The barndoors are both barred from the inside, requiring a successful DC 18 Strength (Athletics) check to break them open.

\subsubsection{Portcullis}
There are two heavy iron portcullis in 5. Each has AC 19, 27 hit points, and a damage threshold of 10, but the margrave gladly orders the gate opened if the PCs wish to enter the arena.

\subsubsection{Illumination}
The margrave has allowed his manse to fall into disrepair. Little light filters through barred windows and many rooms are lit only by sparse candles, providing dim light throughout the manse.

\subsubsection{Margrave Millard}
Margrave Millard (NE human \textbf{druid}) is a tall, pale man in his early 30s. He dresses in baggy formalwear that hangs off his frame, giving him the impression of being far scrawnier than he actually is. He's always been obsessed with physical strength and the freedom and respect he believes it can provide. He also really likes to hear himself talk, which is the foundation of his friendship with the stubbornly quiet Researcher Voss (NE human \textbf{priest}).

While Millard is well-educated and is keen to justify himself with philosophical diatribes about the fundamental meaninglessness of life and the purity of a natural life, he's utterly uninterested in what others have to say. Allies have their words interpreted as supporting him, and everyone else is casually dismissed as worthless.

Millard can speak to anyone in his manor, no matter what room they're in, via the enchanted pyramidal stone in 9; it's one-way communication only, so he can't hear any reply. The margrave uses the opportunity to monologue at his guests while they explore the manor. He can also see what his guests are doing, but only in 1, via the enchanted spectacles in 9.

An unlabelled version of the map can be found here.

\subsection{1. Entryway}
\begin{DndReadAloud}{}
This extravagant foyer may once have been luxurious, but now it's best described as eerie. Half the candles on the chandelier are unlit, casting the room in dim, unpleasant light, while the once-gleaming hardwood floors are now marred with overlapping scratch and claw marks. Two threadbare chairs (with their cushions missing) flank the entryway. Directly ahead there's a door to the north with a large, sneering face carved into the wall above it.

\end{DndReadAloud}

\subsubsection{Creatures}
Margrave Millard watches the entrance hall, using a pair of enchanted spectacles (see 9) to look through the eyes of the sneering face above the door. When he spots the PCs, he introduces himself, his voice reverberating all around the PCs, as if coming from the house itself. He thanks them for volunteering as his final test subjects, then urges them to step through the door to the north, so that the demonstration can begin. He calls them cowards if they go through any other door.

\subsection{2. West Barracks}
\begin{DndReadAloud}{}
This plain room hosts three bunks, each with a pair of small wooden chests at its foot. A polished wooden table sits in the center of the room, with several chairs and stools arranged around it. The room is better maintained than the rest of the manse—clean, comfortable, and free of damage.

\end{DndReadAloud}

\subsubsection{Creatures}
Three \textbf{thug} and three \textbf{veteran} are relaxing in this room. They are dedicated servants of the margrave and attack intruders, fighting to the death. However, if the margrave has died in the arena, they flee the manor.

\subsubsection{Treasure}
Most of the chests contain mundane personal effects, but inside one is a \textit{bag of holding} containing 15 human left ears and 2 lbs. of beef jerky.

\subsection{3. East Barracks}
\begin{DndReadAloud}{}
This ostentatiously painted room hosts three bunks, each with a pair of small wooden chests at its foot. The floor has been covered in wolf pelts that have suffered the wear and tear of having many, many boots trod upon them. Standing in the center of the room is a set of table and chairs and in the northeast corner is a writing desk.

\end{DndReadAloud}

\subsubsection{Creatures}
A \textbf{thug}, two \textbf{bandit captain}, and two \textbf{veteran} are eating a meal in this room. They are dedicated servants of the margrave and attack intruders, fighting to the death. However, if the margrave has died in the arena, they flee the manor.

\subsubsection{Treasure}
Most of the chests contain mundane personal effects, but one contains 1,500 sp and a note detailing payroll for the Manse's guards over the next month.

\subsection{4. Small Creature Holding Pens}
\begin{DndReadAloud}{}
This long, narrow room expands into a small alcove at its northern end, where you can see shreds of half-eaten flesh and fragments of bone scattered. Wooden cages are pressed against the walls here and wherever else there is space, and the area reeks of rotting meat and the acrid stench of a spilled alchemical reagents. This looks as if it might have once have been a kitchen, as you see a chimney and oven against the outer wall, next to it a water basin, and to the northeast a counter—now dedicated to what must be monster food—that may once have been used for food preparation. A barndoor on the western wall, which is barred on this side, appears to lead outside the manse.

\end{DndReadAloud}

Behind the northernmost door leading into 5 is a heavy iron portcullis.

\subsubsection{Creatures}
If they have not yet been released into the arena (5), there are four \textbf{capybear}), four \textbf{razorfeather raptor}, and an animal handler (a \textbf{thug}) in this area. If the PCs are spotted, the thug orders the animals to attack.

If the creatures are in this area, add the following read aloud text:

\begin{DndReadAloud}{}
Near one of the cages, an animal handler looks to be training—or perhaps just playing with—a feathered bipedal dinosaur taller than he is.

\end{DndReadAloud}

\subsubsection{Cages}
Several of the cages are more like dog crates—they're left open, and the creatures just have their own space that they can call their own. Both the \textbf{capybear} and \textbf{razorfeather raptor} are quite intelligent and loyal and have been well-trained.

Smaller crates have typical livestock—chickens and the like—for feeding the margrave's creatures.

\subsection{5. Arena}
\begin{DndReadAloud}{}
This ballroom must once have been grand, its chandeliers glistening from mage-light, its polished floor made for dancing, the walls austere yet elegant with their marble decorations. Now, the marble is bloodstained, shattered furniture is scattered across the floor, and the stench of death suffuses the air. Heavy iron portcullis stand before the two northernmost doors on the east and west walls.

\end{DndReadAloud}

\subsubsection{Creatures}
If the PCs entered from 1, the door behind them closes and magically locks after they enter (via the \textit{arcane lock} spell). The margrave thanks the PCs for entering his arena, then unleashes several waves of monstrosities upon them, grandiosely introducing each and explaining how it will tear them apart. The margrave grants the PCs two rounds of rest between attacks.

\textbf{First Wave:} Four \textbf{capybear} (from 4), one \textbf{dire owlbear} from 7.

\textbf{Second Wave:} Four \textbf{razorfeather raptor} from 4.

\textbf{Third Wave:} Three \textbf{grolar bear} from 7.

\textbf{Final Wave:} Margrave Millard (NE human \textbf{druid}) and Researcher Voss (NE human \textbf{priest}) leave 6 and one of them enters the arena through each portcullis. Upon arriving, the margrave gives the PCs an earnest congratulations, then announces the perfection of his experiments and transforms into a \textbf{grolar bear alpha}. He and Voss then attack; both fight to the death.

The \textbf{thug} in Areas 4 and 7 lead the animals in their charge to the northern entrance of the arena, open the doors, then lift the portcullis and release Millard's beasts and monstrosities inside.

If the margrave is defeated, any of his remaining servants flee.

\subsection{6. Observation Room}
\begin{DndReadAloud}{}
The southern wall of this room is an enormous one-way mirror looking into the bloodstained arena. Arrayed before the mirror are a series of comfortable chairs with small tables between them, clearly seating for guests to watch whatever horrors are to be put on display. A door lies in the center of the northern wall; flanking it are a pair of tables that hold the accoutrement for drinks and hors-d'oeuvres. Towering portraits of sharp-faced men and women line the walls, and a one of the small tables has atop it several empty glasses that have not yet been cleared.

\end{DndReadAloud}

While the PCs fight in the arena, Millard and Voss watch them from these seats.

The mirror is magically protected by a persistent \textit{wall of force} effect.

\subsection{7. Large Creature Holding Pens}
Behind the northernmost door leading into 5 is a heavy iron portcullis. Three very large—and one enormous—animal cages occupy the northern end of this area.

\subsubsection{Creatures}
If the creatures have not yet been released into the arena (5), three \textbf{grolar bear}, a \textbf{dire owlbear}, and an animal handler (a \textbf{thug}) are in this area. If the PCs are spotted, the handler orders the animals to attack.

\begin{DndReadAloud}{}
An enormous cage with an ferocious-looking—and large—owlbear inside is pressed against the northern wall of this sprawling chamber, while three smaller—but still quite large—cages pushed against other walls contain some kind of snarling bear. An animal handler stands next to the southernmost bear cage and pets the creature through the bars.

\end{DndReadAloud}

\subsubsection{Tactics}
On her first turn, the handler opens the first \textbf{grolar bear} cage, then rushes to the next and opens that one as well. On her next turn, she releases the \textbf{dire owlbear}. She then goes to release the final \textbf{grolar bear}, but the \textbf{dire owlbear}—perhaps enraged by the scent of blood in the air—turns on her, and kills her with a swipe of its mighty claws. The third \textbf{grolar bear} therefore remains caged. Before her untimely demise, upon releasing each of the beasts, the handler orders each one to attack.

If possible, the two \textbf{grolar bear} move to clog up the chokepoint in the hallway to allow the handler to release the \textbf{dire owlbear}; if she's killed by the PCs before she can release the \textbf{dire owlbear}, on its next turn it breaks free of its cage.

\subsection{8. Laboratory}
\begin{DndReadAloud}{}
Tables and alchemical equipment and cauldrons fill this room. The sound of slowly roiling liquid is omnipresent, as is the acrid scent of strange reagents burning.

\end{DndReadAloud}

\subsubsection{Creatures}
Two \textbf{slithy tove} have finished gestating in one of the cauldrons. If the PCs start looting the laboratory, spend more than 1 minute inside, or discover the toves, the creatures crawl out of their cauldrons and attack.

\subsubsection{Treasure}
In this room are two sets of \textit{alchemist's supplies}, 4 pieces of fine glassware (worth 100 gp total) and 15 flasks of alchemist's fire.

\subsection{9. The Margrave's Chambers}
\begin{DndReadAloud}{}
Bookcases line the walls of this chamber, and a formal table stands in the middle of the room with paperwork strewn across it. To the east and west of the table are a pair of comfortable-looking chairs, with smaller tables beside them, each of which hold a stack of books. On the central table is a strange, pyramidal stone engraved with runes, with a pair of large spectacles are chained to it.

\end{DndReadAloud}

\subsubsection{Creatures}
If they were not defeated elsewhere, Margrave Millard (NE human \textbf{druid}) and Researcher Voss (NE human \textbf{priest}) wait in this room. Once the PCs are spotted, Millard turns into a \textbf{grolar bear alpha} and attacks.

\subsubsection{Tactics}
Researcher Voss casts \textit{spirit guardians} to protect themself and Millard and heals the margrave if he is reduced to fewer than 15 hit points. Margrave Millard loses all self-control when he transforms and attacks whatever PC is closest at any given time. Both fight to the death.

\subsubsection{Treasure}
The spectacles allow the wearer to spend a bonus action to gaze through the eyes of the carving in 1. Any creature that speaks into the pyramidal stone broadcasts their voice to the entire manse. The bookshelves contain books on biology and monster ecology (worth 700 gp total to the right buyer). On the central table is a scroll of \textit{polymorph} and a small box that contains 4 diamonds (worth 600 gp total).

\section{Concluding the Adventure}
The adventure concludes when the PCs bring an end to Margrave Millard's reign of terror.

If Millard or Voss are left alive, they say whatever they need to in order to convince the PCs that they'll cease their evil ways, but they absolutely will not. Neither is motivated by revenge and so are not likely to seek out the PCs again, but they could become a problem that the PCs have to deal with again in the future.

If both Millard and Voss are killed, but any of the creatures are left alive, the PCs must decide what do with them. The \textbf{razorfeather raptor} and \textbf{dire owlbear} are monstrosities too dangerous to release into the wild, thought the \textbf{razorfeather raptor} is intelligent enough (and has been around humans enough) that it's possible clever PCs—likely with the aid of magic—could convince the creatures to travel somewhere remote and not prey on sentient creatures. Under normal circumstances, the \textbf{grolar bear} and \textbf{capybear} could potentially be released into the wild, but the margrave has brought out the most vicious instincts in them and thus none are likely to make good neighbors.

\chapter{The Green Sanctum}

\begin{DndTable}{c}

\textbf{An adventure for four to five characters of 7th level}\\
\end{DndTable}

\section{Adventure Background}
\DndDropCapLine{T}{}he halfling Zelen Barishka adventured during most of her early adulthood, growing in power as a druid as she traveled the world and fought monsters. Her true passion involved the study of plants, and she had occasion to note unusual species thriving in hidden locations or driven to near extinction because of the danger they posed to nearby settlements. She cataloged these plants' locations and vowed to return to retrieve, rescue, and study them.


When she retired, she used her accumulated wealth to build a hidden sanctuary where she could store and examine these notable plants. She also procured a strange golem capable of gathering and safely storing her prizes. Still afflicted with some wanderlust, Zelen also journeyed to recover plants herself. On her final trip, she fell to her death attempting to collect a cliff-dwelling shrub. Meanwhile, her golem assistant arrived from its latest assignment to obtain the fruit from a plant the halfling knows as an alaguara and now awaits Zelen's return so it can release its bounty.

\section{Adventure Hook}
The following can used to start the PCs on this adventure:


\begin{itemize}
\item \textbf{\textit{Botanist Check-in.}} One of Zelen's botanist colleagues has been expecting to receive an alaguara seed from the druid to assist in the plant's study, and its delivery is now considerably past due. The botanist offers the PCs a reward to travel to Zelen's sanctum to check in with her.
\item \textbf{\textit{Fashioning a Cure.}} The PCs may require the alaguara plant held by the golem to fashion a cure for a disease or other malady afflicting one of their number, an important NPC, or an entire settlement.
\end{itemize}

\section{The Green Sanctum}
\begin{DndReadAloud}{}
Ahead, the trees thin, providing better access to the river flowing through the forest and revealing a stone structure nestled on an island about thirty feet from the bank. Two clear glass domes top the vine-covered building.

\end{DndReadAloud}

Granite forms the basis of the island supporting the stone building, allowing Zelen to comfortably commission the material used for the building's construction without fear of it sinking. Narrow slits in the walls allow ample light to penetrate the building while keeping out all but the smallest creatures. During the day, bright light fills the building except where noted, but, at night, darkness engulfs the building, as there are no continuous light sources. Contrasting with the stone walls, wooden doors allow access to the sanctum's inner chambers.

\subsubsection{Creatures}
Twelve \textbf{dust mephit} serve Zelen by carrying out mundane tasks and warding off intruders while she is away. When the PCs enter any room except 16, they have a chance to encounter the creatures. Roll a d10; on a roll of 1, they encounter two \textbf{dust mephit}, or on a roll of 2–4, they encounter one. Each \textbf{dust mephit} can use its action to activate the entangling roots growing beneath the sanctum. These roots can only be activated once per day in each area and reset at dawn.

\subsubsection{Entangling Roots}
Thanks to the druid's coaxing, dozens of roots grow from the riverbed and symbiotically intertwine with the granite to create one of the sanctum's defense measures. Roots snake through holes in the floor in a 15-foot radius centered on a designated point for 1 minute. For the duration, that area is treated as difficult terrain for creatures other than plants. In addition, each creature in that area must succeed on a DC 12 Strength saving throw or become restrained. A creature can use its action to make a DC 12 Strength check, freeing itself or another entangled creature within reach on a success.

\subsubsection{Fire Suppression System}
Another magical defense system counters the threat fire poses to the sanctum's rare specimens. After a creature uses magic or an effect that deals fire damage, the sanctum—as if it were a creature using its reaction—douses any fire caused by the effect, and the creature responsible, using the following attack:

\subparagraph{Freezing Spray}
\textsl{Ranged Spell Attack:} +4 to hit, one target. \textsl{Hit:} 4 (1d8) cold damage and douses any nonmagical flames on the target.

The fire suppression system can also target an ongoing spell that deals fire damage, such as \textit{flaming sphere}, as if it cast \textit{dispel magic} (+4 bonus on the ability check). The sanctum gets its reaction back on initiative count 20.

An unlabelled version of the map can be found here.

\subsection{1. Beach Landing}
\begin{DndReadAloud}{}
The wide river separating the banks from this modest island in its center proves easy to navigate. Thorny grasses cover much of the sandy soil surrounding the building.

\end{DndReadAloud}

The island stands 30 feet from the nearest bank and the water reaches a maximum depth of 15 feet halfway between the bank and the island. Near the island, the depth varies from 5 to 7 feet. Zelen left the native grasses alone apart from those she cleared prior to the sanctum's construction.

\subsubsection{Creatures}
Two \textbf{vine-covered owlbear}, trained and altered by Zelen, patrol the beach. Without hearing the proper series of whistles known only to the druid, the \textbf{vine-covered owlbear} attack any creatures that set foot on the beach. They do not pursue creatures that move more than 10 feet from the beach.

\subsubsection{Creature Change}
Once per short or long rest, the \textbf{owlbear} can use their action to produce an effect similar to the \textit{entangle} spell (spell save DC 13). It has a range of only 20 feet, but requires no concentration. The vines covering the creatures shoot out to the point within range to create the effect.

\subsection{2. Sanctum Entrance}
\begin{DndReadAloud}{}
As you enter the sanctum, you see a long hallway stretching ahead north; in the immediate vicinity, to the east and west are closed wooden doors. The flooring is cool to the touch and looks to be natural stone, polished mostly smooth.

\end{DndReadAloud}

If the PCs arrive during the day, the light that shines through the dome in 6 provides enough illumination through the Observation Pane (3) to be notable even from the entrance.

\subsection{3. Observation Pane}
\begin{DndReadAloud}{}
The hallway ends at a large, frigid, ice-rimed glass pane through which one can view the plants contained within a glass dome. At the southern end of the hall, closed doors lead to the east and west.

\end{DndReadAloud}

\subsubsection{Glass Pane}
The glass pane has AC 15, 30 hit points, a damage threshold of 10, and immunity to piercing, slashing, poison, and psychic damage. Destroying the pane attracts the attention of 1d4 + 1 \textbf{dust mephit}, which arrive in 2 rounds to attack the offenders.

\subsection{4. Zelen's Living Quarters}
\begin{DndReadAloud}{}
This spare living space contains a neatly made bed, a small wooden table with numerous ring stains from mugs and tankards, a writing desk with an open leather-bound book, and a wardrobe.

\end{DndReadAloud}

The desk is clear of items except for Zelen's journal open to its final entry, dated three weeks ago; the entry reads: "I've sent my assistant to find one of the last known specimens of the alaguara, which has faced eradication from people unkindly disposed to its invasive qualities. While he is out on his mission, I shall undertake my own quest to retrieve a rare flowering plant that climbs the Blackrock Cliffs." The journal spans a month prior to the last date and features discussions about a species of orchid that spits poisonous thorns and a plan to seek out the flower later.

\subsubsection{Treasure}
A set of \textit{goggles of night} hang from a hook behind a heavy cloak in the wardrobe.

\subsection{5. Plant-Lined Hallway}
\begin{DndReadAloud}{}
Numerous flowering shrubs and bushes thrive in this hallway and fill it with sweet fragrances. The temperature drops significantly as the hallway proceeds north.

\end{DndReadAloud}

As the corridor turns east, a set of double doors on the eastern wall stand open (to 6); opposite that, on the western wall, is a closed single door (to 7).

\subsubsection{Creatures}
Two \textbf{alpine creeper}, which typically dwell in the arctic sanctum, lurk in the colder northern section of this hall, just south of the doors leading to 6. The plants ignore the \textbf{dust mephit}, but they attack anyone else, fighting to the death.

\subsection{6. Arctic Sanctum}
\begin{DndReadAloud}{}
A blast of cold hits you as you enter this room. Ice coats most of this domed, circular room's surfaces, but the plants within clearly thrive here. The frost on the dome above dims the light shining through it.

\end{DndReadAloud}

Permanent magic keeps this chamber's temperature at just below the freezing point. A successful \textit{dispel magic} (DC 17) ends the magic; if dispelled, the room warms to match the ambient temperature after four hours, causing most of the plants located here to wither and die.

\subsubsection{Creatures}
An \textbf{ice willow} stands along the room's southeastern arc. It waits until prey draws close enough for it to unleash its spray of ice shards with greatest effect. The willow fights until reduced to 20 hit points, at which point it attempts to retreat as far as it can within 6 and cowers, saying in Sylvan "was scared", "no kill" or the like. If the PCs seem determined to kill it, however, it continues fighting to the death.

\subsection{7. Storage Shed}
\begin{DndReadAloud}{}
Gardening tools—trowels, shovels, rakes, and the like—haphazardly fill this shed and clutter the workbench along the northwestern wall. Though disorganized and well-used, the tools show no sign of disrepair and are free of dirt and rust.

\end{DndReadAloud}

\subsubsection{Treasure}
A successful DC 12 Intelligence (Investigation) check finds 4 flasks packed away under the workbench. Two of the flasks contain \textit{plant eradicator} and two contain \textit{plant revivifier}.

\subsection{8. Tucked-Away Shrine}
\begin{DndReadAloud}{}
The corridor ends at a wall ahead, at the foot of which is a shrine enshrouded in ferns and flowers.

\end{DndReadAloud}

The shrine is anointed to various nature deities to whom Zelen gave prayers while she adventured.

\subsubsection{Secret Door}
A successful DC 17 Wisdom (Perception) check discovers a faint crack in the wall, revealing a door that pushes inward to 15.

\subsection{9. Vine-Covered Trellises}
\begin{DndReadAloud}{}
Wooden frames covered with climbers line the north and south walls of this short hallway.

\end{DndReadAloud}

The vines react to vibrations, including footsteps, by shuddering in response, though they pose no danger.

\subsection{10. Meeting Room}
\begin{DndReadAloud}{}
A grand oak table surrounded by mahogany chairs fills most of this square chamber.

\end{DndReadAloud}

A single sheet of paper as well as a cup that once contained coffee are on the table; the paper seems to be notes from a meeting, from one druid to another. The room is empty apart from the table and chairs.

\subsection{11. Moss-Coated Hallway}
\begin{DndReadAloud}{}
Different strains of moss and lichen fill every inch of this hallway, including the ceiling and floor, deadening sound within it. The plants glow with blue and green iridescence when blocked from sunlight.

\end{DndReadAloud}

These plants are all just common mosses and lichen.

\subsection{12. Dark Room}
\begin{DndReadAloud}{}
A riot of fungi covers a table along the southern wall of this lightless chamber. Shelf fungus covers most of the northern wall.

\end{DndReadAloud}

\subsubsection{Creatures}
A \textbf{shrieker} lurks among the other fungus here. If it raises its alarm, it draws 1d4 + 1 \textbf{dust mephit} to the chamber 1 round later.

\subsection{13. Storage Room}
\begin{DndReadAloud}{}
Neatly stacked piles of books and papers cover the floor.

\end{DndReadAloud}

The material in this chamber contains mundane details about nature and plants.

\subsubsection{Secret Door}
A successful DC 17 Wisdom (Perception) check detects the seams of the door.

\subsection{14. Research Vault}
\begin{DndReadAloud}{}
The secret door opens to reveal a dark, cool chamber. Against the north wall stands a table, against the west wall a bookcase. Numerous tomes fill the bookshelf, and the desk has an array of journals and papers arrayed across it.

\end{DndReadAloud}

Zelen stored important correspondence, research tomes, and her own journal in this secret room.

\subsubsection{Treasure}
Zelen's journals contain an array of information useful for those who study plants (worth 150 gp total to the right buyer). A creature who spends a month studying the research gains proficiency in Intelligence (Nature) checks; if they are already proficient, their proficiency bonus is doubled for any Intelligence (Nature) check.

The journals also contain instructions for operating the \textbf{ice golem} in 16, including the command words necessary to activate and deactivate the construct. A PC can locate that information without spending days reading through all of the journals with a successful DC 15 Intelligence (Investigation or Nature) check; a druid makes this check with advantage.

\subsection{15. Laboratory}
\begin{DndReadAloud}{}
This oddly shaped chamber holds a handful of tables laden with beakers, mortars and pestles, and tools used for the fine manipulation of herbs and reagents. A door provides the only entryway to a circular chamber contained within this laboratory.

\end{DndReadAloud}

\subsubsection{Lab Purpose}
If a PC spends 10 minutes or more examining the notes and equipment in the lab, with a successful DC 18 Intelligence (Investigation or Nature) check, they determine that Zelen used this lab to study harmful and dangerous plants to make them more harmonious with their surroundings, including the alaguara fruit. She wished to preserve its healing properties in a hybrid that wouldn't wipe out competing plants.

\subsubsection{Locked and Trapped Door}
A successful DC 15 Dexterity check using \textit{thieves' tools} picks the lock. Electricity charges the door when opened and every creature within 10 feet of the door must make a DC 17 Dexterity saving throw, taking 21 (6d6) lightning damage on a failed save or half as much damage on a successful one.

\subsubsection{Prevention}
A successful DC 17 Intelligence (Investigation) check while inspecting the door finds nearly imperceptible runes carved in the wood. A spell or other effect that can sense the presence of magic, such as \textit{detect magic}, reveals an aura of evocation magic on the door. A successful \textit{dispel magic} (DC 14) cast on the door deactivates the trap. The trap can be disarmed with a successful DC 17 Dexterity check using \textit{thieves' tools}. Failing this check by 5 or more discharges the trap and causes the PC to have disadvantage on the resulting saving throw.

\subsubsection{Treasure}
Three herbalism kits can be collected from the equipment on the various tables. There are also 3 potions of healing and a \textit{potion of greater healing}.

\subsection{16. Inner Sanctum}
\begin{DndReadAloud}{}
This domed chamber houses several strange trees and flowering plants, some of which seemingly move independently. A metal humanoid standing along the eastern wall seems incongruous to the greenery.

\end{DndReadAloud}

\subsubsection{Creatures}
A \textbf{shambling mound} guards this sanctum and attacks anyone other than Zelen who enters, fighting to the death.

Zelen's \textbf{ice golem} stands beside the eastern wall. It seeks only to protect the plant within its chest cavity, so it only attacks any creatures that move within its reach, targets it with ranged attacks, or uses spells or effects targeting it. If the proper command words found in Zelen's journals in 14 are spoken, they deactivate the \textbf{ice golem} and command it to release the alaguara fruit held in its chest cavity.

Attempting to forcibly remove the alaguara fruit causes the \textbf{ice golem} to go berserk. In this state, the \textbf{ice golem} attacks the nearest creature it can see, excluding the \textbf{shambling mound}, even pursuing targets outside this area. It remains in this state until it is destroyed or deactivated.

If the PCs successfully navigate the \textbf{ice golem} without having to fight it, but you would like there to be more of a challenge to this final encounter, you can add additional plant creatures to this combat—of which there are nearly two dozen in \textit{Tome of Beasts 3}.

\subsubsection{Treasure}
The \textbf{ice golem} holds the alaguara fruit in stasis inside its chest cavity. There are no other objects of value in this room, but the plants growing here are almost all extremely rare or unique specimens of considerable value to those searching for or studying such plants.

\section{Concluding the Adventure}
If the PCs retrieve the alaguara fruit, they receive a reward commensurate with the adventure hook you chose.

Another druid or scholar with an interest in Zelen's studies may eventually take up residence in the Green Sanctum and contact the PCs to share information about the location's denizens and traps.

If the PCs retrieved Zelen's journals, interested groups may contact them and ask about obtaining the journals and their secrets for a price.

The PCs may also receive a request to find Zelen's remains so she can receive a proper burial.

\chapter{The Salons of Mother Celeste}

\begin{DndTable}{c}

\textbf{An adventure for four to five characters of 8th level}\\
\end{DndTable}

\section{Adventure Background}
\DndDropCapLine{T}{}he fur trader \textit{Whippet} has been lost, presumably due to a non-functioning lighthouse and the foggy shoals outside of Wrackwater Cove (see "Incident at Wrackwater" adventure for additional background). Wrackwater locals have recently found confirmation of a nearby wreck location. Undoubtedly, it's \textit{Whippet}, but the ship's investors and underwriters must have incontrovertible proof. They're also eager to recover some specific items that will confirm, document, and perhaps mitigate their financial losses.


\section{Adventure Hook}
There have been many strange occurrences in the region recently, including the appearance of strange, dangerous creatures, and now the merchant schooner \textit{Whippet}, which didn't arrive in port when expected, is no longer merely "overdue" but officially "lost at sea."

The Wrackwater town council, with funding provided by a group of maritime investors and guarantors, are looking for the assistance of "heroes" to recover the lost vessel. Given other recent events, it's entirely too dangerous for a mere salvage crew (though they tried and no one could be convinced to take the job), clearly a job for heroes.

The following hooks can used to start the PCs on this adventure. The PCs:


\begin{itemize}
\item Already helped Wrackwater in the "Incident at Wrackwater" adventure and are recruited by their contacts from that previous adventure.
\item Are asked by one of their allies or patrons—who has a friend, relative, or colleague among the investors and/or town council—to travel to Wrackwater to investigate the strange goings-on there.
\item Are traveling in the vicinity and hear of the plight of the people of Wrackwater.
\end{itemize}

\subsubsection{The Job}
The investors wish to discuss enlisting the PCs for ship identity confirmation (by retrieving the ship's bell), insurance assessment (by retrieving the ship's manifest), and to report any other salvage opportunities at the recently confirmed wreck site. The town council is naturally most concerned with the PCs putting an end to the threat posed by the strange creatures that have been lurking about of late, but the investors are footing the bill and so their concerns are at the forefront of the negotiations.

If the PCs agree, the guarantors propose a total payment of 2,000 gp for the job. With the generous funding of the guarantors, the council offers transportation to and from the wreck site, with necessary diving and salvage equipment provided.

\section{The Good Barge Barnacle}
The PCs are brought to the wide, flat deck of the salvage barge Barnacle. This tattered barge sports a small pilot house, tool racks and cargo nets, and a small winch and a 150-foot boom. The PCs are greeted by Old Wil, \textit{Barnacle's} master and commander, and then are introduced to his crew: Daisy (a geriatric parrot) and Dog (a pet harbor seal).

Wil cheerfully sets a crate down in front of the PCs and says, "The fine folk of the town council says for me to pass this along to ye heroes, so I'm'a doin' it." Inside the crate is a \textit{sealskin coldsuit} (which grants the wearer a swimming speed of 20 feet), 3 glowmold globes, and a \textit{crowbar}. Wil says he'll wait topside with "emergency refreshments and moral support."

\begin{DndSidebar}[float=!b]{Some Nautical Terminology}
\textbf{Fore:} The front (bow) of the ship



\textbf{Aft:} The back (stern) of the ship



\textbf{Starboard:} The right side of the ship if facing the bow



\textbf{Port:} The left side of the ship if facing the bow



\textbf{Deck:} Floor; also used to refer to the levels of the ship—i.e., Deck 1, Deck 2, Upper Deck, etc.



\textbf{Overhead:} Ceiling



\textbf{Bulkhead:} Wall



\textbf{Bilge:} The lowest inner part of a ship, designed to collect excess water



\textbf{Binnacle:} A pedestal mounted at the helm that houses the ship's compass, made of wood and brass



\textbf{Capstan:} A type of large winch used to help the crew lift heavy weights and/or wind the ship's anchor



\textbf{Fo'c's'le:} Forecastle, the forward (main) deck of the ship (pronounced "FOKH-sil")



\textbf{Gangway Ladder:} Steep, ladder-like stairs



\textbf{Ship's Boat:} A smaller boat (typically a \textit{rowboat}) carried by a larger vessel, used to communicate with shore or perform other tasks the larger ship may be otherwise incapable of



\end{DndSidebar}

\section{The Wreck of the Whippet}
\begin{DndReadAloud}{}
You descend more than seventy feet, diving down through water that seems uncharacteristically warm, with a clear current that feels like it zigzags through these waters. Before long, the capsized vessel comes into view and you see a jagged, wide hull breach that exposes the expansive darkness of the ship's bilge.

\end{DndReadAloud}

The PCs may free-dive or can be lowered down to the wreck via Barnacle's boom winch.

\subsubsection{Capsized Vessel}
Since the ship is upside-down, any "floor" (called "deck" on ships) the map depicts is actually ceiling (called "overhead"). Likewise, all descriptions referring to "the deck" mean the surface above and "the overhead" refers to the surface below.

\subsubsection{No Bodies}
The shipwreck is curiously devoid of dead crew, except for Areas 12–14. \textit{Whippet} had a crew of 16.

\subsubsection{Hatches and Doors}
While opening a hatch or door on a normal ship would be a free action, opening them underwater requires an action. Any Dexterity (Stealth) check while opening a door or hatch is made with disadvantage due to the disturbance it causes in the water.

An unlabelled version of the map can be found here.

\subsection{1. Bilge Breach}
The hull breach is 10 feet long and ranges from 1 foot to 3 feet wide, large enough for Small and Medium creatures to squeeze through without too much effort. Looking inside, the bilge is 4 feet high and crisscrossed with ribbing, ship beams, and shoring. The bilge access hatch to/from the cargo hold (3) hangs open 15 feet aft of this hull breach.

\subsection{2. Forward Hull Breach}
A large breach pierces the hull here near the bottom of the bow; it's 2 feet by 7 feet and opens into 5. Inside, a vast amount of cargo is haphazardly piled up on Deck 2 as far as the eye can see.

\section{Deck 2: Dire Depths}
As a reminder, Deck 2 actually rests below the bilges and above Deck 1 due to the \textit{Whippet} being inverted.

\subsection{3. Main Cargo Hold, Aft}
\begin{DndReadAloud}{}
Huge piles of ruined goods and freight crowd the hold's aft end, some of which stretch nearly all the way from the overhead to the deck above.

\end{DndReadAloud}

\subsubsection{Hatch}
A portside hatch in the deck above leads to the bilge (1).

\subsubsection{Vision}
The towers of cargo here may impede a creature's line of sight and/or cover options.

\subsubsection{Creatures}
Two \textbf{dire lionfish} lurk here in the shadows, making their home here as if the piles of cargo were coral. The pair are viciously territorial and have claimed the entirety of Deck 2 as their domain. Unless the PCs succeed on a DC 19 group Dexterity (Stealth) check, the \textbf{dire lionfish} detect the PCs approaching and attempt to ambush them when they come within 15 feet. They fight to the death to protect their territory.

Once combat ensues, there is a 60\% chance of attracting 1d4 + 1 curious \textbf{bilge gremlin} from 9 who come through the cargo hatch (that connects 9 to 4) and join the fray. (The chance increases to 100\% if the PCs use a cacophonous effect like the spell \textit{thunderwave}.) The curious \textbf{bilge gremlin} arrive after 2 rounds and happily attack both the \textbf{dire lionfish} and the PCs.

\subsubsection{Treasure}
If the PCs wish, they can harvest the dire lionfish spines with a successful DC 18 Wisdom (Survival) check; on a failure of 5 or more, the PC attempting the task takes 2d6 poison damage. 30 spines are available to collect; the PC must repeat the check every 10 spines they harvest, and the process of removing them takes 5 minutes of excruciating, cautious work. If another PC wishes to help, the Wisdom (Survival) check can be made with advantage, but the helping PC must make a DC 16 Constitution saving throw or suffer 2d6 poison damage.

\begin{DndSidebar}[float=!b]{Bilge Gremlin Tactics}
\textbf{bilge gremlin} (and \textbf{bilge gremlin bosun}) are wily and slippery foes. They rely on their evasiveness to be effective, and so they take full advantage of their ability to cast \textit{misty step at will} nor \textbf{bilge gremlin bosuns} can do this.}, rarely ending a turn adjacent to an enemy if at all possible. Likewise, whenever they detect danger, they make frequent use of their at-will \textit{invisibility} nor \textbf{bilge gremlin bosuns} can do this.} to spring ambushes.



\end{DndSidebar}

\subsection{4. Main Cargo Hold, Forward}
\begin{DndReadAloud}{}
This area is mostly clear of cargo and detritus, though the mainmast bears a deep indentation as if something forcefully slammed into it. Amidship, you can see the outline of the cargo hatch in the overhead that goes to the next deck. A gaping hull breach in the starboard bow is visible from here.

\end{DndReadAloud}

The \textbf{dire lionfish} in 3 headbutted the mast as part of claiming its territory and/or letting out some aggression.

\subsubsection{Hatch}
The cargo hatch in the overhead (the floor, since the ship is upside-down) leads from Area 4 to 9.

\subsection{5. Forward Stores \& Hatchway}
\begin{DndReadAloud}{}
Several crates and barrels are strewn about this area, the wood on several of them shattered, spilling their contents. Some bits of grain float in the water here, remnants, perhaps, of one of these broken barrels or sacks, the rest already washed out to sea.

\end{DndReadAloud}

\subsubsection{Hatch}
A portside gangway ladder (staircase) leads to a hatch that opens onto Deck 1.

\subsubsection{Breach}
A hull breach (see 2) is located on the starboard bow here. On the edge of the breach, a large spine protrudes from the bulkhead (from the \textbf{dire lionfish} in 3).

\subsubsection{Treasure}
The \textit{dire lionfish spine} can be collected and used as a weapon.

\section{Deck 1: Battle of the Bilgies}
Nine \textbf{bilge gremlin} and a \textbf{bilge gremlin bosun} have taken ownership of Areas 6–9. If combat ensues in any of those locations, the gremlins call out to each other and then converge and attack as a group.

\subsection{6. Boatswain's Locker}
\begin{DndReadAloud}{}
Suspended from the capstan on the deck are great lengths of four-inch-thick rope dangling in the seawater and coiling in tangled snarls below. Yards and yards of tangled mooring and anchor lines likewise litter the overhead.

\end{DndReadAloud}

\subsubsection{Hatch}
A portside gangway ladder (staircase) leads to a hatch that opens onto the fo'c's'le, the forward part of the main deck.

\subsubsection{Rope}
Due to being extremely waterlogged, these hemp ropes are not floating freely as they would otherwise.

\subsubsection{Creatures}
Three \textbf{bilge gremlin} are napping in this room, hidden amongst the ropes. They wake when the hatch or door is opened, unless the creature opening the door or hatch succeeds on a DC 12 Dexterity (Stealth) check. If the PCs enter the room undetected, it requires a successful DC 16 Wisdom (Perception) check to notice the \textbf{bilge gremlin}.

\subsubsection{Tactics}
If the gremlins are alerted to the PCs' entrance, they cast \textit{invisibility on themselves} can't do this.} and wait for the PCs to enter the hallway; once inside, they sing out to their comrades in Areas 7, 8, and 9.

\subsection{7–8. Forward Crew Quarters}
\begin{DndReadAloud}{}
Hammocks and tie-downs float like ghosts between ship beams, articles of clothing bedding and common sailor's effects drift about the flooded compartment.

\end{DndReadAloud}

Both of these rooms are identical crew quarters for 6 individuals.

\subsubsection{Creatures}
One \textbf{bilge gremlin} sleeps in each of these rooms. If any combat ensues anywhere in Areas 6–9, they awaken and join the fray.

\subsection{9. Upper Cargo Hold}
\begin{DndReadAloud}{}
This cargo hold is in as dismal condition as you might expect. All manner of cargo has been cast about, making the space a chaotic mess.

\end{DndReadAloud}

\subsubsection{Creatures}
Four \textbf{bilge gremlin} are scavenging in this cluttered, warehouse-like space. The gremlins converge on the PCs once they become aware of them. If the PCs open the hatch above, the gremlins all cast \textit{invisibility on themselves} can't do this.}.

In the southwest corner, a large bell lies on its side with the top of a crate leaning against its mouth. Once the PCs are detected, a \textbf{bilge gremlin bosun} comes out of the bell that he has claimed as his home.

\subsubsection{Locked Door}
The door leading to 10 is locked, but can be picked with a successful DC 12 Dexterity check using \textit{thieves' tools}.

\subsubsection{Treasure}
Inside the ship's bell is a bottle of vintage brandy (worth 150 gp), a sextant, a leather pouch containing several (human) finger bones, and a brass cabin key (which opens the door to 10).

\subsubsection{The Ship's Bell}
The \textbf{bilge gremlin bosun} "home" is the \textit{Whippet's} ship's bell. Made out of solid bronze, it's 18 inches tall, with a mouth diameter of 20 inches. Salvaging the bell, which weighs 225 lbs., requires \textit{Barnacle's} boom winch, clever planning and teamwork, and/or magic.

\section{Deck 1: Wrath of the Wrackwraiths}
In Areas 12–14 lurk five \textbf{wrackwraith}. If combat ensues in any of those areas, all of the \textbf{wrackwraith} converge and attack as a group.

When the PCs first encounter one of the \textbf{wrackwraith}, read or paraphrase the following:

\begin{DndReadAloud}{}
Out of the corner of your eye, you detect motion, and you turn to see, rising in the water, a humanoid shape formed by torn fabric and sundered wood, shorn hair and shattered bone. Then you hear a sound, almost as if from inside your own head—a sound reminiscent of the last gasp of a person drowning.

\end{DndReadAloud}

\subsubsection{Beach Bound}
The "beach" these \textbf{wrackwraith} are bound to is the seafloor upon which the \textit{Whippet} wrecked. Their spirits are unable to leave the ship, where the callous \textbf{bilge gremlin} that caused the wreck laughed as they watched them drown. The \textbf{wrackwraith} won't enter 9 or beyond.

\subsection{10–11. Galley \& Kitchen}
\begin{DndReadAloud}{}
This area was clearly the ship's galley. Chairs, cutlery, and dishes are scattered about; the tables the crew ate at hanging above you, attached to the deck. Through the open door to port, you can see a smaller room that is clearly the galley's kitchen.

\end{DndReadAloud}

Inside the galley and the kitchen, the PCs find about what you'd expect in such places, nothing more.

\subsection{12–13. Passenger Staterooms}
\begin{DndReadAloud}{}
This cabin was once nicely appointed, seemingly for passengers. Naturally, now it's a total ruin. A bloated corpse floats in one of the corners.

\end{DndReadAloud}

Both of these rooms are essentially identical quarters for 1–2 individuals.

\subsubsection{Door Ajar}
The door to this room is ajar.

\subsubsection{Creatures}
Each room contains one \textbf{wrackwraith}.

\subsubsection{Treasure}
One of the corpses—a woman wearing a fashionable dress—is wearing an onyx ring (worth 90 gp) bearing the inscription "Wrackwater 5 Annual Springfair Queen."

\subsection{14. Aft Crew Quarters}
\begin{DndReadAloud}{}
This room was clearly once used as crew quarters; you can see the torn remnants of numerous hammocks and destroyed footlockers, as well as clothes and other personal effects. Now, however, the room is clearly used for storing bodies. The bloated, floating corpses of the entire crew seem to have been shoved into this room, with no respect for the dead, or any sense of propriety whatsoever. But the crew aren't the only dead present: Floating nearby are three dead gremlins.

\end{DndReadAloud}

There were once nine hammocks hanging in these quarters. Now most of them are on the ground . . . or being used for nefarious purposes by the current occupants.

\subsubsection{Door Ajar}
The door to this room is ajar.

\subsubsection{Creatures}
Three \textbf{wrackwraith} lurk in this room.

\subsubsection{Dead Gremlins}
The \textbf{bilge gremlin bosun} in 9 ordered this area off-limits after the three gremlins were killed. If any living gremlins are questioned, they don't know what's in Area 14, but they know they don't want to mess with it after Scobbe, Votrag, and Blarck didn't come back.

\subsubsection{Bodies}
Further inspection of the area reveals that all of the crew's fingers have been severed and are missing.

\section{Upper Deck}
As a reminder, the "upper" deck is actually the lowest level of the \textit{Whippet}, due to it being inverted.

\subsection{15. Bosun's Hatch}
A hatch on the deck opens from Area 15 to 6. A portside gangway ladder (staircase) is on the other side of the hatch.

\subsection{16. Main Deck}
\begin{DndReadAloud}{}
Looking up, you can see the main deck above you, along with half of the splintered mainmast—the other half sundered on the seafloor. Broken gear, tangled rigging, and shattered spars and yards are littered about this area. In the center of the deck above, you can see a closed cargo hatch. Among the debris on the seafloor, the overturned ship's boat lies capsized, partially embedded in the sand. Of more immediate concern, however, are the figures before you. Festooned in nudibranchs, starfish, and anemones, an amphibious troll leads two large, clawed and finned humanoid creatures with merfolk-like tails toward the sterncastle.

\end{DndReadAloud}

\subsubsection{Ship's Boat}
The ship's boat is approximately 12 feet by 8 feet; large enough to carry six Medium-sized creatures

\subsubsection{Creatures}
A \textbf{breakwater troll} leads two \textbf{merrow} toward the sterncastle (18).

\subsubsection{Treasure}
In the ship's boat is an airtight flask containing 2 doses of a \textit{potion of superior healing}.

\subsection{17. Quarterdeck}
\begin{DndReadAloud}{}
Lengths of anchor shanks and chain lie on the seafloor nearby a barrel of harpoons lying on its side. Portside, the ship's binnacle—a wooden and brass pedestal about three-and-a-half feet tall—protrudes from the seabed, partially embedded in the sand, its brass and glass dome cracked and filled with water. Aftward, you see the inverted sterncastle, a large cabin that stands atop the aft end of the deck.

\end{DndReadAloud}

\subsubsection{The Compass Array}
The binnacle contains the ship's compass array. Salvaging the compass array, which weighs 275 lbs., requires \textit{Barnacle's} boom winch, clever planning and teamwork, and/or magic.

\subsubsection{Treasure}
The PCs can also collect 8 harpoons (use statistics for \textit{spear}).

\subsection{18. Sterncastle}
\begin{DndReadAloud}{}
Rainbow-hued colonies of anemones, sponges, and starfish drape the glistening walls and carpet the deck and overhead in this room that was once the captain's quarters. A grotesque humanoid figure floating in the center of the room—with blue, bloated, dead-looking skin, and hair resembling kelp—looks up at you as you enter. From the north, another figure comes around the half-wall, draped in seaweed and slime, encrusted with barnacles. This one is also humanoid, but has scaled skin and eel-like tail instead of legs. On her face is a terrifying smile, which falters into a momentary look of confusion upon seeing you, before recovering. "You are not who I was expecting. But welcome," she says, moving toward you and gesturing for you to enter. "Come in, come in."

\end{DndReadAloud}

\subsubsection{Magical Wards}
Area 18 can't be scryed upon, teleported into, or otherwise breached by any means short of a \textit{wish} spell. Only Mother Celeste can negate/bypass this effect.

\subsubsection{Creatures}
Mother Celeste, a \textbf{brine hag} and a \textbf{sea hag}, attack once the PCs enter.

\subsubsection{Magic Chest}
Against the aft wall is a magically locked chest (via the \textit{arcane lock} spell); it can be unlocked with Mother Celeste's key or picked with a successful DC 30 Dexterity check using \textit{thieves' tools}. Any inanimate object inside the chest is immune to the effects of water; if the object has been stored in the chest for at least 4 hours, if it is removed, it maintains this immunity for 24 hours. Both it, and the key (which Mother Celeste carries) radiate abjuration magic if \textit{detect magic} is cast.

\subsubsection{Treasure}
Mother Celeste wears an \textit{amulet of proof against detection and location} and, also around her neck, is a key on a chain.

Inside the chest: There is 308 sp, 221 gp; 4 scrolls of \textit{sending}; the \textit{Whippet's} ship's manifest; the ship's log, which contains sketches depicting a number of sunken shipwrecks arranged around \textit{Whippet} in such a way that makes it very closely resemble an underwater fortress; stacks of current nautical charts, including local tidal tables and vessel names—and their sizes and shipping schedules.

\section{Concluding the Adventure}
Returning topside safely, the PCs find Old Wil there waiting for them, as promised—with the promised refreshments. With the Barnacle's winch and Old Wil's know-how, it's a simple enough matter to recover the heavy salvage items.

If the PCs retrieve the ship's manifest and the ship's bell, they receive the agreed upon reward and the eternal thanks of the guarantors and investors. If the PCs also salvage anything else valuable from the wreck other than the manifest and bell, at GM discretion, the guarantors offer an additional reward for items such as the ship's compass and/or log.

If the PCs also put an end to the threat of Mother Celeste and her coterie, the town council is even more grateful, dubbing them "The Heroes of Wrackwater" and insisting on throwing a feast in their honor.

\chapter{Chaos at the Caldera}

\begin{DndTable}{c}

\textbf{An adventure for four to five characters of 8th level}\\
\end{DndTable}

\section{Adventure Background}
\DndDropCapLine{C}{}enturies ago, the \textbf{red greatwyrm} Soriantathalga laired in the heart of Mount Sagur. Mount Sagur was one of two small, adjacent "twin" volcanos at the edge of a large mountain range near the coast, the other being Mount Magra. Adventurers slew the dragon in a cataclysmic battle, which caused Sagur to erupt, killing them and everything else for miles.


A small caldera formed where Mount Sagur had been. Only about 100 feet in diameter, it has been little more than a blasted pit for centuries; meanwhile, Mount Magra has remained dormant.

In the aftermath of the cataclysm, the plains to the west of Mount Magra and the Sagur caldera became incredibly fertile. Over time, plants grew again in the ground and people settled there, founding villages and cultivating farmland. The closest village is Dorma, about 5 miles away, situated on a burbling stream fed by mountain springs.

Over the course of the last week, several farmers in the area heard strange sounds coming from the caldera—and others reported sightings of monsters, though none of the creatures attacked or approached any of the farms.

Upon learning of this, Ferriahn Colvane—a cleric of Volund, god of earth and fire—went to investigate. She inspected the caldera and its surroundings, finding no sign of volcanic activity in the caldera or Mount Magra, nor did she find any creatures, though she did find some scorched trees. Unsatisfied but out of leads, Ferriahn returned to the village to research and await further developments. After hearing more of the same from the farmers, she returned to investigate again only to discover that the caldera was bubbling with magma and Mount Magra was leaking toxic volcanic gas. Using her magic and knowledge gained by her lifelong study of volcanoes, Ferriahn sensed a planar rift in the magma of the caldera and that it was somehow causing its twin volcano to become unstable. Troubled, she returned to Dorma after narrowly escaping an encounter with a \textbf{fire elemental}.

After some research and prayer, Ferriahn determined that the rift was unstable, and that it's likely that if a certain amount of explosive energy—fire, lightning, thunder, or force—is used in the caldera, it may cause Mount Magra to erupt. Using the \textit{scrying} spell, she observed fiends, dragons, and elementals emerging from the rift, any of which could employ such explosive energies. She believes the planar rift will remain open as long as any of the creatures that came through it survive or remain on this plane—and that closing the rift will render the volcano and caldera dormant once again. If that is not done, she fears the volcano will erupt, causing incredible devastation.

\section{Adventure Hook}
The following can used to start the PCs on this adventure:


\begin{itemize}
\item \textbf{\textit{Hired Heroes.}} Ferriahn met the PCs once and has heard much of their accomplishments since. She casts \textit{sending} to them, imploring them to come to Dorma, offering a reward if the PCs slay the creatures and prevent the eruption. If she finds them honest and forthright, she also offers one uncommon magic item of the GM's choice.
\item \textbf{\textit{Local Lineage.}} One or more of the PCs is from the region or even from Dorma itself; or, a PC is descended from the adventurers who slew the dragon and hope to honor their ancestor's sacrifice.
\item \textbf{\textit{Fire Friend.}} If like Ferriahn, a PC is a worshiper, cleric, or paladin of Volund (or a similar faith), they will likely be more motivated to assist her and the town.
\end{itemize}

Ferriahn is a no-nonsense half-elven cleric of Volund. She's capable, smart, and powerful, but few would describe her as warm. No one has ever doubted her commitment to her faith. She loathes anything that upsets the natural order of mountainous regions, and this certainly qualifies.

She meets the PCs at the village's only tavern, a merry place by the name of the Mountain's Hearth. She informs them of what she knows (see "Adventure Background"). If the PCs ask what would happen if the volcano erupts, Ferriahn tells them that everything in a 10-mile radius would be absolutely devastated and thousands of people would likely die.

\begin{DndSidebar}[float=!b]{What if the Volcano Erupts?}
If the volcano erupts, the eruption lasts two hours, devastating everything in a 10-mile radius. Beyond the initial eruption, lava and gasses continue to leak from the volcano for many weeks afterward.




\begin{itemize}
\item The lava takes two hours to reach the village of Dorma, obliterating everything in its path. See 6 for rules regarding contact with the lava.
\item All creatures and objects in the affected area take 14 (4d6) acid damage from acidic rain. Any creature that takes this damage is blinded for 3 (1d6) days unless it succeeds on a DC 15 Constitution saving throw. Creatures inside buildings are unaffected unless the acid destroys the building.
\item After one hour, volcanic hail rains down on the affected area. Treat this as the spell \textit{ice storm} that affects the entire area and deals fire instead of cold damage. Creatures inside buildings are unaffected unless the hail destroys the building.
\item Each PC must make a DC 13 Wisdom (Survival) check. If more than half of the PCs succeed, they are able to find shelter for their whole party. If less than half succeed, the PCs who succeeded found shelter and can choose whether to use it for themselves or for another PC. Any PC left without shelter takes full damage from the eruption's effects or one with shelter takes half as much. If the Wisdom (Survival) check succeeds by 5 or more, any PC with shelter takes no damage. Creatures with the ability to move quickly or interdimensionally might be able to escape and potentially bring their allies with them.
\item Any creatures already on this plane when the volcano erupts are no longer "gated" to the caldera.
\end{itemize}



\end{DndSidebar}

\section{The Trek to the Mountain}
Once the PCs agree to help, Ferriahn remains behind in the town. She intends to stay there to protect it from any marauding monsters that might flee once the PCs begin their efforts. Before they go, she offers to cast \textit{aid} on the entire party (at 4th level) and \textit{death ward} on one PC.

\begin{DndReadAloud}{}
The brief trip from the village of Dorma to the caldera takes you through verdant fields that were enriched by minerals from the volcanic eruption centuries ago. Farmers wave to you as you pass by. One gnome farmer runs up to you with a basket of freshly baked bread and wishes you luck.

\end{DndReadAloud}

\subsubsection{Travel Trouble}
Travel to the mountain takes the PCs 90 minutes on foot or half that if mounted. When they are halfway there, roll 1d10. If you roll a 10, roll another d10 and consult the "Caldera Creatures" table to determine which creatures they encounter. If the PCs drive off—but do not kill—any creatures encountered, those same creatures may show up again as the PCs are closing the gate.

\begin{DndTable}[header=Caldera Creatures]{XX}

d10 & Encounter\\
1 & 2 \textbf{red dragon wyrmling}\\
2 & 4 \textbf{azer}\\
3 & 1 \textbf{chain devil}\\
4 & 3 \textbf{hell hound}\\
5 & 3 \textbf{bearded devil}\\
6 & 2 \textbf{barbed devil}\\
7 & 1 \textbf{bone devil}\\
8 & 1 \textbf{fire giant}\\
9 & 2 \textbf{fire elemental}\\
10 & Reroll\\
\end{DndTable}

\section{The Caldera}
\begin{DndReadAloud}{}
Traveling east from Dorma, the road you're traveling ends at the ring of rocky terrain that forms the outer edge of the Sagur caldera. Three hundred feet further east stands Mount Magra, beyond which larger mountains loom. The temperature had been rising imperceptibly as you approached the caldera, but now that you stand on its precipice the heat is now oppressive.

\end{DndReadAloud}

\subsubsection{Lighting}
The glow of the lava causes the caldera to be dimly lit. Any area within 5 feet of lava are brightly lit.

\subsubsection{Volatile Volcano}
While the PCs are inside the caldera, tally the amount of thunder, lightning, force, or fire damage dealt by any creature in Areas 2–6, ignoring a creature's immunities or resistances for this purpose. Once that total exceeds 200, Mount Magra erupts. The volcano also erupts if an hour passes before the PCs have defeated all the planar creatures. See the "What if the Volcano Erupts?" sidebar for more.

\subsubsection{Poisonous Gases}
The lava in the caldera releases harmful volcanic gases that irritate the eyes and lungs. At the start of each PC's turn that they're inside the caldera, they must succeed on a DC 13 Constitution saving throw or become poisoned until the start of their next turn; if the roll failed by 5 or more, the PC is also blinded. A PC that fails 3 saving throws (these need not be consecutive) is poisoned for 1 hour. A \textit{lesser restoration} spell removes one of the effects and also removes one failure. A \textit{greater restoration} spell removes both effects and all failures.

An unlabelled version of the map can be found here.

\subsection{1. Ground Level}
The area surrounding the caldera is rocky terrain, but not difficult to traverse.

\subsection{2. Caldera Rim}
\begin{DndReadAloud}{}
As you crest the rocky rim, you can see the entirety of the caldera before you. On the ledge below stand four large rams with glowing red eyes and flaming hooves. An enormous lava pool sits at the lowest point of the caldera, atop of which writhes a massive serpent, its glassy, obsidian scales shining in the glow of the lava around it. The rams and the serpent are in the midst of an all-out battle.

\end{DndReadAloud}

The outer rim of the caldera rises up in a 15-foot-tall escarpment at a 45 degree angle that comes to a peak, then another 45 degree escarpment descends 15 feet to 3, back down to the same elevation as 1. Climbing the outer escarpment requires a successful DC 15 Strength (Athletics) check, which is made with advantage if the PC uses a \textit{climber's kit} or gear such as pitons; no skill check is required to climb down the escarpment. The escarpment is considered difficult terrain.

\subsubsection{Creatures}
See the "Caldera Combat" section.

\subsection{3. Upper Caldera Ledge}
This area is a ledge of volcanic rock that encircles the lava pool below. To the north, south, and east, you can see small lava flows that have burst up out of the rock.

\subsection{4. Lower Escarpment}
A 45-degree-angle escarpment descends here, ranging from 10 to 15 feet down, to another ledge below (5). For climbing rules, see 2.

\subsection{5. Lower Caldera Ledge}
A rock ledge that ranges from 3 to 8 feet wide surrounds a pool of lava (see 6).

\subsubsection{Extreme Heat}
The lava creates extreme heat. A creature without resistance or immunity to fire damage must succeed on a DC 15 Constitution saving throw every 10 minutes it spends in Area 5 or suffer one level of exhaustion.

\subsection{6. Lava Pool}
At the center of 5 is a pool of lava 40 long and 25 feet wide. Channels of lava extend away from the pool to other parts of the caldera, with spurs running upward to the northwest, northeast, and south—with the southern spur branching off east and west on the southern section of 3. These channels have a depth of 1 foot or less, but the central pool itself is extremely deep, leading to an underground channel that connects with Mount Magra.

\subsubsection{Traversing Lava}
If a creature enters or starts its turn in contact with lava, it takes 55 (10d10) fire damage. If a creature starts its turn completely immersed in lava, it takes 99 (18d10) fire damage instead.

\subsubsection{Avoiding Lava}
The lava does not cause contact damage to creatures standing within 5 feet of it, but if that creature begins or ends its turn within 5 feet of it, it must make a DC 15 Constitution save or take 1d6 points of fire damage from the extreme heat.

\subsubsection{Touching Lava}
Creatures that touch the lava, even for a moment, take 2d6 fire damage.

\section{Caldera Combat}
\subsubsection{Creatures}
Within the caldera, four \textbf{ibexian} are fighting a pitched battle with an \textbf{obsidian ophidian}.


\begin{itemize}
\item{\textit{Hampered Foes}.}
The \textbf{ibexian} have been badly injured in their battle, so they are only at a quarter of their hit point maximum. Also, none of their fire-based abilities affect the \textbf{obsidian ophidian}, which also regenerates.

The \textbf{obsidian ophidian} is likewise hindered in the battle. Splashing the \textbf{ibexian} with lava is completely ineffective, and they are resistant to its attacks.

\item{\textit{Volatile Energy}.}
When the PCs first arrive on the scene, they see the \textbf{obsidian ophidian} use its Lava Splash ability and see one of the \textbf{ibexian}, with the \textbf{obsidian ophidian} near the edge of the lava pit, charge down the escarpment and use its Pyroclasm ability. Seeing this use of fire damage should force the PCs to intervene immediately. Don't count this damage in the "Volatile Volcano" tally and allow the \textbf{obsidian ophidian} to automatically succeed on recharging its Lava Splash this once.
\end{itemize}

\subsubsection{Additional Creatures}
Roll 1d10 at the start of each round of combat. On a 10, roll another d10 on the "Caldera Creatures" table, rerolling if the result is a creature that the PCs have already encountered during this adventure. The creature bubbles up from the lava and acts at the end of the round. It instinctively knows that its ability to continue roaming this world is dependent on killing the PCs. It cannot leave the area surrounding the caldera unless the volcano erupts.

\subsubsection{Driven-off Creatures}
If the PCs fought one of the creatures on the way to the caldera, but drove it off rather than killing it, that creature (at full health) could appear on a roll of a 10 instead of gating a new creature in.

Also, if all other creatures are slain, a creature the PCs drove off on their way automatically shows up. It's possible the PCs forgot about this creature, leading them to wonder why the volcano is still spouting gases, still ready to explode. But they won't have long to wonder as the previously repelled creature shows up in 1d6 rounds.

\subsubsection{Leaving the Caldera}
If the PCs leave the caldera, the creatures within follow them, except for the \textbf{obsidian ophidian}. This is a sound strategy because taking the fight outside the caldera gives them more room to spread out, potentially dividing and conquering their foes. In addition, any energy effects they use outside the caldera that deal fire, lightning, thunder, or force damage do not further destabilize the volcano. Finally, the PCs are not vulnerable to the harmful gases of the caldera while outside it.

\subsubsection{Enemy Enmity}
It's important to remember that most of these creatures are not particularly intelligent and probably don't like each other. If another creature gets in a monster's way or harms them, they have no compunctions about attacking that creature. They also take opportunity attack against any creature that leaves their reach, except creatures of their same kind.

\begin{DndSidebar}[float=!b]{Encounter Difficulty}
The encounter at the caldera can prove extremely deadly depending on the encounter table results. Here are some ways to reduce the threat level, should you feel the battle has become too lethal.




\begin{itemize}
\item If you roll a 10 on the encounter table during the battle, simply ignore it.
\item Have the creatures kill each other. Do this sparingly, but the creatures present might gang up on one of the other creatures. Or it could be that two of the creatures particularly hate each other.
\item Make the volcanic eruption require more damage to trigger it.
\item Give the PCs a hint, suggesting they might benefit considerably from drawing the battle away from the caldera. This is particularly reasonable if a PC has the History skill or a military or tactical background.
\end{itemize}



If, however, the encounter proves too easy, here are some ways to amp up the difficulty:




\begin{itemize}
\item Change the die you roll on the encounter table to determine if there's an encounter from a d10 to a d8, or lower the die size by one step each round: d10, d8, d6, d4, d2. You would still roll a d10 to determine which creature appears.
\item Instead of rolling on the encounter table, pick one of the higher CR options.
\end{itemize}



\end{DndSidebar}

\section{Concluding the Adventure}
Once the PCs have killed or otherwise removed all the creatures in the caldera from this plane, the planar rift closes and the volcanic threat is no more. Ferriahn thanks them profusely, gladly handing over any promised reward.

The PCs have no time to explore the mountain during the adventure, but they may do so afterward. Mount Magra is not detailed here, but if the PCs go there they might encounter more of the creatures described in this adventure. They may find the corpses of the ancient heroes and any treasure they had on them. They might also find the bones of the ancient dragon and possibly its hoard. There might even be a new dragon lairing there, attracted by the brief awakening of the volcano.

\chapter{The Eclipsed Chapel}

\begin{DndTable}{c}

\textbf{An adventure for four to five characters of 9th level}\\
\end{DndTable}

\section{Adventure Background}
\DndDropCapLine{T}{}he Eclipse of the Sun Chapel is a legend shared across generations. In the story, a priesthood devoted to the sun god Khors (see \textit{Midgard Worldbook}, or use another god of the sun or light appropriate to your setting) lost their home when ooze cultists attacked during an auspicious solar eclipse. The cult defiled the chapel, creating a planar gate from which oozing horrors spilled forth. Amidst the massacre, a cleric sacrificed themself to seal the portal and block part of the chapel with a magical wall, but the damage was done and the chapel was lost.


Now, a third-generation painter named Pol is visited by an angel who shows him visions of his ancestor at the chapel's fall. The angel then anoints Pol, blessing him with divine authority and proclaiming him an apostle—chosen to reconsecrate the chapel through the working of a miracle. As outlandish as this may seem, Pol accepts his charge with conviction and sets forth on his quest.

\section{Adventure Hook}
Pol (CG human \textbf{apostle}), who works miracles through his art, is guided by visions to seek out the PCs and ask for their help. He has nothing with which to pay the PCs, but claims they have been—like he himself—chosen to save the lost chapel. If the PCs demand payment, Pol reluctantly offers them his painting studio, a property in a mountainside town worth approximately 8,000 gp.

When the PCs agree to help, Pol takes them to where he has a horse-drawn cart laden with supplies and declares he is ready to set forth. The ensuing trek takes Pol and the party six days through mountainous terrain before finally reaching the chapel. Complications during the trek are minimal. Pol always seems to know where to go to avoid danger or hostile creatures. When asked about his prescience, Pol simply responds with "I had a vision."

\section{The Chapel of the Sun}
The Chapel of the Sun faces east on a mountaintop overlooking Pol's hometown. The chapel was built to observe the rising sun and, with its brilliant stained glass windows, was a shining beacon for travelers. The mountain range around the chapel grew rich with alpine berries and flowers which the priests used to make dyes and pigments for blessed works of art.

Pol and the PCs arrive first thing in the morning, just as the sun is rising behind the chapel.

\subsubsection{The Nave Beckons}
Upon arriving at the chapel, Pol leads the PCs directly to 3, so he can immediately begin his work. He is certain—thanks to his visions—that the nave doesn't contain any creatures, traps, or hazards inside it, and, unwilling to waste any time, he is unwilling to wait for the PCs to clear any rooms. He goes directly to the altar and starts his work on his miracle.

\subsubsection{The Work Begins}
After surveying the nave for several minutes, Pol explains that he will paint a great mural on the wall behind the altar. It will be a triptych, he says, and points out the three central sections of wall between the window panes that he intends to use as his canvas. He believes that the act of observing affects the act of creation, so, through prayer, he learned a subtle magic that he uses to obscure his works until they're finished. Thus, as he paints, the PCs cannot see what he's creating. He insists that he will reveal the work only when it is complete.

\subsubsection{Desecration}
When the cultists attacked the chapel, they ravaged the area; they burned the fields, razed the rectories to the ground, destroyed precious artworks and sacred relics, and smeared slime and ooze into the priests' sacred pigments.

The chapel's desecration causes the following effects:


\begin{itemize}
\item Any creature that doesn't have the elemental, fiend, ooze, or undead type is prevented from traveling within 100 feet of the chapel by teleportation or extradimensional/interplanar means, though short-range teleportation within the grounds, such as via \textit{misty step}, still works.
\item While on the chapel grounds, whenever a creature's hit points are restored by healing magic, the number of hit points restored is halved. (Hit dice recovery is unaffected.)
\item Plants cannot grow without magical means within 1 mile of the chapel.
\end{itemize}

\subsubsection{Chapel Ceilings}
The ceiling of the nave reaches 60 feet in height while all the rest of the ceilings in the chapel are 15 feet high.

An unlabelled version of the map can be found here.

\subsection{1. Vestibule}
\begin{DndReadAloud}{}
Gobs of ooze, thick with grit, coat most of this vestibule's walls. A table with a crushed alms box stands against the north wall. Opposite a set of ornate double doors is a shallow alcove, grime and muck obscuring it almost entirely.

\end{DndReadAloud}

In addition to the double doors, there's also a single door on the western wall.

\subsubsection{Mural}
In the alcove on the southern wall is a mural, though it is nearly completely hidden behind layers of grime and muck. A pair of small windows shine down from above the mural—or would if they too weren't entirely obscured.

If cleansed (see the "Removing Ooze Grime" sidebar), the mural is revealed to be a triptych that depicts: (1) a priest building a mechanical moth; (2) the priest releasing the moth to fly toward the sun; and (3) the priest showing pride at a job well done.

\subsubsection{Chapel's Defenders}
If a PC examines and/or attempts to glean the meaning of the mural, they will be able to repair any deactivated \textbf{stained-glass moth} in the nave (3, see "Dawn Encounters" section).

\begin{DndSidebar}[float=!b]{Removing Ooze Grime}
The ooze-based grime used by the cultists to defile the chapel is strangely resistant to prestidigitation and other magic that would otherwise easily remove such filth. But any grime in the chapel can be cleansed by spending 30 minutes continuously using such magic or with one hour of intense scrubbing and/or cleaning (or half as long if a PC is proficient with \textit{painter's supplies} or \textit{alchemist's supplies}).



\end{DndSidebar}

\subsection{2. Pastor's Chambers}
\begin{DndReadAloud}{}
Front and center in this room, a lone skeleton sits slumped over a desk, a large tome under its hand and skull. A closed wardrobe leans against one wall, a fouled wash basin by another. The room is covered in dust, but seems to have escaped the defacement seen in the vestibule.

\end{DndReadAloud}

A successful DC 13 Wisdom (Perception) check discovers a (plain, non-magical) steel \textit{dagger} under the desk. If a successful DC 15 Wisdom (Medicine) check is made before anyone touches or jostles the skeleton, it determines, based on the physical evidence, that this person had been stabbed in the back. If the skeleton is jostled at all, it falls apart.

\subsubsection{Creatures}
Interacting with the tome on the desk awakens a \textbf{gryllus swarm} from inside it. If the PCs defeat or calm and befriend the swarm, they gain access to its tome, an illustrated history of the chapel and its clergy. The information within could be useful to Pol, and/or the swarm could be convinced to help the PCs. The swarm does not pursue PCs beyond room 2.

\subsection{3. Nave}
\begin{DndReadAloud}{}
A ragged rug stretches along the length of this room, flanked by twin rows of moldering pews. Three devotional alcoves line the eastern and western walls, inside of which stand pedestals with shattered stained glass littering the floor around them. Two large pulpits at the north end of the enormous chamber stand before the raised altar at the room's far end. Natural light from above filters into the room through in a kaleidoscopic haze from the few intact stained glass windows that remain.

\end{DndReadAloud}

\subsubsection{The Alcoves}
The pedestals in the alcoves on the eastern and western sides of the nave are all bare or smashed, hinting at missing works of art each crafted by a priest.

\subsubsection{The Altar}
The solar burst \textit{holy symbol} is broken and covered in grime. A PC who spends an hour cleansing the altar (see the "Removing Ooze Grime" sidebar) and succeeds on a DC 17 Intelligence check using \textit{tinker's tools} or \textit{mason's tools} can reassemble the \textit{holy symbol}, which causes it to burst into flame as the \textit{continual flame} spell.

\subsubsection{Shattered Glass}
The piles of shattered glass on the floor are all difficult terrain; a creature can choose to ignore the difficult terrain, but if so it takes 1d4 slashing damage each 5 feet it moves.

\subsubsection{Treasure}
A shelf in the western pulpit holds a ruined worship text. Inside the base of the pulpit is a hidden compartment (that is impossible to see or detect with divination magic) which holds a \textit{mace of disruption}; the mace has an additional blessing from the sun god that causes its feature that affects fiends and undead to also affect oozes. The only way to open the compartment is depicted on the alcove wall in 6; a creature must: (1) light a \textit{torch} with the altar's flame; (2) hold the \textit{torch} aloft, at which point the light cast by the \textit{torch} allows the creature to see the hidden compartment; (3) touching the compartment door while holding the \textit{torch} opens it.

\subsection{4–5. Narthex}
\begin{DndReadAloud}{}
This room is separated from the main chapel by a door and decorative half wall. Religious instruments and objects of devotion lie forgotten in corners and cabinets. One cabinet—a wardrobe—holds priestly vestments that crumble if disturbed. In the corner of the room, a staircase leads down into the undercroft, the steps covered in slime trails and footprints from another age.

\end{DndReadAloud}

Areas 4 and 5 are functionally the same. The half wall is 4 feet tall.

\subsubsection{Treasure}
The narthices are mostly empty, but a successful DC 15 Intelligence (Investigation) check in each room finds up to 1,000 gp in material components for divine ritual spells (at the GM's discretion); on a failure, only up to 250 gp is found in each room.

\subsection{6. Undercroft Hallway}
\begin{DndReadAloud}{}
A foul stench permeates this underground corridor, but most striking is the wall of shimmering light that seals closed an archway in the southern wall. The dim light cast by the wall shines on several crates and barrels, a toppled shelf in an oily puddle, and an empty alcove with its walls covered in grime.

\end{DndReadAloud}

\subsubsection{Mural}
A mural is painted along the alcove of the northern wall, though it is obscured behind layers of muck and grime.

If cleansed (see the "Removing Ooze Grime" sidebar), the mural is revealed to be a triptych that depicts: (1) a priest kneeling before an altar to light a \textit{torch} on its aflame \textit{holy symbol}; (2) the priest standing with the \textit{torch}—which has now become a mace—raised high, burning with holy fire; and (3) the priest valiantly fighting creatures of the dark with the weapon.

\subsubsection{Shimmering Wall}
The wall is a \textit{wall of force} (as the spell), which can be identified with a successful DC 15 Intelligence (Arcana) check; a creature that can cast \textit{wall of force} automatically succeeds on the check.

\subsection{7. Sanctum of Saints}
\begin{DndReadAloud}{}
Beyond the impassable light barrier is a large chamber shrouded in darkness. Silhouettes in the soft glow of the barrier suggest a circular font in the center, half a dozen pillars, and the chiseled edges of large stone sarcophagi.

\end{DndReadAloud}

A stone sarcophagus sits in each of the alcoves to the east and west—and in the larger alcove to the south. A large circular font occupies the center of the chamber.

\subsubsection{Font}
Currently, the font is an oozing cesspool that subtly roils and is viscerally upsetting to look upon. A successful DC 15 Intelligence (Religion) check notices some char marks on the edge of the font, along with the faint, lingering scent of incense in the air, which suggests that it was once used was a focal point of remembrance and contemplation.

\subsubsection{Sarcophagi}
All of the lids of the sarcophagi have been covered in grime, though it is clear that they bear bas-relief carvings on them. If cleansed (see the "Removing Ooze Grime" sidebar), the lids of the sarcophagi in the east and west alcoves are revealed to feature bas-relief carvings that depict the sun god's apostles and their good deeds. The lid of the one to the south is more elaborate, with its bas-relief carving also being painted with gold leaf and bright colors.

\subsubsection{Creatures}
This catacomb is now both lair and prison to a \textbf{sinoper ooze} due to the \textit{wall of force} blocking egress into the room (see "Pol's Miracle" section).

\subsubsection{Cesspool}
Any creature with the ooze type that starts its turn in the cesspool regains 11 (2d10) hit points.

\subsubsection{Consecration}
Once the \textbf{sinoper ooze} is defeated, the cesspool bubbles violently and completely evaporates to reveal a bas-relief carving of the sun. Once all of the sarcophagi are cleansed, the sun carving immediately fills the room with holy light.

\section{Standing Vigil}
Pol's work takes an entire day to complete, during which time the PCs must stand vigil—protecting and providing for Pol—and clear the rooms of the chapel to ensure no unseen dangers lurk.

Mechanically, the PCs' vigil is divided into three phases: Dawn, Day, and Dusk. Each phase lasts 3 hours, during which the PCs experience a set of three encounters. Allow the PCs to take a short rest in between phases and time to explore the chapel as well.

\subsection{1. Dawn: Enchanted Pigments}
Pol points out that there is evidence of enchanted pigments on the walls in the nave and asks the PCs to look around the chapel for more. Each PC must make a DC 17 Intelligence (History, Investigation, or Religion) or Wisdom (Perception) check, made with advantage if the PC is proficient with \textit{painter's supplies}. Tally the number of successes and consult the "Enchanted Pigments" table to determine what the PCs learn; a higher result finds the items from lower results as well.

\begin{DndTable}[header=Enchanted Pigments]{XX}

Successes & Information Learned\\
1–2 & Some of the paints used—antimony, cinnabar, and sinoper—are acidic and poisonous substances that can be fatal if ingested or handled improperly.\\
3 & The cultists that attacked the chapel stole many of the pigments and spilled the rest in the undercroft.\\
4 or more & An ooze absorbed the spilled, toxic pigments, potentially making it more dangerous.\\
\end{DndTable}

\subsubsection{Reward}
Any PC that succeeds on this roll gains a blessing of inspiration, which allows them to add 1d4 to one ability check using Charisma.

\subsection{2. Dawn: Malfunctioning Machines}
\subsubsection{Creatures}
Six \textbf{stained-glass moth}, high in Area 3's rafters, activate with the rising sun and immediately attack the PCs. The moths recognize Pol's ancestry and ignore him. After 2 rounds of combat, on initiative count 20, a \textbf{gearmass} enters from the undercroft. There are moth parts stuck inside the gearmass, causing any remaining moths to ignore it when it arrives. The gearmass is drawn to the PCs and also to any incapacitated constructs.

\subsubsection{Moth Repair}
After the moths are reduced to 0 hit points, if a PC has studied the mural in 1, they can spend 1 hour (or 10 minutes if they're proficient with \textit{tinker's tools}) repairing a moth (they're too complicated to repair with the \textit{mending} spell), using spare parts scrapped from two other moths. The moths have been corrupted by the ooze defilement; otherwise they wouldn't have attacked the PCs. The only way to cleanse them of the defilement is to take them outside and release them toward the sun, as in the mural. Any repaired moths remain non-hostile for 1 hour, after which, if they haven't been released toward the sun, they turn violent once again. After being cleansed, any repaired moths return to the rafters in 3 to act as chapel defenders. If the PCs struggle with a combat encounter in 3, any repaired moths will come to their aid and attack any hostile creatures.

\subsection{3. Dawn: Mural Modeling}
Pol asks the PCs to act as models for his mural. Each PC must make a DC 15 ability check using the skill of their choice. GMs should ask the players to describe what their pose looks like, such as a heroic pose for Strength (Athletics) or a magical flourish for Intelligence (Arcana). He sketches the poses in a sketchbook to later incorporate into the mural.

\subsubsection{Reward}
Each PC that participates gains the benefit of the \textit{aid} spell for the duration of the vigil.

\subsection{7. Day: Crisis of Faith}
Pol asks the PCs if they believe in the gods. Unsure of where he himself stands even as he works to paint his miracle, Pol is suffering a crisis of faith and needs guidance. The painter is open to hearing all opinions, not just those of trained cleric or paladin, and wants to hear from each PC.

\subsubsection{Reward}
Any PC that engages with this question with Pol receives the benefit of the \textit{shield of faith} spell for 1 hour.

\subsection{8. Day: Painter's Block}
Pol has hit an artistic block and is desperate for inspiration; he asks the PCs to search the premises for some creative insight. Each PC must make a DC 17 Intelligence (History, Investigation, or Religion) or Wisdom (Perception) check, made with advantage if the PC is proficient with \textit{painter's supplies}. Tally the number of successes and consult the "Creative Insights" table to determine what the PCs learn; a higher result finds the items from lower results as well.

\begin{DndTable}[header=Creative Insights]{XX}

Successes & Information Learned\\
1–2 & There is imagery of warrior angels in the narthices and undercroft. The angels appear harsh, meting out judgments of holy flame and condemnation.\\
3 & The angels are painted with long dried cinnabar and sinoper pigments and are oriented toward the basement.\\
4 or more & The smearing and acidic pockmarking around the images suggests something crawled across the images, taking some of the toxic paint with it.\\
\end{DndTable}

\subsubsection{Reward}
Any PC that succeeds on this roll gains a blessing of inspiration, which allows them to add 1d4 to one ability check using Wisdom.

\subsection{9. Day: Pseudopod Predator}
\subsubsection{Creatures}
A \textbf{copperkill slime} enters 1, spreading itself against the alcove wall. It tries to draw attention to the room by using its pseudopod to grab and throw the nearest heavy object, possibly the alms table if it is still in the room. If this fails to bring anyone within reach of it, the ooze tries to sneak through the chapel to get close enough to a creature to engulf them. If it is still alive after 4 hours or having successfully eaten, it flees the area.

Meanwhile, a \textbf{leavesrot ooze}, which manifested on the grounds as a result of the chapel's defilement, has entered the chapel through the broken window in 2. It lies in wait, oozing under the door to attack once the \textbf{copperkill slime} springs its ambush. The \textbf{leavesrot ooze} is indistinguishable from a pile of leaves when unmoving, but a PC with a passive Perception of 17 or higher notices that some leaves seem to have blown into 2 and are now just visible under the door.

\subsection{10. Dusk: Defiled Defender}
\subsubsection{Creatures}
Four \textbf{gray ooze} come up the stairs of Areas 4 and 5, two on each side.

Once the oozes emerge into 3 to attack, a spirit rises through the altar and steps forward: a \textbf{ghost}, the spirit of one of those interred in the undercroft, awoken and corrupted by the ooze defilement. It seeks to defend the chapel from Pol and the PCs and sees the oozes as its allies. It does not use Possession, as it seeks to oust, not control, the intruders. It doesn't know who it was anymore and can't be reasoned with.

After 2 rounds of combat, on initiative count 20, four more \textbf{gray ooze} come up the narthex stairs and join the fray.

\subsection{11. Dusk: Burial Grounds}
Pol has another vision, this time of hungry undead clawing and biting him after bursting up from below. He asks the PCs to search the chapel to find if there are dead buried below or on the chapel grounds. Each PC must make a DC 17 Intelligence (History, Investigation, or Religion) or Wisdom (Perception) check, made with advantage if the PC has a religious background or one associated with the dead. Tally the number of successes and consult the "Chapel Search" table to determine what the PCs learn; a higher result finds the items from lower results as well.

\begin{DndTable}[header=Chapel Search]{XX}

Successes & Information Learned\\
1–2 & There is a graveyard nearby, and the enchantments in place to prevent undead from rising there are still intact.\\
3 & The room behind the light barrier is called the Sanctum of Saints.\\
4 or more & The Sanctum of Saints contains six dead priests and a saint and seems to be the primary source of evil and corruption defiling the chapel.\\
\end{DndTable}

\subsubsection{Reward}
Any PC that succeeds on this roll gains a blessing of inspiration, which allows them to add 1d4 to one ability check using Intelligence.

\subsection{12. Dusk: Fear of Fate}
Pol takes one last break to have a small meal with the PCs. At the meal, he confides that he's afraid of what will happen next, what is contained below the chapel, and what his fate will be when his quest concludes. At the end of the meal, Pol has another vision, warning of an undulating horror waiting beneath the chapel that must be destroyed.

After the meal, the PCs have 1 hour to prepare for the final encounter while Pol finishes his work.

\subsubsection{Reward}
The PC that reassures, supports, and engages with Pol the most during this discussion (as adjudicated by the GM) receives the benefits of the \textit{protection from evil and good} spell in perpetuity whenever they are on the chapel grounds; this blessing adds oozes to the creature types affected by the spell.

\section{Pol's Miracle}
After all four encounter phases are complete, Pol stands and triumphantly states that he is finished and reveals to you his miracle.

Pol's mural is truly an astonishing work, employing vibrant colors and remarkable skill. The triptych, from left to right, depicts the adventure: The PCs arriving at the chapel, a scene showing whichever encounter the players enjoyed the most, and the culmination of their work. The GM is encouraged to try to make it feel meaningful for the players.

After showing off the mural, Pol falls unconscious. Simultaneously, the chapel fills with a brief flash of light as the light barrier in the undercroft is dispelled.

\subsection{Creatures}
After the light barrier is dispelled, the \textbf{sinoper ooze} and two \textbf{ochre jelly} use their Spider Climb and Amorphous abilities to ooze up through the floorboards in the nave and attack. After 5 rounds or when the \textbf{sinoper ooze} is reduced to half its hit point maximum, it Disengages and retreats, again using its Amorphous ability to slip through cracks in the floorboards to return to 7, to heal itself within the cesspool.

\section{Concluding the Adventure}
When the \textbf{sinoper ooze} is defeated, all of the remaining evil is driven out, and the chapel's reconsecration is complete. Pol regains consciousness and makes good on any deal agreed upon at the beginning of the adventure. He also blesses some of his leftover paint, gifting it to the PCs. The paint is 1 pot of marvelous pigments or \textit{interplanar paint} (see \textit{Vault of Magic}). Pol remains in the chapel, and within a few months he recruits followers who work to repair the chapel and the priesthood.

\chapter{The Twisted Sanctuary}

\begin{DndTable}{c}

\textbf{An adventure for four to five characters of 10th level}\\
\end{DndTable}

\section{Adventure Background}
\DndDropCapLine{T}{}he Veiled Mannequin is a tailor's shop for discerning customers run by the gnome couturier Yofi Whorthaus. Unbeknownst to the inhabitants of the city, Yofi and her halfling assistant Naver guard the secret entrance to the Twisted Sanctuary, the safehouse of a rakshasa cult hidden in the heart of the city.


The Cult of the Claw secretly corrupts and manipulates city officials as part of a wide-ranging plot designed to enhance their own power and influence. Narga Kan-Sorkis, the leader of the cult, needs several key components to complete a magical pool that will give him the ability to dominate minds from a distance. The cult's rakshasas (and their other minions) have been ordered to acquire the final components at all costs. Over the past three weeks, a spree of vicious crimes has rocked the city, shrouding the otherwise thriving community in a fog of frightened paranoia.

\section{Adventure Hook}
The following can used to start the PCs on this adventure:


\begin{itemize}
\item \textbf{\textit{Missing Child.}} Looking into the kidnapping of a young child, the PCs follow the evidence to the door of the Veiled Mannequin.
\item \textbf{\textit{Grisly Murder.}} While investigating the crimes in the city, or while involved in other city matters, the PCs come upon the scene of a grisly murder of a city guardsman. A fragment of silk in the dead man's hand is of a type only sold in a few establishments in the city, including the Veiled Mannequin.
\item \textbf{\textit{Find the Halfling.}} Desperate city officials offer the PCs 500 gp to put a stop to the rash of thefts, murders, and outright terrorism plaguing the city. The only lead they have is the description of a halfling spotted by guardsmen at the scene of three of the incidents. An hour or two spent investigating points the PCs at three halflings fitting the description, including Navar, the apprentice at the Veiled Mannequin.
\end{itemize}

\section{The Veiled Mannequin}
An elegantly carved sign above the door reads, "The Veiled Mannequin. Fashions for Discerning Lords, Ladies, and Lieges. Y. Whorthaus, Couturier/Designer."

An unlabelled version of the map can be found here.

\subsection{1. Showroom}
\begin{DndReadAloud}{}
Fanciful dresses and sharp-looking coats hang on racks along the east wall, and three mannequins along the north wall are draped in partially constructed garments. On the long wooden counter to the west are several garments laid out, as well as a book open to a page showing a sketch of a dress design.

\end{DndReadAloud}

This room is the primary workspace and showroom for Yofi Whorthaus.

\subsubsection{Creatures}
The proprietor, Yofi (NE \textbf{gnome cult fanatic}) and her assistant Naver (NE \textbf{halfling cultist}) can usually be found here during the day. At night, Naver sleeps in a hammock in the Northwest corner while Yofi retires to 2. Both Yofi and Naver are fanatically loyal to the cult and would never willingly betray any of its secrets.

\subsubsection{Alarm Stone}
If Yofi decides a PC is a threat, she touches a gemstone enchanted with the \textit{alarm} spell that is affixed to the underside of the counter. This silently alerts Narga Kan-Sokris and his right-hand Aliza Sahar to the danger (see Areas 5 and 13).

\subsection{2. Yofi's Quarters}
\begin{DndReadAloud}{}
This is a small but elegant bedroom with a large wardrobe against the north wall. The bed is covered with a very fine embroidered quilt and feather pillows.

\end{DndReadAloud}

\subsubsection{Secret Door}
A successful DC 15 Wisdom (Perception) check spots faint scratches on the wooden floor where the wardrobe has slid from left to right. A secret switch inside the wardrobe allows the bulky furniture to roll, revealing a hidden door. A successful DC 17 Intelligence (Investigation) check finds the switch. The secret door opens to a spiral staircase going down (3).

\subsection{3. Spiral Stairs}
\begin{DndReadAloud}{}
A curving marble staircase winds down the inside curve of a twenty feet wide, eighty feet high, cylindrical shaft that descends deep into the ground. The domed ceiling of the chamber is fashioned to resemble a human face with an open screaming mouth.

\end{DndReadAloud}

\subsubsection{Rolling Fire Trap}
This magical trap activates when one or more creatures steps on the center of the first step. Once triggered, a ball of fire emerges from the mouth on the ceiling and begins rolling down the stairs. Each creature on the stairs must make a DC 13 Dexterity saving throw; the creature takes 3d6 fire damage on a failed save or half as much damage on a successful one. The mouth spits additional balls of fire as the PCs descend the stairs; to reach the bottom, each PC must make 3 separate saves. The balls of fire vanish when they reach the foot of the stairs.

A successful DC 15 Wisdom (Perception) check notices faint glyphs on the stair, which allows the PCs to avoid the trap if they step around it. A successful \textit{dispel magic} (DC 13) removes this enchantment.

\section{The Twisted Sanctuary}
The cultist's lair is richly appointed and elegant, constructed of fine marble, decorated with murals and mosaic tile. Unless noted in the text, none of the doors of the Twisted Sanctuary are locked, and the chambers are brightly illuminated by glowing hanging orbs.

\subsection{4. The Entry Hall}
\begin{DndReadAloud}{}
Just beyond the landing of the spiral staircase is a short, wide flight of stairs that leads to grand hall lined with carved stone pillars. The floor is tiled in mosaic patterns, and rich tapestries hang from the walls. At the far side of the room, the hall ends in a large metal door set into an imposing wall carved with tiger heads and images of ritual slaughter.

\end{DndReadAloud}

\subsubsection{Creatures}
When the PCs arrive in this chamber, they discover an \textbf{apostle} overseeing the training of three \textbf{cult fanatic} and four \textbf{cultist}. If the foes are alerted to the PCs' intrusion (either via the \textit{alarm} in 1 or triggering the trap in 3), the \textbf{apostle} and cultists attack immediately. Otherwise, the PCs have surprise.

\subsubsection{Treasure}
The cultists are armed with daggers and scimitars. Two of the \textbf{cult fanatic} wear pendants made of gold, shaped like a clawed hand, holding a large ruby (worth 1,000 gp each).

\subsection{5. Audience Hall}
\begin{DndReadAloud}{}
A colorfully painted gazebo with four pillars supporting a domed roof sits in the center of the chamber. To the east, a few short steps lead up to a raised nook where a figure in a hooded robe sits atop plush cushions while smoking on a brass hookah. Standing behind and to her left is a tall, tigerfaced woman in scale mail.

\end{DndReadAloud}

Aliza Sahar (LE human \textbf{apostle}) is tasked with meeting with members and handling the day-to-day operations of the cult. Aliza uses this room for that purpose and also to collect information from the cult's spies. Almost no one is even aware of Narga's existence, believing that Aliza runs the cult. The pair is happy to keep up this ruse, which allows Narga plenty of freedom to carry out his schemes. Unwaveringly loyal to Narga, Aliza attempts to ferret out as much information as she can from any intruders and fights to the death to protect the door that leads deeper into the sanctuary.

\subsubsection{Creatures}
Aliza Sahar is a human \textbf{apostle} with five \textbf{myrmidon rakshasa}, one standing behind her, and four lurking inside hidden alcoves in each corner of the chamber; the alcoves allow a creature inside it to see and hear into the room.

\subsubsection{Locked Door}
The door to the north is locked, but can be picked with a successful DC 17 Dexterity check using \textit{thieves' tools}.

\subsubsection{Treasure}
Beneath the cushions in the raised nook is a small coffer containing a \textit{potion of invisibility}.

\subsection{6. Long Hall}
\begin{DndReadAloud}{}
The walls of this hallway feature beautiful—but upsetting—artwork depicting tiger-headed humanoids oppressing and torturing all manner of other races. A carved pillar stands in the center of an intersection to the north, the floor around it decorated with exquisite mosaic tilework.

\end{DndReadAloud}

\subsubsection{Pillar Trap}
Snarling tiger heads decorate each face of this diamond shaped pillar. Any non-rakshasa that steps on the mosaic floor around the pillar (10 foot radius) must succeed on a DC 14 Wisdom saving throw or suffers the effects of the \textit{confusion} spell for 1 minute; any creature wearing a ruby claw pendant (see 4) automatically succeeds on the saving throw. A successful DC 17 Wisdom (Perception) check notices glyphs hidden on the mosaic tiling on the floor. A successful \textit{dispel magic} (DC 14) removes this enchantment for 1 hour.

\subsection{7. Storage Area}
\begin{DndReadAloud}{}
Boxes and crates line the walls of this chamber, as well as a row of crates down the center of the room.

\end{DndReadAloud}

\subsubsection{Creatures}
Three \textbf{servitor rakshasa} labor here. They attempt to flee to 8 if confronted, taking the Disengage and/or Dash action if necessary.

\subsubsection{Treasure}
The crates contain supplies of food, wine, mundane gear, and construction supplies. A successful DC 12 Wisdom (Perception) check uncovers a leather bound \textit{spellbook} (see Development) in an unlocked chest. The incantations in the book appear to be formulae for the following spells: \textit{dominate person}, \textit{disintegrate}, \textit{finger of death}, \textit{dominate monster}, \textit{power word kill} and \textit{wish}.

\subsubsection{Development}
The tome is actually a \textbf{pustakam rakshasa} who was napping in book form when the PCs entered. The fiend waits until it is in the possession of a PC, then, at an opportune moment, it uses its \textit{suggestion} and \textit{dominate person} spells to sow chaos among the party.

\subsection{8. Lakeside Quarters}
\begin{DndReadAloud}{}
The damp air in this room is heavy with the heady scent of incense, strange spices, and the metallic tinge of blood. Embroidered fabrics line the walls, arching up to gather near the ceiling, giving the interior a pavilion-like facade. Piles of rich cushions festoon the area, with elegant side tables groaning beneath the weight of full wine pitchers, platters of sweetmeats, scattered playing cards, candied fruits, and curried meats on long skewers. The northern wall of the chamber is a series of arches and pillars, draped with diaphanous silks, opening onto the shore of an underground lake.

\end{DndReadAloud}

When not disguised as mortals and advancing their foul plots or participating in evil rituals, the rakshasa member of the cult lounge here, waited on by lesser rakshasas and \textbf{imp}.

\subsubsection{Creatures}
Six \textbf{myrmidon rakshasa} lounge here, waited on by three \textbf{servitor rakshasa} and two \textbf{imp}. If warned by the servitors in 7—or by any active fighting in 6—the rakshasas lurk along the southern wall, ready to ambush the PCs. Otherwise, they relax in the chamber, feasting and drinking.

\subsubsection{Treasure}
A search of the area finds 5 bottles of very fine wine (worth 500 gp total). 1d6 of the playing cards found on the various tables are random cards from a \textit{deck of illusions}.

\subsection{9. Underground Lake}
\begin{DndReadAloud}{}
The midnight water flows from the northeast into a huge cavern, pooling in a deep, murky lake, before sweeping outward to the southwest.

\end{DndReadAloud}

Apart from the tiny rocky shore along the northern edge of 8, the walls of the cavern drop vertically to the bottom of the lake.

\subsubsection{Creatures}
Five \textbf{bearded devil} and three \textbf{imp} inhabit the lake and attack any creature foolish enough to enter their domain.

\subsubsection{Illusion}
The cavern wall to the northwest is an illusion that hides the dock in 14. A creature that uses its action to examine the image can determine that it is an illusion with a successful Intelligence (Investigation) check against DC 16.

\subsubsection{Development}
A successful DC 15 Wisdom (Perception) check notices a golden glow emanating from something at the bottom of the lake 15 feet from the shore of 8.

\subsubsection{Treasure}
The rakshasas threw a \textit{gem of brightness} into the lake after they failed to corrupt it. The item lies at a depth of 15 feet.

\subsection{10. Stone Bridge}
\begin{DndReadAloud}{}
A graceful bridge arches across a subterranean river here. Crystals atop a pair of pillars on each side of the span shine with an orange light, casting distorted reflections on the black waters below.

\end{DndReadAloud}

The crystals shine with a \textit{continual flame} spell but are otherwise merely decorative.

\subsection{11. The Chiropteran Stepwell}
\begin{DndReadAloud}{}
A huge, inverted step pyramid descends three levels down to a circular dais topped with two rings of freestanding, ornamentally decorated pillars. Each step is 15 feet below the one above it, the final level resting 50 feet beneath the ceiling of the chamber. An archway on the northern wall is carved to resemble a snarling tiger's head with a large metal door set into its mouth. The air in the chamber bears the musky scent of incense and guano.

\end{DndReadAloud}

The cult's spymasters, Chiaana and Sunana (NE \textbf{half-elf atavist}), keep their colony of bats in this large chamber. Together they raise and care for the bat colony, hand-feeding and training the bats daily. As such, the bats understand and follow basic orders, such as to attack a specific target, from the twins, who need only speak single-word commands (no action required) to direct the bats.

\subsubsection{Creatures}
Chiaana and Sunana hide on the ceiling alongside their bat companions (twelve \textbf{swarm of bats}), using their Malleable Physiology trait to sprout fleshy, membranous wings from their shoulders.

\subsubsection{Dais}
The mosaic tiles at the center of the dais form a sliding tile puzzle. A successful DC 12 Wisdom (Perception) check spots the missing tile that begins the puzzle. Solving the puzzle requires a successful DC 15 Intelligence (Investigation) check, and the revealed image depicts a tiger's head. When the puzzle is solved, a niche opens on one of the pillars, revealing a lever that unlocks the door in the northern wall.

\subsubsection{Metal Door}
This door is secured with an \textit{arcane lock} spell (DC 25) but can be unlocked by pulling the levers revealed when one solves the puzzle.

\subsection{12. Hallway}
\begin{DndReadAloud}{}
Mosaic tiles cover the floor of this long hallway in intricate patterns and colors while glowing orbs shine down from recesses in the arched ceiling and stone reliefs of tiger's heads snarl out from the walls.

\end{DndReadAloud}

\subsection{13. The Tiger's Eye}
\begin{DndReadAloud}{}
Tall pillars support the cavernous ceiling of this ornate chamber. An oval pool of viscous, reddish brown liquid glows in the center of the room, at the foot of a wall of rich crimson velvet drapery bisecting the chamber. Steep mosaicked stairs lead 15 feet up to a platform covered with glowing runes. A semicircular dais overlooks the glowing liquid from the other side of the pool.

\end{DndReadAloud}

Narga Kan-Sorkis labors to construct his domination pool in this chamber. It is incomplete still and he fights to the death to protect it.

\subsubsection{Creatures}
\textbf{Narga Kan-Sorkis} (LE \textbf{slayer rakshasa}), waits here to attack the interlopers—invisible if the PCs have been detected. Two \textbf{myrmidon rakshasa} lurk behind the curtains, attacking if they hear their master in danger.

\subsubsection{Legendary Actions}
Narga can take 3 legendary actions, choosing from the options below. Only one legendary action can be used at a time and only at the end of another creature's turn. Narga regains spent legendary actions at the start of his turn.


\begin{itemize}
\item \textbf{\textit{Teleport.}} Narga magically teleports to an unoccupied space he can see within 60 feet of him.
\item \textbf{\textit{Restorative.}} Narga recovers 7 (2d6) hit points.
\item \textbf{\textit{Fetch (Costs 2 Actions).}} Narga summons a \textbf{servitor rakshasa}. It appears in an unoccupied space Narga can see within 30 feet and acts on his initiative count.
\item \textbf{\textit{Whelm (Costs 3 Actions).}} A pseudopod of magical liquid lashes out from the pool. A target within 30 feet of the pool must make a DC 15 Strength saving throw. On a failure, the target takes 13 (3d8) bludgeoning damage. If it is Large or smaller, it is also grappled (escape DC 15). Until this grapple ends, the target is restrained and unable to breathe unless it can breathe water. The pool can only grapple one creature at a time.
\end{itemize}

\subsubsection{Curtain}
Beyond the velvet curtain is Narga's palatial living quarters. A secret door in the eastern wall—which is detected with a successful DC 15 Wisdom (Perception) check—leads to 14.

\subsubsection{Treasure}
Narga keeps his greatest treasures on display in his living quarters including: a \textit{wand of magic detection}, an ornate \textit{dagger of venom}, and a \textit{circlet of blasting}. An open coffer holds 4,230 gp in assorted coins. The various works of art and rich furnishings could bring 7,810 gp to the right buyer. A table covered with papers could offer important information on local nobles or officials (at the GMs discretion).

\subsection{14. Hidden Dock}
\begin{DndReadAloud}{}
This small natural grotto features a small wooden dock that reaches out onto the dark water. A 15-foot-long wooden skiff floats next to it.

\end{DndReadAloud}

Narga keeps this vessel here in case he needs to make a hurried escape from the lair.

\subsubsection{Treasure}
The boat is actually a \textit{folding boat}, though it only produces a 15 foot craft without a deck cabin.

\section{Concluding the Adventure}
The adventure ends when the PCs defeat Narga. They are given the appropriate reward for their efforts, depending on the hook used to begin the quest. However, thwarting Narga brings the PCs to the attention of the other rakshasa cults, and if there is one thing the fiends know, it's revenge.

\chapter{A Midnight Ride}

\begin{DndTable}{c}

\textbf{An adventure for four to five characters of 11th level}\\
\end{DndTable}

\section{Adventure Background}
\DndDropCapLine{A}{} strange, black horse draped in silver chains has been seen in the woods and fields around the small farming village of Derrymead. The locals claim it is a \textbf{púca}, and they're certain it will run off with their children. More concerning, however, is the group of cultists trying to get their hands on the \textbf{púca} and use it to gain access to the other planes. The cultists have claimed an old barrow near Derrymead for themselves and scour the roads and forests day and night in hopes of finding the beast.


\section{Adventure Hook}
The following can used to start the PCs on this adventure:


\begin{itemize}
\item \textbf{\textit{Heroes Wanted.}} A flyer requests brave heroes to help rid the town of Derrymead of a frightening creature that endangers the town's children.
\item \textbf{\textit{Hail Traveler.}} A farmer from Derrymead hauling a cartful of goods meets the PCs on the road and passes on the information about the \textbf{púca}.
\item \textbf{\textit{Just Passing Through.}} The PCs are passing through Derrymead and stop at the Fallows, Derrymead's inn.
\end{itemize}

No matter how the PCs arrive in Derrymead, they're directed to the Fallows to speak with the innkeeper, Felicity Dunn.

\section{Derrymead}
Derrymead is a small town surrounded by farms and fields that are wild and often . . . strange. The farmers do their best to push back the forest and cultivate the land, but danger constantly lurks in those woods. The folk are hardy and strong, but hold tight to superstitions and folktales to keep themselves and their children safe. The town itself consists mainly of a town square with a few small shops and the Fallows Tavern and Inn.

\subsection{The Fallows Tavern and Inn}
The Fallows Tavern and Inn is a modest establishment in the tiny town square of Derrymead. It boasts but six tables and has three rooms for rent. A notice board just inside the door sports scraps of paper looking for help with hauling wood or sowing grain—and the flyer described in the "Heroes Wanted" adventure hook; it directs interested parties to speak to the proprietor of the inn, Felicity Dunn (NG human \textbf{commoner}), who is a handsome woman of middle years with dark skin and curly black hair.

\subsection{What Felicity Knows}
A strange black horse draped in silver chains was first spotted in the woods outside of town four weeks ago. The eerie, cold light in the beast's eyes hinted that it is not of this world, and some of the farmers are convinced it's a \textbf{púca}—a corrupted fey steed that lures innocents onto its back and then rides off with them.

A week ago, several strange, robed figures appeared nearby town and began roaming the woods day and night as if hunting something. People say that their eyes are pure black but with specks of light in them, like a starry sky. Felicity doesn't know who they are, but she's sure they're here because of the creature—though whether they're working with the beast or against it, she doesn't know. Even more troubling: Three days ago, the blacksmith's son, Auri Rorsk, saw one of the robed figures while walking the road to town with his sister. The fool boy confronted the stranger, and, well . . . his mind hasn't been the same since. She doesn't know more than that about Auri, but suggests that perhaps the boy's sister Leena might be able to tell them more.

The farmers have managed to pool together a pouch of 500 gp as payment for anyone able to kill or drive off these frightening strangers. Additionally, one of the farmers is offering up his grandfather's dwarven plate armor as a reward if the adventurers can rid the town of the \textbf{púca}—or whatever the creature is—as well.

\section{The Cult and the Púca}
The cult is enthralled by an entity known as the Astral Rider, who is said to ride a majestic black steed across the Astral Plane, driving all he encounters mad. They're convinced that capturing the \textbf{púca} will please their patron.

\subsection{Encountering the Cultists}
Any time the PCs venture into the woods or fields around Derrymead, there is a chance they encounter wandering cultists searching for the \textbf{púca}. The cultists are hostile and cannot be reasoned with, immediately attacking the PCs. Their faces are all pale and drawn, and their eyes are as Felicity described.

\subsubsection{Creatures}
If the PCs are out during the day, roll 1d10. On a result of 1–5, the PCs encounter that number of \textbf{Astral Rider cult fanatic}. If the PCs are out at night, they encounter a more robust group pursuing the \textbf{púca} (see the "Encountering the Púca" section).

\subparagraph{Creature Change}
Any \textbf{cult fanatic} the PCs encounter in this adventure have darkvision (60 ft.) and the following changes to their spell list: replace \textit{sacred flame} with \textit{eldritch blast}; replace \textit{shield of faith} with \textit{bane}; and replace \textit{spiritual weapon} with \textit{darkness}.

\subparagraph{Questioning the Cultists}
Interrogating a \textbf{Astral Rider cult fanatic} is tricky business. For every minute the PCs attempt to question one, roll a d20. On a result of 1–5, the \textbf{Astral Rider cult fanatic} spends the next minute lost in their madness; they cannot answer questions (instead, they babble incomprehensibly or wail or stare into the distance), and spells like \textit{detect thoughts} reveal nothing but otherworldly shrieking. If a spell such as \textit{calm emotions} is used, reroll the d20.

\subsection{Encountering the Púca}
If the PCs venture into the woods and fields at night, they encounter the \textbf{púca}. A successful DC 18 Wisdom (Perception) check allows a PC to hear the clinking of the chains before the creature emerges from the trees.

\subsubsection{If the Púca Spots the PCs}
It approaches them at a walk, waiting for one of the PCs to climb onto its back. It only attacks the PCs if they are hostile.

\subsubsection{If the PCs Hide From the Púca}
It continues slowly on its way. Three rounds later, five \textbf{Astral Rider cult fanatic} wearing iron spurs arrive in the company of a \textbf{cult thaumaturge} (use statistics for \textbf{mage} with the changes noted below) atop a \textbf{leashed lesion}. The cultists attack the PCs if they intervene, but otherwise ignore them and chase after the \textbf{púca}.

\subsubsection{Creature Change}
The \textbf{cult thaumaturge} has darkvision (60 ft.), AC 16 (\textit{breastplate}), and the following changes to its spell list: replace \textit{fire bolt} with \textit{eldritch blast}; replace \textit{mage armor} with \textit{hellish rebuke}; replace \textit{misty step} with \textit{darkness}; replace \textit{fly} with \textit{hypnotic pattern}; and replace \textit{cone of cold} with \textit{hold monster}.

\subsubsection{If the PCs Do Not Intervene}
The \textbf{cult thaumaturge} uses \textit{hold monster} on the \textbf{púca}, then commands one of the \textbf{Astral Rider cult fanatic} to mount it. They then proceed back to their home base (see "The Barrow" section) with the creature.

\subsubsection{If the PCs Attack the Cultists}
The \textbf{púca} fights alongside the PCs (unless they attack it as well). If any of the \textbf{Astral Rider cult fanatic} mount the púca during the combat, they flee with the creature back to the barrow, having gotten what they came for.

If, at any time while fighting the cultists or the PCs, the \textbf{púca} drops to one-quarter of its health (42 hit points) or lower, it flees.

\subsubsection{Treasure}
Along with the items assigned to them by their statblock, the cultists each wear a pair of iron spurs.

\section{Visiting the Farms}
The PCs can visit any of the local farms to attempt to gain more information on the \textbf{púca} and/or the cultists. There are twenty farms within a day's travel of Derrymead. Any of the farmers questioned only know what Felicity knows, though they direct the PCs to Mell Rorsk—the town blacksmith, whose son encountered one of the cultists on the road.

\subsection{The Blacksmith}
Mell Rorsk is Derrymead's blacksmith. She lives not far from the town center, in a small house with her husband Jorun, and her 16-year-old son Auri and 14-year-old daughter Leena. Auri was struck with madness when he confronted one of the wandering cultists on the road three days ago. He hasn't recovered, and he now spends all his time in his bedroom, staring blankly out the window. If the PCs wish to ask any questions about Auri's encounter, Mell directs them to Leena, who was with the boy when he met the cultist.

A successful DC 12 Wisdom (Medicine) check determines that the madness afflicting Auri can be cured by the \textit{lesser restoration} spell and that, otherwise, he should recover after 7 days have passed.

\subsection{What Leena Knows}
Leena is an impulsive teenager, with a stubborn streak and a lingering sense of guilt over her brother's current condition. If asked about the encounter with the cultist, she tells them the following:

Leena and Auri were walking from their home toward town for market day when they spotted a cloaked, hooded figure in a copse of trees by the road. Having heard the fearful rumors about these hooded strangers lurking in the woods, Auri drew his sword and advanced on the figure. The man turned toward him, and a moment later Auri fell to the ground, clutching his head and screaming. Leena only got a glimpse of the man's face—it was pale and drawn, and his eyes were like a starry night sky. Rivulets of blood ran from his eyes and down his cheeks. The experience was terrifying, but she is angry and ready to fight—though her parents won't let her.

A successful DC 16 Wisdom (Insight) check reveals that Leena is not being completely honest about the situation. If pressed, a successful DC 14 Charisma (Intimidation or Persuasion) check forces her to break down and admit she's the one who advanced on the figure initially—her brother wanted to turn back and get their parents, but Leena wanted to confront the individual—and her brother paid the price for it. She knows nothing else.

\subsection{Iron Spurs}
If the PCs tell Mell they are hoping to drive off the \textbf{púca} as well as the cultists, she provides them with two sets of iron spurs. She tells them her grandmother always said the creature despised iron, and, that by wearing the spurs, the PCs might be able to control the \textbf{púca}. She apologizes—she'd offer more, but someone recently ransacked her forge and stole the other five sets of spurs she'd already made.

\section{The Barrow}
The barrow is an old burial site located roughly an hour west of Derrymead. It has been mostly forgotten by the locals, as it was sealed generations ago, but the cultists discovered it and have claimed it as their own.

The barrow has several features that are common throughout:

\subsubsection{Stone Doors}
The doors within are thick, gray granite. Each door has 17 AC, 27 hit points, and a damage threshold of 10. They can be pushed open with a successful DC 15 Strength (Athletics) check.

\subsubsection{Domed Ceilings}
The roof of each chamber is curved like the inside of a dome. Each ceiling is made of stacked stone, is 10 feet tall where it meets the wall, and rises 15 feet tall at its apex.

\subsubsection{Earthen Walls}
The walls are made of stacked stone, behind which lies the dirt of the mound itself. The walls and ceiling are reinforced with wooden beams.

\subsubsection{Torchlit}
Each of the three main rooms has torches mounted in sconces along the wall, every 15 feet or so. All the torches are lit, causing the main three rooms to be brightly lit. They are easily removed from the sconces.

An unlabelled version of the map can be found here.

\subsection{1. Entryway}
This hallway is 20 feet long and 10 feet wide and is bracketed by double doors. There is no light source here.

\subsubsection{Creatures}
If any of the \textbf{Astral Rider cult fanatic}, the \textbf{cult thaumaturge}, or the \textbf{leashed lesion} escaped from the earlier encounter (see the "Encountering the Púca" section), the PCs encounter them as they exit the barrow.

\subsection{2. First Ritual Chamber}
This round chamber measures roughly 40 feet in diameter. In the center of the room is a stone dais upon which sits a 3-foot-tall, 8-foot-diameter stone brazier. The bowl of the brazier looks to be lined with iron, and it is cold and unlit. Double doors to the north lead to 3, and four single doors lead to Areas 5, 6, 7, and 8.

\subsection{3. Second Ritual Chamber}
This room is nearly identical to 2, with one exception: A portion of the southeastern wall and ceiling has caved in, exposing the dirt of the mound above. Thin roots poke through the dirt like withered, pale fingers. Piles of stone and dirt lie beneath the holes; these areas are considered difficult terrain. Double doors to the north lead to  4, and four single doors lead to Areas 9, 10, 11, and 12.

\subsection{4. Third Ritual Chamber}
In this room, a portion of the ceiling to the east has fallen in, causing that area to be difficult terrain. Five single doors lead to Areas 13, 14, 15, 16, and 17.

\begin{DndReadAloud}{}
This circular chamber is nearly identical to the others, ringed by single stone doors. To the east, a portion of the ceiling has fallen in, and debris lies in a mound beneath the hole. Whereas in the other chambers, the braziers were cold, this one burns with a searing-bright green and purple flame, lighting several figures behind it in stark relief.

\end{DndReadAloud}

\subsubsection{Creatures}
Standing before the brazier on the far side of the room are three \textbf{Astral Rider cult fanatic} and the cult's leader Cae Kel Dranus (LE human \textbf{psychophant} cultist).

If the cultists captured the \textbf{púca}, it is present as well; if so, add the following to the read-aloud text above:

\begin{DndReadAloud}{}
And standing before you, bound by its own chains to the brazier, is the púca, its head hanging low.

\end{DndReadAloud}

\subsubsection{Cultist Conversation}
When the PCs enter the room, Dranus raises one hand to keep her cultists from attacking, and invites the PCs in to talk. If they do, she explains the following:


\begin{itemize}
\item She and her people are only here for the \textbf{púca}.
\item They have no quarrel with the farm folk or the PCs.
\item If the PCs killed any of her cultists, she magnanimously forgives them such transgressions.
\item She and her people will leave as soon as they have the \textbf{púca}—or, if they already have the \textbf{púca}, they'll leave as soon as she has it fully under her thrall.
\item If the PCs either (a) bring the creature to her if the cultists don't have it already, or (b) if they'll stop interfering with cult business, she offers them their pick of the grave goods within the barrow. If the PCs seem reluctant to agree, she casts \textit{charm person} or \textit{suggestion} on whoever seems to be the leader of the group; the \textbf{psychophant} spellcasting doesn't require any V, S, or M components, so it cannot be detected without magical means (such as an active \textit{detect magic} spell).
\end{itemize}

\subsubsection{Cultist Combat}
If the PCs attack, the \textbf{psychophant} and her cultists fight to the death.

\subsubsection{Chained Púca}
If the \textbf{púca} is present, it fights against the cultists if combat ensues. Due to the cultist's ritual, the \textbf{púca} cannot use its Beguiling Aura until it takes a long rest, and, because its own chains have bound it to the brazier, it can't use its Chain Whip attack. It can use its action to make one of its Hooves attacks, but it's chained tightly enough that its speed is 0, and its attacks are made with disadvantage.

\subsubsection{Freeing the Púca}
If a PC attempts to free the \textbf{púca}, combat ensues if it hasn't already. If it's freed, it attacks the cultists and is able to use its full array of attacks.

The shape of the brazier prevents the chains from being slid up over the top, and attacking the chains injures the \textbf{púca}. But the brazier itself is a viable target: It has 17 AC, 27 hit points, and a damage threshold of 10. If it is destroyed, the \textbf{púca} is freed. It's almost impossible to lift and/or flip the brazier over, but a PC can do so with a successful DC 30 Strength (Athletics) check.

The easiest way to free the \textbf{púca} is to quench the ritual flame in the brazier. The flame is producing a magical effect, but is not itself magical—so any methods that would normally quench fire successfully put it out; if a PC uses its action to use any reasonable method to quench the flame, they can do so with no rolls. Once the flame is quenched, the \textbf{púca} rears back and regains control over its chains, freeing itself of the brazier.

\subsection{5–17. Burial Chambers}
These chambers are all 15 feet long and 10 feet wide, with no light inside. At the far end of each is a simple stone sarcophagus that holds human remains and grave goods.

The sarcophagus lids are made of thick stone. They have 17 AC, 27 hit points, and a damage threshold of 10; they can be pushed aside with a successful DC 15 Strength (Athletics) check.

The contents of each sarcophagus follows:


\begin{itemize}
\item A skeleton wrapped in decayed leathers, a \textit{longbow}, a \textit{quiver} containing 10 +1 arrows, and a \textit{silvered dagger}.
\item A skeleton wrapped in moldering cloth and a ruby ring worth 300 gp. The ceiling in this chamber has mostly collapsed, and the room is filled with rubble. It is considered difficult terrain, and the debris must be removed in order to access the sarcophagus.
\item A skeleton wrapped in decaying furs and a silver decanter carved with intricate scrollwork of grapevines worth 150 gp.
\item A skeleton wrapped in stained crimson silk robes and a \textit{staff of healing}.
\item A mostly-decayed body wrapped in a \textit{rope of entanglement}. If the PCs remove the rope, a \textbf{wraith} emerges from the sarcophagus and attacks. A successful DC 18 Intelligence (Investigation) check reveals scuffs on the interior of the sarcophagus lid, indicating that individual was interred alive.
\item A mostly-decayed, small body and a wooden toy.
\item A mostly-decayed body in corroded \textit{chain mail}, a kite \textit{shield}, and a \textit{longsword of life stealing}.
\item A mostly-decayed body in gray and while silk robes and a pouch containing 25 pp.
\item A mummified corpse in decaying clothes, and a +1 cloak of protection draped over the corpse.
\item A mummified corpse in pristine \textit{+1 leather armor}.
\item A mummified corpse in musty furs and a pouch of gems worth 300 gp.
\item A mummified corpse in decaying cotton robes wearing a \textit{medallion of thoughts}.
\item A mummified corpse wearing a \textit{robe of useful items} that is missing one \textit{dagger} patch and one hempen rope patch, but also has the following additional patches: 10 gems worth 100 gp each, bag of 100 gp, iron door, \textbf{mastiff} (2), pit, a scroll of \textit{magic missile} (3rd level), \textit{rowboat}, silver coffer, \textit{potion of healing} (4), and wooden ladder.
\end{itemize}

\section{Concluding the Adventure}
The conflict with the cultists and the \textbf{púca} can be resolved in a number of ways.

\subsection{The Fate of the Cultists}
If the PCs kill the cultists, Felicity rewards them with the promised coin and turns a blind eye if they took any of the grave goods, though it would be exceedingly unlikely she would know that they did unless the PCs tell her.

If the PCs help the cultists capture the \textbf{púca}, the cultists remain near Derrymead for another week—the length of time needed for Cae Kel Dranus to fully charm the \textbf{púca}. Once they are gone, Felicity rewards the PCs with the promised coin but does not hide her displeasure. If the PCs took any grave goods, she works to turn the town's favor against them. After a few months, Derrymead falls into shadow and darkness (see "The Fate of Derrymead").

\subsection{The Fate of the Púca}
If the PCs kill the \textbf{púca}, the cultists leave Derrymead, as they have no reason to remain.

If the PCs redeem the \textbf{púca} and kill the cultists, the \textbf{púca} remains on the outskirts of Derrymead, helping the locals and any lost travelers. It is seen less often, but its appearance is thereafter considered good luck instead of something to fear, and the locals leave milk and honey out for it in hopes of luring it (and its good favor) to their door.

If the PCs free the \textbf{púca} but do not kill the cultists, the cultists remain in the barrow, determined to capture the creature. They are eventually successful, leading to Derrymead's fall into shadow and darkness (see "The Fate of Derrymead" below).

If the PCs did not kill or redeem the \textbf{púca}, the creature continues to haunt Derrymead. It ensnares several farmers and travelers, taking them on nightmarish rides but returning them unharmed (albeit exhausted and terrified). Eventually, Leena Rorsk tracks down the \textbf{púca}, rides it, and redeems it herself.

If the PCs kill or free the \textbf{púca}, Felicity provides them with the dwarven plate armor, regardless of the fate of the cultists.

\subsection{The Fate of Derrymead}
If the cultists manage to gain control of the \textbf{púca}, they leave Derrymead—for a time. With her newfound power, Cae Kel Dranus is able to travel to other twisted planes of existence and eventually returns to Derrymead to "enlighten" its populace. After six months, all the locals are either dead or transformed into \textbf{Astral Rider cult fanatic}.

\chapter{Down in the Dark of The Deep King's Domain}

\begin{DndTable}{c}

\textbf{An adventure for four to five characters of 12th level}\\
\end{DndTable}

\section{Adventure Background}
\DndDropCapLine{D}{}own deep in the depths of the world, among twisting natural caverns where bioluminescent fungi shine as motes of light in the darkness, lies the lair of the Deep King. A pervading gloom lingers there, an unseen weight pressing down upon mind and soul. Long ago, some forgotten crafters carved pillars into the walls of a great cavern as the heart of a new subterranean hold. The pounding of their mallets echoing through the dark brought forth the cave sovereign, a fell alien intelligence that consumed their spirits and transformed their bodies into mindless husks. Today, the Deep King hunts the world below, feeding on the spiritual essence of its victims and dragging their bodies back to its lair.


\section{Adventure Hooks}
The following can used to start the PCs on this adventure:


\begin{itemize}
\item \textbf{\textit{Legendary Blade.}} After famed adventurer Harnac Stormcrow failed to return from his apparently doomed quest to defeat the Deep King, the race is on for adventurers to recover Stormcrow's legendary blade and seize the hero's towering reputation for themselves.
\item \textbf{\textit{Cure the Prince.}} Every attempt to cure the king's child of a magical illness has failed, but there is a glimmer of desperate hope; following auguries and divinations, the PCs are hired to recover the only thing that will save the child, the heart-blood of the Deep King, a strange, powerful entity that lurks deep underground.
\item \textbf{\textit{Planewalkers.}} The PCs need to access a hard-to-reach plane. Luckily, they have learned the forgotten builders slain by the Deep King were planewalkers who concealed a still-active portal to the plane the PCs must reach in one of the pillars in the sovereign's throne room.
\end{itemize}

\section{The Deep King's Domain}
The fungi-streaked caverns that comprise the domain of the Deep King (CE \textbf{cave sovereign}) may at first appear benign, but the subterranean complex presents grave dangers to any who venture within.

\subsubsection{Temperature}
The temperature is warm and humid, perfect conditions for the bioluminescent fungus dotting nearly every surface. Water drips from the ceiling and tufts of slippery moss coat the floor, adding an additional obstacle to traversing the area.

\subsubsection{Light}
The caverns are not illuminated and are filled with darkness unless the PCs provide their own light.

\subsubsection{Ceilings}
The ceiling in the domain are 15 feet high, unless otherwise noted.

\subsubsection{Fungi}
The specimens of fungi found throughout the caverns are generally harmless, though some areas do have specific fungal dangers. The bioluminescent qualities of the fungi and their spores do not produce enough illumination to provide even dim light; instead, glowing motes hang on surfaces or drift in the air.

\subsubsection{Unsettling Atmosphere}
All living creatures feel an oppressive psychic weight, as if under constant observation while some unseen presence presses upon their minds. Due to this, completing a long rest while in the Deep King's domain is impossible, even with the aid of spells such as \textit{leomund's tiny hut} or \textit{mordenkainen's private sanctum}.

\begin{DndSidebar}[float=!b]{Domain Actions}
The Deep King can use the standard cave sovereign lair actions in 7. But it also has special "domain" actions that function in areas beyond which it can see; they are similar to the cave sovereign's lair actions, but work a little differently as noted below.



This entire cavern complex is considered the domain of the Deep King. It is aware of the presence of any living creature within its domain unless the creatures are protected by spells that block divination magic such as \textit{nondetection} or \textit{mind blank}. Otherwise, the Deep King detects the PCs as soon as they enter 1.



\subsubsection{Domain Actions}
The Deep King can use a domain action on any creature it is aware of in its domain. When the PCs enter each area (except 5), select or randomly choose a PC to experience one of the Deep King's domain actions.


\begin{itemize}
\item \textbf{\textit{Inciting Madness.}} The target must succeed on a DC 18 Wisdom saving throw or suffer a random short-term madness for 1d4 rounds.
\item \textbf{\textit{Altered Perception.}} The target must succeed on a DC 18 Wisdom saving throw or perceive the Area it is presently in as difficult terrain for 1d4 rounds.
\item \textbf{\textit{Dimmed Vision.}} The target must succeed on a DC 18 Wisdom saving throw or treat all areas of bright light as dim light, areas of dim light as darkness, and areas of darkness as magical darkness for 1d4 minutes. This doesn't affect the light shed by the Deep King's Deathlights.
\end{itemize}



\subsubsection{Not Illusory}
The domain actions alter the target's understanding and perception—they aren't illusory effects.



\end{DndSidebar}

An unlabelled version of the map can be found here.

\subsection{1. Entrance}
\begin{DndReadAloud}{}
As you descend, conditions in the tunnels grow steadily warmer and more humid, until the water drips from the rough ceiling and glowing motes from bioluminescent fungal spores hang in the pungent air.

\end{DndReadAloud}

\subsubsection{Tracks}
A successful DC 15 Wisdom (Survival) check discerns a number of tracks leading in and out of this area. The most numerous footprints appear to belong to Medium-sized humanoids that lead toward the northwest passage. Similar footprints mark the way to the eastern passage but are fewer in number. More disturbing are the Huge-sized, segmented claw marks that appear to go back and forth in both passages.

\subsection{2. Illusionary Terrain}
As the PCs travel through this area, the Deep King assaults their minds, forcing each PC to make a DC 18 Intelligence saving throw.

From the point of view of each PC, the terrain of the tunnel or cavern changes, becoming filled with dangerous hazards or that which the PC fears the most. Each PC sees a different change to the terrain. Moreover, from each PC's point of view, their companions either vanish or immediately perish from the changed terrain.

On a failed save, the target believes the changed terrain is real, and it takes 2d10 psychic damage each turn for 1 minute. The target can repeat the saving throw at the end of each of its turns, ending the effect on a success.

On a successful save, the target takes half as much psychic damage, and they discern the illusion for what it is, thereafter seeing it only as a vague image superimposed on the real terrain.

While a creature is under the sway of the illusion, the actions or influence of other PCs appear to be aspects of the illusionary terrain. Shouts become the roar of beasts or the clatter of stones, physical interactions become the attacks of foes, etc. If a PC under the influence of the illusion casts an offensive spell, or makes an attack, they randomly target the nearest creature, believing them to be part of the hazard they combat.

The illusion includes audible, visual, tactile, and olfactory elements. Suggestions include (but are not limited to) the following:


\begin{itemize}
\item The cavern collapses or the walls begin closing in.
\item Hordes of monsters attack.
\item Poisonous gas seeps from cracks in the walls.
\item Burning lava pours from the walls filling the area.
\item Stinging insects swarm from all directions.
\item The tunnel become the fang-filled mouth of some titanic beast.
\item The walls fall away to reveal an airless void.
\end{itemize}

\subsubsection{Truesight}
A creature with truesight can see through the illusion.

\subsection{3. Fungus Cavern}
\begin{DndReadAloud}{}
The tunnel spills into a circular cavern with a high, curved ceiling held aloft by a large, natural stone column in the center of the area. The walls are choked with glowing fungi, and the ground bears a carpet of moss studded with large clumps of fungus. A fog-like miasma of spores that carries the odor of rotting, moldy vegetation fills air. At the far side of the chamber, another passage leads north.

\end{DndReadAloud}

\subsubsection{Terrain}
Due to the slippery moss on the floor, the area is difficult terrain. During combat, a creature that moves more than half its speed on a turn must make a DC 12 Dexterity saving throw. On a failed save, the creature falls prone and their speed drops to 0 until the beginning of their next turn.

\subsubsection{Fire and Water}
This area is very humid and damp, so any attempt to burn out the fungus will not work without liberally covering the area with a flammable accelerant such as oil or naphtha.

\subsubsection{Creatures}
Among the clumps of fungus are four fungal mounds (use statistics for \textbf{shambling mound}). They are indiscernible from normal fungus until they move.

\subsubsection{Treasure}
Among the other fungi in the cavern are three witch's gift mushrooms, which are gold-colored sac fungi with black streaks; eating a \textit{witch's gift} within 1d4 days of its harvesting provides the same effect as a \textit{potion of superior healing}.

\subsection{4. Ossuary}
\begin{DndReadAloud}{}
The passage opens into a cavern where a carpet of bones arranged in decorative patterns covers the floor and walls. The remains range in size from tiny finger bones to huge skulls and show varying states of decay. A dozen arches formed of skulls and set into the walls around the cavern hold the decomposing forms of large bull-headed humanoids. The opening of another passage leads off toward the northeast.

\end{DndReadAloud}

In the shadows of the northwestern corner of the chamber is a \textbf{duskwilt}. The \textbf{duskwilt} is the willing thrall of the Deep King and tends to the ossuary like an overprotective artist.

The \textbf{duskwilt} pride and joy are the displayed corpses in the skull arches. The dozen minotaur were the original builders who accidentally awakened the Deep King during their construction.

\subsubsection{Creatures}
Aside from the \textbf{duskwilt}, the skull arches hold eight \textbf{minotaur skeleton} who animate and move to attack as soon as the PCs enter the chamber. The skeletons act on their own but will not attack the \textbf{duskwilt}.

\subsubsection{Treasure}
On the mummified finger of one of the minotaurs is a \textit{ring of shooting stars}.

\subsection{5. Overlook Tunnel}
\begin{DndReadAloud}{}
This cavern tunnel appears less trafficked than the other areas you've encountered. Thin, gray metallic lines run through the stone floor and walls while glowing fungi cling to various surfaces as the passageway slopes gradually upward.

\end{DndReadAloud}

The Deep King's domain actions do not function within this area.

This tunnel leads to a small ledge high above the Deep King's throne room in 7. The Deep King does not use this passageway often as the veins of lead running through the walls partially block its psionic awareness.

\subsubsection{Lead-Lined}
The metallic lines in the floor and walls appears to be a heavy, shiny gray metallic ore. A successful DC 10 Intelligence (Nature) check allows a PC to identify the substance as lead.

\subsubsection{Overlook}
The curving passageway leads to a small ledge 25 feet above the floor of Area 7. PCs observing the room from the overlook easily spot the Deep King (see 7). It does not notice the PCs until they do something to attract its attention.

\subsection{6. Larder}
\begin{DndReadAloud}{}
The stench of rotting flesh fills this long cavern as more than a dozen animated corpses shuffle around the southern end of the room. An untold number of remains lie trampled on the floor, including a few that still occasionally twitch. Beyond, a golden glow emanates from the short passage to the south.

\end{DndReadAloud}

Here, the past victims of the Deep King gather to await their master's orders.

\subsubsection{Terrain}
The crushed remains of no-longer animated zombies makes for treacherous footing and the entire floor is considered difficult terrain.

\subsubsection{Domain Actions}
Any PC targeted by one of the Deep King's domain actions in this location has disadvantage on the saving throw.

\subsubsection{Creatures}
Guarding the passageway to the south are five \textbf{ogre zombie}, five \textbf{zombie}, and five \textbf{ghoul}. The undead try to keep any foes from entering 7, and all fight to the death. These creatures all wander about this area, and at any given time there's at least one located in the passages west and north of Area 6, where they would see PCs passing by.

\subsection{7. The Deep King's Hall}
\begin{DndReadAloud}{}
A knee-high swirling mist of water vapor and fungal spores lies thick upon the ground in this wide, 30-foot-high cavern. The walls have been cleverly carved into closely spaced towering columns. Two floor-to-ceiling pillars are carved in the shape of stacked bulls' heads, though the features have been defaced, scratched, and marred. A soft, golden light fills the cavern from a longsword stuck point first into a heap of gold and other treasures in the middle of the chamber.

\end{DndReadAloud}

\subsubsection{Creatures}
The Deep King (a \textbf{cave sovereign}) lurks on the ceiling between the two northeastern entrances to its lair. A successful DC 18 Wisdom (Perception) check spots the creature. If the PCs enter from 6, the Deep King waits until the majority of the PCs have entered the room, then drops next to the last one to enter and attacks. When it appears, or when combat begins, read or paraphrase the following:

\begin{DndReadAloud}{}
Suddenly, a hulking creature leaps into view, landing on massive insectoid limbs. It is a huge serpentine horror that is some aberrant hybrid of crustacean, arachnid, and eel. The soft radiance of its numerous antennae intensifies as its deadly maw opens to reveal multiple rows of semi-translucent teeth, teeth as large and sharp as swords.

\end{DndReadAloud}

The Deep King fights to the death.

\subsubsection{Glowing Sword}
The glowing sword is a sentient magic weapon (see "Akefalos"). Presently, it wishes nothing more than to kill the Deep King, who slayed and devoured his beloved former owner, the great hero Harnac Stormcrow; as such, it begins vibrating slightly with excitement when the PCs enter Area 7 and can attune instantly to any creature it feels is capable of wielding it to destroy the Deep King. If none of the PCs are likely to want to use a \textit{longsword}, change its initial form to some other bladed melee weapon and adjust the read-aloud text accordingly.

\subsubsection{Sentience}
Akefalos is a sentient neutral good weapon with an Intelligence of 12, a Wisdom of 14, and a Charisma of 17. It has hearing and normal vision out to a range of 60 feet. The weapon cannot speak, but can convey emotions to the creature attuned to it. Its special purpose is to destroy aberrations.

\subsubsection{Treasure}
The Deep King's hoard is a strange, eclectic assortment of coins and items gathered from its many victims. It includes: 20,142 cp, 8,321 sp, and 12,095 gp; 12 containers of various trade goods (1 ton, worth 1,500 gp total); various art works (worth 2,500 gp total); an \textit{arrow of slaying}; a scroll of \textit{dominate monster}; and a \textit{potion of vitality}.

\begin{DndSidebar}[float=!b]{Managing Difficulty}
The encounter with the Deep King is very dangerous (particularly if the party only has 4 PCs), so it's essential to have the 6 encounter before the PCs encounter the Deep King, as having those additional foes would make the combat too deadly. If the PCs manage to get to Area 7 without dealing with the 6 creatures, perhaps the creatures are the thralls of the Deep King but will not step foot in Area 7.



If, on the other hand, it seems like the PCs are doing extraordinarily well against the Deep King, you can make some of the remains in 6 reanimate as one or more of those creatures to join the fray.



\end{DndSidebar}

\section{Concluding the Adventure}
The adventure ends when the PCs defeat the Deep King. Victory is sweet but may prove fleeting as transporting the hoard of treasure to the surface may attract many unwelcome encounters. If the PCs were hired to defeat the monster or bring back a portion of it, they are likely due a significant reward for their daring-do. If this adventure is placed in the context of a larger campaign, a scroll found among the treasures of the Deep King may point the PCs toward the next portion of their quest.

\end{document}
